% !TEX program = lualatex
% Continuous Chinese translation without facing-page English source plates.
\documentclass[UTF8,11pt,a4paper,openany,fontset=fandol]{ctexbook}
\usepackage[margin=2.35cm,headheight=25pt]{geometry}
\usepackage{amsmath,amssymb,bm,mathtools}
\usepackage{fancyhdr,microtype,xcolor,hyperref,xurl,needspace}
\setlength{\parindent}{2em}
\setlength{\parskip}{0.55em}
\setlength{\emergencystretch}{3em}
\linespread{1.22}\selectfont
\clubpenalty=10000
\widowpenalty=10000
\displaywidowpenalty=10000
\raggedbottom
\pagestyle{fancy}
\fancyhf{}
\fancyhead[C]{\small\nouppercase{\leftmark}}
\fancyfoot[C]{\thepage}
\renewcommand{\headrulewidth}{0.3pt}
\hypersetup{hypertexnames=false,colorlinks=true,linkcolor=black,urlcolor=blue!60!black,pdftitle={计算流体动力学前沿（2006）},pdfauthor={D. A. Caughey、M. M. Hafez 主编}}
\title{计算流体动力学前沿（2006）\\[0.45em]\Large 中文译本}
\author{D. A. Caughey、M. M. Hafez 主编\\\small 中文初译经 CFD 专业术语表统一}
\date{}
\begin{document}
\frontmatter
\maketitle
\chapter*{译本说明}
本译本按原书章节顺序翻译并连续排版，正文不附英文原文，供中文阅读和检索使用。作者姓名、文献题名、专业缩写以及数学符号按学术写作惯例保留必要的拉丁字符。

中文初稿由离线神经机器翻译批量生成，并通过 CFD 专业术语表统一 Navier-Stokes 方程、有限体积法、伴随方法、多重网格、预处理、湍流、磁流体动力学等词汇。由于全书篇幅大且原 PDF 文本层存在断行与公式混排，本文件属于可继续校订的学习译本，不替代正式授权译本。
\tableofcontents
\mainmatter
前沿F 计算流体动力学2006

此页面故意留空

由D.A.编辑的COFm p u ta t i o na I Fluid Dynamics 2006 考伊康奈尔大学理学硕士 哈菲兹加州大学戴维斯分校rp世界科学新泽西州*伦敦新加坡BELJlNG上海-香港台北*金奈

由世界科学出版公司出版。 重新。 有限公司 5 Toh Tuck Link, Singapore 596224 美国办公室：27 Warren Street, Suite 401-402, Hackensack, NJ 07601 英国一次：伦敦考文特花园谢尔顿街57号 WC2H 9\% 大英图书馆编目出版数据 本书的目录记录可从大英图书馆获得。 前沿计算流体动力学 2006 版权所有 0 2005 世界科学出版公司。 重新。 有限公司 保留所有权利。 未经出版商书面许可，不得以任何形式或任何方式（电子或机械）复制本书或其部分，包括复印、记录或任何已知或将要发明的信息存储和检索系统。 如需复印本卷材料，请通过版权清关中心支付复印费，地址为222 Rosewood Drive, Danvers, MA 01923, USA。 在这种情况下，不需要出版商的复印许可。 国际标准书号981-256-527-2由世界科学打印机（S）私人有限公司在新加坡印刷

\chapter*{奉献}
奉献 本卷由献给David A教授的论文组成。 Caughey在他60岁生日之际。 大多数论文是在2004年6月22日至24日在纽约伊萨卡康奈尔大学举行的题为“计算未来IV”的研讨会上发表的。 作者是David的朋友和同事，很高兴这本书献给他，以感谢他对该领域的贡献。 大卫于1944年3月5日出生在密歇根州的大急流城。 他于1965年在密歇根大学获得了航空和航天工程理学学士学位，以及他的硕士学位。 （1967）和博士学位。 （1969）普林斯顿大学的航空航天和机械科学学位。 1969年10月至1970年8月，他是莫斯科苏联科学院计算中心的交流科学家。 他在圣的麦克唐纳道格拉斯研究实验室工作。 1971年至1975年，密苏里州路易斯担任科学家。 他于1974年加入康奈尔大学，担任西布林机械与航空航天工程学院的客座助理教授，次年转为正式终身助理教授，并于1984年成为正教授。 从1993年到1998年，他担任西比利学校的院长。 在过去的30年里，他教授了本科和研究生的几门课程，包括流体力学、航空学、可压缩和不可压缩空气动力学、航空航天推进、飞行动力学和计算空气动力学。 他监督了18名博士生、5名硕士生和23名工程硕士生的设计专案。 他在普林斯顿大学（1981年）、美国宇航局艾姆斯研究中心（1989年）、莱特-帕特森空军基地的空军研究实验室（1998年）和英国威尔士斯旺西的威尔士大学度过了休假。 （1999年和2004年再次）。 大卫很幸运能够与一系列才华横溢的教师和研究人员进行学习和合作。 他是普林斯顿大学华莱士·海耶斯教授的学生，他与奥列格·S教授一起工作。 Ryzhov在莫斯科，与博士。 McDonnell Douglas研究实验室的Raimo Hakkinen，普林斯顿的Antony Jameson教授，最近康奈尔大学的Stephen Pope教授。 在休假期间，他访问了

Ui奉献博士 美国宇航局艾姆斯的Gil Chyu博士 莱特·帕特森的约瑟夫·尚和斯旺西的肯·摩根、欧贝·哈桑和奈杰尔·韦瑟里尔教授。 迄今为止，他已经发表了约90次受邀讲座和研讨会，撰写了135篇论文，并且他仍然非常积极地从事研究和教学。 他曾担任过许多公司和机构的顾问。 1989年至1991年，他是AIAA杂志的副编辑，1988-1991年是《国际杂志：大气与海洋物理学》英文翻译的技术编辑。 他曾是国家科学基金会、能源部以及该领域许多期刊的审稿人。 作为一名学生、一名教师和一名研究人员，David获得了许多奖项和荣誉。 1965年，他获得了密歇根大学航空航天工程系颁发的杰出成就奖。 他于1967-1968年获得美国宇航局的研究生实习，并于1965-1967年和1968-1969年获得NSF研究生奖学金。 1977年，他获得了康奈尔工程师协会和康奈尔分会Tau Beta Pi颁发的工程教学卓越奖。 1979年，他获得了美国航空航天研究所（AIAA）颁发的劳伦斯·斯·斯佩里奖，“以表彰他对跨声速流动场关于实际感兴趣的复杂配置的高效数值计算做出的杰出贡献。” 他还于1980年获得了美国宇航局兰利研究中心颁发的优秀证书，以表现在美国宇航局技术简报出版物的主题，该出版物题为“FLO-22-跨声速流动经过扫翼的数值计算”。 他于1994年当选为AIAA院士，“以表彰对理论的开创性贡献和 计算空气动力学的应用。” 1998年，他获得了美国土木工程师协会颁发的阿奇·卡特出版奖，“以表彰《流体力学：互动文本》的发展”，该奖利用最先进的技术教授大学生。 同年，他还获得了美国协会高管协会颁发的技术创新卓越一等奖，“以表彰“开创了工程流体力学第一本电子教科书的开发”。 2003年，David获得了AIAA的特别服务奖章，“以表彰他在康奈尔大学担任教师顾问，以及对AIAA和行业的技术贡献。” 最后，David特别自豪地在2005-2010年期间被任命为威尔士大学斯旺西工程学院的名誉教授， 并感谢Ken Morgan、Nigel Weatherill和Oubay Hassan多年来的友谊和支持。 本书包括在研讨会上发表的论文：计算未来IV。 这本书还包括David无法参加研讨会的朋友和同事的几篇论文。 第一章专门介绍David Caughey对CFD的贡献。 其余章节分为五个部分。 第一部分涵盖了设计和优化的主题，包括：对工程系统（道路隧道和硬盘驱动器外壳）的应用；

奉献Uii在气动外形优化中取得了进展；螺旋桨叶片的设计和优化；以及边界条件用于模拟自适应和优化配置上的不稳定流动。 第一部分涵盖了基本算法和准确性的主题，包括：隐式、基于残差、紧凑方案的稳定性和效率；用于动态、非结构网格模拟的高阶时间积分；电磁学的显式、时域、有限元方法；CFD解中网格诱导误差的估计；以及使用涡量约束处理涡流。 第11部分涵盖了流动稳定和控制的主题，包括：使用合成和脉冲喷射的流动控制；使用电磁场的分离控制；跨声速流动在扁平侧翼上的分叉行为；涡流对在纤细的锥体上的稳定性；上游条件对湍流边界层分离的影响；以及高超音速、磁流体动力学相互作用。 第四部分涵盖了多相和反应流的主题，包括：使用AUSM+-up方案计算多相流；复杂几何中多相流的有限体积界面追踪法；以及湍流火焰的计算建模。 最后一部分，第五部分，包含一章关于教育，副标题是教育改革在空气动力学中的影响，反映了编辑和整个社群所欣赏的这一主题的重要性。 最后，David Caughey的朋友和同事们祝愿他继续他成功的职业生涯，充满贡献和成就。 穆罕默德·M。 哈菲兹·戴维斯，加利福尼亚州，2005年3月

大卫·A。 康奈尔大学机械与航空航天工程教授Caughey

\chapter*{内容}
内容奉献V 1。 David Caughey对计算M的贡献。 M。 哈菲兹流体力学 1 1.1 1.2 1.3 1.4 1.5 1.6 1.7 1.8 1.9 1.10 1-A 1-B 介绍................................ 1 激波结构和声爆............ 2个势流模拟...................... Euler 方程的3个解决方案.................. Navier-Stokes 方程的4个解决方案.............. 10 湍流反应流的模拟............ 13个专题............................... 14篇评论文章....................... 15 流体力学：交互式文本............. 16 结束讲话....................... 17 博士 由David A监督的学生。 Caughey...... 19 David A的出版物。 Caughey............... 21 我。 设计和优化2。 计算流体动力学工程系统分析和设计37 M。 达莫达兰。 S。 阿里和S。 Dayanandan 2.1 介绍................................ 37 地下道路隧道规划............ 38 硬盘驱动器外壳中的流程建模 43 2.4 结论性评论...................... 46 2.5 参考书目.............................. 46 2.2 2.3 消防控制策略和场景的流程建模......

X内容3。 气动外形优化 A的进步。 詹姆森 3.1 3.2 3.3 3.4 3.5 3.6 3.7 3.8 3.9 3.10 介绍............................ 49 优化程序的制定......... 52 3.2.1 梯度计算.................. 52 使用Euler 方程的设计............... 55 减少梯度配方.............. 61 优化程序..................... 63 3.5.1 梯度定义中索博列夫内积的必要性............ 63 3.5.2 外形优化的索博列夫梯度...... 65 3.5.3 设计程序大纲............ 66个案例研究............................. 67 3.6.1 跨声速翼型设计的二维研究 67 3.6.2 B747欧拉平面结果.............. 69 3.6.3 超级B747...................... 71 超级P51赛车................................ 71 结论............................ 75 确认....................... 76 参考书目........................... 76 3.7.1 外形优化用于Transonic商务喷气式飞机... 73 4. 螺旋桨叶片L的设计优化。 Martinelli和J.J. Dreyer 4.1 介绍................................ 81 4.2 制定作为控制问题............... 82 4.2.1 螺旋桨叶片的成本函数........ 83 4.2.2 搜索程序.................. 84 4.3 实施................................ 86 4.4 低空穴叶片部分的优化...... 86 4.4.1 与水隧道测量的比较...... 90 4.5 结论................................ 94 4.6 参考书目.............................. 95 5. 4D H中的流程边界条件建模。 优化的Sobieczky。 自适应和不稳定配置 97 5.1 介绍............................ 97 5.2 CDimensional问题的几何概念...... 98 5.3 优化............................... 101 5.4 自适应配置..................... 101 5.5 不稳定 边界条件................. 102

内容 Xi 5.6 Bic流体力学应用............:..... 102 5.7 结论............................ 104 5.8 参考书目.............................. 104 I1。 算法和准确性6。 基于隐性残基的紧凑方案的稳定性和效率107 C。 Corre和A。 Lerat 6.1 介绍................................ 107 6.2 隐含计划描述.................. 108 6.3 直接求解器 效率...................... 112 6.4 隐性治疗描述................. 114 6.6 结束讲话...................... 123 6.7 参考书目.............................. 126 6.5 迭代求解器 效率和稳定性............ 120 7. 高阶时间积分方案€或动态非结构网格 CFD模拟129 D。 J。 Mavriplis和Z。 杨7.1摘要............................... 129 7.2 介绍................................ 130 7.3 控制方程在任意-拉格朗日-尤拉（ALE）形式和基础流动求解器................... 131 7.4 高阶时间积分和离散几何守恒律..................... 132 7.5 网格运动策略..................... 134 7.5.1 张力弹簧类比................. 135 7.5.2 线性弹性类比................. 135 7.6 加速策略...................... 137 7.7 网格运动结果...................... 网状运动方程的137 7.7.1 收敛...... 139 7.8 使用反向差分方案的非定常流动模拟.............................. 141 7.8.1 多重网格收敛 效率............ 141 7.8.2 时间准确度验证................ 144 7.9 动态网格问题的隐式Runge-Kutta方法。 147 7.10 结论............................ 151 7.11 致谢...................... 152 7.12 参考书目.............................. 152 7-A BDF3的几何收敛定律........... 155

Xii内容8。 电磁学K的显式时域有限元方法。 摩根。 M。 El Hachemi。 0。 Hassan和N。 Weatherill 8.1 8.2 8.3 8.4 8.5 8.6 8.7 8.8 介绍............................ 电磁散射................... 8.2.1 控制方程................... 8.2.2 边界条件................... 网格生成......................... 数值解算法................ 8.4.1 时间离散化.................... 8.4.2 空间中的离散化.................. 8.4.3 计算细节..................... 8.5.1 PEC球体........................ 8.5.2 PEC杏仁....................... 处理更大的电子散射器......... 8.6.1 更高阶的Taylor-Galerkin 时间推进方案。 8.6.2 高阶空间离散化............. 结论............................ 数字示例....................... 参考书目.............................. 9。 估计CFD解决方案183 T中的网格诱导误差。 1. P。 Shih 9.1 介绍................................ 183 9.2 方法分类.................. 184 9.3 离散误差传输方程概述...... Euler 方程的FV解决方案的186 9.4 DETE...... 解决方案的188 9.4.1 有限体积法.......... 189 9.4.2 有限体积法的DETE......... 191 9.5 DETE的有用性.................. 192 9.5.1 测试问题1：无黏流动在翼型上...... 192 9.5.2 测试问题2：黏性流动在冰翼型上.. 193 9.6 最后讲话...................... 195 9.7 参考书目.............................. 196 10. 使用涡量约束 199 J处理伏特流。 Steinhoff和N。 Lynn 10.1 摘要.............................. 199 10.2 介绍................................ 200 10.2.1 基本概念..................... 201 10.3 说明性一维示例............. 204 10.4 涡量约束...................... 207

目录 10.5 10.6 10.7 10.8 10.9... xaaa 10.4.1 基本配方.................. 210 10.4.2 VC2配方与传统不连续性增深方案的比较。 214 10.4.3 结果.............................. 217 10.5.2 气缸唤醒...................... 218 其他研究............................ 220 10.6.1 导弹基础阻力计算............ 220 10.6.2 叶片涡流相互作用（BVI）............ 220 10.6.3 结论............................ 221 致谢...................... 222 参考书目.............................. 223 VC2配方的计算细节... 215 10.5.1 翼尖................... 217 10.5.3 动态摊位...................... 219 湍流 特效模拟...... 221 I11。 流量稳定和控制11。 合成和脉冲喷射器R的流量控制应用。 K。 阿加瓦尔。 J。 瓦迪洛。 Y。 棕褐色。 J。 崔。 D。 郭。 H。 耆那教 A。 W。 Cary和W。 W。 鲍尔11.1摘要.............................. 241 11.2 介绍................................ 242 11.3 CFD 流动-Solvers 采用..................... 243 11.4 结果和讨论..................... 使用合成喷射执行器的翼型的244...... 244 用合成喷射器对初级喷射器进行矢量控制............................ 247 使用合成喷射器控制向后台阶后面的再循环流区域...... 250 合成喷射与平板湍流边界层的相互作用............... 250 使用脉冲吹制控制亚音速腔剪切层......................... 254 11.5 结论................................ 260 11.6 致谢......................... 260 11.7 参考书目.............................. 262 使用电磁场在圆形圆柱体上控制流动分离：数值模拟B。 H。 丹尼斯和G。 S。 Dulikravich 12.1 命名法........................... 265 11.4.1 虚拟空气动力学形状修改 11.4.2 11.4.3 11.4.4 11.4.5 12。

Xiv 目录 12.2 导言.............................. 266 12.3 EMHD的二阶分析模型.......... 268 12.4 最小二乘有限元法............. 269 12.4.1 简化EMHD的非维一阶表格...................... 270 12.4.2 准确性验证................ 273 12.5 数值结果...................... 274 12.6 结论................................ 277 12.7 致谢...................... 278 12.8 参考书目........................... 278 跨声速流动在扁平翼型 285 A上的分叉。 Kuz'min 13.1 介绍................................ 285 13.2 问题陈述和数值方法......... 286 13.3 升力系数作为M函数的分析，...... 286 13.4..... 289 13.5 结果总结..................... 290 13.6 结论............................ 290 13.7 参考书目............................... 291 尤拉计算对涡流对在纤细圆锥体上的稳定性研究 297 J。 蔡，H-M。 蔡，S。 罗和F。 刘13。 14变化的稳定性分析。 14.1 14.2 14.3 14.4 14.5 14.6 14.7 14.8 摘要.............................. 297 介绍................................ 298 欧拉求解器和流动模型............. 302计算网格和边界条件...... 303 静止对称和不对称解决方案 14.5.1 时间不对称扰动......... 306 14.5.2 静止对称涡流........... 307 14.5.3 静止对称涡流的稳定性。 308 14.5.4 静止不对称涡流的稳定性 310 14.5.5 不对称涡流的镜像... 311 14.5.6 当前尤拉求解器的对称性质...... 313 14.5.7 与稳定性理论预测的比较 314 14.5.8 与稳定性实验数据的比较。 314 涡核的结构..................... 315 14.6.1 计算结果................. 315 总结和结论...................... 322 参考书目........................... 323及其稳定性...................... 306 14.6.2 与实验数据的比较......... 320

内容xu 15。 上游条件对全球分离0时湍流边界层速度缺陷曲线的影响。 S。 Ryzhov 15.1 介绍................................ 329 15.2 奇异隐形压力梯度............... 330 15.3 控制方程....................... 331 15.4 隐形子层1........................ 332 15.5 外部 湍流 子层 2................... 333 15.6 外层湍流 子层 3................. 333 15.7 压力主导的流动模式................ 334 15.8 与实验的比较.................. 336 15.9 结论................................ 336 15.10 参考书目.............................. 337 16. 高超音速磁流体动力学相互作用341 J。 S。 Shang 16.1 16.2 16.3 16.4 16.5 16.6 16.7 16.8 16.9 16.10 摘要.............................. 341 命名法.............................. 342 介绍................................ 342 控制方程...................... 344 等离子模型...................... 346 电-流体-动态相互作用............... 349 磁流体动力学 互动.............. 354 结束讲话....................... 359 致知................................ 361 参考书目.............................. 361 IV。 多相和反应流17。 使用AUSM+-up方案367 M.-S.计算多相流量。 Liou和C.-H. Chang 17.1 摘要.............................. 367 17.2 介绍................................ 368 17.3 控制方程（多相流模型）...... 369 17.3.1 热力学平衡模型 [14]...... 369 17.3.2 双流体模型.................... 373 17.3.3 多相分层流体模型.......... 376 17.3.4 对流通量，Ff3kl/2............... 379 17.3.5 压力通量，F:j+l和FP............ 380 17.3.7 时间积分.................... 381 2.3 17.3.6 界面压力修正期限...... 381

Xvi 目录 17.4 计算示例和讨论............... 383 17.4.1 赎金的水龙头问题................ 383 17.4.2 空气-水激波管问题............. 385 17.4.3 休克-泡泡相互作用问题.......... 386 17.5 结束讲话....................... 389 17.6 致谢...................... 391 17.7 参考书目.............................. 391 17-A 数字通量公式...................... 394 17.4.4 冲击水柱相互作用问题...... 389 18. 用于计算M的有限体积界面追踪法。 Muradoglu 18.1 介绍................................ 395 18.2 数学表述.................... 397 18.3 数值法...................... 399 18.3.1 流动方程的积分............ 400 18.3.2 界面追踪法................. 402 18.3.3 整体解决方案程序............ 404 18.4.1 振荡跌落..................... 406 18.4.2 直通道中浮力驱动的下降下降 407 18.4.3 连续收缩通道中浮力驱动的上升下降...... 410 18.4.4 一滴混乱的混合......... 413 18.5 结论................................ 414 18.6 参考书目........................... 415 18-A 斯托克斯极限的最佳人工可压缩性.... 419 复杂几何中的多相流动 395 18.4 结果和讨论...................... 405 穿过蜿蜒通道19。 湍流 Flames S.B.的计算建模 教皇 19.1 介绍............................ 421 19.2 湍流火焰的PDF计算.............. 422 19.2.1 驾驶喷气火焰................... 423 19.2.2 在Vitiated Co-Flow中解除喷气式火焰...... 423 19.3 湍流混合的建模................. 425 19.4 致知...................... 427 19.5 参考书目.............................. 427

目录xvii V。 教育20。 教育未来：教学改革的影响D。 L。 空气动力学中的Darmofal 433 20.1 20.2 20.3 20.4 20.5 20.6 20.7 20.8 20.9 介绍................................ 433 课程概述......................... 434 概念理解和主动学习...... 435 理论的整合。 计算。和实验...... 438 Project-baed Learning...................... 438 20.5.1 军用飞机设计项目............ 439 20.5.2 混合翼体设计项目.......... 440个结果.............................. 440 20.6.1 教学的有效性............... 440 20.6.2 课前作业的影响............ 442 20.6.3 学生评论................... 444 展望.............................. 445 致谢........................ 446 参考书目.............................. 446

此页面故意留空

\clearpage
\chapter{第1章 David Caughey对计算流体动力学 Mohamed M的贡献。 哈菲兹}
\section{1.1 介绍}
第1章 David Caughey对计算流体动力学 Mohamed M的贡献。 Hafez' 1.1 介绍 在本章中，我们简要讨论了David Caughey教授对计算流体动力学（CFD）的贡献。 他的出版物涵盖了该领域的各个方面：网格生成、离散化方案、快速求解器以及并行计算。 他开发了解决外部和内部空气动力学的二维和三维、稳定和非稳定、可压缩和不可压缩流动的电位、欧拉和Navier-Stokes 方程的方法。 他模拟了机翼、翼身和翼尾组合、机舱、旋翼和风扇上的超音速流。 他对应用和理论工作都感兴趣；例如，虽然他的代码在工业中使用，但他研究了Euler 方程在跨声速系统中非唯一解决方案的稳定性！ 他还研究过超声速流动、声爆传播和冲击噪声。 除了空气动力学之外，David还研究了空气弹性，以及求解概率密度函数（pdf）方程的技术，以模拟湍流，反应流动。 他在关于计算流体和空气动力学的书籍中写了几篇评论文章和章节。 我们还将涉及他的教学活动，包括使用CFD作为流体力学教学的教育工具。 '加州大学戴维斯分校机械与航空工程系，加利福尼亚州戴维斯，95616。

\section{1.2 激波 结构和 声爆}
2 Caughey对CFD的贡献 本章末尾包含David出版物的完整列表。 讨论中只提到少数几个。 本章分为几个部分，涉及他的研究活动，一个部分涉及他的教学和他介绍的教学工具。 1.2 激波结构和声爆 Caughey教授的博士论文是1969年在Wallace D教授的监督下，题为“超音速流中的二阶波结构”。 海耶斯在普林斯顿大学。 它作为美国宇航局承包商报告（CR-1438）发布。 该论文涉及非黏性超声速流动经过薄体的二阶非线性扰动理论，重点是波系统在大距离上的行为。 研究表明，一阶非线性理论是统一有效的渐近展开中的第一个项，并证实了由Lighthill修正和扩展的平面流的Friedrichs结果。 然而，在后一项工作中指出了一个错误，其中在非常远的距离上的激波位置被错误地给出。 对于有限物体的流，匹配的渐近展开与区域性解一起使用，以提供内部边界条件。 事实证明，完整的二阶理论在圆锥流的情况下比一阶理论有了明显的改进（其中重要的是包括在领先激波附近的非线性效应）。 Caughey教授和他的顾问就二阶波结构和声爆相关的理论问题写了几篇论文。 1971年，他用俄语写了一篇关于“根据非线性理论超音速理想气体流动中的波阻”的论文。 1968年，他还在田纳西州图拉霍马的田纳西大学航天研究所的“激波的生成和传播与声爆的应用”的短期课程中发表了演讲，在那里他探索了黏度在定义激波在大气中传播的结构中的作用。 他将射线几何视为已知的，并研究了通过广义形式的Burgers 方程在不同面积的导管下传播的弱冲击问题。 在特殊情况下，Hopf-Cole变换将Burgers 方程还原为线性热方程，这导致许多感兴趣的情况的闭合形式解。 研究表明，对于压缩脉冲，波的形状保持不变，而对于整个对称N波，形状确实随着衰减而改变。 还讨论了其他耗散机制。 虽然这些论文不是以CFD为导向的，但很明显，它们背后的分析有助于为以后的发展做好准备。 例如，对超声速流动使用高阶非线性扰动电位补充了Hayes的二阶跨声速流动理论。 速度势的存在以及基于克罗科关系的熵和涡量的大小，以及远场特征的使用在Caughey教授后来的作品中重新出现。 特别是，在波系统中使用质量和动量通量守恒，以及

\section{1.3 势流模拟}
M。 M。 波阻的哈菲兹3方程是非常有用的练习。 1.3 势流模拟Caughey教授写了大约50篇论文，涉及基于（非线性）势流方程的跨声速流动模拟。 1972年，他引入了“新的跨声速小扰动理论”，他使用（复杂）不可压缩势函数的实部分和虚部分作为独立变量。 小扰动理论是通过扰动电位的多重极限过程开发的，测量与不可压缩势流的偏差。 该理论结合了翼型的共形映射到更简单的规范（矩形）域的几何优势，允许在物体表面应用精确的非线性边界条件，同时保持小扰动公式的简单性和计算效率。 超音速翼型问题获得的解决方案在准确性上与全势方程的解决方案相当。 基于此近似值的计算机程序记录在麦克唐纳道格拉斯公司报告中。 1976年，Caughey教授在《飞机杂志》上发表了一篇关于“跨声速板条机翼的非黏性分析”的论文，该论文进一步扩展了这一理论，与实验数据的一致性相当好（对于对板条的黏性影响不太重要的情况）。 1974年，Caughey教授开始与Antony Jameson教授就FLO-22进行富有成效的合作，这是一个基于非保守形式的全势方程的跨声速流动通过扫翼的数值模拟的计算机代码。 该代码在研究实验室和工业界被广泛使用（今天它仍然用于初步设计的某些方面）。 因此，Caughey教授的努力在1979年和1980年得到了AIAA和美国宇航局兰利研究中心的认可。 为了更准确地用激波计算跨声速流动，Caughey教授和Jameson教授开发了有限体积方法，并在一系列论文中解决了围绕复杂几何形状的流动的守恒形式全势方程。 与此同时，David发表了一篇关于“生成有用共形映射的系统程序”的有趣论文和另一篇关于“跨声速翼体流场计算的近共形网格生成方法”的论文。 此外，1985年在《飞机杂志》上发表了一篇关于“翼尾部-机身配置的网格生成”的论文。 此时，Caughey教授将注意力转向了多网格加速技术，并应用于三维势流。 他在1982年的《流体力学年报》和1987年的《科学和技术百科全书》中总结了跨声速势流计算的工作。 他还在《流体和流体机械手册》（1984年）中写了一个关于“有限体积方法”的部分。

\section{1.4 Euler 方程的解决方案}
Caughey对CFD 1.4 Euler 方程解决方案的贡献 本节有25篇由Caughey教授撰写或合著的论文。 第一篇论文是关于“Euler 方程解决方案的自适应网格技术”。 继Brackbill和Saltzman之后，使用变分公式来生成网格，其中平滑度、正交性和浓度的函数组合被最小化。 在计算平面中引入了一个新的定向集中函数，并实施了一个映射边界程序，以帮助保持边界的网格正交性。 为了推导网格自适应，在压力梯度方面形成了新的泛函。 该方法应用于计算跨声速流动过去的2D机翼，结果展示了拟议策略和网格自适应过程的收敛的优势，提高了解决方案在适应网格上的准确性。 下一篇论文（与Eli Turkel教授合作）是关于“数值耗散对可压缩流动问题的有限体积解的影响”。 以细胞为中心的空间离散化通过第二和第四差异的自适应混合来增强，以引入数值耗散。 离散方程使用詹姆森教授开创的显式Runge-Kutta时间步进方案或Caughey的对角化交替方向隐式（DADI）方案求解。 这两种计算都是使用多网格技术加速的。 第二个差项只需要控制激波附近的0-scillations。 因此，在低速区域，它乘以马赫数的函数，以减少其贡献。 幽灵细胞用于强制执行翼型表面边界条件。 第一个差异在区域外的幽灵细胞表面设置为零。 幽灵细胞中的压力是使用正态动量方程找到的。 法向速度被反对称地反射，速度的切分量被外推。 密度由反射或线性外推决定。 另一个选择是从总焓中计算密度，总焓反射到幽灵细胞中。 在他们推荐的程序中，总焓和熵都是从内部推断出来的。 前者用于确定总能量，而后者用于确定密度。 因此，确定了幽灵细胞中速度的大小，而选择其角度是为了在翼型表面强制执行无通量条件。 此外，还研究了能量中的两种形式的人工耗散：一种是总焓，另一种是内能。 第一个允许求解恒定总焓，第二个是对角化ADI配方所必需的。 所有这些变化的结果都记录在论文中。 数字边界条件及其对结果准确性的影响将在关于使用欧拉和Navier-Stokes代码评估阻力的论文中再次讨论。 与Culbert B一起撰写了三篇关于“极端控制”的论文。 拉尼，

M。 M。 1990年和1991年，Hafez 5是康奈尔大学应用数学研究生。 （后来，在1998年，博士 Laney写了一本关于计算气体动力学的广泛使用的书。） 本研究的目标是设计具有理想冲击捕获能力的求解守恒定律的高阶方法，即没有虚假振荡和超标。 为对无界域上的标量守恒定律进行半离散近似，并推导出了防止极值增长和/或产生新极值的最佳条件。 一类后处理器（滤波器），确保单调和无超标解，同时均匀地保持二阶精度，应用于一维模型问题的标准完全离散近似。 该方法可以扩展到多维矢量问题。 本节中的大多数论文都涉及稳定Euler 方程的快速解法。 我们将在这里讨论隐式LU方案、对角线隐式多网格和预制（LU-SGS）多网格方法。 1986年，Caughey教授共同撰写了一篇发表在AIAA上的论文，标题为“Euler 方程应用于任意级联的隐性LU方案”。 继詹姆森和图尔克尔之后，雅各布人被分为两部分，仅使用两个因子的方案甚至可以用于三维问题。 通过使正差重建雅各布矩阵的特征值非正，使后重建雅各布矩阵的特征值非负，从而确保稳定性。 二阶和四阶耗散项仅添加到残差中。 下块和上块对角线系统是通过逐点进行域解决的。 对于周期性问题，系统是循环的，它简化为通常的行进技术，操作量是两倍。 结论是，对于稳定流问题，LU方案比所测试的二维问题的AD1方案快2-3倍。 三年后，出现了对三维非稳态流动的扩展和多网格技术对稳态流动问题的应用。 隐式运算符中使用了一个额外的常数系数二差人工耗散样项。 隐式方案允许采取比显式方案通常允许的更大的时间步骤，而多重网格法被纳入以加速收敛速率以进行稳定的计算。 （在另一篇论文中研究了采用多阶段Runge-Kutta时间步进的多网格算法，通过利用分析确定的控制参数组合来增强。） 继变暖、光束和海耶特关于气体动力学矩阵的对角化和同步对称化的工作之后，开发了一种对角倒置LU隐式方案。 对于三维欧拉方程，矩阵系统通过局部对角变换来处理，将它们解耦为标量方程组。 时间守恒不受解耦合过程的影响，因为对角化变换没有从隐性空间差异中排除。 这种解耦合减少了解决LU近似所需的计算工作量，并且可用于

6 Caughey对CFD的稳定和非稳态方程的贡献。 Caughey教授在1988年的论文《Euler 方程的对角线隐式多网格算法》是众所周知的。 为了在高拉伸网格上构建一个有效的多网格加速度平滑算子，包括人工耗散项的准确表示，特别是没有块五对角线系统反转的四阶项，方程首先在每个点使用Chaussee和Pulliam之后的局部相似性变换对角化。 因此，解耦方程只需要每个因子的阶乘五对角方程的求解。 由此产生的方法具有良好的高波数阻尼，因此它是一个很好的平滑算法。 此外，本研究中使用的耗散性术语对每个方向的缩放方式不同。 这种策略允许使用最小耗散来稳定每个坐标方向的一维问题（而不是像詹姆森公式中那样使用最大值），从而在计算中产生更少的虚假熵。 在模拟跨声速流经过机翼和超音速流经过二维入口和在二维入口内的应用，证实了该方法的效率。 被证明，在扫翼上扩展到三维超音速流也是成功的。 Caughey教授将他的DADI方案与Jorgeson和Turkel提出的对称总变异减耗散应用。 这种对称的TVD耗散方案中的基本元素是：（1）改进的冲击传感器，（2）耗散项的系数是矩阵。 前者允许数值方案在几乎任意强度的冲击附近简化为简单的、上风的一阶近似值，而后者允许对系统中每个单独的方程的耗散项进行适当的缩放。 矩阵耗散使得在没有振荡的情况下捕获非常强和非常弱的冲击成为可能，以进行稳定流动模拟。 此外，收敛速率对耗散常数的选择不敏感。 使用并行计算机在块状结构网格上高效实施DADI方案是几篇论文的主题。 研究了推进多网格周期的模式：多网格周期在所有区块中同步推进的水平模式，允许在周期中它们之间进行数据交换，以及多网格周期在每个区块中独立推进的垂直模式。 事实证明，由于接口边界条件的频繁更新，水平模式的收敛速率比垂直模式的速率快。 跨声速流经过机翼的结果验证了两种模式的准确性，在块间边界处没有引入虚假误差。 在平等位器计算机上可以获得接近理论极限的加速。 在随后的一篇论文中，对垂直模式的改进版本进行了检查。 使用缓冲阵列允许在粗网格上异步更新接口边界条件，消除了边界条件在整个过程中冻结时遇到的收敛问题

M。 M。 哈菲兹7周期。 跨声速流动过去机翼的结果验证了这个版本的垂直模式表现出与水平模式相同的收敛特性，但不需要频繁同步。 上述方法仅限于具有完全连续性的接口（即网格线跨块边界是连续的）。 对于具有不连续界面（即网格密度和间距不连续）的更通用的非重叠块结构网格，必须构建准确和保守的界面方案。 为了将以细胞为中心的有限体积离散应用于界面两侧的细胞，有必要在界面上的细胞面上将不黏性和人工耗散通量进行复合。 因变量的双线性插值用于界面上的细胞面的中点。 只需要从界面隔开的两层单元格中获取信息。 界面的通量和所需的几何量存储在特殊的表面射线中。 这些数组的数据结构遵循通常的多块/多网格框架。 块间边界的完全保守处理允许跨块边界的不连续性。 三维外部和内部流动的计算结果证明了该方法的能力。 接下来的两篇论文涉及预处理，即一维和二维低马赫数稳流量计算的可压缩欧拉方程。 这个想法与Chorin的伪可压缩性密切相关，将压力的时间导数添加到连续性方程中，并调整该项的系数，以优化收敛到稳态解的速率。 Turkel通过将压力时间导数添加到连续性和动量方程中，从而引入了第二个自由参数，从而扩展了Chorin的方法。 在Caughey教授的工作中，压力的时间导数被添加到连续性、动量和能量方程中，从而引入了三个自由参数。 选择这些参数是为了避免Euler 方程在低马赫数下的特征值刚度，同时确保预制系统保持良好定位。 可以选择参数的方式是，当马赫数变为零时，所有特征值都变得与局部流速有序。 预处理修改了人工耗散术语以及边界条件的处理，这是基于修改后系统的特点。 收敛-发散喷嘴的结果证实了分析。 对于翼上的二维流，同样的想法也适用于低马赫数流量。 为了降低超音速和超音速制度的刚度，预处理矩阵经过修改，使预条件系统具有一组不同的特征值，因此雅各布可以（单独）对角化，并且可以选择自由参数，以最小化马赫数全谱的特征条件数。 使用DADI和多网格方法将空间离散方程集成到稳态中。 数值结果

8 Caughey对CFD的贡献表明，收敛速率在lop5 < M < 0.6范围内独立于马赫数，而修改后的预处理矩阵显著改善了收敛速率。 结果还表明，收敛速率在M <时会退化，计算在更小的马赫数下分解。 有人认为，能量通量计算产生的四舍五入误差显著增加了“机器零点”，因此残差只能减少到有效机器零点，然后将其调平。 尽管拟议的方法对实际应用非常有用，但仍有一个理论问题需要回答。 根据Kreiss的说法，当马赫数消失时不应该有奇点，零马赫数流量应该是一个特殊情况。 在这里，拟议的方法并没有消除用于数值计算的公式引入的人工奇点；它只是使先入条件的公式比非先入条件的公式更适合较低的马赫数流量计算。 同样的评论也适用于Turkel、van Leer和Merkle的预处理策略。 本节中的第三个算法是最近的预制条件LU-SGS多网格突破。 为了准备这个阶段，我们讨论了第一个Yoon和Jameson LU实现的对称高斯-塞德尔算法，作为1987年发布的多重网格法的平滑方案。 稳态Euler 方程的线性隐式方案基于LD-IU形式的三个因子。 雅各布矩阵A和B分别分为A+和A-以及B+和B-，其中B+和B-的表达式相似。 在这里，PA和PB分别是A和B的光谱半径。 在这种情况下，D减少到（PA + PB）I，因此是对角线。 矩阵A+和A-的特征值分别为非负和非正。 反向和正向差异用于构建方案的L和U因子，确保只需要标量对角反转的对角线主导地位。 数值耗散项被明确添加到右侧的残差中（因为单边差分方案自然是耗散的）。 Yoon和Jameson也使用相同的方案来解决层流和湍流流的Navier-Stokes 方程。 考虑到基于上风方案（特别是通量分裂）的Euler 方程解的对称高斯塞德尔迭代的成功，正如Chakravarthy、Walters、MacCormack等人所显示的那样，雅各布人现在分裂了，因为2001年，Jameson和Caughey引入了他们的算法。以下：其中IAI = QIAIQ-\textasciitilde{}

M。 M。 Hafez 9 这里/A/是对角矩阵，其条目是Jacobian 矩阵 A和Q的特征值的绝对值，Q-l是A及其逆的模态矩阵。 Bf和B-也使用类似的表达方式。 校正现在用（IAI + IBI）进行预处理，而残差以F+和F-以及G+和G-计算，表示在相应网格坐标方向上，细胞面积乘以通量向量的逆变分量的分割近似值。 残留通量使用詹姆森对称限正（SLIP）近似的标量或对流上风分裂压力（CUSP）版本来近似。 通过将残差局部转换为与原始变量中写成的方程相对应的残差，然后将校正转换为保存变量，使该程序的实现在计算上非常高效。 每个周期都会对超音速区进行额外的校正，通过冻结矩阵系数IAl和IBI来减少这些多次扫描所需的CPU时间。 通过仅存储雅各布的对称形式，将存储这些系数矩阵所需的额外内存最小化。 该方法已针对准一维喷嘴中的流动和经过边界符合“O”型网格的机翼的二维跨声速流动实施。 事实证明，对于这一类问题，该方法明显比任何可用的显式或隐式方法更快，在三到五个多网格周期内获得精确的解决方案，达到精细网格上的截断误差水平！ 在后续论文中，提出了边界符合“C”型网格上通过翼型的无黏性和黏性（层）流动的结果。 虽然不黏性溶液需要三到五个多网格周期，但黏性溶液仍然需要多达二十个多网格周期。 在他们的黏性流动计算中，他们使用Abarbanel和Gottlieb描述的变换同时对称非黏性和黏性Jacobians，并只考虑Navier-Stokes 方程中出现的非混合二次导数的隐式处理。 收敛速率随着网格的高度拉伸而恶化，导致长宽比更大的网格单元。 在一篇刚刚提交出版的论文中，Caughey教授将SGS算法扩充套件到非结构化网格上的跨声速翼型计算。 与基于原始显式（Runge-Kutta）的代码版本相比，收敛的渐近率明显改善了，每个多网格圆盘，但非结构网格上耗散项的计算效率远不如原始显式方案，因此每CPU秒的收敛速率只能比较。 完整的多电网程序仍有待实施。 Caughey教授在结构化网格上使用快速算法，在自由流Mach num-ber mathbfM, = 0.78上演示了Jameson 5-78 翼型的非唯一解决方案。 解决方案收敛到机器零。 他还通过调查不稳定的行为来研究这些解决方案的稳定性

\section{1.5 Navier-Stokes 方程的解决方案}
10 Caughey对CFD的数值贡献，使用他的对角线隐式多网格方案的双时间推进（临时子迭代）版本。 非稳定Euler 方程写入一般曲线坐标系，其中通量使用二阶精确有限体积方案近似，并添加对称的TVD耗散，形成解的第二和第四差的混合。 引入了非稳态方程的二阶精确隐式（向后向）方案，方程在每个时间步骤中通过迭代过程求解，其中校正使用Caughey的对角线隐式方案计算，由最初为稳定流开发的多网格加速。 从各种频率下正弦变化的俯仰角的升力和阻力系数的归一化时间历史可以看出，在足够低的降低频率值下，非独特性导致高度非线性滞后效应，解从稳态解的一个分支跳到另一个分支。 对非定常流动经过Korn 翼型的类似计算表明，这种迟滞不会发生在流过稳定的解决方案独特的机翼中。 在一篇后续论文中，Caughey教授在当前作者研究的课堂上仔细研究了跨声速流动过去的机翼。 这些机翼是对称的，有接近抛物线的前缘，然后是平坦（或几乎平坦）的表面，然后用尖顶的后缘闭合。 对于平面扁平的机翼，其振动与詹姆森的翼型相似。 另一方面，当引入轻微的正弦波时，这些机翼似乎在零角攻击附近具有独特的解决方案，这与现任作者和其他调查人员的说法相反。 在本文中，使用新的LU-SGS多网格求解器的双时间推进版本计算了非稳定流的时间准确解。 1.5 Navier-Stokes 方程的解决方案 本节将回顾十篇论文，包括使用并行计算机模拟层流和湍流流，以及气动性应用。 在第一篇论文中，对角化交替方向隐式对角线多网格算法扩展到黏性、可压缩流的Navier-Stokes 方程的解。 注意力集中在将黏性贡献纳入隐性因素上，以增强稳定性的方式，但不会干扰对角线算法的效率。 计算了经过二维机翼的层流，以证明方案的稳定性和效率。 从隐性因素中完全忽视黏性项限制了方案的稳定性。 遵循Pulliam，通过将黏性项的对角近似添加到适当的隐式因子中，为薄层方程构建了一个方案。 另一个选项是使用额外的隐式运算符，其中包含黏性术语的独家贡献。 然而，第二个选项延长了稳定性极限

M。 M。 哈菲兹11只有条件。 在另一篇论文中，该方法已扩展到湍流流，由于Baldwin和Lomax，它整合了简单的代数模型。 如果流是附加的，包括湍流跨声速流边界层在内的结果与实验数据很一致，而分离流的一致性很差。 收敛率与使用显式Runge-Kutta 多重网格法获得的比率相比是有利的。 然而，黏性术语被明确对待。 用于欧拉和Navier-Stokes 方程解的隐式（对角化ADI）多网格算法是在多个块结构网格的框架内实现的，其中物理域在空间上分解为几个块，解决方案在每个块和分布式计算环境中并行推进。 同步和非同步多电网策略使用显著的加速和收敛速率与单块方案相似。 复合块结构网格允许根据流动的物理特性，在不同的块上解决不同的控制方程；在流动几乎不黏性的子域中，Euler 方程被解决，而在黏性效应很重要的子域中，使用Navier-Stokes 方程。 在一篇有趣的论文中，Caughey教授研究了数值耗散对Navier-Stokes解和阻力计算的影响。 重要的是，添加的数值耗散不会压倒真正的黏性耗散，特别是在后者很重要的区域。 开发了一种估计数值耗散对Navier-Stokes 方程溶液的综合效应的方法。 该方法基于动量方程的积分和数值耗散对阻力积分的校正的计算。 （在使用修正后的外积分计算升力系数时，数值耗散的修正可以忽略不计--所考虑的案例不到0.1\%。） 从基于动量积分方程的分析中可以清楚地看出，固体表面数值边界条件的特定选择将使数值耗散对固体边界积分的贡献归零。 对于混合的第二和第四差分，数值耗散通量--第一和第三差分--在固体表面上设置为零。 通过将通量乘以局部马赫数的某个函数，在正向方向上缩放数值耗散也很常见；因此，在表面附近，实黏性项占主导地位。 由于添加了数值耗散，对阻力的校正可用于定量估计解决方案的质量。 总数值误差可以使用理查德森外流法从表面的压力和皮肤摩擦阻力值中估计。 基于此分析，从计算模拟中计算的阻力是相同的，与用于评估它的等高线无关，如果解已经收敛到稳态，并且如果使用的公式

12 Caughey在评估阻力积分（包括数值分散的贡献）时对CFD的贡献与用于确定解过程中通量平衡的贡献一致。 此外，与不同轮廓相对应的阻力值的差异可用于判断迭代解决方案的收敛程度。 在另一篇论文中，Caughey教授将同样的想法应用于Euler 方程的解决方案。 应该注意的是，只有明确识别数值耗散时，才能进行此分析。 Caughey教授认为，“如果在最关键的区域，即靠近身体表面的溶液中存在大误差（由数值耗散或其他不准确引起），似乎很难证明人们可以通过任何技术获得对阻力的准确预测。” 使用这个论点，人们可能会质疑数值耗散在连续性和能量方程中的影响。 Caughey教授只考虑了动量方程，但数值质量通量也被数值耗散污染了。 速度乘以数值耗散对质量通量的贡献，代表了数值动量通量的修改，并影响阻力计算--这可能是一个更高阶效应，但这并不明显。 数值耗散在能量方程中的影响更加微妙，因为它隐含地污染了压力。 我们以两篇关于非稳流的隐式多网格计算的论文结束本节，该论文适用于空气弹性；一篇涉及流动过流圆柱体，另一篇涉及方形横截面的流动过流圆柱体。 对固定圆柱体计算了周期性涡流脱落的观测模式，并将斯特鲁哈尔数与实验确定的中等雷诺数值进行比较，并有显著的一致性。 耦合问题的解决方案，其中气缸的运动由空气动力学力决定，并与文献中的结果进行了比较。 使用的运动方程是简单的弹簧-品质-阻尼器。 假设结构运动的x和y分量是独立的，它们与流动方程同时集成。 方程被写成一阶系统，并使用与流动方程相同的三级隐式时间离散化进行集成，具有升力和阻力系数的当前值。 流体方程和结构方程在相位收敛时完全耦合。 对于方形横截面的圆柱体，与经验的一致性很差。 事实上，实验数据本身存在相当大的分散。 安装在弹簧阻尼器悬架上的气缸的计算极限周期振幅显示出一种迟滞现象，随着结构频率的增加，导致振荡振幅更大，而结构频率范围比气缸固定时涡流脱落频率略大。

\section{1.6 湍流反应流的模拟}
M。 M。 Hafez 13 1.6 湍流反应流的模拟 最近（1998年至今）Caughey教授与康奈尔大学的同事Stephen Pope教授共同撰写了四篇论文，发表在《计算物理学杂志》上，关于一种新的“湍流反应流PDF方程的混合有限体积/粒子方法”。 在这种方法中，平均质量、动量和能量的守恒方程由有限体积法求解。 这基本上是一个欧拉求解器，用于添加源项的可压缩流。 雷诺应力、斯卡通量和反应项是从通过粒子法计算的粒子场中提取的，并输入到有限体积计算中。 这种方法最重要的问题之一是两种方法的耦合。 有限体积法为粒子系统提供了平均速度场和平均密度。 事实证明，在有限体积计算中满足粒子质量密度与平均可感内能以及相应量之间的一致性条件对耦合的稳定性和准确性至关重要。 此外，需要一个准确的方案，将平均速度场从有限体积数据插值到粒子位置。 新的混合方法有几个优点：建模方程是一致的；而且，由于偏置误差非常小，与纯基于粒子的方法相比，每个单元格的给定精度水平可以使用的粒子少得多。 通过使用有限体积代码和时间平均计算的平均域来减少统计误差。 该方法具有非常好的全局收敛速率。 在1999年发表的一篇论文中，在1D反应随机理想流的简单设置下研究了混合方法开发中出现的算法和数值问题。 采用中点和梯形规则的组合来积分时间中的粒子演化方程，Caughey的对角化隐式有限体积算法用于求解场方程。 集成了局部预处理，以消除由于低马赫数下特征波速之间的巨大差异而导致的特征值刚度造成的困难。 还开发了一种简单的算法来消除场方程中大源项引起的化学刚度。 统计误差、偏差、空间截断误差和时间截断误差在预期速率下收敛。 还表明，提供最佳全局收敛速率的时间平均策略是将粒子算法在每个外部周期中要采取的时间步数增加一倍。 第二篇论文涉及一种紧密耦合算法，其中一个有限体积方案时间步骤和一个粒子方法时间步骤来完成一个迭代。 在这项工作中，有限体积方案基于二阶上风求解器来计算隐性通量，并应用具有局部CFL数的显式Runge-Kutta 时间推进方法与预处理技术一起克服低马赫数刚度。 关于“悬崖体稳定流动的PDF模拟”的第三篇论文包括

\section{1.7 特别主题}
14 Caughey对CFD的贡献，比较了混合方法和独立粒子网格方法。 考虑了混合方法的两个版本。 第一篇论文的松耦合算法和第二篇论文的紧密耦合算法，扩展到二维（轴对称）问题。 就网格加密和粒子数而言，收敛解决方案是用所有三种算法获得的。 理查森外推用于获得大量网格单元格极限的解。 推断的解决方案彼此一致，并且与可用的实验数据非常一致。 在第四篇论文中，确定了在数值溶液水平上完全一致性需要满足的条件，并开发了强制执行这些条件的校正算法。 此外，还提出了能量方程和状态方程的新公式。 混合方法适用于非预混合的引导式喷气火焰。 这些结果与其他早期的计算和实验数据很一致。 第五篇关于“使用具有详细化学成分的联合PDF模型计算悬崖体稳定火焰的论文”正在出版。 我们以以下引文结束本节：“由于数值误差大幅降低，对于给定的网格大小和粒子数，目前的混合方法代表了粒子/网格方法求解PDF方程的计算效率的重大进步。” 1.7 特别主题 本节讨论了Caughey教授共同撰写的不同领域的几篇论文。 1983年，康奈尔大学撰写了一份关于“稳定二维跨临界分层流动问题的数字技术（应用于海洋热能转换厂的中间场动力学）”的报告。 基于浅水方程和可压缩气体动力学的类比，使用势流方程的Jameson \& Caughey 有限体积法模拟了带有水力跳跃的跨临界流，而使用MacCormack的方案模拟了具有摩擦效应的超临界流。 该报告基于康奈尔大学环境工程博士论文，由Caughey教授部分指导。 1993年，Caughey教授参与了一项关于“A 激波声音放大的数值调查”的研究。 使用二阶MacCormack和高阶EN0方案对收敛-散射喷嘴中准一维冲击流场的时间依赖性数值计算，在入口边界引入声波受到干扰，很好地预测了喷嘴中的扰动振幅和相速度，使用线性理论验证了小扰动的结果。 在另一篇论文中，对计算缓慢移动的激波的难度进行了调查。 发现数值误差主要表现为伪熵波，其结构是所用算法的函数。 福图-

\section{1.8 评论文章}
M。 M。 Hafez 15 nately，这一类问题的声压波不受伪熵产生的影响。 1996年，AIAA发表了一篇关于“冲击运动和变形在冲击噪声生成中的作用”的论文，其中模拟了轴对称流，允许涡量与声波和熵波相互作用。 冲击运动是通过声波与冲击的相互作用来建模的，而冲击变形是通过涡旋环与冲击的相互作用来建模的，同时产生声波和接触面。 作者声称，这两种机制在产生冲击噪声方面都很重要。 Caughey教授共同编辑了三卷《计算流体动力学的前沿》。 在1994年的第一卷中，他共同撰写了一篇关于“詹姆森对CFD的贡献的回顾”的论文；在1998年的第二卷中，他共同撰写了一篇关于“Earl1 Murman对跨声速流动和CFD的贡献的回顾”的论文；在2002年第三卷中，“Robert MacCormack对CFD的贡献”。 作者很荣幸和高兴地回顾David Caughey在本系列的当前和第四卷中对CFD的贡献。 1.8 评论文章 在过去的三十年里，Caughey教授写了几篇评论文章。 除了对前面提到的势流计算的审查外，1980-2003年期间还有10篇论文。 AIAA的一篇名为“计算空气动力学的（有限）视角”（1980年）的论文涵盖了面板方法、跨声速势方法（包括几何考虑、网格生成、有限体积制定和迭代方案）、欧拉方程方法（Lax Wendroff和MacCormack的显式方法、隐式方案、ADI和LU分解），并以一些应用结束。 在1994年康奈尔工程季刊上，Caughey教授发表了一篇关于“空气动力学设计中的计算方法”的文章。 他在1995年计算流体动力学评论中写了一篇关于“跨声速流动的计算”的文章，包括算法开发、大规模并行分布式内存系统的实现、空气动力学、航空声学和设计应用。 1995年，他还写了一篇关于“可压缩空气动力学的多网格方法”的文章，其中他回顾了对角线ADI作为多网格的平滑方案，以解决基于有限体积空间离散化并添加二阶和四阶耗散项的尤拉和Navier-Stokes 方程。 他表明，隐式因子的对角化，产生标量五对角线系统，大大节省了计算劳动力。 1997年，Caughey教授在《CRC计算机科学与工程手册》中写了一章关于“计算流体动力学”，1998年，在《CRC流体工程手册》中写了一章关于“收敛加速”。 在后者中，他回顾了显式和隐式

\section{1.9 流体力学：交互式文本}
16 Caughey对解决欧拉和Navier-Stokes 方程的CFD方法的贡献。 他涵盖了四阶Runge-Kutta方法、焓阻尼、残差平滑、ADI和LU因子化方案、LU-SGS，它结合了LU因法和对称高斯-塞德尔松弛的优势，强隐式程序（SIP）和多网格方法，以及预处理技术。 在《物理科学与技术百科全书》（2002年）中，Caughey教授写了一篇关于“计算空气动力学”的文章，最近在2003年，他与Jameson教授共同撰写了一篇关于“跨声速流动计算技术的发展：历史视角”的论文。 1.9 流体力学：互动文本Caughey教授是一位优秀的老师。 他对使用计算机教授流体力学特别感兴趣。 1998年，他与James Liggett教授合著了一本关于这个主题的书。 这是第一本流体力学入门电子教科书。 这本基于计算机的书籍在CD-ROM软盘上分发，旨在在个人电脑上学习。 这种演示模式启用的功能包括：与超文本相关的索引、搜索和交叉引用工具；显示动态现象的动画和视频的能力；以图形和数值形式访问各种流体属性数据；以及访问解决线性和非线性问题的各种计算工具，允许呈现各种参数值的结果，而无需阅读器的表格查找或迭代的繁重。 这本互动教科书为用MATLAB编写的许多例程提供了图形用户界面（GUI），允许学生解决各种问题，包括：0单位转换问题0标准大气的属性0函数和数据的一维和二维绘图，具有确定表面图下体积的数值积分例程0线性和非线性或超验方程组0维度分析问题0稳态管道网络流量0不可压缩管道流，用于确定湍流流的管道尺寸、流量和摩擦损耗。 0 基本解的SU-perposition的平面和轴对称不可压缩势流问题

\section{1.10 结束说}
M。 M。 Hafez 17 混合诺伊曼和狄利克雷边界条件的任意二维几何中的潜在流动 具有面积变化的等向异性可压缩流动和常数激波 可压缩流动在恒定面积管道中流动，具有热量加成和/或摩擦 开放通道流动问题 0 弹性管道中的水锤问题 这本书是一个很好的计算机辅助教学工具，可以对下一代流体力学教学产生重大影响。 有了包含的计算实用工具，学生有能力在很短的时间内对复杂现象进行参数化研究，而他或她需要更少的繁琐努力，例如，冲击损失或加热和摩擦如何影响管道中的流动，湍流管道流量中的头损失与流量的关系等。 图形能力也非常令人印象深刻，根据几个变量（压力、温度、密度、速度、涡量、熵、熵等）对流场的细节，以及分析这些巨大数据的工具（例如，释放粒子并追踪它们）。 据说，一张图片胜过千言万语（而一部电影也许价值数百万）。 流动可视化本身就是一个非常有趣的研究课题，使用计算机研究流体动力学的领域确实很有前途。 过渡的困难问题和湍流已经从当今计算机能力的进步中受益。 希望Caughey教授在不久的将来会为我们的本科和研究生计算机图书馆增加更多电子书。 注意：虽然出版商不再支持原始CD教科书，但Caughey教授和Liggett教授刚刚发布了MATLAB实用程序的更新版本；有关详细信息和可用性，请参阅http://FluidMechSolns.com。 1.10 结论 Caughey教授在CFD上工作了近35年。 他为该领域的各个方面做出了原创贡献。 他的工作对航空航天业的影响是不可否认的，不仅在美国，而且在国际上。 在康奈尔大学担任教授的30年里，他监督了18个博士和5个理学硕士学位，以及23个工程硕士设计项目。 在过去的10年里，他一直是AIAA的研究员。 他有135篇出版物，其中35篇发表在档案期刊上。 他是第一本电子形式的流体力学教科书的合著者（与Jim Liggett），整个文本中都使用了动画、视频剪辑和计算工具，以明确

18 Caughey对CFD的贡献解释了基本的物理现象。 他与A教授在CFD的合作。 詹姆森值得注意。 1975年，他们一起成功使用旋转差分方案生产了FLO-22--这是第一个在工业中广泛用于跨声速流动机翼模拟的全潜力代码！ 后来，他们使用保守的有限体积方案，引入了FLO-27，取代了FLO-22。 最近，在2002年，他们生产了Euler 方程最快的预制多电网解决方案，用于跨声速翼的稳定流量。 与Stephen Pope教授的合作也很富有成效。 多年来，他们共同研究了用于模拟湍流反应流的混合有限体积/粒子方法，并取得了显著成功。 Caughey教授曾与工业界和学术界的许多科学家合作过，他的同事们（无一例外）都钦佩他；他非常聪明、勤奋、公平、诚实，毕竟是一个正派的人。 我们只是祝愿他一切顺利。 作者一直高度尊重Caughey教授，在他完成这篇评论后，他确信Caughey教授应该得到这个领域任何人的最高尊重。

\section{1-A博士 由David A监督的学生。 考伊}
附录1-A博士 由David A监督的学生。 Caughey教授Caughey指导了康奈尔大学航空航天工程、机械工程、应用数学以及理论和应用力学等研究生领域的18名学生的博士论题。 在本节中，我们列出了在Caughey教授的指导下完成博士学位论文的学生，以及他们的研究生领域、毕业日期和论文标题。 1. 乔杰·S。 杜利克拉维奇，航空航天工程，1979年1月。 通过转子和风扇进行Inviscid 跨声速流动的数值计算2。 爱德华·K。 Buratynski，理论与应用力学，1983年5月。 应用于任意级联3的Euler 方程数值解的下-上因子隐式方案。 Arvin Shmilovich，航空航天工程，1984年1月。 使用多重网格法 4计算跨声速势流经过翼尾-机身配置。 Frederik J. deJong，航空航天工程，1985年1月。 跨声速势流计算中的浮动激波拟合 5. Minwoo Park，航空航天工程，1985年8月。 超音速势流过去倾斜椭圆锥6的多网格计算。 Jang-Hyuk Kwon，航空航天工程，1986年6月。 带和不带解决方案适应7的级联数字网格生成。 Murali Damodaran，航空航天工程，1987年1月。 使用Euler 方程 8对非稳无黏性旋转跨声速流动过去机翼进行数值计算。 杰弗里·W。 横田，航空航天工程，1987年5月。 在旋转涡轮机械通道9中解决跨声速流动的Euler 方程的L-U隐式多网格算法。 韦恩·A。 史密斯，机械工程，1987年8月。 Euler 方程 10的多网格溶合。 邓·C。 刘，航空航天工程，1987年8月。 解决Euler 方程 11的自适应网格技术。 Yoram Yadlin，航空航天工程，1990年8月。 Euler 方程的块隐式多网格解决方案

20 Caughey对CFD 12的贡献。 库尔伯特·B。 Laney，应用数学，1991年5月。 守恒定律数值近似的单调性和超射条件13。 托马斯·L。 泰辛格，航空航天工程，1992年5月。 可压缩Navier-Stokes 方程的隐式多网格解决方案，适用于分布式并行处理14。 Ram Varma，机械工程，1993年5月。 可压缩lhrbulent流的多网格对角隐式解决方案及其评估15。 王丽霞，航空航天工程，1993年8月。 复杂几何中可压缩空气动力学问题的隐式多网格解决方案16。 Guang Chang，航空航天工程，1995年1月。 蒙特卡洛PDF/有限-湍流火焰卷研究17。 Kristine Meadows，航空航天工程，1995年5月。 基本冲击噪声机制的数值研究18。 Metin Muradozlu，航空航天工程，2000年1月。 湍流、反应流的PDF方程的一致混合有限体积/粒子方法

\section{1-B David A.的出版物。 考伊}
附录1-B David A.的出版物。 Caughey 在本附录中，我们包括了Caughey教授出版物的完整列表。 1. 二阶波结构：平面流动，（与W. D。 Hayes），1968年5月20日至21日，多伦多大学航空航天研究所，AFOSR-UTIAS空气动力学噪声研讨会。 2. 回顾二阶波结构，（与W. D。 Hayes），在声爆的第二次会议上，美国宇航局SP-180，pp. 129-132，1968年。 3. 1968年12月9日至14日，田纳西州塔拉霍马市田纳西大学空间研究所激波的生成和传播短期课程的讲义中，对声爆问题的应用。 4. 超声速流动中的二阶波结构，博士 论文，普林斯顿大学，新泽西州普林斯顿，1969年。 此外，美国宇航局CR-1438，1969年9月。 5. 与声爆相关的理论问题，（与W. D。 Hayes等人），在第三次关于声爆的会议上，美国宇航局SP-255，pp. 27-31，1971年。 6. Volnovoe Soprotivlenie v Sverkhzvukovom Potoke Ideal'nogo Gaza v Ne-lineinoi Teorii，（根据非线性理论的理想气体超声速流动中的波阻），Zhurnal Vychislitel'noi Matematiki i Matematicheskoi Fiziki，卷。 11，pp. 982-991，1971年7月至8月，（俄语）。 7. 用于数值分析的新Tkansonic 小扰动理论，Mc-Donnell Douglas报告MDC Q0478，1972年12月。 8. POTMAP用户指南：复杂势平面的翼型映射，麦克唐纳道格拉斯报告MDC Q0505，1973年12月。 9. CPROX用户指南：Transonic 翼型分析程序，Mc-Donnell Douglas报告MDC Q0506，1973年12月。 10. NACEL程序：用于对称分析的计算机程序，Tkan-sonic Nacelles，McDonnell Douglas报告MDC Q0564，1975年9月。 11. 跨声速的隐形分析，板条机翼，飞机杂志，卷。 13，pp. 1976年1月29-35日。

22 Caughey对CFD 12的贡献。 纽约詹姆森-凯的简要描述。 U。 计算机程序\textasciitilde{}FLO-22，（与安东尼·詹姆森，佩里·A。 纽曼和鲁比·M。 戴维斯），美国宇航局TMX-73996，1976年12月。 13. 跨声速势流场的复杂三维配置的计算，（与安东尼·詹姆森一起），在涡轮机械中的跨声速流动问题中，T。 C. 亚当森和M。 F。 Platzer，Eds.，pp. 274-291，半球出版公司，1977年。 14. 跨声速流动经过扫翼的数值计算，（与安东尼·詹姆森合作），纽约大学ERDA报告COO 3077-140，1977年6月。 15. 跨声速势流计算的有限体积法，（与安东尼·詹姆森一起），AIAA第三届计算流体动力学会议记录，新墨西哥州阿尔伯克基，pp. 35-54，1977年6月27-28日。 16. 跨声速声纳赛尔流场的加速迭代计算，（与安东尼·詹姆森合作），AIAA杂志，卷。 15，pp. 1474-1480，1977年10月。 （也是AIAA论文76-100，第14届航空航天科学会议，华盛顿，D。 C.，1976年1月26日至28日。） 17. 最近三维Tkansonic流计算的经验，（与Perry A. 纽曼和安东尼·詹姆森），美国宇航局CTOL运输技术会议纪要，1978年2月28日至3月3日，美国宇航局TM-78733。 18. 跨声速电位流动的数值计算，在计算流体动力学进步短期课程的讲义中，田纳西大学空间研究所，田纳西州图拉霍马，1978年12月4日至8日。 19. 生成有用共形映射的系统程序，国际工程数值方法杂志，卷。 12，pp. 1651-1657，1978年。 20. 三维跨声速流动的有限体积方法的开发，（与安东尼·詹姆森和D. 尼克松），流动研究报告编号 134，1979年2月。 21. 跨声速势流关于翼体组合的数值计算，（与安东尼·詹姆森一起），AIAA杂志，卷。 17，pp. 175-181，1979年2月。 （另外，AIAA论文77-677，在1977年6月27日至29日在新墨西哥州阿尔伯克基举行的第10届流体和等离子动力学会议上发表。）

M。 M。 哈菲兹 23 22。 有限体积方法在机翼机身计算中的应用进展，（与安东尼·詹姆森合作），AIAA论文79-1513，AIAA第12届流体和等离子动力学会议，弗吉尼亚州威廉斯堡，1979年7月24日至26日。 23. 关于复杂三维几何的跨声速流场的高阶有限差分方案，（与L. T。 陈），AIAA第4届计算流体动力学会议记录，1979年7月23日至24日，弗吉尼亚州威廉斯堡。 24. 使用广义坐标计算跨声速入口流场，（与L.T. 陈），飞机杂志，卷。 17，pp. 167-174，1980年3月。 25. 通过转子和风扇对Tkansonic 势流进行有限体积计算，（与D.S. Dulikravich），Sibley机械和航空航天工程学院，流体动力学和空气动力学专案，FDA-80-03报告，康奈尔大学，纽约伊萨卡，1980年3月。 26. 计算空气动力学的（有限）视角，西伯利机械和航空航天工程学院，流体动力学和空气动力学项目，报告FDA-80-07，康奈尔大学，纽约伊萨卡，1980年7月。 27. 三维跨声波流的有限体积方法的开发，（与John E. Mercer、Wen Huei Jou、Antony Jameson和David Nixon），Flow Research Report No. 166，1980年8月。 28. 'Ikansonic 进气道流动使用通用网格生成方案的现场计算，（与L. T。 陈），流体工程杂志，卷。 102，pp. 309-315，1980年9月。 29. 关于轴对称体的三维跨声速流动域的高阶有限差分方案，（与L. T。 陈），飞机杂志，卷。 17，pp. 668-676，1980年9月。 30. 有限体积方法在机翼机身计算中的应用进展，（与安东尼·詹姆森合作），AIAA杂志，卷。 18，第1281-1288页，1980年11月。 31. Tkansonic势流的加速有限体积计算，（与Antony Jameson，W. H。 Jou、John Steinhoff和Richard Pelz），在激波，A中计算隐性跨声速流动的数值方法。 Rizzi和H。 薇薇安，编辑，pp. 11-27，Vieweg，不伦瑞克/魏斯巴登，1981年。

24 Caughey对CFD 32的贡献。 关于冲击时势流方程的各种治疗，（与L. T。 陈），Proc。 数字边界条件程序研讨会，美国宇航局艾姆斯研究中心，加利福尼亚州莫菲特菲尔德，1981年10月19日至20日，33。 超音速电位流动的计算，流体力学年度回顾，卷。 14，pp. 261-283，1982年。 34. 超音速计算，势流过去三维配置，在超音速、冲击和多维流动中：科学计算的进步，Richard E. Meyer，Ed.，pp. 77-105，学术出版社，纽约，1982年。 35. 一种用于跨声速翼体流场计算的近共形网格生成方法，（与L. T。 陈和A。 Verhoff），AIAA论文82-108，第20届航空航天科学会议，1982年1月11日至14日，佛罗里达州奥兰多。 36. 多网格方法在机翼-机组合的跨声速电位流动计算中的应用，（与Arvin Shmilovich合作），计算物理学杂志，卷。 48，pp. 462-484，1982年12月。 37. 有限体积法用于跨声速势流计算的基本进展，（与安东尼·詹姆森一起），空气动力学流动的数值和物理方面，T。 Cebeci，编辑，pp. 445-461，Springer-Verlag，纽约，1982年。 38. 稳定、二维、跨临界、分层流动问题的数值技术，应用于海洋热能转换厂的中间场动力学，（与Janet M. 琼斯和格哈德·H。 Jirka），康奈尔大学土木与环境工程学院，1983年3月。 39. 网格生成用于机翼-尾-机身配置，（与Arvin Shmilovich合作），在网格生成的进步中，K。 N。 Ghia和U。 Ghia，编辑，pp. 189-197，美国机械工程师学会，1983年。 40. 直线和扫尖型转子叶片上测量和计算压力的比较，（与M. E。 陶伯，I-C。 Chang和J。 J。 Philippe），NASA TM 85872，1983年12月。 41. 三维、跨声速、势流的多网格计算，应用数学和计算，卷。 13，pp. 241-260，1983年12月。

M。 M。 哈菲兹25 42。 跨声速计算程序，势流机翼-尾-机身配置，（与Arvin Shmilovich合作），西布利机械和航空航天工程学院，流体动力学和空气动力学项目，报告FDA 83-08，康奈尔大学，纽约伊萨卡，1983年12月。 43. 计算通过机翼和机翼机身组合的超音速电位流动的多网格程序，西布利机械和航空航天工程学院，流体力学和空气动力学项目，报告FDA 84-02，康奈尔大学，纽约伊萨卡，1984年1月。 44. 有限体积方法，第10章第10.3.3-1节，流体力学中的计算方法，S。 中村，Ed.，在《流体和流体机械手册》中，John Wiley \& Sons，1984年出版。 45. 使用多网格方法（与Arvin Shmilovich一起）计算跨声速势流经过翼尾部机组配置，Proc。 第九届流体力学数值方法国际会议，物理学讲义，卷。 218，pp. 508-513，Springer-Verlag，柏林，1985年。 46. 网格生成用于机翼-尾-机身配置，（与Arvin Shmilovich合作），飞机杂志，卷。 22，pp. 467-472，1985年6月。 47. 通过机翼-尾-机组合组合的纳米流的多网格计算，（与Arvin Shmilovich合作），飞机杂志，卷。 22，pp. 581-586，1985年7月。 48. 多网格计算？ 2-泛音速，电位流，（与Arvin Shmi-lovich一起），在流体中数值方法的最新进展：计算跨声速的进步，W。 Habashi，Ed.，Pineridge Press，Swansea U. K.，pp. 83-108，1985年。 49. 在跨声速流动关于nee Air和Wind第11IMACS世界隧道中的翼翼的计算（与Arvin Shmilovich一起），在Proc。 大会，挪威奥斯陆，1985年8月5日至9日。 50. 稳定、二维跨临界流动问题的数值方法，（与Janet M. Jones和Gerhard Jirka），斯旺西，U。 K.，1985年。 51. 关于空气和风洞中的机翼、轴对称项目和机芯的跨声速流动计算（与Arvin Shmi-lovich一起），报告编号。 MDC 53852，道格拉斯飞机公司，加利福尼亚州长滩，1985年10月。

26 Caughey对CFD 52的贡献。 适用于任意级联的Euler 方程的隐式LU方案，（与Edward K. Buratynski）AIAA杂志，卷。 24，pp. 39-46，1986年1月。 （也是AIAA论文84-0167，第22届航空航天科学会议，内华达州里诺，1984年1月9日至12日。） 53. 关于直升机转子叶片的丹索克流的计算，（与Rimon Arieli，M. E。 Tauber和D。 A. Saunders），AIAA杂志，卷。 24，pp. 722-727，1986年5月。 54. 双曲喷嘴中的Zlansonic 势流，（与Minwoo Park合作），AIAA杂志，卷。 24，pp. 1037-1039，1986年6月。 55. 计算空气动力学，《科学技术百科全书》文章，卷。 3，pp. 285-304，学术出版社，1987年。 56. 通过旋转通道的Inviscid Dansonic流的多网格解决方案，（与Wayne A. 史密斯），AIAA论文87-0608，第25届航空航天科学会议，内华达州里诺，1987年1月。 57. 用于可压缩流动计算的对角线隐式多网格算法，在偏微分方程计算机方法的进步中-VI，R。 Vichnevetsky和R. S。 Stepleman，编辑，pp. 270-277，IMACS，新泽西州新不伦瑞克，1987年。 58. 适用于可压缩流的Euler 方程的高效隐式多网格算法，Proc。 国际流体力学会议（1987年北京），pp. 392-397，北京大学出版社，中国北京，1987年7月1日至4日。 59. 数值耗散对可压缩流动问题的有限体积解决方案的影响，（与Eli Turkel合作），AIAA论文88-0621，第26届航空航天科学会议，内华达州里诺，1988年1月11日至14日。 60. Euler 方程的对角线隐式多网格算法，AIAA杂志，卷。 26，pp. 841-851，1988年7月。 （另外，AIAA文件87-0354，在1987年1月12日至15日在内华达州里诺举行的第25届航空航天科学会议上发表。） 61. Zlan-sonic 势流方程解法人工可压缩性方法的评估，（与Matthew Wynn合作），西布利机械和航空航天工程学院，流体动力学和空气动力学项目，报告FDA 88-12，康奈尔大学，纽约伊萨卡，1988年7月。 62. 三维Euler 方程的对角线隐式多网格解，（与Yoram Yadlin）Proc。 第11届流体力学数值方法国际会议，6月，弗吉尼亚州威廉斯堡

M。 M。 Hafez 27 27 - 1988年7月1日，物理学讲义，卷。 323，D。 L。 德沃耶，M。 Y。 Hussaini和R. G。 Voigt，Eds.，pp. 597-601，Springer-Verlag，柏林，1989年。 63. Tkansonic 势流经过扫翼的数值计算-Flo-22程序的改进版本，（与王立霞一起），西布利机械和航空航天工程学院，流体动力学和空气动力学项目，报告FDA 88-17，康奈尔大学，纽约伊萨卡，1988年8月。 64. 三维欧拉方程的L-U隐式多网格算法，（与Jeffrey W. 横田），AIAA杂志，卷。 26，pp. 1061-1069，1988年9月。 65. Inviscid Tkansonic 翼型-Vortex Interaction的有限体积计算，（与Murali Damodaran），AIAA杂志，卷。 26，pp. 1346-1353，1988年11月。 （还有AIAA论文87-1244，第19届流体力学、等离子体动力学和激光会议，1987年6月8日至10日，夏威夷檀香山。） 66. 进气道流动场的对角线隐式多网格计算，（与Ravi Iyer一起）AIAA杂志，卷。 27，pp. 110-112，1989年1月。 67. 三维Euler 方程的对角线隐式多网格解，（与Yoram Yadlin一起）第六届陆军应用数学和计算会议交易，科罗拉多州博尔德，1988年5月31日至6月3日，ARO报告89-1，pp. 1041-1049，陆军基础研究办公室，1989年2月。 68. Euler 方程解决方案的自适应网格技术，（与D. C. 刘）西比利机械和航空航天工程学院，流体力学和空气动力学项目，报告FDA 89-02，康奈尔大学，纽约伊萨卡，1989年3月。 69. 可压缩流体流动Euler 方程的块多网格隐式解决方案，（与Yoram Yadlin一起）西布利机械和航空航天工程学院，流体动力学和空气动力学项目，报告FDA 89-05，康奈尔大学，纽约伊萨卡，1989年6月。 70. 下上隐式方案的对角反转，（与J. W。 横田和R。 V。 Chima）AIAA杂志，卷。 28，pp. 263-266，1990年2月。 （另外，第一届全国流体力学会议的Proc.，第1部分，AIAA，华盛顿特区。 C.，1988年7月，pp. 104-111，还有美国宇航局技术。 备忘录。 100911，1988年7月。）

28 Caughey对CFD 71的贡献。 极端控制：人工黏度的影响，（与Culbert Laney合作）在1990年6月19日至22日，纽约伊萨卡举行的第八届陆军应用数学和计算会议记录中。 72. 可压缩层流的对角线隐式算法，（与Thomas L. 泰辛格）在1990年6月19日至22日，纽约伊萨卡举行的第八届陆军应用数学和计算会议记录。 73. 可压缩湍流流的对角线隐式多网格解决方案，（与Rama R. Varma）在1990年6月19日至22日在纽约伊萨卡举行的第八届陆军应用数学和计算会议记录中。 74. 可压缩Navier-Stokes 方程的多网格ADI算法，（与Thomas Tysinger合作）AIAA论文91-0242，第29届航空航天科学会议，内华达州里诺，1991年1月。 75. 极端控制II：对守恒定律的半离散近似，（与Culbert Laney一起）AIAA论文91-0632，第29届航空航天科学会议，内华达州里诺，1991年1月。 76. 可压缩流体流动Euler 方程的块多网格隐含解决方案，（与Yoram Yadlin一起），AIAA杂志，卷。 29，pp. 712-719，1991年5月。 77. 极端控制III：对保护定律的完全离散近似，（与Culbert Laney一起）AIAA论文91-1534，Proc。 AIAA第10届计算流体动力学会议，pp. 81-94，夏威夷洪奥卢鲁，1991年6月。 78. 可压缩湍流流的对角线隐式多网格解决方案，（与Rama R. Varma）AIAA论文91-1571，AIAA第十届制度流体力学会议，pp. 487-500，夏威夷檀香山，1991年6月。 79. 可压缩Navier-Stokes 方程的并联块多网格解决方案，（带Y。 Yadlin和T。 Tysinger）Proc。 AIAA第十届制度流体力学会议，pp. 965-966，夏威夷檀香山，1991年6月。 80. 数值耗散对Navier-Stokes解决方案的综合效应的估计，（与R. 瓦尔玛）程序。 第四届计算流体动力学国际研讨会，卷。 11，pp. 1161-1166，加州大学戴维斯分校，加利福尼亚州戴维斯，1991年9月9日至12日。

M。 M。 哈菲兹 29 81。 欧拉/Navier-Stokes 方程在具有不连续接口的块状结构网格上的隐式多网格算法，（与Lixia Wang）Proc。 ICFD流体动力学数值方法会议，雷丁大学，1992年4月。 82. 在并行计算机上阻止Euler 方程的隐式多网格解，（与Y. Yadlin）在平行计算流体动力学中：实现和结果，H。 Simon，Ed.，麻省理工学院出版社，马萨诸塞州剑桥，pp. 127-145，1992年。 83. 并行计算 Euler 方程的块状多网格隐式解决方案策略，（与Yoram Yadlin合作），AIAA杂志，卷。 30，pp. 2032-2038，1992年8月。 84. 可压缩Navier-Stokes 方程的多网格ADI算法，（与Thomas Tysinger合作）AIAA杂志，卷。 30，pp. 2158-2160，1992年8月。 85. 跨声速势流经过扫翼的改进计算，（与L. 王）J。 飞机，卷。 29，不。 5，pp. 961-964，9月至10月。 1992年。 86. 分布式并行处理应用于隐式多网格欧拉/Navier-Stokes算法，（与T. L。 Tysinger）AIAA论文93-0057，第31届航空航天科学会议，内华达州里诺，1993年1月。 87. 模拟20和3D可压缩流的多块/多网格尤拉方法，（与王立霞一起）AIAA论文93-0332，第31次航空航天科学会议，内华达州里诺，1993年1月。 88. 使用数值耗散的综合效应对Navier-Stokes解决方案的评估，（与Rama Varma一起）AIAA论文93-0539，第31届航空航天科学会议，内华达州里诺，1993年1月。 89. 可压缩流动、计算机和流体的隐式多网格技术，卷。 22，不。 2/3，pp. 117-124，1993年。 90. 激波声音放大的数值调查，（与K. R。 Meadows和J。 卡斯帕）在计算航空和水声学中，ASME FED卷。 147，pp. 47-52，美国机械工程师学会，纽约，N。 Y。 在1993年6月20日至24日的ASME流体工程会议上发表。 91. 具有对称总变化的隐式多网格欧拉解决方案，减少耗散，Proc。 AIAA第11届计算流体动力学会议，pp. 676-684，佛罗里达州奥兰多，1993年7月6日至9日。

30 Caughey对CFD 92的贡献。 计算用于空气声学应用的非稳态激波，（与Kristine R. Meadows和Jay Caspar）AIAA论文93-4329，第15届AIAA航空声学会议，加利福尼亚州长滩，1993年10月25日至27日。 还有AIAA杂志，卷。 32，pp. 1360-1366，1994年7月。 93. 使用数值耗散的综合效应对Navier-Stokes解决方案的评估，（与Rama Varma一起），AIAA杂志，卷。 32，pp. 294-300，1994年2月。 94. 康奈尔大学航空航天工程，康奈尔工程季刊，pp. 2-3，1994年冬天。 95. 空气动力学设计中的计算方法，康奈尔工程季刊，pp. 23-28，1994年冬季。 96. 1994年6月20日至23日，在内华达州太浩湖举行的ASME流体工程会议上，使用计算机教授流体力学。 97. 计算流体动力学的前沿-1994年，与M共同编辑。 M。 哈菲兹，约翰·威利和儿子，奇切斯特，1994年。 98. 安东尼·詹姆森对计算流体动力学的贡献，（与M. M。 哈菲兹），在计算流体动力学-1994年的前沿中，D。 A. Caughey和M。 M。 哈菲兹，编辑，约翰·威利和儿子，pp. 1-42，1994年。 99. 跨声速流的计算，在1995年计算流体动力学评论中，M。 哈菲兹和K。 大岛，编辑，John Wiley \& Sons，pp. 567-579，1995年。 100。 空气动力学计算中的阻力和数值误差，Proc。 第六届计算流体动力学国际研讨会，第四卷，pp. 23-28，内华达州太浩湖，1995年9月。 101. 可压缩空气动力学的隐式多网格方法，在计成流体动力学技术中，W. 哈巴希和M。 Hafez，编辑，Gordon \& Breach，pp. 677-697，1995年。 102. 计算空气动力学，流体力学研究趋势，J. L。 Lumley等人，编辑，pp. 55-60，美国物理学会，1996年。 103. 冲击运动和变形在产生冲击噪声中的作用，（与K. R。 Meadows），AIAA论文96-1777，第二届AIAAICEAS航空声学会议，宾夕法尼亚州立大学，1996年5月6日至8日。

M。 M。 哈菲兹31 104。 计算流体动力学，CRC计算机科学与工程手册，第37章，pp. 873-891，CRC出版社，博卡拉顿，1997年。 105. 可压缩Euler 方程的收敛在低马赫数，（与M。 Muradoijlu），在《流动模拟技术的进步》中，纪念Joseph L.的会议记录。 Steger，加利福尼亚州戴维斯，1997年5月2日至4日。 106. 预制Euler 方程的隐式多网格解决方案，（与M. Muradoijlu），Proc。 AIAA第13届计算流体动力学会议，pp. 648-658，科罗拉多州斯诺马斯，1997年6月29日至7月2日。 107. 收敛加速，在CRC流体工程手册中，第32章，pp. 32-1-32-23，CRC出版社，博卡拉顿，1998年。 108. 预制多维欧拉方程的隐式多网格解，（与M. Muradoglu），AIAA论文98-0114，第36届航空航天科学会议，内华达州里诺，1998年1月12日至15日。 109. 流体力学入门计算机教科书，（与J. A. Liggett），在Proc。 1998年ASEE年会，美国工程教育学会，华盛顿州西雅图，6月29日至7月2日。 1998年。 110. 流体力学：互动文本，（与James A. Liggett），美国土木工程师学会，1998年8月。 111. 计算流体动力学的前沿-1998年，与M共同编辑。 M。 哈菲兹，世界科学出版公司，新加坡，1998年。 112. Earl1 Murman对Tkansonic Flow和计算流体动力学的贡献回顾，（与M. M。 哈菲兹），在计算流体动力学的前沿-1998年，D。 A. Caughey和M。 M。 哈菲兹，编辑，世界科学，pp. 1-28，1998年。 113. 一本基于计算机的流体力学教科书，（与James A. Liggett），在计算流体动力学的前沿-1998年，D。 A. Caughey和M。 M。 哈菲兹，编辑，世界科学，pp. 465-481，1998年。 114. 流体力学计算机辅助教学：电子教科书，论文FEDSM99-6824，1999年ASME/JSME联合流体工程会议记录，工程教育进步论坛，美国机械工程师学会，加利福尼亚州旧金山，1999年7月18日至22日。

32 Caughey对CFD 115的贡献。 使用交互式流体力学教科书，（与J. A. Liggett），1999年国际水资源工程会议记录，美国土木工程师学会，华盛顿州西雅图，116年8月。 解决湍流反应流的完全联合PDF nansport方程的一致混合算法，（与P. 珍妮，S。 B. 教皇和M。 Muradoglu），被接受在燃烧研究所第二十八次燃烧研讨会（国际）上发表演讲。 117. 湍流反应流的PDF方程的一致混合有限体积/粒子方法，（与M. 穆拉多卢，P. Jenny和S。 B. 教皇），J. 补偿。 物理，卷。 154，pp. 342-371，1999年。 118. 如果飞机飞入太空会发生什么？，问一个科学家专栏，伊萨卡日报，2000年4月19日。 8-11，1999年。 119. 使用计算机教授流体力学，第二部分：电子教科书，Comp。 流体Dyn。 J.，卷。 9，不。 3，pp. 222-230，2000年10月。 120. 湍流反应流联合PDF方程的混合算法，（与P. 珍妮，S。 B. 教皇和M。 Muradoglu），J。 补偿。 物理，卷。 166，第218-252页，2001年。 121. 流体动力学和环境中的非稳态流动的隐式多网格计算与空气弹性的应用：动力学方法，J。 L。 Lumley，Ed.，Springer-Verlag，pp. 35-62，2001年。 122. 解决Steady的Euler 方程需要多少个步骤，可压缩流动：寻找快速解决方案算法，（与A. 詹姆森），AIAA论文2001-2673，AIAA第15届计算流体动力学会议，2001年6月11日至14日，加利福尼亚州阿纳海姆。 123. 平方截面圆柱体的非稳态流动的隐式多网格计算，计算机和流体，卷。 30，编号。 7-8，pp. 939-960，2001年9月至11月。 124. 悬崖体稳定火焰的JPDF计算，（与K. 刘，P。 珍妮，M。 Muradoslu和S。 B. 教皇）摘要提交到湍流火焰工作室。 125. 湍流反应流PDF方程的混合方法：一致性条件和校正算法，（与M. Muradoslu和S。 B. 教皇），J. 补偿。 物理，卷。 172，第841-878页，2001年。 126. Frontiers of 计算流体动力学 - 2002，与M共同编辑。 M。 哈菲兹，世界科学出版公司，新加坡，2002年。

M。 M。 哈菲兹33 127。 Robert W的贡献。 MacCormack到计算流体动力学，（与M. M。 哈菲兹），在计算流体动力学-2002年的前沿中，D。 A. Caughey和M。 M。 哈菲兹，编辑，世界科学，pp. 1-25，2002年。 128. 计算空气动力学，《物理科学与技术百科全书》第三版，卷。 3，pp. 469-485，学术出版社，2002年。 129. Euler和Navier-Stokes 方程的快速预制多网格解决方案，用于稳定、可压缩的流动，（与A. 詹姆森），AIAA论文2002-0963，第40届航空航天科学会议和展览，2002年1月14日至17日，内华达州里诺。 130。 大涡模拟 工具，（与 G. Jothiprasad），使用数值模拟数据库研究湍流-九：2002年暑期课程论文集，P. 布拉德肖，编辑，pp. 117-127，湍流研究中心，美国宇航局艾姆斯研究中心和斯坦福大学，2002年12月。 131. 非结构化网格上非稳定Navier-Stokes 方程的高阶时间积分方案，（与G. 乔蒂普拉萨德和D。 J。 马夫里普利斯），J。 补偿。 物理学，卷。 191，pp. 542-566，2003年11月。 132. 不稳定的跨声速流动过去“非独特”机翼，在Transsonicum IV研讨会中，Kluwer Academic，哥廷根，pp. 41-46，2003年。 133. 跨声速流动计算技术的发展：历史视角，（与A. 詹姆森），在Transsonicum IV研讨会中，Kluwer Academic，哥廷根，pp. 183-194，2003年。 134. Euler和Navier-Stokes 方程的快速预制多网格解决方案，用于稳定、可压缩流动，国际。 J。 数字。 液体中的冰毒，卷。 43，pp. 537-553，2003年。 135. 非定常流动过去机翼的稳定性表现出泛音速非独特性，CFD杂志，卷。 13，pp. 427-438，特别今井周年纪念版，2004年。 136. 具有高黏度的均质湍流的直接数值模拟，（与A. G。 拉摩根塞斯和S。 B. 波普），被接受发表在《流体物理学》上。 137. 为什么回旋镖会回来？，问一个科学家专栏，伊萨卡杂志，2005年3月24日。

34 Caughey对CFD 138的贡献。 使用带有详细化学成分的联合PDF模型计算虚张声势体稳定火焰，（与Kai Liu和Stephen B. 教皇），燃烧与火焰，卷。 141，pp. 89-117，2005年。 139. 关于结构化和非结构化网格的Euler 方程的对称高斯-塞德尔多网格解决方案，提交给国际。 CFD的J.，Antony Jnmeson的70岁生日特别刊。

\clearpage
\part{我。 设计和优化}
第一部分 设计和优化

此页面故意留空

\clearpage
\chapter{第2章计算流体动力学工程系统分析和设计M。 达莫达兰，S。 Alil和S。 达亚南丹}
\section{2.1 介绍}
第2章计算流体动力学工程系统分析和设计M。 达莫达兰'，S。 Alil和S。 Dayanandan' 2.1 介绍 计算科学和工程的技术发展，如高性能计算平台、网格计算基础设施、用于解决工程科学数学模型的快速并发算法和设计优化，正在为开发有效的模拟和设计优化工具以解决工业问题铺平道路。 工业中的问题往往是复杂和多学科的，其中许多问题需要使用计算流体动力学（CFD）模拟流体流动。 Davidson [3]概述了CFD在21世纪工业问题上日常使用的巨大潜力。 在本文中，简要概述了并讨论了一些需要广泛使用CFD并由作者最近解决的复杂工业问题。 一个涉及开发用于通风空气动力学、交通堵塞污染控制和火灾建模的计算模型，该模型正在新加坡最古老的地下隧道之一“南洋理工大学机械与生产工程学院”，南洋大道，新加坡639798。

\section{2.2 地下道路隧道火灾控制策略和情景规划的流程建模}
38 CFD工程系统（a）北行地下道路隧道的草图显示滑道（b）显示通风口、风扇和环境传感器位置的地下道路隧道示意图图2.1：新加坡最长的地下高速公路道路隧道，以及该计算模型的潜在应用，使当地工程师能够使用隧道中的各种流量调节设备模拟各种消防和污染控制策略的情景规划和决策。 考虑的另一个例子是开发一个计算模型，以研究硬盘外壳中的气流特性和粒子轨迹，其中所有组件都到位，以及该模型在设计和原型研究中的潜力。 这两个应用程序本质上都是计算密集型的，因此Fluent [5]套件中的商业CFD软体已与其他使用者编写的例程一起使用，以使用平行计算机获得洞察力和工程估计。 这些选定的例子讨论如下。 2.2 地下道路隧道消防策略和情景规划的流程建模图2.1（a）显示了新加坡目前正在建设的地下高速公路，以缓解交通拥堵。 隧道长约12公里，其中地下约8.5公里。 拟议的隧道设计包括一个主隧道，有15条滑道，使驾车者能够进入和/或离开隧道。 地下道路隧道建设的一个重要方面涉及通风的空气动力学，旨在将车辆尾气中的有害污染物排放保持在受控水平以下，并有效地

M。 达莫达兰，S。 阿里和S。 在消防部门到达现场扑灭火灾之前，Dayanandan 39控制车辆火灾。 图2.1（b）显示了排气口、供应通风口和喷气扇的拟议位置。 还规定了这些流动设备的运行模式，由沿隧道放置的环境传感器（用于处理污染检测）和烟雾探测器（用于处理火灾检测）的读数决定。 烟雾探测器的数量和位置尚未确定。这项工作的主要目标是开发一个有效的计算模型，该模型可用于预测地下隧道中所有可能的三维通风空气动力学流、污染和火灾情况，并验证监管机构提出的污染和火灾控制建议。 地下道路隧道通风研究的流动建模、模拟交通堵塞期间的污染控制和隧道内车辆火灾的模拟，以及通过操作流量调节和通风系统对隧道中模拟火灾的管理和控制，可以作为工程师执行场景规划的工具，以有效设计和控制隧道内通风。 在隧道的各个位置施加适当的边界条件，以模拟所需的通风空气动力学流动。 隧道入口和出口向大气敞开，并被处理为0帕表压力边界条件。 隧道内的条件决定了空气是进入还是离开隧道。 温度的环境值在门户边界上指定。 隧道壁假定在恒定的当地温度下。 喷射风扇被视为无穷小圆盘（即作为细胞之间的接口，而不是物理实体），每个风扇的压力跳跃是根据风扇和风扇区域的设计推力估计的。 在固流体界面，壁面函数方法用于消除解决线性子层的需要。 对于动量方程，无滑边界条件强加在固体表面。 通过考虑涉及所有流量调节单元处于活动模式的隧道中全通风空气动力学的场景，证明了该计算模型在分析隧道中气流的有效性。 这个场景提供了隧道最大可能通风的图片。 计算的量规压力等高线、速度等高线和等高线速度向量场沿着滑道D附近的水平平面绘制（如图所示）。 2.1）如图所示。 2.2（a）-2.2（b）。 有关此模拟和验证的各个计算方面的更多详细资讯，请参阅Dayanandan [4]。 为了模拟道路隧道中火灾的影响，并影响隧道中流量调节装置的激活，不模拟实际的燃烧过程。 相反，火被建模为热源和质量源。 火灾区域释放的热量作为能量方程中的源项引入，火灾区域释放的烟雾或物质作为物种运输方程中的源项引入。 正如Karki和Patankar[9]中描述的那样，燃料的对流热释放率Qc由以下方程给出

工程系统中的40 CFD（a）沿滑道附近水平平面的压力曲线D D（b）沿滑道附近水平平面的速度向量图2.2：新加坡最长的地下高速公路隧道Qc = kfuHfuq（l - XR），其中rhf.是燃料消耗率，Elfu是燃料的加热值，q是燃烧效率，XR是辐射分数。 假设完全燃烧，q的值=1。 辐射分数XR考虑到了来自火的热量，这些热量辐射到隧道的墙壁上，并且不会在隧道内对流。 从Markstein [lo]对扩散火焰进行的实验来看，XR通常取值在0.2到0.4之间。 如Karki和Patankar[9]所述，燃料的对流热释放率QC由以下方程k,,\&e = kfu(l + s)给出，其中msm\&e是烟雾的产生率，s是燃料的化学计量比，即燃烧过程所需的公斤空气与公斤燃料的比率。 隧道中的车辆火灾是用化学反应方程定义的甲烷燃烧建模的，即CH4 + 202 + 2H20 + COz。 根据工业供暖设备协会的手册[7]，甲烷燃烧的化学计量比为10，这相当于产生的烟雾率是甲烷消耗率的11倍的情况。 该模型已被用于模拟初始局部火灾的火灾增长，如地下隧道中的孤立汽车火灾，通过在火灾发生期间将火灾模型与非定常流动方程耦合，可以跟踪隧道内火灾的增长和蔓延，作为时间的函数。 发生火灾时，热烟的浮力使其在火上升。 当烟羽到达天花板时，它会分流成喷射物，在天花板上蔓延。 一旦羽流覆盖了天花板，烟雾层的深度就会增加，直到隧道部分充满烟雾，或者直到烟雾产生的速度等于烟雾排出的速度。 根据国家消防协会[ll]的说法，火灾管理的主要目标是保持一个环境，其中烟雾的影响和

M。 达莫达兰，S。 阿里和S。 Dayanandan LRDl -A JRlcl - 喷气扇（E）Slrp公路喷气扇（J-R）排气口（EX）Saccardo喷嘴（SN）图2.3：隧道段的典型火灾情况对乘员的热量，不会危及生命。 这意味着有效的火灾管理要求烟雾层的高度必须在一定时间内保持在隧道内人类占用的最高水平以上，这比疏散所需的时间要长。 沿着KPE隧道已经确定了一些可能的火灾情景，并就如何调整流量调节装置来控制火灾提出了建议。 同样，这些建议是基于简化的经验模型，必须使用更现实的三维数字流动模型来验证它们的准确性。 图2.3显示了隧道一段五起局部火灾的情景。 模拟了几个火灾场景，但在本文中只展示了与场景2相对应的油轮火灾的模拟，以说明火灾模拟能力。 关于处理第二种情况的流量调节设备运行模式的建议，要求排气EX1在供应模式（297 m3/s）下运行，排气EX2在排气模式（495 m3/s）下运行，SN1（供应通风口）在供应模式（495 rn3/s）下运行。 这是为了确保来自EX1的供应将沿着隧道推动烟雾，烟雾将被EX2的废气吸走，然后通过SN1的供应将新鲜空气泵入隧道。 根据Joyeux [8]的研究，油轮火灾的平均散热率约为100兆瓦。 假设完全燃烧，并假设辐射分数为0.3。 因此，甲烷的所需消耗率约为2.4千克，消耗此量甲烷时的烟雾释放率为26.4千克/秒。 热量释放被引入为能量方程中的源项，烟雾释放被建模为物种运输方程中的源项。 最初，隧道内的所有流量调节装置都被关闭，并解决了此条件的稳态流量。 使用该解决方案作为初始条件，在4米乘2米乘1.5米的范围内引入了100兆瓦的热量释放和26.4千克的烟雾释放。 随后，非定方程被求解100秒。 正如Dayanandan [4]中描述的那样，烟羽的生长用7Ooc等温度可视化。 烟雾

工程系统中的CFD（a）起火后25秒（b）起火后50秒（c）起火后100秒图2.4：起火后起火后25秒、50秒和100秒的烟柱生长如图所示。 2.4（a）-2.4（c）。 据推测，传感器会在火灾启动后检测到100秒。 因此，在这一点上，火灾情景2的建议适用于管理火灾。 建议开始实施后，50s、loOs、150s和200s的烟羽流如图所示。 2.5（a）-2.5（d）。 从这些数字可以看出，消防控制建议的实施有效地驱散了烟雾。 该模型最终将与道路交通模型相结合，以预测隧道中的车辆集中度，这些信息也可用于隧道中车辆排放的污染控制。 （A）激活后50s tivat ion tivation tivation tivation（b）ac后100s-（c）ac-后150s-（d）ac-后200s-图2.5：根据拟议策略控制烟雾

\section{2.3 硬盘驱动器盘柜中的流程建模}
M。 达莫达兰，S。 阿里和S。 Dayanandan 43（a）硬盘外壳中的粒子轨迹（b）硬盘的简化区域图2.6：硬盘外壳2.3硬盘驱动器外壳中的流量建模作者目前正在调查的一个有趣的行业问题涉及硬盘驱动器（HDD）外壳内的气流特性对操作过程中HDD组件的颗粒运输和气流诱导振动的影响。 目前正在解决的问题包括流体-结构相互作用（FSI）和粒子动力学建模，因为气流引起的振动和粒子污染物运输是硬盘故障的主要原因。 建模的目的是深入了解气流特性，从而可能改进HDD外壳的设计。 目前对FSI的大部分调查都集中在Imai [6]中概述的气流对磁盘颤动的影响，以及Watanabe等人[12]中对手臂的振动的影响，因为这两个区域是硬盘驱动器（HDD）中的关键组件。 工业标准硬盘外壳内气流特征和粒子轨迹的详细计算研究，如图所示。 2.6（a）在Ali等人[2]中报道。 在这里进行的初步调查中，如图所示，一个带有单个磁盘的简单硬盘外壳。 2.6（b）用于研究圆盘主轴单元与周围空气之间的耦合流体结构相互作用。 简化的机箱边界处的边界条件来自整个硬盘机箱内的完整流程模拟。 假定主轴由铝合金制成，密度为2.7克/立方厘米，E = 70 GPa和= 0.33。 圆盘由96\%的二氧化硅玻璃制成，密度为2.18克/立方厘米，E = 68 GPa和=0.19。 结构问题被归结为线性弹性问题，并研究了其在流体压力下的线性变形。 同样，还研究了圆盘主轴单元变形引起的流体位移的影响。 方法

工程系统中的44 CFD（a）模式1（f = 1608.5）（b）模式2（f = 1614.9）（c）模式3（f = 1622.3）（d）模式4（f = 2141.9）图2.7：本次调查的圆盘主轴单元的归一化特征模式和特征频率是流体和结构区域之间的弱耦合，其中每个区域的控制方程分别解决，仅通过共享边界交换。 共享的边界条件是作用在圆盘主轴单元上的流体压力以及流体侧和结构侧之间的速度一致性。 使用非线性有限元软件Abaqus 111作为结构求解器，Fluent [5]作为流体求解器，弱耦合是通过Fluent的动态网格功能完成的，该功能模拟了从Abaqus获得的结构位移。 使用Abaqus计算的圆盘主轴单元的自然模式和频率如图所示。 2.7. 从流体侧获得的稳态不可压缩流动溶液，用于以5400转/分旋转的圆盘主轴单元，作为本次调查的起点。 椎间盘旋转在2秒后缓慢减少到停止，并研究了流体结构相互作用。 图中得到的结构变形。 2.8代表作用在圆盘上的压力差，如图所示。 2.9. 通过比较速度向量的图可以看出，10纳米的变形太小了，无法影响流动模式

\section{2.4 结束讲话}
M。 达莫达兰，S。 阿里和S。 Dayanandan 45 (a) t = 0.1 s (b) t = 0.3 s (c) t = 0.5 s (d) t = 0.8 s 图2.8：不同时间点的位移大纲图2.9：从t=O.ls到t=0.8s的圆盘周长压力差。 起始位置在（0.0466，0）和径向位置是逆时针的。在臂附近，对于有和没有结构变形的情况，如图所示。 2.10. 对于粒子动力学，HDD研究侧重于张和博吉[14]提出的拉格朗日模型，该模型跟踪单个粒子。 Whitby等人[13]概述的通过细颗粒模型初步使用气溶胶动力学来模拟颗粒污染似乎很有希望，从使用0.2微米至0.5微米的颗粒的简单硬盘外壳中获得的结果所示。 图中的气流特性。 如图所示，2.11可以看到影响粒子分布。 2.12，正如理论上所预测的那样，随着分子集中向壁附近的更高压力区域移动。 由于离心力的作用，壁附近的粒子浓度由较大的粒子组成。 2.4 结束词 从这些例子中可以看出，CFD在广泛的工业问题和商业的日常使用中发挥着重要作用

\section{2.5 参考书目}
工程系统中的46 CFD图2.10：0.5秒内圆盘边缘周围速度的向量图。 CFD软件以及这些软件与其他用户编写的代码的交互，以解决多学科工程设计问题，将继续为工程系统的原型设计以及场景规划提供更丰富的见解，大大提高生产力和成本节约。 应该承认，虽然许多商业CFD求解器在湍流建模领域存在局限性、网格质量的影响和其他计算问题，但各个小组在该领域正在进行的研究将产生新的湍流模型和修复，最终将进入这些商业CFD求解器，从而减少这些限制并扩大CFD的应用范围。 参考书目ABAQUS软件用户指南6.3-1，2003年。 阿里，S.，宋，H. B.，Damodaran，M.，\& Ng，Q。 Y.，几何结构对硬盘驱动器外壳中气流特性和粒子轨迹的影响，2005年2月12日至15日在美国佛罗里达州奥兰多举行的SIAM计算科学与工程会议上接受演讲。 Davidson，D。 L.，计算流体动力学在工艺工业中的作用，The Bridge，美国国家工学院的出版物，华盛顿特区。 C.，美国，第32卷。 第4期，2002年冬季。 Dayanandan，S.，地下道路隧道通风空气动力学的计算建模，MEng论文，机械学院和

M。 达莫达兰，S。 阿里和S。 Dayanandan 47（a）仪表压力轮廓（b）速度向量图2.11：HDD外壳内的气流特性。 生产工程，新加坡南洋理工大学，2004年。 [5] FLUENT 6.1，流程模拟软件用户指南，Fluent Inc，美国新罕布什尔州，2003年。 [6] Imai，S.，通过测量圆盘之间的压力来进行圆盘颤动的流体动力学机制，IEEE磁学交易，第37卷。 不。 2，第837-841页，2001年。 [7]工业加热设备协会，燃烧技术手册-甲烷燃烧，第4版，1998年。 [8] Joyeux，D.，封闭停车场的自然火灾，报告编号 INC-96/294d-DJ/NB，1997年。 [9] Karki，K.C. \& Patankar，S.V.，CFD喷气风扇通风系统模型，Proc，第10届汽车隧道空气动力学和通风国际研讨会，BHRA流体工程，美国波士顿，2000年。[lo] Markstein，G.H.，烟雾点与浮力湍流和层流扩散火焰辐射发射之间的关系，第20届（国际）燃烧研讨会，燃烧研究所，匹兹堡，第193-329页，1984年。 [Ill NFPA，NFPA标准92B，心房、有盖的商场和大区域的烟雾管理系统指南。 马萨诸塞州昆西：国家消防协会，1995年。

48 CFD in Engineered Systems NWOI xxu \%"Ewe I(sd -,.a SA, CmPiunolw\textasciitilde{}"l --r Nam 2ay FUlEM li 1 (sd 51) MPinolw\&drn-m 图2.12：基于气溶胶动力学的粒子分布 [12] Watanabe, T., Gross, H. M.，和Bogy，D. B.，手臂厚度对同旋转圆盘之间流动诱导头部振动的影响，CML研究报告，部门。 机械。 工程，加州大学伯克利分校，2001年。 [13] Whitby, E., Stratmann, F., \& Wilck, M., FLUENT手册的细颗粒模型（FPM），FPM 1.0.4，2003年。 [14]张，S. \& Bogy，D。 B.，滑块/磁盘界面中粒子的运动，包括升力和壁效应，IEEE Trans。 Mag.，第33卷，第 5，第3166-3168页，1997年。

\clearpage
\chapter{第3章 气动外形优化安东尼·詹姆森的进展}
\section{3.1 介绍}
第3章 气动外形优化安东尼·詹姆森的3.1导言 自1970年左右作者首次参与计算流体动力学以来，情况发生了巨大变化。 当时，面板方法刚刚投入使用，世界上最快的超级计算机Control data 6600只有131000字的记忆体（约1兆位元组）。 在Murman和Cole [1970]的突破之前，没有发现用激波计算跨声速流动的可行算法。 到1980年，超级计算的标准是Cray 1，它实现了约100兆浮点的性能，但至少在最初很难获得内存超过128兆字节的Cray。 目前，许多笔记本电脑的处理速度为2-3千兆赫，内存为1千兆字节，远远超过了80年代中期的Cray XMP。 事实上，从1986年的80386到目前的奔腾4，英特尔微处理器的速度在17年内提高了一千多倍。 这些事态发展在1970年是难以想象的。 至少在一些空气动力学问题上，算法也取得了几乎同样巨大的进步。 部分源于Godunov[4]的开创性工作，开发了许多有效的冲击捕获算法。 此外，虽然1980年解决稳态Euler 方程的可用方法需要5000-10000次迭代才能达到合理的“斯坦福大学航空和航天系，斯坦福，CA 94305-

收敛的50 气动外形优化水平，没有一个会完全收敛到机器零[34]，现在可以通过3-5个步骤获得Euler 方程围绕翼型流动的解[lS]。 作者在《计算力学百科全书》[14]的一篇文章中回顾了这些发展。 一些问题，如过渡和分离的预测，或通用湍流模型的制定，仍然顽固。 然而，软件和硬件的综合进步使得解决比30年前想象的复杂程度高很多的问题成为可能。 即使在一开始，智能使用计算流体动力学（CFD）也可能对设计产生重要影响，而本作者一直认识到，真正的挑战不仅仅是预测给定形状的流动，而是找到一个优越的形状，根据一些有用的标准进行优化。 事实上，作者的第一个CFD程序Synl（1970年7月）为设计理想（非自转和不可压缩）流动的翼型的逆问题提供了一个完整的解决方案，该流将产生指定的目标压力分布。 源于与道格拉斯飞机的马尔科姆·詹姆斯的讨论，该方法找到了将圆转换为所需的翼型的共形映射。 这是Lighthill方法的延伸，Thwaites [36]将其描述为一个不完整的解决方案，因为它需要在圆平面上指定目标速度。 Synl的输入是作为弧长s的函数的目标压力。 然后，由于沿着轮廓的电位在圆平面上是和4，圆平面上的角度0可以通过牛顿迭代被定为s的函数。 如果目标压力不可实现，Synl会找到产生最接近可实现压力分布的形状。 如今，它在笔记本电脑上运行的时间不到八十年代末，在Hicks和Henne [8]使用翼型和机翼设计的数值优化的一些早期实验之后，解决气动外形优化（ASO）的一般问题的时机似乎已经成熟。 在1980年2月参加了ICASE关于流量控制的研讨会后；作者认为，控制理论提供了一条通往AS0的间接路线，这条路线可能比之前尝试过的方法更有效。 作者随后发现，Pironneau也为椭圆方程探索了使用外形优化控制理论的想法[31]。 偏微分方程的控制理论，其中控制采取边界运动的形式，是变分法的自然扩展，它使代价函数相对于形状的无限维（Frechet）导数能够通过并行方程的解来确定。 然后，这个梯度信息可用于改进形状，这个过程可以重复，直到形状收敛。 有了这种方法，一个人做第二。

A. 詹姆森51不考虑10-100范围内的一些设计参数。 相反，该形状被视为自由曲面，在离散模型中可能由曲面网格点表示，或适当基础函数集的扩展。 例如，在翼型设计的情况下，可以用与洛朗级数相对应的傅里叶系数来描述轮廓，该系数定义了圆的共形映射。 线性PDE的控制理论是在Lions的经典著作中制定的[28]。 扩展到具有可能不连续解决方案的非线性PDE提出了一些困难的问题，其中一些问题仍未解决。 然而，作者在1988年为跨声速势流和Euler 方程推导出了必要的伴随方程[9]，并在同年晚些时候开发了跨声速势流的翼型设计软件。 第一个数值结果于1989年公布[lo]。 还开发了一个初步的欧拉辅助代码（Syn82），并获得了AFOSR的支持，以进一步追求该概念。 需要探讨的问题之一是，从PDE中以连续形式推导出连续形式的伴随PDE，然后离散化，（“连续伴随”方法），还是先离散化流动方程，然后直接推导离散伴随方程（“离散伴随”方法）更好。 在作者看来，推导连续伴伴方程很重要，以深入了解方程的属性和适当的边界条件。 但适当的离散化当然应该反映流动方程的离散化。 例如，如果使用上风方案，伴风离散化就会显示为下风方案（实际上，伴随方程中向相反方向传播的波浪的上风）。 当使用带有非线性限制器的冲击捕获方案时，离散伴随方法会产生非常复杂的离散方程。 在实践中，连续伴伴方法已被证明非常有效，但有时用离散方法治疗边界条件更容易[30]。 本文主要侧重于连续伴伴方法。 伴随方程组的形式与流动方程相似，因此为流动方程[21、15、21）开发的数值方法可以用于伴随方程。 虽然从伴随解获得的梯度信息可以输入到任何基于梯度的搜索程序中，但事实证明，在实践中，通过二阶微分方程隐式平滑梯度来在定义的方向上重复小步骤是非常有效的。 这个过程保证了重新设计的形状序列的平滑性，相当于在索博列夫空间中重新定义梯度，它作为一个有效的预调节器，通常在10-20个设计周期内产生最佳值。 还表明，使用连续伴伴方法（但不是离散方法），可以直接从伴伴解和表面运动中推导出梯度，独立于网格修改。 这消除了评估取决于网格扰动的体积积分的需要。 如果一个人希望使用非结构获得点梯度-

\section{3.2 优化程序的制定}
\subsection{3.2.1 梯度计算}
52 气动外形优化 tured网格，这些积分变得非常昂贵，因为网格变形的传播必须为每个表面网格点的偏转单独计算。 因此，它们从梯度中消除为使用非结构网格的外形优化开辟了道路。 在过去的十年里，基于控制理论的优化技术已经发展成为一种强大的产品工具。 它们已成功应用于超音速和超音速设计，并在美国宇航局的HSR计划中发挥了重要作用[33、37、23、321。 它们已被扩展到包括使用雷诺平均Navier-Stokes 方程的黏性效应[19，241。 这很重要，因为无黏性优化可能导致陡峭的不利压力梯度，从而导致分离。 此外，延迟跨声速阻力上升所需的机翼截面修改通常与边界层位移厚度相同，因此适当的设计必须允许边界层的效果。 最近，机翼平面参数被纳入设计变量，斯坦福大学航空航天计算实验室成功设计了一个机翼，该机翼以最小阻力产生指定的升力，同时满足其他标准，如低结构重量、足够的燃料量以及稳定性和控制[25]。 使用非结构化网格技术对动态设计具有相当大的希望，它促进了复杂配置的处理，而不会在网格生成中产生令人望而却步的成本和瓶颈。 使用非结构网格进行设计的计算可行性基本上是通过使用连续伴从方法和“减少梯度”式式[22]实现的。 完整配置的代表结果显示在最后一部分。 3.2 优化程序的制定 3.2.1 梯度计算 对于所考虑的航空力学优化问题，设计空间本质上是无限的。 假设系统设计的性能可以通过代价函数 I来测量，该代价函数 I依赖于描述形状的函数F(x)，在设计SF(z)的变化下，成本的变化是SI。 现在假设SI可以表示为SI = G(z)SF(z)dz s，其中G(z)是梯度。 然后通过设置SF(2) = -XG(z)

A. 詹姆森53获得改进61 = -A @(x)ds s，除非G(x) = 0。 因此，梯度的消失是局部最小值的必要条件。 计算复杂系统的代价函数梯度可能是一项数值密集型任务，特别是如果设计参数数量大，并且代价函数是一个昂贵的评估。 最简单的优化方法是通过一组设计参数来定义几何形状，例如，这可能是应用于一组形状函数\&(x)的权重ai，因此形状被表示为然后选择代价函数 I，这可能是阻力系数或升力与阻力比；I被视为参数ai的函数。 现在可以通过在每个设计参数中依次进行小幅变化6ai，并重新计算流量以获得I的变化来估计灵敏度。 那么，这种有限差分方法的主要缺点是，估计梯度所需的流量计算数量与设计变量的数量成正比[7]。 同样，如果诉诸于直接代码微分（ADIFOR [3，5]）或复变量扰动[I]，则确定梯度的成本也与用于定义设计的变量数量成正比。 一种更具成本效益的技术是通过解决伴随问题来计算梯度，例如参考文献[18，12，111。 基本想法可以总结如下。 对于任意物体的流动，定义代价函数的空气动力学属性是流场变量（w）和物体物理形状的函数，可以用函数F表示。 然后I = I（w，F）和变化。 F导致代价函数 l3IT dIT 61 = -6w + -63的变化。aW a3使用从控制理论中提取的技术，将流场的控制方程引入为约束，这样梯度的最终表达不需要重新评估流场。 为了

54 气动外形优化实现这一点，6w必须从上述等式中剔除。 假设控制方程R表示流场域D中w和3的依赖性，可以写成R(w,3) = 0。 （3.1）然后从方程中确定6w 接下来，引入拉格朗日乘数？ 我，我们有一些重新安排选择11，以满足伴随方程 [ElT\$= dIT "I> aF 63。 （3.3）在代价函数的变化中可以消除乘以6w的术语，我们发现61 = 663，其中dIT a3 [ 3 6=--11,T 优点是代价函数的变化与6w无关，因此I相对于任意数量的设计变量的梯度无需额外的流场评估即可确定。 如果（3.1）是偏微分方程，则伴随方程（3.3）也是偏微分方程，必须确定适当的边界条件。 事实证明，合适的边界条件取决于代价函数的选择，并且可以很容易地推导出涉及表面压力积分的成本函数。 涉及场积分的成本函数导致伴随方程中出现源项。 求解伴随方程的成本与求解流动方程的成本相当。 因此，无论设计空间的维度如何，获得梯度的成本都与两个函数评估的成本相当。

\section{3.3 使用Euler 方程进行设计}
A. Jameson 55 3.3 使用Euler 方程的设计 控制理论在空气动力学设计问题中的应用在本节中说明了使用可压缩Euler 方程作为数学模型的三维机翼设计的情况。 参考文献[19、17、131）介绍了治疗Navier-Stokes 方程的方法的扩展。 事实证明，用21、52、23和u1、u2、u3表示笛卡尔坐标和速度分量很方便，并使用重复指数i = 1到3的总和的惯例。 然后，三维Euler 方程可以写成dw dfi在axi -+-=O inD,, fi = (3.4) (3.5)，Sij是Kronecker delta函数。 此外，pH=pE+P（3.7），其中y是比热的比率。坐标为el，E2，[3）的固定计算域的解，为了简化伴随方程的推导，我们映射了K..- 23 - [:I]，- J = det (K)，s = JK-\textasciitilde{} S的元素是K的辅因子，在有限体积离散中，它们只是投影在xi、22和23方向的计算单元的表面区域。 使用排列张量Eijk，我们可以将S的元素表示为

56 气动外形优化 在随后分析形状变化的影响时，需要注意的是，ax, ax, s,j = Ejpg - -, 862 at3 ax, ax, s2j = €jpg - -, at3 at1 ax, ax, s3j = €jpq--. at1 at2 现在，将方程(3.4)乘以J并应用链式规则，其中aw J- + R(w) = 0 afj - a R(w) = s..- - - (Sijfj), a3 a\& a\& at使用（3.9)。 我们可以将遍体速度分量写成（3.10）（3.11）（3.12）变换通量，以缩放的con- r puiu1 + SilP 1为方便起见，选择描述固定计算域的坐标ti，以便每个边界都符合这些坐标之一的常数值。 然后，形状的变化会导致Ki定义的映射导数的相应变化。 假设性能由包含边界和场贡献的代价函数来测量，其中dBc和dDc是计算域中的表面和体积元素。 一般来说，M和P将取决于流变量w和定义计算空间的度量S。 设计问题现在被当作控制问题

A. 詹姆森57，其中边界形状代表控制函数，选择控制函数将I最小化，受流方程定义的约束（3.11）。 形状变化产生流动解6w和指标6s的变化，这反过来又产生代价函数的变化（3.13）的变化，这可以分为SI = SIl + SIII，（3.14）6M = [\&]I\textasciitilde{}W+\textasciitilde{}MII，6P = [PW]，6W+6P1I，（3.15），我们继续使用下标I和II来区分与流动解决方案6w变化相关的贡献和与度量变化相关的贡献。 因此，[\&I和[P,]代表和\$\$与固定的指标，而SMII和6P11代表指标变化6s对6M和6P的贡献。 在稳态中，约束方程（3.11）指定了状态向量6w的变化（3.16）在这里，6R和6Fi可以使用符号a at2 6R = -6Fi = 0分割为与6w和6s相关的贡献。其中afi [Fa,], = saj-. aW（3.17）乘以共状态向量\$，它将发挥与方程（4.4）中引入的拉格朗日乘法相似的作用，并在域上积分产生（3.18）假设\$是可微的，下标I的项可以按部分积分以给出

58 气动外形优化 该方程直接来自流动方程的弱形式的变化，其中1c被视为一个任意的可微测试函数。 由于左手表达式等于零，它可以从代价函数（3.13）的变化中减去，以给出Now，因为1c是一个任意的可微函数，因此可以选择SI不再明确依赖状态向量6w的变化。 然后可以直接从度量变化中评估代价函数的梯度，而不必重新计算由每个设计变量的扰动引起的变化Sw。 比较方程（3.15）和（3.17），可以通过将所有域项与下标“I”等同，从（3.20）中消除变化Sw，从而产生一个管理\$（3.21）的差异伴随系统。取方程（3.21）的转置，在代价函数中没有场积分的情况下，非黏性伴随方程可以写成T C，- = 0在D，xi（3.22）中，其中变换空间中的非黏性雅各布矩阵由方程（3.20）中的子脚本“I”边界项等同来产生相应的伴伴边界条件，在B上产生niQT [\%I，= [\&]I。 （3.23）然后，方程（3.20）中的其余项产生了代价函数变化的简化表达式，该表达式定义了梯度（3.24），该梯度纯粹由包含度量变化的项组成，流动解是固定的。 因此，一旦定义了网格扰动和形状变化之间的关系，就可以推出梯度的显式公式。

A. Jameson 59 梯度公式的细节取决于边界形状作为设计变量的函数的参数化方式，以及边界修改时网格变形的方式。 利用网格变形和曲面修饰之间的关系，通过沿着表面发出的坐标线积分，将场积分简化为曲面积分。 因此，61的表达式最终简化为SI = \textasciitilde{}g6的形式。 FdB\textasciitilde{}其中3表示设计变数，G是梯度，这是在边界表面上定义的功能。 流动方程满足的边界条件限制了伴随边界条件左侧的形式（3.23）。 因此，对代价函数 M的边界贡献不能任意指定。 相反，它必须从允许取消方程（3.20）边界积分中包含Sw的所有项的函数类别中选择。 另一方面，对代价函数 P的场贡献的规范没有这种限制，因为这些项可能总是作为源项被吸收到伴联场方程（3.21）中。 为了简单起见，将假定下进行形状修改的边界部分仅限于坐标曲面\& = 0。 然后，方程（3.20）和（3.23）可以通过纳入条件n = 1来简化，因此只需要在墙边界考虑变化6F2。 在E2 = 0时没有流经壁边界的条件等同于u, = 0, SU2 = 0 n1 = 723 = 0, d\& = d\&d[3，因此当边界形状被修改时。 因此，边界的黏性通量的变化减少到s23，因为6F2只取决于压力，现在很明显，边界M（w，S）上的性能测量可能只是压力的函数和

60 气动外形优化公制术语。 否则，边界积分中包含6w的项完全取消是不可能的。 例如，人们可以包括代价函数中力和力矩的任意测量，因为这些是表面压力的函数。 为了设计一个将导致所需压力分布的形状，一个自然的选择是设置pd是所需的表面压力，并在实际表面积上评估积分。 在计算域中，这被转换为数量表示与计算域中面面积的单位元素相对应的面面积。 现在，为了取消边界积分对bp的依赖性，伴随边界条件减少到nj是曲面法线的分量，这相当于与动量分量相对应的共态变量上的蒸腾边界条件。 请注意，它对边界处\$J的切向分量没有施加限制。 最后，我们发现，在这里，成本变化的表达式取决于域积分中出现的整个域的网格变化。 然而，形状变化的真正梯度不应该取决于网格变形的方式，而只取决于真正的流动解决方案。 在下一节中，我们展示了如何消除场积分，以产生仅取决于边界运动的简化梯度公式。

\section{3.4 减少梯度配方}
A. 詹姆森61 3.4 减少梯度表述考虑具有固定边界的网格变化的情况。 然后，61 = 0，但变换通量存在变化，6Fi = Ci6\textasciitilde{} + 6Sij fj。 在这里，真正的解决方案是不变的。 因此，变化6w是由于每个网格点的网格运动6倍。 因此，dW 6w = vw.6x = -sxj (= SW*) dXj，由于d -6Fi = 0，因此（3.28）下面验证了这种关系在边界运动的一般情况下成立。 现在在墙边界Czbw = 6F2 - 6S2j fj。 （3.30）因此，通过选择q5来满足伴随方程（3.22）和伴随边界条件（3.23），我们将成本变化降低到只取决于表面位移的边界积分：（3.31）

62 气动外形优化 为了完整，这里介绍了方程（3.28）的一般推导。 使用公式（3.8）和属性（3.9）a la - 1 862，ax，+--)！ \& ax, 862, ZEjpqEirs (z Rr ats 1 ax 8.f. - - ZEjpqEirs \{ \& (\textasciitilde{}xpgs)\} a Now express 62，就原始计算坐标ax 62的偏移而言，= \$gk。 然后我们得到fpqj Gsi - -- a a\& atr - (bS2j fj) = -（\&中的术语在这里乘以6\&是ax，ax，a fj ax，ax，a fj at1 at2 at3 at1 at2 --- - ---根据公式（3.10），这可以被识别为af 1 8.f 1 s2j - + s3j - at2 at3 8f 1 -4j -. a51或，使用方程的准线性形式（3.12）进行稳定流动，因为乘以6\&和St3乘以的项是ax，ax，a fj ax，ax，a fj KzK - \%KG（3.32）（3.33）

\section{3.5 优化程序}
\subsection{3.5.1 梯度定义中对索博列夫内积的需求}
A. 詹姆森63和axp ax，a fj axp ax，d fj --- - \textasciitilde{}\textasciitilde{}- a\textasciitilde{} at2 at3 at3 x3 at2 因此，\&中的术语被简化为最后，\&和\&中的术语的类似简化，我们得到了有待证明的。 3.5 优化程序 3.5.1 梯度定义中需要索博列夫内积 成功实施连续伴随方法的另一个关键问题是为梯度定义选择合适的内积。 事实证明，使用修改后的索博列夫梯度具有巨大的好处，它能够生成一系列光滑的形状。 这可以通过考虑变分微积分中问题最简单的案例来说明。 假设我们希望找到具有固定端点y(a)和y(b)的最小化b I = / F(Y, i)dx a的路径y(x)。 在（x）的变化下，因此将梯度定义为aF d dF g=dy-zd？ /'

内积为b 气动外形优化 (u,w) = 1 uvdz a，我们发现，如果我们现在设置6y= -xg，x >o，我们得到一个改进6I = -X(g,g) 5 0，除非g = 0，必要的条件@最小。 请注意，g是y、y、y、9 = dY、Y、Y')的函数，例如，在众所周知的Brachistrone问题中，当粒子落入重力下时，它要求确定两个侧向分离点之间的最快下降路径，以及1 + y'2 F(Y,Y') = y J 1 + y'2 + 2yy" 2 (y(1 + y'2))3'2 g=-可以看出，每个步骤yn+l = yn - xngn将y的平滑度降低两类。 因此，计算的轨迹变得越来越不流畅，导致不稳定。 为了防止这种情况，我们可以引入加权索博列夫内部产品[20]（u，w）= /（uu + Eu“Ul）dz，其中E是控制导数权重的参数。 我们现在定义梯度jj，使61 =（Q，6Y）

\subsection{3.5.2 外形优化的索博列夫梯度}
A. Jameson 65 然后我们有SI = (ijSy + c<Sy')dx J，其中3 = 0在终点。 因此，ij可以通过平滑方程从g中得到。 现在，该步骤有所改进，但yn+'具有与yn相同的平滑度，因此过程稳定。yn+l = yn - Anijn SI = -Xn(ijn,gn) 3.5.2 外形优化的索博列夫梯度 在将控制理论应用于气动外形优化时，使用索博列夫梯度对于保持重新设计的表面的平滑度等级同样重要。 因此，使用上面定义的加权索博列夫内积，我们定义了一个修改后的梯度4，使61 =< Q,63 >。 在一维情况下，Q是通过求解平滑方程（3.34）获得的。在多维情况下，平滑以乘积形式应用。 最后，我们设置SF = -AS（3.35），结果是除非S = 0，否则G = 0。给定的节点i可以表示为SI = -A < 8，Q > < 0，当二阶中心差分应用于（3.34）时，Gi - E（Bi+1 - 2Bi + Qi-1）的方程= \textasciitilde{}i，1 5 i 5 n，其中Gi和Qi分别是平滑前后节点a的点梯度，n是在这种情况下，设计变量的数量等于网格点数。 然后，B = AG，

\subsection{3.5.3 设计程序大纲}
66 气动外形优化 Flow Solution djoint Solutio 梯度计算 Q I Sobolev Gradient I 图3.1：设计周期，其中A是n x n三对角矩阵，使1+2€ -€ 0。 0 A-1=[...: -€ 1. 1 + 2E 在每次设计迭代中使用最陡峭的下降方法，一步，6。 F，被采取这样6。 F = -XAB。 （3.36）从这种表达式的形式可以看出，隐式平滑可以被视为一种预处理，它允许在搜索过程中使用更大的步骤，并导致收敛所需的设计迭代次数大大减少。 3.5.3 设计程序最终可以总结如下：设计程序大纲1。 求解p，UI，'112，213，p的流动方程。 2. 在适当的边界条件下解决4个伴随方程。

\section{3.6 案例研究}
\subsection{3.6.1 跨声速翼型设计的二维研究}
A. 詹姆森67英尺3。 评估G并计算相应的索博列夫梯度c。 4. 将应变。投射到一个满足任何几何条件的允许子空间中。 根据最陡峭的下降方向更新形状。 6. 返回1，直到达到收敛。 设计方法的实际实现在很大程度上依赖于状态（w）和共状态（+）系统的快速和准确的求解器。 第3.6节中获得的结果是使用经过验证的欧拉和Navier-Stokes 方程解决方案经过多年开发的软件获得的[21、29、351。 逆向设计，升降机由目标压力固定。 在阻力最小化中，修复升力系数也是合适的，因为诱导阻力是总阻力的很大一部分，只需通过降低升力就可以减少阻力。 因此，在每个流动解决方案期间调整攻角，以迫使达到指定的升力系数，并且攻角变化的影响包含在梯度的计算中。 涡流阻力也取决于跨度载荷载，跨度载可能受到其他考虑因素的限制，如结构载荷载或冲击。 因此，通过调整扭转分布以及流动解决方案期间的攻角，提供了强制跨度负载的选项。 3.6 案例研究 3.6.1 跨声速翼型设计的二维研究 当使用非黏性Euler 方程对流动进行建模时，阻力的来源是激波引起的工资阻力。 因此，如果形状在固定升力下被优化为最小阻力，那么最好的结果是零阻力的无冲击翼型。 根据这个标准，最佳形状是完全不独特的，因为所有无冲击的轮廓都同样好。 作者在过去15年的经验证实，无冲击型材可以从各种初始形状中获得，同时保持固定的升力系数和固定的厚度。 最近，作者的二维尤拉设计代码Syn83被用来探索马赫数和升力系数的可实现极限，在该极限下可以实现给定厚度的无冲击机翼[6]。 当设计目标是两个极端时，性能往往会迅速下降到设计点，强烈的双重冲击通常出现在设计点下方。 因此，Cl-Mach空间中无冲击翼型的边界有些模糊。 然而，这项研究证实，对于十分厚的机翼，人们可以

68 气动外形优化 图3.2：各种形状优化的机翼的可实现无冲击解决方案沿着穿过Cl.6和Mach.78和Cl.7和Mach.77的边界实现良性无冲击形状。 其中第二点如图所示。 3.3. 随着厚度的减少，边界向上移动。 事实上，跨声速相似性规则可用于找到逐渐变薄的轮廓，这些轮廓在增加马赫数时没有冲击。 此外，与典型的平顶和后部装载的超临界截面没有相似之处的轮廓可以实现无冲击流动。 然而，似乎后部装载，也许在发散后缘的帮助下，可以帮助将无冲击流动扩展到更高的升力系数。 优化前优化后图3.3：DLBA-243 翼型的压力分布和马赫轮廓

\subsection{3.6.2 B747欧拉平面图结果}
A. Jameson 69 3.6.2 B747 Euler Planform 结果 提高跨声速性能所需的机翼部分变化非常小。 然而，为了获得真正的最佳设计，应考虑更大的规模变化，如机翼平面的变化（背、跨度、弦、截面厚度和锥度）。 由于这些直接影响结构重量，只有考虑同时考虑空气动力学特性和重量的代价函数才能获得有意义的结果。 在参考[25、27、261中，代价函数被定义为CW 3-是机翼重量的无维度量，可以从统计公式或对代表性结构的简单分析中估计，允许面板屈曲等失效模式。 引入系数a2是为了给设计者提供对压力分布的一些控制，而阻力和重量的相对重要性由系数a1和a3表示。 通过改变这些，可以计算出给定阻力系数重量最小的设计的帕累托正面，或给定重量的最小阻力系数的设计。 这些相对重要性可以从宝璏范围方程中估计；4m ST，f图3.4显示了从波音747机翼研究中获得的帕累托前线[26]，其中流动由Euler 方程建模。 机翼平面和截面同时变化，平面由六个参数定义；回扫、跨度、三个跨度站的弦和机翼厚度。 重量是根据对结构箱所需的截面厚度的分析估算的。 该图还显示了帕累托前部的点，当选择2，使飞机的航程最大化。 最佳机翼，如图所示。 3.5，在机翼的内侧部分具有更大的跨度、较低的扫角和较厚的机翼部分。 跨度的增加导致诱导阻力的减少，而截面形状的变化使冲击阻力保持在低水平。 与此同时，较低的扫角和较厚的机翼部分减少了结构重量。 总体而言，最佳机翼提高了空气动力学性能和结构重量。 阻力系数从108计数减少到87计数（19\%），而重量系数CW从455计数减少到450计数（1\%）。

0 052 0 05 0 0 048 0 046 0 044 0 042 0 04 0 0 038 气动外形优化 帕累托前部优化截面，固定平面基线-7drange / \textbackslash\{\} x =优化截面和平面图 85 90 95 100 105 110 80 图3.4：帕累托前部截面和平面修改图3.5：波音747的基线（较短跨度）和优化截面和平面（跨度较长）几何形状的叠加。 重新设计的几何形状具有更长的跨度、较低的扫角和较厚的机翼部分，提高了空气动力学和结构性能。 优化在0.87马赫和固定CL.42下进行，其中选择2来最大化飞机的航程。

\subsection{3.6.3 超级B747}
\section{3.7 超级P51赛车}
波音747 CL CD cw计数计数.45 141.3 499 3.6.3 超级B747 超级B747 为了探索可实现性能的极限，B747机翼已被一个全新的机翼取代，以生产“超级B747”。 通过将从其他优化中获得的超临界机翼部分与上一节描述的平面研究中发现的最佳平面形式混合在一起，创建了初始设计。 然后，使用RANS优化代码SynlO7在4.8、0.85、0.86和0.87的4个设计点上获得最小阻力，如图所示。 3.6（a）-（d）裸露的机翼固定升力系数为0.45，当包括机身升降机时，相当于约0.52的升力系数。 由于新机翼部分明显更厚，新机翼估计比基线B747机翼轻700磅，如表3.1所示。 与此同时，从0.78马赫到.92的整个范围内，阻力都减少了，在0.87马赫时，最大收益为25次，如图所示。 3.7（a）。 图3.7（b）和表3.2显示了0.86马赫的升力阻力极点。 超级B747的阻力系数在升力系数为0.5时为135次计数，而基线B747在升力系数时具有相同的阻力略低于.44。 这代表了LID的改善超过11\%。 结合机翼重量的减少和由于机翼部分较厚而增加的燃料量，这应该会导致航程的大幅增加。 （107.0压力，34.3黏性）（96.0压力，39.2黏性）（82,550磅）（81,870磅）.50 135.2 495 3.7 超级P51赛车 在此应用中，自动设计方法重新设计了P51“Dago Red”的机翼，这是一架参加里诺空中竞赛的飞机。 由于激波的出现，飞机的速度达到500英里/小时以上，并遇到压缩阻力。 目标是在不改变机翼结构的情况下延迟阻力上升。 因此，形状修改仅限于在机翼表面增加一个凸起，只允许向外移动。 此外，变化仅限于部分和弦范围，如图所示。 3.8。 这个凸起产生的扰动沿着特征传播，并从声线反射回来，以削弱冲击。 改进如图所示。 3.9.

72 气动外形优化 cp=-2'o. -. -.-. （A）马赫.78 \_\_-- -。 \_ -.... I cp=-20 -\_.. -\_.- (b) 马赫.85 (c) 马赫.86 (d) 马赫.87 图3.6：(a)-(d)：超级B747分别为马赫.78、35、36、.87。 破折线表示初始配置的形状和压力分布。 实线表示重新设计的配置。

\subsection{3.7.1 Transonic公务机的外形优化}
A. 詹姆森阻力上升在nxed CL.45（a）机翼CD g E CD（munt-I（b）阻力极点图3.7：（a）：超级B747在固定C，.45时阻力上升，（b）：基线和超级B747的阻力极点在26马赫。 （实线代表超级B747。 破折线代表基线B747。） （O'O'（图3.8：为Transonic公务机添加了凸起，以实现无冲击机翼3.7.1 外形优化。非结构化设计方法也已应用于几个完整的飞机配置。 商务喷气式飞机的结果如图所示。 3.10（a）和（b）。 初始配置的外翼部分存在强烈的冲击，这基本上被重新设计所消除。 在大约8个设计周期内，阻力从235计数减少到215计数。 通过扰动攻角，升力被限制在0.4。 此外，在设计过程中保持了机翼的原始厚度，以确保重新设计的形状将保持燃料量和结构完整性。 机翼的厚度限制施加在沿机翼跨度的切割平面上，并通过将受限的形状运动传递回表面三角测量的节点。

74 气动外形优化（a）重新设计之前（b）重新设计后图3.9：通过在原始P51机翼上增加凸起而获得的改进。 （a）C，原始机翼在0.78马赫的分布，（b）C，重新设计的机翼的分布。 颠簸消除了冲击，（c）Co vs. 固定CL 0.1的马赫数，（d）LID与 固定CL 0.1的马赫。

\section{3.8 结论}
A. 詹姆森75表3.2：拖曳极波音747 0.0045 0.0500 0.1000 0.1501 0.2002 0.2503 0.3005 0.35007 0.4009 0.4512 0.5014 0.5516 94.3970 82.2739 74.6195 72.1087 73.9661 79.6424 88.7551 101.5293 118.0487 141.2927 177.0959 228.1786 B747与 超级B747在0.86 超级B747 0.0005 0.0504 0.1505 0.2005 0.2506 0.3006 0.3507 0.4008 0.4508 0.5009 0.7509 74.5328 66.5664 64.1332 65.3071 69.0516 75.2481 83.4291 93.2445 103.2 186 116.3848 135.2461 163.8937 0.6016 298.0458 I 0.6011 204.3010（计数）注意CL.45的基线B747和CL.5的超级B747的等阻力。 3.8 结论 二维和三维设计研究的一个重要结论是，降低抗震强度或产生无冲击Aow所需的翼部分不需要像熟悉的平顶干旱后部装载的超临界剖面。 今天几乎任何飞行的飞机，如波音747或麦克唐纳-道格拉斯MD 11，都可以调整，以在选定的设计点产生无冲击的流量。 过去十年的积累经验表明，大多数以超音速巡航的现有飞机都可以减少3\%至5\%的阻力，或者马赫数的阻力上升至少增加0.02。 这些改进可以通过非常小的形状修改来实现，这些修改太微妙了，不允许通过试错方法来确定它们。 当允许进行更大规模的修改，如平面变化或新的机翼部分时，可以在5-10\%的范围内获得更大的收益。 考虑到整个航空公司的燃料成本，潜在的经济效益是巨大的。 此外，如果在设计过程中充分利用升力与阻力比的增加，可以设计出更小的飞机来执行相同的任务，从而进一步降低成本。 这种类型的方法肯定会为未来的空气动力学设计提供基础。

\section{3.9 确认}
\section{3.10 参考书目}
76 气动外形优化（a）基线（b）重新设计图3.10：M = 0.8，a = 2" 3.9的公务机的密度轮廓 确认 本文描述的研究从空军科学研究办公室在赠款下的继续支持中受益匪浅。 AF F49620-98-1-2004。自1990年以来。 3.10 参考书目[1] 安德森，W. K.，纽曼，J. C.，惠特菲尔德，D. L.，和尼尔森，E. J。 使用复杂变量对非结构化网格的Navier-Stokes 方程的敏感性分析。 AIAA论文99-3294，弗吉尼亚州诺福克，1999年6月。 [2] Barth，T. J。 欧拉和纳维尔·斯托克斯方程的非结构化网格和有限体积求解器的方面。 AIAA论文91-0237，AIAA航空航天科学会议，内华达州里诺，1991年1月。 [3] Bischof, C., Carle, A., Corliss, G., Griewank, A., \& Hovland, P. 从Fortran程序生成衍生代码。 内部报告MCS-P263-0991，计算机科学部，阿贡国家实验室和平行计算研究中心，莱斯大学，1991年。 [4] 戈多努诺夫，S. K。 流体动力学方程不连续解的数值计算的差分法。 垫子。 Sbornik，47：271-306，1959年。 由美国翻译为JPRS 7225。 商务部，1960年。 [5] 格林，L. L.，纽曼，P. A.，和海格勒，K. J。 通过自动微分，高级CFD算法和黏性建模参数的灵敏度导数。 AIAA论文93-3321，第11届AIAA 计算流体动力学会议，佛罗里达州奥兰多，1993年。

A. 詹姆森 77 [6] 哈贝克，M. 探索无冲击跨声速翼型设计的局限性。 航空和航天问题报告，斯坦福大学，斯坦福，加利福尼亚州，2004年7月。 [7] Hicks，R. M。 \& Henne，P. A. 通过数值优化设计机翼。 飞机杂志，15407-412，1978年。 [8] 希克斯，R. M。 \& Henne，P. A. 通过数值优化设计机翼。 AIAA论文79-0080，1979年。 [9]詹姆森，A。 通过控制理论进行空气动力学设计。 科学计算杂志，3:233-260，1988年。[lo] Jameson，A. 飞机设计的计算空气动力学。 科学，245\textasciitilde{}361-371，1989年7月。[ll]詹姆森，A。 使用CFD和控制的最佳空气动力学设计。 AIAA论文95-1729，AIAA第12届计算流体动力学会议，加利福尼亚州圣地亚哥，1995年6月。 [12]詹姆森，A。 使用控制理论进行最佳空气动力学设计。 玉米输液流体动力学评论，第495-528页，1995年。 [13]詹姆森，A。 气动外形优化使用伴随法。 2003年2月3日至7日，比利时布鲁塞尔冯·卡曼流体力学研究所冯·卡曼研究所的2002-2003年系列讲座。 [14]詹姆森，A。 空气动力学。 在Erwing Stein、Rene de Borst和Thomas J.R. 休斯，编辑，《计算力学百科全书》，第11章。 John Wiley \& Son，2004年。 1SBN：O-470-84699-2。 [15]詹姆森，A。 和贝克，T。 J。 飞机尤拉方法的改进。 AIAA论文87-0452，AIAA第25届航空航天科学会议，内华达州里诺，1987年1月。 [16] 詹姆森，A. 和Caughey，D。 A. 解决稳定可压缩流动的Euler 方程需要多少个步骤：寻找快速解决方案算法。 AIAA论文2001-2673，第15届AIAA 计算流体动力学会议，加利福尼亚州阿纳海姆，2001年6月11日至14日。 [17]詹姆森，A。 \& 马丁内利，L. 基于控制理论的气动外形优化技术。 数学讲义\#1739，由工业问题驱动的计算数学进展，CIME（国际数学之夏（中心），意大利Martina Franca，1999年6月21日至27日。

78 气动外形优化 [18] 詹姆森，A.，马丁内利，L.，阿隆索，J. J.，瓦斯伯格，J. C. \& Reuther，J. 基于模拟的空气动力学设计。 IEEE航空航天会议，密苏里州大天空，2000年3月。 [19]詹姆森，A.，马丁内利，L.，和皮尔斯，N. A. 最佳空气动力学设计理论。 计算机。 流体力学，[20] Jameson, A., Martinelli, L., \& Vassberg, J. 连续极限中伴随分级公式的减少。 AIAA论文，第41届AIAA航空航天科学会议，内华达州里诺，2003年1月。使用Navier-Stokes 方程。 10:2 13-237，1998年。[所有Jameson, A., Schmidt, W., \& Turkel, E. 使用Runge-Kutta 时间推进方案的有限体积方法求欧拉方程的数值解。 AIAA论文81-1259，1981年1月。 [22] Jameson, A., Sriram, \& Martinelli, L. 一种用于跨声速流动的非结构化伴伴方法。 AIAA论文，第16届AIAA CFD会议，佛罗里达州奥兰多，2003年6月。 [23]詹姆森，A。 \& 瓦斯伯格，J. 用于动态设计的计算流体动力学：其当前和未来的影响。 AIAA论文01-0538，AIAA第39次航空航天科学会议，2001年1月，内华达州里诺。 [24]金，S.，阿隆索，J. J.，和詹姆森，A. 设计优化使用黏性连续伴联方法的高升力配置。 AIAA论文2002-0844，AIAA第40届航空航天科学会议和展览，内华达州里诺，2002年1月。 [25] Leoviriyakit，K. \& 詹姆森，A. 翼的气动外形优化 AIAA纸2003-0210，第41届航空航天，包括平面变化。 科学会议和展览，内华达州里诺，2003年1月。 [26] Leoviriyakit，K. \& 詹姆森，A. 空气结构机翼计划优化-内华达州里诺，2004年1月。 第42届航空航天会议记录。 科学会议和展览。 [27] Leoviriyakit, K., Kim, S., \& Jameson, A. 翅膀的黏性气动外形优化，包括平面变量。 A IAA论文2003-3498，第21届应用空气动力学会议，佛罗里达州奥兰多，2003年6月23日至26日。 [28]狮子，J。 L。 由偏微分方程支配的系统的最佳控制。 Springer-Verlag，纽约，1971年。 由S.K.翻译 米特。 [29] 马丁内利，L. \& 詹姆森，A. Reynolds平均方程的多重网格法验证。 AIAA论文88-0414，1988年。

A. 詹姆森 79 [30] 纳达拉贾，S。 K。 气动外形优化的离散伴随方法。 博士论文，斯坦福大学航空航天系，斯坦福，加利福尼亚州，2003年1月。 [31] 皮罗诺，0。 椭圆系统的最佳形状设计。 Springer-Verlag，纽约，1984年。 [32] Reuther，J. 超音速飞机的气动外形优化。 IAA论文2002-2838，第32届AIAA流体动力学会议，St. 密苏里州路易斯，2002年6月。 [33] Reuther，J.，Jameson，A.，Alonso，J. J.，Rimlinger，M. J.，和Saunders，D. 使用伴随公式和并行计算机约束多点气动外形优化。 Aircrajl杂志，36，1999年。 也在AIAA 97-0103，第35届航空航天科学会议和展览中，内华达州里诺，1997年。 [34] Rizzi，A. 和Viviand，H。 用激波计算隐性跨声速流动的数值方法。 在Proc中。 GAMM研讨会，斯德哥尔摩，1979年。 [35] Tatsumi, S., Martinelli, L., \& Jameson, A. 一种新的高分辨率方案，用于带冲击的可压缩黏性流动。 AIAA论文出庭，AIAA第33届航空航天科学会议，内华达州里诺，1995年1月。 [36] Thwaites，Bryan，编辑。 不可压缩的空气动力学。 牛津大学出版社，1960年。 由S.K.翻译 米特。 [37]瓦斯伯格，J。 C. \& 詹姆森，A. 里诺赛车飞机的气动外形优化。 《国际车辆设计杂志》，28:318-338，2002年。

此页面故意留空

\clearpage
\chapter{第4章 螺旋桨叶片的设计优化 Luigi Martinellil和James。 德雷尔}
\section{4.1 介绍}
第4章 螺旋桨叶片Luigi Martinellil和James 3的设计选择。 Dreyer2 4.1 介绍不可压缩黏性流动的形状设计有几种应用，包括飞机机翼、潜艇附属物、船体形式、游艇帆和海洋螺旋桨的优化，仅举几例。 传统上，选择设计变体的过程是通过反复试验进行的，依靠设计师的直觉和经验。 计算流体动力学（CFD）和计算机硬件的进步使今天能够将数值模拟与自动搜索和优化程序相结合，这有可能提高大量空气动力学设备的性能。 一种越来越流行的方法是通过遗传算法对大量变异进行搜索。 这可能允许在非常复杂的多目标问题中发现（有时是意外的）最佳设计选择，但当代价函数的每次评估都需要密集计算时，就会变得非常昂贵，就像空气动力学问题一样。 由于这种限制，我们选择了一种基于控制理论的替代方法，该理论在八十年代末由皮罗诺[7]和詹姆森[3]独立开创。 在这种方法中，梯度（灵敏度）的评估只需要解决普林斯顿大学机械和航空航天工程系的伴手系统，普林斯顿，计算力学系，ARL，宾夕法尼亚州立大学，州立大学，新泽西州08544。 宾夕法尼亚州16804

\section{4.2 制定作为控制问题}
82 螺旋桨叶片 设计方程即使在无限多的设计变量的极限情况下。 在九十年代，主要由于可用的计算资源增加和CFD演算法的进步，基于辅助的设计方法已成功应用于几个二维和三维问题。 特别是，Martinelli和Cowles [5]开发了一种基于人工可压缩性方法的伴随方法，并将其应用于湍流 黏性流动中三维机翼和帆的逆向设计。 Dreyer和Martinelli首先将这种方法扩展到三维定子和转子的逆设计[6]，最近他们将这种方法扩展到螺旋桨叶片部分的设计优化。 在本文中，我们总结了我们设计方法的关键特征，并报告了使用我们方法设计的螺旋桨叶片特性的实验验证结果。 实验工作是在宾夕法尼亚州立大学应用研究实验室（ARL）的水隧道中进行的[a] 4.2 制定为控制问题 任何与气动外形优化相关的适当代价函数都可以写成I = I（W，F），其中w是流场变量，3是边界的位置。 例如，对于逆叶片设计，I是叶片表面集成的设计压力和目标压力之间的差值。 边界F的变化导致代价函数的变化。 流场的控制方程 R(w,F) = 0。 （4.2）表达了w和的依赖性。 F被引入为流场域V内的约束。 然后从方程（4.3）3中确定bw，下标I和11用于区分流解中变化6w的贡献与与修改6直接相关的变化。 形状为F

\subsection{4.2.1 螺旋桨叶片的成本函数}
L。 Martinelli和J。 J。 Dreyer 83 接下来，我们使用拉格朗日乘法4用流动方程的变化来增加代价函数。 选择+来满足伴随方程第一个项被消除了。我们发现61 = 663，（4.4）其中g=--*T aiT一旦方程（4.6）建立，可以通过形状变化6T = -xg进行改进，其中A是正的，并且足够小，第一个变化是61的准确估计。 然后代价函数的变化变为61 = -xgT6 < 0。 对梯度最终形式的检查揭示了基于控制理论的方法的强度。 由于Sw没有出现在方程中，因此评估梯度的成本相对于流动解来说是微不足道的。 设计周期中唯一的额外费用是相伴系统的解决方案，它所需的工作量与流程解决方案大致相同。 通过基于有限差分的优化，对每个设计变量进行流场评估。 对于由大量设计变量（如螺旋桨叶片或帆船龙骨）组成的几何形状，控制理论方法比有限差分方法的成本优势令人震惊。 4.2.1 螺旋桨叶片的成本函数 对于实际应用，必须使用多目标函数进行优化，以延迟包括设计条件在内的一系列操作点上的空化开始。 因此，一个多点代价函数

\subsection{4.2.2 搜索程序}
84 螺旋桨叶片设计可以通过单个设计点的成本函数的简单线性组合来形成。 Nower n=1 为了尽量减少空化，我们选择了以下代价函数；其中和While 代价函数基于轴向力的最小化，选择优化效率。 这允许实现改善设计外空化的目标，同时最大限度地减少设计性能的损失。 4.2.2 搜索程序 通常，对本地最小值的搜索从基线设计开始，并通过搜索方法遍历设计空间。 搜索的效率取决于找到本地最小值所需的步骤数量以及每个步骤的成本。 为了加快搜索速度，人们可能会诉诸于使用牛顿方法。 在这里，搜索方向基于梯度消失表示的方程，G(F) = 0，并通过非线性方程的标准牛顿迭代求解。 如果黑森法可以准确且廉价地评估，牛顿法通常非常有效。 不幸的是，气动外形优化的情况并非如此。 准牛顿方法从连续步骤中记录的6个变化中估计黑森数或其逆数。 对于N设计的离散问题

L。 Martinelli和J。 J。 Dreyer 85个变量，需要N个步骤才能获得黑森的完整估计值，这些方法的属性是它们可以在正好N个步骤中找到二次形式的最小值。 因此，一般来说，准牛顿搜索的成本随着设计空间的维度而扩展。 由于高效的空气动力学形状主要是平滑的，因此可以对搜索方法采取合理的替代方法。 为了确保优化序列中的每个新形状保持平滑，可以平滑梯度，并在下降过程中用其平滑值替换G。 这也起到了预制剂的作用，允许使用更大的步骤。 这个程序可以正式化如下。 定义一个修改后的索博列夫内积（u，w）= l（uv + EVU。 Vw)dR，那么（u，v）=（u，u）+（EVU，Vw），其中（u，w）是Lz中的标准内积。 按零件积分产量使用内积符号，代价函数 I的变化可以表示为，因此我们可以隐式求解G B - V（EVG）= 6。 然后，如果设置6F = -XG，dI = -A@，G）= -A（G，G），如果X足够小和正，则保证改善，除非过程已经达到= 0的静止点（因此6=0）。 此外，保留了边界的原始平滑度。 事实证明，这种方法对使用梯度的近似值非常宽容，因此在改变形状之前，流动解和伴随解都不需要完全收敛。 这大大节省了完整优化过程的计算成本。 此外，观察到平滑下降法收敛到固定的步数，与设计变量无关。

\section{4.3 在离子上实施}
\section{4.4 离子下低Cav-it的叶片部分的优化}
86 螺旋桨叶片 设计 4.3 在离子实施 一个设计周期所需的步骤如下：1. 求解u1、uz、u3、p的流动方程。 2. 解决伴随方程。 II，受制于适当的边界条件。 3. 评估Q 4。 将Q投影到一个满足任何几何约束的允许子空间中。 5. 根据最陡峭的下降方向更新形状。 6. 返回1，直到达到收敛。 黏性设计方法的实际实现在很大程度上依赖于状态（w）和共状态（11）系统的快速和准确的求解器。 在这项工作中，我们使用了一个Flo103版本，最初由Jame-son和Martinelli [4]为可压缩流动编码，该版本经过修改，允许使用Chorin的人工可压缩性方法解决RANS及其伴随系统的不可压缩形式[l，81。 该代码基于完整的多网格时间推进方案，该方案将快速收敛产生到稳定状态。 此外，叶片部分的优化是使用级联配置进行的。 4.4 离子低空穴叶片截面的优化 这项工作的主要动机是从基线NACA 65410开始定义一个优化的叶片截面，并展示水隧道中一系列入射角的预测性能改进。 然而，在进行优化之前，必须解决一个微妙的问题。 在固定的隧道速度下，在测试部分的垂直壁范围内改变水翼的入射角度并不等同于改变水翼级联上入射流动的角度--这种情况与推进器设计师更相关。 然而，在这种情况下，人们可以表明，水隧道中的0.403英寸f5英寸发病率变化大约相当于非交错的20英寸间距级联中的2英尺418.5英寸流动角度变化[a]。 这是一个重要的结果，因为它表明，我们可以根据推进器设计所需的不同流动条件优化级联配置的水翼，并在具有不同叶片方向的固定流动条件下验证水隧道的性能。

L。 Martinella和J。 J。 Dreyer 87 基于这些发现，设计优化继续进行，从NACA 65410水翼的非交错级联开始，叶片间距为20英寸。 这种俯仰的选择不是任意的：这个尺寸与ARL直径12英寸的水隧道的矩形测试部分的宽度一致，这种选择允许我们为设计优化以及隧道内性能的后期设计分析使用类似的计算网格。 外形优化的意图是提高基线NACA 65410水翼的稳健性，主要是在非设计流量入射条件下的空化开始，同时保持设计点的部分提升。 在非入射条件下，二维水翼将在前缘吸峰处开始空化：流量入射越大，起始指数越大。 因此，增加水翼的空化起始稳健性等于在更大的流量入射角度下延迟起始，或者，另一种方式是扩大其空化桶。 使用的外形优化方法是一种多点、多目标的方法。 生成了两个替代的优化设计。 两者都在+/- 5O入射角下最小化了基于吸力峰值的代价函数，并在设计入射点最大化了基于效率的代价函数。 这可以实现改善设计外空化的目标，同时最大限度地减少设计性能的损失。 这两种情况之间的唯一区别是复合成本函数梯度构建中成本函数的相对权重。 在第一个优化案例中，案例1，设计点效率代价函数的权重为5，相对于非设计空化代价函数的权重为1。 图4.l（a）是本案例通过50个设计迭代的关键参数演变的总结。 这种权重选择的效果主要是限制刀片厚度的增长，当试图减弱非设计前缘吸峰的形成时，这是不可避免的。 这种特殊的权重选择产生了一个优化的截面，该截面在设计入射率下增加了4\%的阻力，并且在分别为0.26和0.61的高入射和低入射条件下改善了空化初始。 在这种情况下，与基线水翼相比，最终的水翼形形状显示出大致不变的最大厚度--尽管有些明显改变的弧度线。 为了保持设计点提升，最终部分的交错角度也有所增加。 在第二个优化案例中，案例2，复合材料代价函数的效率最大化部分被完全删除，其权重为0，相对于非设计空化成本函数的权重为1。 这种加权的效果是允许水翼厚度增长，而不考虑其对设计点阻力的影响--然后设计点流仅用于确定保持所需升力所需的交错角度。 图4.l(b)（与图4.l(a)相同的比例绘制）总结了演变

88 螺旋桨叶片 设计 weignts=5.,1.,1. 0 0061 1.4 0 OOSI 0 11 L: 10 0 W51 0.6 0.7 05' NDES 重量=O., l.,l. 0 0060 1.3 OW58 IJ 1.1 G 10 0 0056 0.7 O -o NDES ((a)) 案例 1 ((b)) 案例 2 图 4.1：此案例关键性能参数的设计演变。 很明显，受此 外形优化 影响的变化明显大于案例 1 中观察到的变化。 具体来说，不受此权重选择的制定，设计点截面阻力比设计演变增加了近12\%（而案例1中只有4\%）。 虽然这在设计点上是一个不利的影响，但它仍然创造了条件，在这种条件下，非设计空化性能可以得到更大的改进。 高发病率条件显示预测的起步空化指数改善为0.69（与病例l的0.26相比），而低发病率的改善为0.75（与0.61相比）。 这个新刀片的厚度约为13\%。 有趣的是，最佳截面的倾角和交错角与NACA 65410基本无法区分，厚度的增加似乎是一种全局现象。 后一种观察特别值得注意，因为优化器关注的流动现象非常局部于前缘区域，但前缘半径在优化过程中基本上保持不变。 与NACA 65410相比，第二次优化在空化性能方面产生了更大的三角洲。 因此，我们的第二个设计预计在水隧道环境中的实际性能会更容易观察到，最终被选中进行测试。 进行了最后一系列CFD预测，以预测水隧道20英寸宽测试部分范围内选定部分的性能差异。 在实验配置中，水翼围绕其中弦点相对于隧道的静止垂直壁旋转。 图4.2（a）和4.2（c）显示了在两种入射条件（0英寸和-7英寸）下优化的水翼的计算域。 笔记

L。 Martinelli和J。 J。 Dreyer 89（（a））计算域：Oo。 （（B））静压轮廓：0'。 （（c））计算域：-7O。 （（d））静压轮廓：-7O。 图4.2：优化水翼（案例2），没有尝试解决隧道墙边界层；它们被视为简单的无通量表面。 图4.2（b）和4.2（d）显示了图4.2（a）和4.2（c）所示网格的相应静压场。 应该指出的是，计算流入平面的位置大致对应于物理测试部分的入口，所有压力都参考该平面上的中间通道点。 为NACA 65410和跨越f9'的入射角度的优化水翼生成了一系列网格。 所有模拟都是在2,200,000的和弦雷诺数下运行的。 在大约f8O的入射角之外，流动求解器开始显示出伪不稳定的迹象，在某些情况下，完全发散。 推测这表明了水翼部分出现的开始停滞现象。 图4.3（a）显示了NACA 65410和优化水翼的入射角度范围内的升力和阻力变化。 在f8O范围内，升力曲线几乎是线性的，两个水翼的升力曲线基本上在彼此的地块上。 这是对设计点升力施加约束的直接后果（也就是说，两个水翼的设计升力保持不变）。 两个部分的阻力曲线都显示了特征性空化桶，两个水翼的最小阻力发生在大约-lo的入射点。 请注意，在0.4O设计点附近，优化的水翼的阻力比基线水翼高出约12\%。 图4.3（b）显示了NACA 65410和最佳化水翼的空化结果。 这个图中的几个特征值得注意。 在空化铲斗相对平坦的底部附近，优化过程实际上

\subsection{4.4.1 与水隧道测量的比较}
90 螺旋桨叶片 设计基线 NACA 65410 6 优化截面性能 wh 入率 Varlatioon 在 12in WT 矩形测试 Senion Bassline NACA 65410 6 优化 Senbn 性能与 bctdence Varition 在 12in WT 矩形测试截面 a ((a)) 提升和阻力 a ((b)) 空化参数 图 4.3：NACA 65410 和最佳化水翼的计算性能的比较使 NACA 65410 的空化性能略有下降。 这是由于这些操作点的最小压力不是发生在前缘吸峰，而是靠近吸侧的中弦，在流动的更良性变化区域。 此外，将较厚的水翼放在相同的流动条件下自然会导致较低的吸面压力。 更令人感兴趣的是水桶相对陡峭的侧面。 这些是随着水翼入射角远离设计点而形成前缘吸峰的直接结果。 正是在这里，外形优化对性能的影响最大。 在15英寸入射区域，最佳水翼在给定的入射设计角度下，在la的顺序上产生了开始空化指数的改善。或者从另一个角度来看，在给定的空化数下，与NACA 65410相比，最佳水翼为无空化操作提供了约2O的入射范围。 4.4.1 图4.4（a）显示了基线NACA 65410和一系列入射角度优化截面的截面升力变化。 线条是2D RANS预测值，符号是测量值。 应该注意的是，对于基线NACA 65410部分，升力系数是通过测量的静压分布的数值积分推断的；对于优化的部分，升降机也从压力分布中整合，但除此之外，它是直接使用称重单元测量的。 在后一种情况下，很明显，推断和直接测量的升力之间的差异与水隧道测量比较

L。 Martinelli和J。 J。 Dreyer 91 Bareline NACA 65410 \& OpUmked Seclion Pellolmance：比较01 2D RANS与测量a（（a））升力。 Easeline NACA 65410 h Opumaed Section Performance Comparison of 20 RAIlS Wlh Measurements 0 11) - s I7 6 54 I2 10 I 23 1 I6 7 a P 0 a ((b)) 空化参数。 图4.4：由于模型硬件上对压力抽头放置的明智选择，测量和计算性能的比较很小。 在图4.4（a）中立即显而易见的是预测和观测行为的提升曲线的不同坡度。 在负入射角度下，预测和观察之间的一致性相当好；然而，随着入射率变得越来越正，它会下降。 标有“Baseline-3D corr.”的行试图通过使用校正来解释一些三维效果。 事实上，这种三维“校正”确实减少了升力曲线的坡度，并优先影响了更大的入点角度；然而，冲击的大小不足以解释所有差异。 尽管升力曲线坡度之间存在差异，但在图4.4（a）中，也许更重要的是预测和测试硬件的基线和优化叶片部分之间的性能相似性。 当然，在入射角变化上保持横向力变化对于实用的叶片截面外形优化工具至关重要，图4.4（a）表明已经实现了这一点。 图4.4（b）总结了基线NACA 65410和优化叶片部分在入射角度范围内的空化性能。 线条显示二维预测空化性能，符号显示测量性能。 图4.4（b）表示了三个不同的区域。 负入射区（a）以压力侧前缘表面空化为特征，正入射区（c）以吸侧前缘表面空化为特征，设计点（b）周围的小区域以吸面行气泡空化为特征。 这些是非常典型的空化

92 螺旋桨叶片设计桶，即设计点周围相对平坦的底部，由相对陡峭的侧面支撑，以获得更大的正负入射角。 平底是前缘吸峰没有比叶片吸面上最小压力更深的入射范围；阶梯侧是由于前缘吸峰一旦建立就占主导地位。 2D RANS预测与（a）和（b）区域的水隧道观测之间的一致性相当好。 即使是优化叶片部分的幷斗底部空穴性能的微妙下降，也会在数字和空穴呼叫中捕捉到。 此外，对于负入射条件，计算和实验之间的一致性相当好，发现两条曲线的偏移量大约为1'。 然而，在（c）区域，结果不那么明显。 对于（c）区域较大的正入射条件，预测再次显示空化曲线以大约一度的固定角度偏移。 这似乎与在这个区域的实验中看到的偏移量差不多，尽管这里的观测似乎有更多的散点。 然而，显而易见的是，两个叶片部分的预测和观测到的气化行为之间的垂直偏移。 换句话说，在给定的入射角下，2D RANS预测的空化指标始终大于在同一入射角下观察到的空化；它们之间的差异似乎在0.5之间。 这可能与4.4（a）中显示的大正入射角的升力曲线的发散增加有关。 事实上，如果预测在给定的入射点下叶片负载更高，人们会期待更深的吸力峰值。 尽管如此，空化桶从基线NACA 65410到其优化形式的整体扩大，这是形状优化练习的目标，并由2D RANS隧道内分析预测，在非常令人鼓舞的程度上，通过图4.4（b）中总结的观察得到了验证。 图4.5（a）-4.5（f）分别显示了基线NACA 65410和优化叶片部分的几个具有代表性的弦静压分布。 每个图都比较了负设计外入射条件、近设计点条件和正设计外入射条件的预测和测量的表面压力。 图4.5（a）和4.5（b）都显示了2D RANS预测与基线和优化叶片的测量静态压力分布之间的良好一致性，每个叶片的入射率约为-5英寸。 对于这些情况，即使是前缘吸峰也能很好地捕捉到。 由此，人们期望2D RANS可以准确地预测该区域的负载和空化指数。 对-5英寸左右区域的4.4（a）和4.4（b）的检查确实表明了这种情况。 图4.5（c）和4.5（d）显示了预测和mea-之间的比较

L。 Martinelli和J。 J。 Dreyer 93基线NACA 65410，隧道内：具有优化测量的2D RANS的COmWriwn，Ih隧道：2D RANS与测量x（（a））测量和计算（（b））测量和计算的比较压力系数：NACA 65410。 压力系数：已优化。 Basellne NACA65410，隧道内：2D RANS wtm测量的Comlladson Optlmbed，mnnel内：Cornpadson 01 2D RANS与测量（（c））测量和计算（（d））测量和计算压力系数：NACA 65410。 压力系数：已优化。 基线NACA65410，In-Nnnel：Comlladwn 01 2D RANS wim测量a Opumked，In-tunnel：Comllallson of 2D RANS wim Measurements x x（（e））测量和计算（（f））[测量和计算压力系数：NACA 65410。 压力系数：已优化。 图4.5：测量和计算压力分布的比较

\section{4.5 结论}
94 螺旋桨叶片 设计确保靠近基线和优化刀片的设计点。 在这里，我们处于每个叶片的空化桶的底部，因此没有吸峰。 最小压力发生在中弦、吸侧附近，每个情况都非常准确地预测。 然而，对于每种情况，计算出的压力侧表面压力都有些预测过高，导致叶片负载明显增加超过测量值。 基于这些观察，人们期望2D RANS可以准确地预测设计点附近的空化，但过度预测了那里的叶片加载。 图4.4（a）和4.4（b）再次表明，这正是设计点附近的情况。 最后，图4.5（e）和4.5（f）显示了基线和优化刀片部分的2D RANS预测和测量之间的比较，每个截面的入射率约为+5O。 在这里，很明显，吸力峰值和叶片载荷都没有捕捉到其他两个入射角度的最佳元素。 具体来说，在这种情况下，前缘吸峰的深度似乎被两个叶片的2D RANS过度预测。 对于每个刀片，刀片的负载似乎也有些过度预测。 对+5'入射率附近的4.4（a）的检查表明，2D RANS确实预测了比在水隧道中观察到的更大的升力。 此外，在同一区域，图4.4（b）显示了比隧道中测量的更大的预测空化指数。 4.5 结论：基于控制理论的不可压缩湍流的设计方法已经进一步开发并应用于螺旋桨叶片截面的设计。 最近在宾夕法尼亚州立大学应用研究实验室（ARL）进行的针对气化起始优化的叶片截面的实验研究]进一步验证了我们的方法。 特别是，我们已经令人信服地证明，我们的外形优化方法可以成为为设计叶片部分构建实用的多操作点优化程序的可行基础]，该方法在一系列入射角度上保持流体动力学横向力，同时表现出大大改进的空化桶。 然而，实验测量确实表明，对于大正入射角来说，未完全解释的三维效应可能很重要。 我们方法的三维版本已经在并行计算机上开发和实现。 我们的快速多网格流动求解器的结合，以及我们的控制理论方法提供的计算梯度的包含成本]使全三维螺旋桨的优化完全可行。

\section{4.6 参考书目}
L。 Martinelli和J.J. Dreyer 95 4.6 参考书目[l] Chorin，A。 （1967）。 解决不可压缩黏性流动问题的数值方法，计算物理学杂志，2，pp. 12-26。 [2] Dreyer，J. J。 （2005）。 非设计性能的叶片塑造：二维级联中的空穴和效率。 第四届海洋流体力学国际会议记录，英国南安普敦。 [3]詹姆森，A。 （1988）。 通过控制理论进行空气动力学设计。 Scaentafic计算杂志，3，pp. 233-260。 [4] 马丁内利，L. （1987）。 用多重网格法计算黏性流量，博士 论文，MAE 1754-T，普林斯顿大学。 [5] 马丁内利，L. \& Cowles，G。 W。 （2003）。 基于控制理论的形状为不可压缩的Navier Stokes方程Int进行。 《计算流体动力学》杂志17（6）页。 1415-1432。 [6] 马丁内利，L. \& 烘干机，J. J。 （2001）。 使用连续伴随方法的推进器配置的流体力学外形优化。 第15届AIAA CFD会议记录，加利福尼亚州阿纳海姆。 [7] 皮罗诺，0。 （1984）。 椭圆系统的最佳形状设计，Springer-Verlag，纽约。 [8] Rizzi，A. 和埃里克森，L。 （1985）。 旋转隐形不可压缩流动的计算，流体力学杂志，153，pp. 275-312。

此页面故意留空

\clearpage
\chapter{第5章 用于优化、自适应和非稳定配置的4D流式边界条件建模Helmut Sobieczky}
\section{5.1 介绍}
第5章 用于优化、自适应和非稳定配置的4D流式边界条件建模 Helmut Sobieczky 5.1 简介 用于分析机械产品组件的各种负载的边界条件建模已成为许多学科的成熟技术。 结构、空气/水力和热力学荷载及其多学科耦合尚未通过标准程序解决，但许多商业软件包已经可用于解决一些实际问题。 在这种情况下，学术机构改变了课程，以便将教学材料的负荷适应新的技术和技巧。 学生面临挑战，要了解应用案例研究，并链接到针对特定问题的现成软件。 我们观察到，许多基础学科，如数学和力学，显然需要剥离一些更基本的材料，以利于此类特殊问题领域和应用课程工作。 这种发展在分析建模使用DLR德国航空航天中心的学科中尤其如此，Bunsenstr。 10, D-37073 哥廷根, 德国

\section{5.2 四维问题的几何概念}
98 Flow 边界条件 40中的建模对于理解机械现象至关重要：数学函数被用来解释一些复杂效应的特殊性，这些效应在产品操作中造成实际问题。 构建知识库的一个主要工具是分析几何学的基本学科。 其数学背景的后续应用是现代商业计算机辅助设计（CAD）的基础。 发展得不那么好，但也导致了商业软件，是可压缩的流体动力学（CFD）的计算分析，黏性流动增加了特殊的大气或热力学条件等复杂性。 因此，花了很多精力将CAD和CFD用于分析方面，使用几何软件包来定义边界条件来模拟复杂的给定形状，如整个飞机和复杂的涡轮机械部件。 然而，在这种情况下，我们似乎失去了一些基本理解，这是利用新颖想法来影响流动物理学创造新产品组件的设计方面非常必要的。 富有创造力的教师现在挑战使用信息技术工具箱向学生介绍力学和流体力学方面的一些经典知识库。 基础学科的基于计算机的教科书是保持这些基础知识的一种选择，例如，请参阅Caughey和Liggett的开创性工作[11。 这里说明的一些工作侧重于上述耦合CAD-CFD的几何方面，试图考虑可压缩流动的气动力学知识库：发现微小的形状变化是导致大流量质量变化的原因，因此在高速流动中组件操作的实际效率水平发生了巨大变化。 作者在[6]中描述了几何学的作用，以及创建具有明确定义的几何形状定义的测试案例研究的重要性。 在这里，我们专注于改变这些形状的必要性，以便在优化程序中创建整个配置系列可供选择，或模拟适应不同操作条件的智能形状控制，或者最后创建像我们在自然界中观察到的不稳定的机械系统：以真实或虚拟时间为附加维度，我们试图创建4维形状。 5.2 4维问题的几何概念 3D形状定义的概念之前已经解释过，见[7,8]，因此我们专注于扩展到第四维。 该方法的基本特征摘要将表明，扩展到4D是直截了当的，并将使用与我们发现的相同技术，通过改变2D曲线定义的参数来创建3D曲面。 在这种以（流动-）现象为导向的方法中，创建初始配置几何形状，我们遵循了我们所说的气体动力学知识库：

H。 Sobieczky 99逆设计概念是根据可压缩流动运动的基本方程开发的，它教会了我们识别沿着表面prtions的坡度和曲率比表面坐标更重要的敏感区域，这表明通过定义支持数据以及坡度、曲率或奇点指数来有力地控制形状生成函数。 此外，保持对笛卡尔定义的明确遵守，并建立直接函数评估，x3 = Fct（x1，x2），从而避免插值、迭代和积分来获得单个编度数据。 这种哲学似乎适合一大类形状设计任务，用于需要自由形式定义和一组相对较少的参数的航空和流体力学配置。 凭借之前在高速设计方面获得的经验，例如在跨声速马赫数制度中，机翼等基本组件应该从可能很少但特定的参数中定义。 这已在PARSEC系列部分[S]中实现。 与此同时，几位作者利用这些机翼部分获得了令人印象深刻的空气动力学设计结果，因此我们最近开发了交互式、基于网络的软件，让学生熟悉所用参数的作用。 可以定义两种不同类型的机翼，图。 5.1显示了涡轮机械叶片和飞机机翼截面的交互生成数据的选项：PARSEC- 09 这些是涡轮机械叶片部分，有9个参数。 热流体中的涡轮机或压片提出了许多强大的多学科设计任务，Dennis等人提出了使用这些形状的第一个方法[2]。 图5.1：涡轮机械和飞机机翼的交互式2D“PARSEC”翼型定义。

40 PARSEC-I 8中的100 Flow 边界条件建模是飞机机翼的原始PARSEC（- 1 1）套装的精炼版本：具有发散后缘（DTE）的超临界机翼部分可以通过支持高效的Kutta条件来定义在降低的激波附近进行黏性相互作用控制和循环控制。 下一步是第3维翼型参数的适当变化，对于机翼是跨度，对于叶片是径向方向。 基本上，同样的想法将3D扩展到第4维：新维度t要么是时间，要么是类似时间的虚拟坐标，只是为了定义一整组形状变化可供选择。 仅用于机翼或3D特征数据的少数形状生成参数将在给定的间隔内进行修改。 为了获得变体，我们只需要沿着0 I t I 1扫描，t = 0产生最初的3D设计研究，t = 1产生一个新的、极其变形的形状。 要修改的参数的分布函数允许单个影响控制，从线性到斜坡状函数，就像用于支撑站之间的几何曲线部分一样。 如果t是实时的，那么斜坡函数由正弦平方所定制，我们已经拥有了谐波周期周期的一半。 当然，其他周期性函数允许对参数振荡进行单独控制。 0 80 0 70 0.60 048m oooi 0002 0003 0004 0005 * Paretofmnt - - PolarAirfoil f - Polar 翼型 2. j 0. W6 0.007 cdwave) 1.2 1.1 Euler b 00 00 001 c, 002 003 图 5.2：17\%厚机翼的遗传算法优化，马赫=0.7：无黏流动帕累托前和2个机翼的拖曳极，（a）。 黏性流动优化，Re=4OMill，（b）。

\section{5.3 优化}
\section{5.4 自适应配置}
H。 Sobiecchy 101 5.3 优化 通过一组可用的变体，分析一些机械性能参数目标函数，并使用优化策略，人们可以选择一个优越的形状进行应用。 在Klein [4]的工作中，空气动力学性能由升力比（L/D）来定义，他的案例研究范围从机翼到机翼。 图5.2a显示了无黏流动中大量最佳翼型的结果，其最大L/D，位于帕累托前部。 然后选择两个机翼，并分析它们在各种攻击角度的性能。 这产生了它们的阻力极点，这些极点在设计点处接触到帕累托前部。 图5.2b显示了隐形（仅波阻）和黏性（包括摩擦阻力）设计之间的比较。 流动被认为没有冲击。 5.4 自适应配置 通过优化软件选择最合适的配置，如果形状变化是由机械设备有效改变流动边界获得的，那么在硬件上也有对应物。 迄今为止，这实际上可以通过弹性或气动元件来实现，由伺服电机或压电设备控制，这些设备由微型计算机控制，而压电设备又经过预编程，并使用流量质量传感器，将表面压力与预定义的目标值进行比较-图5.3：通过两个独立参数控制卡纳德根位置和二面体来改变卡纳德配置

\section{5.5 不稳定 边界条件}
\section{5.6 生物流体机械应用}
102 Flow 边界条件 在几个表面位置的40个ue建模。 由于需要将自适应设备限制在小面积区域，我们选择通过在机翼表面的跨声速膨胀和再压缩区域添加颠簸来选择使用曲率变化：发现膨胀肩部碰撞（ESB）和重新压缩肩部碰撞（RSB）对跨声速流动状态下翼型或机翼的空气动力学效率有积极影响，请参阅Geissler和Koch [3]的工作。 变形飞机 观察到计算工具的不断增加，以及用于表面制造的新智能材料，我们确认了创建配置的重要性，这些配置可以改变其形状的部分，而不仅仅是非常本地的补丁：变形飞机的想法诞生了，并将采用“智能技术，可以实现机上配置变化，以获得最佳飞行特性”（美国宇航局）。 在这种想法成为现实之前，我们尝试用我们的4D工具模拟形状的变化，并应用CFD来研究由此产生的流动结构，以更好地了解预测的空气动力学性能。 飞机工业最近的一项研究，遇到了将canard概念应用于跨声速运输机的困难，促使作者创建了一个可变的canard形状及其与飞机机身的整合，（图。 5.3），用于为进一步研究创建一个测试案例。 5.5 非稳态边界条件 上述关于自适应翼装置的工作实际上是为稳定和非稳态形状进行的：研究凸起振荡，以影响跨声速翼型上动态失速的开始。 非定常流动边界出现在旋翼飞机上。 刚性转子叶片的设计需要在叶片的跨声波推进阶段和低速后退阶段之间做出妥协。 随着对恒定升力的要求，需要改变攻角，在撤退阶段可能会导致动态失速。 在Trenker [9]的工作中，最初选择的PARSEC-1 1 翼型首先针对推进相位超音速峰值马赫数进行了优化。 随后，机头下垂和后缘密封襟翼模型允许创建可变反射型翼型形状，以适应通过可变增加的攻角的低速流动。 其结果是完全消除了动态失速，这肯定会发生在低马赫数的刚性超音速翼型上。 直升机行业似乎正在采取这个概念，实际案例研究正在进行中。 5.6 生物流体力学应用优化、适应和不稳定的形状变化：所有这些都是

H。 Sobieczky 103在自然界中表演并不断发生，鸟类、昆虫和鱼类历代都优化了它们的形状，不断适应飞行或游泳条件，并需要不稳定的谐波运动，通过拍打翅膀或鳍来产生升力和推力。 因此，通过创建此类动物的参数化模型，应用CFD并了解现象和最佳形状，向自然学习是一个挑战。 图5.4：不稳定的配置示例：蜻蜓拍打翅膀 一个特别具有挑战性的例子是蜻蜓具有很高的机动性，这是由于其独立操作每个翅膀的能力。 因此，这种昆虫已经在风洞中进行了研究，并对拍打运动进行了摄影评估[5]。 我们了解到，在一个周期内，翅膀被不同的周期性功能和前翼和后翼之间的特征相位移所拍打、扫荡和扭曲。 对于我们的几何工具来说，这种不稳定的边界条件是一个具有挑战性的测试案例：从照片中模拟刚性昆虫形状（图。 5.4）是第一步。 如果选择适当的模型功能来创建类似于飞机串联机翼体配置的合适数据集，很少有支撑点就足够了。 随后，测量的拍打振荡[5]定义了一旦数据集输入非稳态CFD模拟，让我们的模型昆虫适当拍打翅膀所需的非稳态参数，并足够接近以观察合理的空气动力学特征。

\section{5.7 结论}
\section{5.8 参考书目}
I04 Flow 边界条件 40 5.7 结论 使用我们的几何预处理器定义各种空气动力学或流体动力学结构的定义变得更加具有挑战性，包括第四维的形状变化。 虽然3D表面网格目前为CAD处理提供了合适的输入，以为各种网格生成、CFD分析和设计以及模型制造软件生成接口数据，但移动或变形表面的扩展邀请为代码开发和现象研究创建新的非稳态测试用例。 5.8 参考书目[l] Caughey，D. A.，Liggett，J. A.：流体力学入门计算机教科书。 在：Caughey，D.，Hafez，M. （编辑），计算流体动力学前沿，世界科学（1 998），pp. 465-48 1 [2] Dennis, B. H.，埃戈罗夫，I. N., Sobieczky, H., Dulikravich, G.S., Yoshimura, S.：涡轮叶片中三维蛇形冷却通道的平行热弹性优化。 ASME GT2003-38180（2003）[3] Geissler，W.，Koch，S.：自适应翼型。 在：Sobieczky，H.，（编辑）：Transsonicum IV研讨会。 IUTAM研讨会Transsonicum IV会议记录，Kluwer（2003），pp. 303-3 10 [4] Klein, M.: Entwurfsaerodynamische Studien an Tragfliigelkonfigurationen im Hochgeschwindigkeitsbereich mit evolutionaren Algorithmen. 论文大学。 哥廷根，（2000）。http:/lwebdoc.sub.gwdg.de/diss/2000/klein/ index. htm [5] Saharon, D., Luttges, M. W.：蜻蜓非稳态空气动力学：机翼相关系在控制产生流量中的作用。 AIAA-89-0832（1989）[6] Sobieczky，H.：理论、应用和教育流体力学的几何学。 在：Caughey，D.，Hafez，M. （编辑），计算流体动力学的前沿，世界科学（1998），pp. 45-55 [7] Sobieczky，H.：CFD和应用空气动力学的几何生成器。 在：H。 Sobieczky（编辑），高速航空运输的新设计理念。 CISM课程和讲座卷。 366. 纽约维也纳：Springer（1997），181 Sobieczky，H.：参数化机翼和机翼。 在：K。 Fujii和G。 S。 Dulikravich（编辑）：数值流体力学笔记，卷。 68，Wies-baden：Vieweg（1 998），pp. 7 1-87 [9] Trenker，M.：具有动态跨声速流动控制的自适应翼翼的设计理念。 J。 飞机卷。 40，第4期，（2003），第734-740页，第137-158页

\clearpage
\part{二。 算法和准确性}
第一部分 算法和准确性

此页面故意留空

\clearpage
\chapter{第6章基于隐性残基的紧凑方案的稳定性和效率C。 科雷和A。 勒拉特}
\section{6.1 导言}
第6章基于隐性残基的紧凑方案的稳定性和效率C。 科雷'和A。 Leratl 6.1 介绍使用隐式方案将收敛速率增加到稳态，无论是物理时间还是双时间，都很普遍，但对最佳隐式处理的追求仍在进行中161。 的确，只有当隐式方案由于需要在每个时间步骤中解决线性系统而导致的每次迭代的潜在高成本，由于使用大时间步骤而无法抵消达到稳态所需的迭代次数减少时，隐式方案才是真正有效的。 在提高隐式方案效率时，可以使用并结合两种主要策略：一种旨在提高方案的内在效率，即减少达到稳定状态所需的迭代次数，而另一种侧重于降低隐式方法每次迭代的成本，有时以牺牲内在效率为代价。 例如，在程序中使用对角化程序来降低块线性系统解决方案的成本时，遵循后一种方法[la，31。 对于大规模计算，内存要求也是有效使用隐式方案的重要制约因素，除了降低每次迭代的成本外，块隐式阶段的对角化[14]或所谓的无矩阵方法[lo]等技术的另一个好处是它们的低内存-'SINUMEF Lab., ENSAM, Blvd de l'H6pita1, 75013 Paris, FRANCE。

\section{6.2 隐含计划描述}
108 隐含的基于剩余的紧凑计划要求。 正是如此，也是为了降低成本和内存存储，高阶上风方案通常与简单的一阶上风隐性阶段耦合[13，21，尽管这种选择引起了内在效率损失。 在目前的工作中，我们有兴趣研究最近为可压缩流动计算开发的基于残差的紧凑（RBC）方案的效率[7，91，并被证明为传统的上风方案提供了替代方案，该方案依靠大型定向模版达到高阶[8]。 我们计划证明，在不调整涉及冲击问题的情况下，准确性和稳健性不是基于残余方法的唯一优势，但它也比相同顺序的传统逆风方案提供了更好的效率。 在本研究中，我们考虑了三阶定向非压实上风显式阶段和基于残差的紧凑显式显式级，两者都与一阶上风隐式级相结合，如本文第6.2节所述。 在第6.3节中，对这些方案的放大系数进行了分析和数值研究，以便在精确解决时评估两种方案各自的内在效率。 然而，由于在实践中，与隐性阶段相关的线性系统是使用因数分解或松弛方法大致解决的，因此第6.4节给出了这些各种解法的统一公式，并提供了内在效率方面隐性阶段处理的第一个层次结构。 第6.5节对最有效处理的放大系数的行为进行了数值分析，用于常规和基于剩余的方案，并就这些方案各自的效率得出结论。 6.2 隐式方案描述 让我们考虑以下守恒定律df P的双曲系统，其中A--是空间方向-dw xp中物理通量fp的雅各布。 在研究在冯·诺伊曼类型分析框架内近似（6.1）的方案的效率时，我们将专注于系统dW d在p=l时的线性化形式，以及围绕（6.1）线性化的物理状态。 此外，在不更改符号的情况下，也可以假设最后一个方程可以简化为一个简单的对流方程，w是对流的标量和

C. Corre \&' A. Lerat 109 标量系数A，与每个空间方向相关的波速。 系统（6.1）由以下守恒格式 Awjn (S,h,)q -+-=o At 62近似，其中J = (jl,...,jd)表示与正法笛卡尔网格的点xj = (\textasciitilde{}16x1,...,jdSxd)相关的多整数，wy = w(xj, nAt)，在时间步，h，是近似物理通量f,的数值通量。 空间方向x,, S,,的一个网格单元格的差分算子是这样(S,W),+;,,, = w\textasciitilde{}+\textasciitilde{}, - wj，如果q \# p，则ep是一个多整数的分量epq等于0，如果q = p，则为1。 如果数值通量h，是与一阶上风方案相关的，它读作h，= ppfp - \$A，IS，w，其中p，表示空间方向x中一个网格单元的平均算子，: (p,\textasciitilde{})\textasciitilde{}+;\textasciitilde{}, = z(w\textasciitilde{}+\textasciitilde{}, + w,)。 在线性或线性化的情况下，方案（6.2）也可以表示为Aw，“= -K'w，”，其中K'是显式阶段运算符，由一阶上风方案K' = CApSppp - 51AplS;, p=l与A, = a,A,的以下表达式给出。 如果一阶上风方案（完全）隐式，则通过FL = Id + KI给出的隐式阶段运算符求解Awjn，其中Id是恒等式运算符。 众所周知，当使用大时间步骤时，由此产生的方案效率很高：首先，（6.4）具有内在效率，当在400时，放大系数可以非常快速地阻尼所有误差模式；其次，涉及每个空间方向的3个点的隐性阶段可以通过一系列对角线系统解决方案有效地解决实际应用（本研究中考虑的隐式处理将在第6.4节中描述）。 不幸的是，一阶精度对于实际目的来说是不够的，人们需要将显式阶段的精度顺序提高到二阶或三阶。 一个标准的方法是使用Van Leer的MUSCL方法，在这种情况下，对于线性问题，导致（6.1）的三阶上风离散化的数值通量为FLAW，" = -K'wjn（6.4）与Ic"' = c Ap(I - p= I)形式的关联显式阶段运算符

I10 基于隐性残差的紧凑方案 如果该方案完全隐式，它将导致一个涉及每个方向5个点的隐性阶段，因此其处理将需要一系列五对角线系统解决方案；这在实践中是不能接受的，因为五角线求解器与简单的三对角线求解器相关的CPU和内存存储的额外成本是不可接受的。 因此，即使显式阶段是更高阶的，也习惯上保留一阶隐式阶段。 因此，线性对流问题的隐式三阶上风方案的通常形式为（Id + KI）Awl= -K'IIwn 3（6.7）从现在开始，隐式方案（6.7）将表示为DNC（3），其中DNC代表定向非紧凑，（3）是指稳态解精度顺序；显式运算符Ic”也将表示为KDNC。 显式阶段（6.6）被称为定向，因为它专门使用xp方向上的网格点来近似同一方向的物理通量fp。 至于方案不紧凑，据说是这样，因为它利用每个空间方向的5个点来达到三阶，而基于残差的方案在3 x 3点模版上产生相同的精度顺序。 正好，现在让我们回顾一下后一种方案的特点--读者可以参考[7，91，在欧拉和Navier-Stokes 方程的二阶和三阶近似的情况下，对方案设计原则进行全面描述，并参考[5]，以获取高阶残差方案的一般推导。 基于残差的方案近似（6.1），其中点j，而Fp（p = 1，d）是r在中点j + iep的中心近似。 系数4p的设计是为了确保（6.8）的RHS确实产生一个耗散方案，正如[7]中详细解释的那样，必须满足\$q = aqpq5p与aqp = IA，I/IA，I的关系。 由于保持这些系数（bp接近单位也是可取的，以避免数值耗散过量，因此在之前的约束下最小化C:=，-11导致选择是稳定残差r = Cp=，d az，a.f，在\$p = min(l,min(l/aqp)) = min(I,min(-)）的中心近似。 IAPI 42P q\#p I A，I 自然，表达式（6.9）仅针对标量对流问题定义良好；然而，它们以直截了当的方式扩展到系统情况，假设矩阵系数4与雅各布矩阵A具有相同的特征向量，并将其特征值表示为（bg）= min（1，minp\#p（lA\textasciitilde{}'I/m（A，）））

C. Corre 63 A。 Lerat 111在哪里A！） 表示A的ith 特征值，而m(A,) = minj(1Af)I）。 现在返回残差的近似值，很容易检查以下选择是否在稳态下产生r“0的四阶非耗散近似值和三阶耗散，即3d点模版上（6.1）的全局三阶近似值。 在线性情况下，与（6.8）相关的显式阶段运算符，（6.10）读取d d KRBC = c [A, (1 + c; 6: 1 6PPP - 5 dpsgn(Ap) (APG + c A46PPP64P4 11. p= 1 q=l,qfp q=Lq\#p (6.11) 为了尽量减少隐式相位处理的成本（就CPU时间和内存要求而言），这个显式阶段也与简单的一阶上风隐式阶段耦合，因此所谓的RBC(3)隐式方案读作（Id + K')Awj” = -KRBCw? 3'（6.12）我们现在将专注于DNC（3）（6.7）和RBC（3）（6.12）方案的放大因子分析。 可以从（6.3）中立即推断出，对于线性问题，两个方案的通用隐式运算元（使用每个方向三点）读作p= 1，1 q = ApSppp - zlAplS；或者，它可以表示为H=D+CSP p=l，p- P P PP'，其中p方向上的移位算子，\&;是\&;Awjn = \textasciitilde{}Iwjn+\textasciitilde{}，非对角系数由A定义：=;（Ap+ IAPI）和A；=；（Ap - lApl）；隐级的对角系数由D = Id+）lApI给出。s - A-\&+l- A+\&-1 p=l

\section{6.3 直接求解器效率}
112 基于隐性残基的紧凑方案 6.3 直接求解器 效率 通过研究其相关的放大系数，对方案DNC（3）和RBC（3）的稳定性和效率属性进行了分解。 分别表示H和K是方案\%AwY = -KwY的隐式和显式阶段的傅里叶符号，放大矩阵G，读作G，= Id - H-' ' K。 进一步指出Tp（resp. S，）运算符lp的傅里叶符号（resp. Sp)，DNC(3)和RBC(3)的H由d d p = 1 p=l给出，Tp = 21AJ,z, + iApsZn(Q), sp = -Apcos(\textasciitilde{}p) + iAps2n(tp)，其中Q是与空间方向p和zp = +(l - COS([\textasciitilde{}))相关的还原波数。 分别与DNC(3)和RBC(3)相关的傅里叶符号读取KDNC= C,=l d,IA,I.; 4 + i C\&l(Id +;.p)sin(lp)Ap KRBC = C,",=, (24,1A,l.p + b\#)Psgn(AP)sin(tP) c:=l,q\#p AqSW)) +iCp=1(1 - \$ c:=l,q\#p z,)A,sin([p)。 （6.13）在对DNC（3）和RBC（3）扩增因子进行数值分析之前，可以通过分析评估由于选择一阶隐式阶段而导致的内在效率损失，该阶段与三阶显式阶段不一致，但受到全局效率（CPU成本）和内存需求考虑的激励。 在最短波长的特殊情况下，对应于减少波数平面中的Q = n，KDNC = f C:=, /A,(而KRBC = c:=, 24,IA,I和H = Id + Cp=l 21ApI对于两种方案。 由于G，=（H-K）/H在标量情况下，因此，当使用大型时间步骤时，GfNC-用于最短的波长，而GflBC-1-'\&''plAp'；让我们回顾一下，当隐式阶段与显式阶段一致时，G，-+ 0。 对于RBC（3），很明显，如果系数\&都等于1，那么最短波长的误差阻尼将是理想的。 然而，由于4和q\$必须满足q\& = aqp4，aqp = lAql/JApl，选择dP = lb'p不能保留在一般c;=，IAPI中

C. 科雷和A。 Lerat 113案例。 如果GfZBC使用以下身份表示，可以评估选择（6.9）对RBC（3）效率的影响。事实上，因为在这种情况下，GfZBC - 1-x（aqp）与x E [0,1]，x取的最小值对应于最短波长的最坏预期内在效率。 很容易检查，对于20对流问题，对于Q = 0.414或Q = 2.414，x（a = u）达到等于0.8284的最小值；因此，最短波长方案RBC（3）的放大系数可以高达0.1716，这当然比一阶隐式方案（G，-0对于大CFL数）效率低，但仍然比方案DNC（3）与GfNC-i的效率更高。 对于30个问题，最短波长最差的阻尼系数是GFBC-0.268，对于RBC（3）来说，这仍然比GfNC-？ j由DNC（3）提供。 DNC(3)和RBC(3)的放大因子可以使用应用于20个标量对流问题或20线性化Euler 方程的数值方法进一步分析。 在第一种情况下，放大系数G取决于与每个空间方向相关的CFL数，A，= CFL，以及减少的波数<，而在第二种情况下，放大系数，即放大矩阵G的光谱半径，取决于马赫数，流动方向u/u（u和u是速度的笛卡尔分量），长宽比AR = 6xZ/bx1和CFL数乘以物理物理时间步骤的CFL数。 在每种情况下，将对G，（el，62）或p（G，（G，（2）））进行全局和局部分析。 在全局分析中，波数平面中放大系数的最大值和平均值是为参数CFL或（M，CFL，u/u，AR）的每个组合计算的：如果所有参数组合的放大系数的最大值不大于1，这意味着正在研究的方案是无条件的稳定的；在这种情况下，通过计算每个给定参数的波数平面中放大因子的平均值G或p（G，），并绘制该平均值的轮廓来评估方案的内在效率是有意义的 参数平面。 分别由（6.7）和（6.12）定义的DNC（3）和RBC（3）方案，以及由K = K'和“H: = Id + K”定义的一阶隐式方案，当其隐式阶段被精确解决（使用所谓的直接求解器，以下简称DS），对于标量对流问题和线性化Euler 方程来说，是无条件稳定的。 透过查看图的左侧，可以一目了然地评估这些方案在标量情况下的效率。 6.1. 绘制G的等高线，确实可以观察到，对于大型CFL数，G，对于一阶隐式方案倾向于0，这是G,(G,\&) + 0 for ([I, \&) \# (0(27r),0(2\textasciitilde{}))这一事实的直接后果。 当我全部

\section{6.4 隐含的治疗描述}
114 基于隐性残差的紧凑方案 使用了三阶显式阶段，如DNC（3），效率损失是显而易见的，因为在大型CFL数字G中，在0.44到0.50之间变化，取决于对流速度相对于网格的方向。 当使用显式阶段RBC（3）时，G在大CFL数中变化在0.09到0.25之间，这说明了RBC（3）比DNC（3）提供更好的误差阻尼，这与之前对G的分析一致，对于最短的波长。 图的右侧绘制了一组CFL数（A1 = 100，A，= 100）的放大系数的典型轮廓。 6.1：完全隐式一阶方案对所有错误模式的阻尼几乎是完美的；使用三阶显式阶段会严重降低这些阻尼属性，然而，在使用基于残差的契约方案时，这些属性的保存效果要好得多。 当这些方案应用于二维线性化Euler 方程时，同样的结论成立。 保持流动方向v/u和长宽比AR固定为1，并扫过Mach和CFL数字，可以再次观察到RBC（3）方案提供的更好效率，尽管当然不如完全隐式一阶方案（见图 6.2）。 6.4 隐性治疗描述 第6.3节中获得的结果并没有说明DNC（3）和RBC（3）隐性方案的全部情况。 事实上，如前所述，由于成本原因，直接求解器很少在实践中使用；相反，已经开发了各种因数分解和放松方法，以大致解决隐性阶段，以获得更合理的成本。 因此，也有必要评估隐性阶段解决方案对方案效率的作用。 在本研究中，我们将考虑以下众所周知的隐式处理方法：近似因子化（AF）和修改近似因子化（MAF）[13]，两者都是非迭代方法；这些因子化方法的迭代版本，其中因数分解错误通过亚迭代过程校正，分别表示为AF（p）和MAF（p）[ll]；Jacobi（ALJ（p））和对称高斯-Seidel（ALGS（p））类型的交替线松弛处理[4]。 正如[2，11）中已经指出的，分解和放松方法之间的这种区别往往具有误导性；因此，在分析这些应用于一阶隐式方案或DNC（3）和RBC（3）的各种隐式处理的稳定性和效率属性之前，我们将首先在统一的框架中表达它们。 涉及子迭代过程的隐式方案可以放在一般形式Aw(0) = 0 1 = 0,p - 1 cu=l,N (6.14) X a. aw(lN+a) = AWe'Pl - (X - Xa)。 Aw(lN+"-l) Awn = Aw(P*)

C. 来\&A。 Lerat 115 图6.1：2D对流。 左：CFL数字平面中JG（A1，Az）l平均值的轮廓。 右：扩增系数的轮廓（G（G，\&）l在A的波数平面上，= A2 = 100。 顶部：一阶隐式方案；中间：DNC（3）；底部：RBC（3）。

116 隐含的基于残基的紧凑方案d 60 Y 05 I I5 Mach Mach 图6.2：2D线性化Euler 方程。 左：（M，CFL）平面中p（G（A1，Az））平均值的等高线，w/u = 1和AR = 1。 右：M = 0.7，CFL = 100和w/u = 1，AR = 1的波数平面中放大系数p（G（fl，\&））的轮廓。 顶部：一阶隐式方案；中间：DNC（3）；底部：RBC（3）。

C. 科雷和A。 Lerat 117，其中Aw(l) = w(l) - wn和Ha是样本隐式算子，在子迭代过程的每个阶段连续替换精确算子'FI；sim-ple意味着与'FI相关的线性系统通常是三对角线的，这反过来意味着'Ha是给定空间方向上的三点算子。 在隐式运算符'FI'的p应用之后，隐式增量的结果估计，Aw(pN)，保留为隐式增量Awn = wn+'-wn的近似值。 对于20中的线-Jacobi松弛过程，N = 2和'HI = D+S1，7-12 = D+S2，这导致i (D + s2)\textasciitilde{}wy+2) = aWeXp1 3 - s1aw,(2z+1) (D + Sl)Aw:2z+') = AweXPz 3 - S2Awy with 1 = 0,p- 1。 同样，通过取N = 4并连续应用以下三点隐式算子7i1 = D+S1，'H2 = D+\& +A,\&\$1, 7\& 7\& D +S2 -A:\&','FI4 = D +S2 得到交替线对称高斯-塞德尔松弛过程。一般公式（6.14）也允许描述AF和MAF处理；事实上，取N = 1与'HI = 'HAP = (Id + 7i)。 （Id + 12）并使用p = 1，z.e.单个内部迭代，在（6.14）中产生标准的AF处理（Id + '\&）。 （Id + Z>Qw,” = Aw,”"\textasciitilde{}'分两个步骤解决，每个步骤都涉及三对角线系统的求解（6.15）同样，MAF由N = 1，p = 1和'HI = 'HMAF = (D + S1)获得。 D-'。 （D + Sz），导致（Id + '\&）Aw; = Aw,eXp1 \{ (Id + = Aw;。 （D + S1）。 D-l(D + S2)Aw; = AweXpz 3，分三个步骤解决，涉及两个三对角线系统解和矩阵向量积（D + s1）a.w; = AWj expl AW;* = DAw; (D + S2)a\textasciitilde{}j" = Aw;*。 （6.16）AF和MAF的迭代版本用p>1获得。 请注意，由于WAF = \textasciitilde{} + 'PAF，PAF = 5给出的因数分解误差。 12，7-11的方案（6.14）='HAF也可以以7-1的形式重新铸造。 啊！' +'）+pAF。 啊！' +') = AWexP1 + PAP。 AW;.') 3 3

118个隐含的基于残差的紧凑方案，清楚地表明亚迭代过程旨在消除因数分解错误[ll]。 相同的公式适用于PMAF = S1的MAF(p)。 D-l。 Sz。 自然，如果公式（6.14）可以统一分析AF、MAF、AF（p）、MAF（p）、ALJ（p）和ALGS（p）的放大系数，必须指出，这个表达式不太适合实际目的，因为它使用了与因数分解误差相关的运算符，计算成本可能很高，因为它是3点方向运算符的乘积。 为了摆脱运算符P，使用增量6w(') = w(') - w('-')作为未知值来表达AF(p)和MAF(p)更方便。 事实上，在这种情况下，（6.14）可以等价地表示为w（0）= wn l=O，p-l，由于上述公式仅使用XI = 7-t\textasciitilde{}\textasciitilde{}或XI = ZMAF和X，因此可以按照AF（p）的（6.15）和MAF（p）的（6.16）启发程序，以较低的成本计算。 使用（6.14）已经证明，[4]与迭代方案（6.14）相关的放大因子G由G = G, + VP(Id - G,) (6.18) V = II:==,(Id - HglH)给出，其中G，是与直接求解器\%Awn = -Kwn相关的放大因子。 这个有趣的公式清楚地显示了高效隐式方案的两个成分：首先，确切的隐式方案，即直接求解器，必须提供高内在效率，该属性仅取决于定义G的傅立叶符号K和H；其次，实践中用于求解隐式阶段的迭代方法必须收敛到参考放大因子G，成本低，即乘积N xp的小值，此收敛速度取决于因子V，其本身是傅立叶符号H和Ha的函数。 由于IG - G,I = lVlPll - G,I在标量情况下，很明显，高效的迭代方法将为大多数错误模式提供接近零的IVI。 另请注意，如果使用单个内部迭代（p = 1）处理AF（p）或MAF（p），公式（6.18）生成G = Id - HikK或G = Id - H\&iFK，这当然是基本AF或MAF方法的放大因子。 众所周知，AF不是有效的隐式处理；假设它用于求解一阶隐式方案，二维标量案例中的放大因子的形式为GAF = (HAF - K)/HAF，K = TI + T2和HAF =

C. 来t3 A。 Lerat 119 H + T1T2 = 1 + K + T1T2，因此，当时间步骤变大时，GAF + 1，因为因数分解误差PAF = T1T2 = O(At2)相对于H = O(At)变得显性。 使用MAF可以在大时间步骤下改进收敛，因为在这种情况下，因式分解误差是PMAF = SID-\textasciitilde{}S\textasciitilde{} = @At），因此，在At 4 00时不再占主导地位[13]。 然而，利用因数分解错误消除过程来恢复接近直接求解器的内在效率更有利。 因此，我们现在将重点分析与一阶隐性阶段的每个迭代隐式处理相关的因子1V）。 AF（p）对于一个标量问题，VAF可以写成VAF = H\&:FH = H:yA,。 因为，对于\&维度问题，PAF = O(Atd)，而H = O(At)，因此\textasciitilde{}VAF\textasciitilde{} --t 1当At变大时，因此将非常缓慢地收敛到G，当使用较大的时间步进时，不太可能高效率。 MAF(p)，其中PMAF = S1D-lS2 对于标量问题，VMAF = "\textasciitilde{}\textasciitilde{}\textasciitilde{}" = HT\%gF，因此，在2D 同样，在3D的情况下，VMAF由（6.19）（6.20）给出。对于最短的波长，对应于波数Q = T，对于大CFL数，VMAF是2D中的VMAF N \$，VMAF N \&在3D中，这确保了与使用AF时更快的直接求解器的收敛方案。 ALJ（p）对于Jacobi类型的交替线松弛，WP = D+Sp，p = 1，d，因此对于二维标量问题VAL J = yA:\&y x 7A:;f'；考虑到H = D + S1 + S2，此表达式可以重新表述为（6.21）从公式（6.19）和（6.21）中可以看出，MAF（p）和ALJ（p）处理在2D中是相同的，因为它们对V的表达式具有相同的表达式。 然而，请注意，这在3D中不再正确，因为在这种情况下（6.22）

\section{6.5 迭代求解器效率和稳定性}
120个隐含的基于剩余的紧凑计划，与（6.20）不同。 对于最短的波长和大CFL数量，VALJ - 20分之9，当然像VMAF一样，但VALJ N \& = 4 in 30，这意味着当使用ALJ(p)而不是MAF(p)时，可以预期收敛到G略快。 ALGS（p）四个以下运算符 对称线高斯-塞德尔松弛在2D中定义为Hl = D + S1 - A,fe-i<2 = H - A- H2 = D + S1 + A; e+ZC2 = H + A; e-aCz H3 = D + S2 - AFe-iE1 = H - A-e+i<l H4 = D + Sz + A,e+acl = H + A+ -6 \{ le so VALG\textasciitilde{}读出vALG\textasciitilde{} = (I - HT\textasciitilde{}H) x (I - H;\textasciitilde{}H) x (I - H,-H,-\textasciitilde{}H) x (I,-\textasciitilde{}H) (6.23) A+A-A-A-，在标量情况下减少到VALGS = E;1\$2HtH:. 对于标量一阶隐式阶段，A：= 0或A；= 0，因此所有波数的VALGS = 0，因此GALGS（，）= G，这并不奇怪，因为在这种情况下，ALGS（p）准确地解决了与隐式阶段相关的线性系统。 公式（6.21）和（6.23）将在下一节中用于研究由ALJ（p）和ALGS（p）解决的DNC（3）和RBC（3）方案的行为，用于20个标量对流问题或20个线性化Euler 方程。 6.5 迭代求解器 效率和稳定性 我们现在转向DNC（3）和RBC（3）方案的放大因子的数值分析，用于二维标量对流问题，以及使用上一节中描述的方法之一解决隐性阶段时的二维线性Euler 方程。 如第6.3节所述，当隐性阶段由Jacobi或Gauss-Seidel类型的线松弛程序解决时，DNC（3）和RBC（3）方案的放大因子的分析将依赖于对（G（G，\&？）（）的最大值的研究。 （resp（G（\&，\&）））在CFL中，平面（resp.（M，CFL）平面）评估稳定性，并研究放大因子的平均值来评估效率；还将检查参数CFL或（MI CFL）固定值的放大系数的轮廓。 标量对流扫荡（CFLl，CFL2）发现ALJ（p）以及AF，AF（p）只有在应用于DNC（3）或

C. 海湾63 A。 Lerat 121 RBC（3）；ALGS（p）在这个标量情况下不考虑，因为它等同于第6.3节中研究的DS。 显然，这种条件稳定性是使用与高阶显式阶段不一致的隐式阶段的后果；当与一阶隐式方案结合使用时，这些方法确实变得无条件稳定。 如图所示。 6.3，增加ALJ（p）的内部迭代次数可以减少数值不稳定性，直到获得无条件稳定性（在研究的域中，每个CFL从-100扩展到100），在DNC（3）的情况下，p = 10，RBC（3）的p = 8。 显然，DNC(3)和RBC(3)的ALJ(p)是JVI < 1，这允许在足够数量的内部迭代中恢复直接求解器效率。 当方案稳定时，绘制其放大系数的平均值是有意义的：由于G + G，当p增加时，找到图6.3中显示的用ALJ（10）求解的DNC（3）和用ALJ（8）求解的RBC（3）的G的等高线图与这些方案的直接求解器版本非常相似（见图6.1），因此，在实践中求解的RBC（3）产生比DNC（3）更好的效率是合乎逻辑的。 线性化Euler 方程扫过马赫数和CFL数（同时保持w/u = 1和AR = l），发现ALGS（p）在条件上是稳定的，不仅在应用于DNC（3）和RBC（3）时，而且在与一阶隐式方案一起使用时，尽管该方案是完全隐式的。 如图所示。 6.4，使用ALGS（p）来解决后一种方案，并将p从1增加到4，增加了放大系数的最大值；这种行为是IVI > 1的直接后果。 在DNC（3）和RBC（3）中也观察到了类似的行为。 通过绘制一组给定参数Mach和CFL数字的放大系数的轮廓，即（A4 = 0.7，CFL = 100），可以更深入地了解这种不稳定性的本地化。 研究这些图（见图6.4）可以观察到ALGS（p）到直接求解器的所有波数的非常快的收敛，但那些靠近O（27r）的波数发生了之前观察到的强烈不稳定。 情况确实随着ALJ（p）的改善而改善，因为这种处理对系统案例的行为仍然与标量案例中记录的相似：对于足够大的p，ALJ（p）对DNC（3）和RBC（3）都变得稳定，表明IVI < 1对于所有波数。 当w/u = 1，AR = 1时，发现p = 8确保了M E [0.3,1.5]和CFL E [lo, 1001的无条件稳定性。 在相同条件下绘制放大系数的平均值（见图6.5）表明，ALJ（8）也为DNC（3）和RBC（3）提供了接近直接求解器的效率（见6.2）。 从全局到本地，并绘制两个方案的放大系数的等高线值（M = 0.7，CFL = 100，v/u = 1和AR = 1）可以检查\textasciitilde{}（GALJ（\textasciitilde{}））确实非常接近DNC（3）和RBC（3）的p（G，），（放松）RBC方案再次提供了更好的效率

122 基于隐性残基的紧凑方案图6.3：二维对流。 顶部和中间：用ALJ（4）（左）和ALJ（8）（右）解决的方案DNC（3）（顶部）放大系数的最大值的轮廓，以及用ALJ（4）（左）和ALJ（6）（右）解决的方案RBC（3）（中间）。 底部：方案DNC（3）/AL J（10）（左）和RBC（3）/AL J（8）（右）的放大系数平均值的轮廓。

\section{6.6 结束语}
C. 海湾63 A。 Lerat 123常规DNC方案。 6.6 总结性言 在构建传统非紧凑三阶上风方案DNC（3）的高效隐式版本时，习惯上使用简单的一阶上风隐式阶段，以尽量减少每次迭代的成本。 这种选择会影响高阶隐式方案的内在效率，特别是由于每方向5点显式阶段和每方向3点隐式阶段之间的强烈不一致，高频错误模式的阻尼速度要慢得多。 如果基于三阶残差的方案RBC（3）使用相同的一阶上风隐式阶段进行隐式，其模板的紧凑性（2D中的3 x 3点）可以显著降低内在效率的损失，正如使用直接求解器求解的DNC（3）和RBC（3）的放大系数的理论和数值分析所证明的那样。 然而，要真正感兴趣，当隐性阶段像实践中所做的成本和内存要求降低时大致解决时，这一观察也必须成立。 隐性阶段的常见近似解方法是因式分解和松弛技术。 为这些各种隐式处理提出了一个统一的公式，并允许将备用线Jacobi和Gauss-Seidel方法确定为特别有效的处理方法，在几次内部迭代后恢复直接求解器内在效率；[ll]中提出的带有迭代错误校正的修改近似因数分解也可以被列为效率，因为它被证明等同于2D中的ALJ（p）。 然后，通过计算使用ALJ（p）求解的2D标量对流问题的DNC（3）和RBC（3）的放大系数和2D线性化Euler 方程进行检查，在足够数量的内部迭代（通常p小于10）下，这两种方案都变得无条件稳定，基于残差的方案提供了更好的误差阻尼。 当使用ALGS(p)时，在线性化Euler 方程的情况下，直接求解器放大因子在一次内部迭代后几乎完全恢复，但如之前在[a]中指出的那样，接近0的波数会出现强烈的不稳定性。 虽然受到一些限制性假设（线性问题、周期性边界条件）的限制，但Von Neumann分析的发现在DNC（3）和RBC（3）隐式方案的实际经验中得到了很好的反映：在[8]中使用这两种方法进行的模型问题和可压缩流的计算确实系统地显示了扩增因子分析预测的RBC（3）的更好效率。 然而，在计算p（G）时，在实践中没有观察到上述ALGS（p）处理的（非常局部的）不稳定性，因为p = 1的替代线高斯-Seidel在[8]中成功使用，但也在[7，91中，以确保RBC（3）的快速收敛达到稳态；将对ALGS（p）进行进一步分析，以了解这种差异。

12.4 基于隐性残数的紧凑方案 Mach 图 6.4：2D Euler 方程。 顶部：使用ALGS（p）求解的一阶隐式上风方案p（G）最大值的轮廓，p = 1（左）和p = 4（右）。 中间：使用ALGS（1）求解方案DNC（3）（左）和RBC（3）（右）的波数平面中的p（G）等高线。 底部：沿着波数平面对角线的上述等高线图的切线。

C. 核心和A。 Lerat 125 1.1.." 马赫'0.9 0.8 0.7 Go。 @ \$0.5 z 0.4 0.3 0.2 0.l 0.5 1.5 1.".. I Mach ' 0.9 0.8 0.1 z0.6 \$0.5 c1 -0.4 0.3 0.2 a. I 1 1.5 2 0 0.5 S，图6.5：2D Euler 方程。 顶部：使用ALJ（8）求解的p（G..方案DNC（3）（左）和RBC（3）（右）的平均值的轮廓。 中间：使用ALJ（8）求解方案DNC（3）（左）和RBC（3）（右）的波数平面中的p（G）的连接。 底部：沿着波数平面对角线的上述等高线图的切线。

\section{6.7 参考书目}
126 隐含的基于残基的紧凑方案 6.7 参考书目 [I] Briley, W. R。 \& 麦当劳，H. 隐式Navier-Stokes算法和近似因式分解的概述和概括，计算机[2]Buelow，P. E。 O., Venkateswaran, S., \& Merkle, C. L。 隐式逆风方案的稳定性和收敛分析，计算机与流体30，pp. 961-988（2001）。 [3] 考伊，D. A.. 尤拉方程和流体的对角线隐式多网格算法，30，pp. 807-828（2001）。tions，AIAA杂志26，pp. 841-851（1988）。 [4] Corre, C., Khalfallah, K., \& Lerat, A. 一类中心方案的线松弛方法，计算流体动力学杂志，5，pp. 213-246（1996）。 [5] 科尔，C. \& 勒拉特，A. 基于高阶剩余的契约计划，2004年7月，加拿大多伦多计算流体动力学第三届国际会议，将于2004年计算流体动力学 D. Zingg（编辑），Springer（2005）。 [6]詹姆森，A。 和Caughey，D。 A. 解决Steady的Euler 方程需要多少个步骤，可压缩流动：寻找快速解决方案算法，AIAA 2001-2673（2001）。 [7] 莱拉特，A。 \& Corre，C。 可压缩Navier-Stokes 方程的A 基于残差的紧致格式，J。 计算机。 物理，170，pp. 642-675（2001）。 [8] Lerat, A., Corre, C., \& Hanss, G. 高阶可压缩流动计算的基于残差的紧凑性与方向性，数值流体力学笔记，卷。 78，pp. 61-76（2001）。 [9] Lerat，A. \& Corre，C。 多维双曲守恒定律系统的基于残基的紧凑方案，计算机与流体，31，第31页。 639-661（2002）。[lo] Luo, H., Baum, J., \& Lohner, R. 用于非结构化网格上可压缩流的快速、无矩阵的隐式方法，计算物理学杂志146，pp. 664-690（1998）。[ll] MacCormack，R. W。 迭代修改近似因式分解，电子和流体，第30页。 917-925（2001）。 [I21 Pulliam，T. H。 \& Chaussee，D。 S。 隐式近似因子化算法的对角形式，J. 计算机。 物理。 39，pp. 347-363（1981）。

C. 科雷和A。 Lerat 127 [13] Pulliam，T。 H.，MacCormack，R. W.，\& Venkateswaran，S. 收敛近似因数分解方法的特征，第16届流体力学数值方法国际会议，法国阿卡雄，1998年7月，物理学讲义，515（1998）。 [14]横田，J。 W.，Caughey，D. A.，和Chima，R. V。 下上隐式方案的对角反转，卷。 28，第2号（1990）。

此页面故意留空

\clearpage
\chapter{第7章动态非结构网格 CFD模拟的高阶时间积分方案Dimitri J. Mavriplis和Zhi Yang}
\section{7.1 摘要}
第7章动态非结构网格 CFD模拟的高阶时间积分方案Dimitri J. Mavriplis' and Zhi Yang' 7.1 摘要 具有相对物体运动的非定常流动模拟的解决方案通常需要使用动态变形网格。 为了为此类问题构建有竞争力的模拟能力，必须设计有效积分流动方程以及网格运动方程的技术。 在这项工作中，我们演示了使用多网格方法来解决非定常流动方程和网格运动方程。 在以前的工作中，发现对静态网格上的非定常流动方程使用高阶时间积分技术优于低阶时间积分方法。 在这项工作中，这些技术被扩展到包括动态变形的网格。 为了保持准确性和稳定性，这需要构建满足离散几何守恒律的方案。 给出了DGCL的两种不同结构，数值结果用于验证三维黏性流动的这些方案。 怀俄明大学机械工程系。 1000 E。 大学Av-enue 拉勒米，怀俄明州82071

\section{7.2 介绍}
130 动态非结构网格的驯化集成 7.2 导言 随着从业人员寻求对更复杂的物理问题进行建模，计算成本不断降低，高保真时间准确流体流动模拟变得越来越重要。 然而，时间精确的计算流体动力学（CFD）问题仍然比等效的稳态问题贵几个数量级，从而促使人们寻找更高效、更准确的算法。 对于在时间和空间尺度上明显分离的问题，如非稳态空气动力学、空气弹性和其他移动物体问题，其中物体运动的时间尺度与特征流体时间尺度相去甚远，实际解决方法需要隐式时间积分方案。 当需要小时间误差时，高阶时间积分（高于二阶）已被证明比低阶时间积分更有效率。 Bijl、Carpenter和Vatsa [3]在结构化网格上调查和比较了高阶隐式Runge-Kutta方案和向后微分方案，而Jothiprasad、Mavriplis和Caughey[7]展示了四阶Runge-Kutta方案（RK64）如何在非结构化网格上优于二阶后向后方差分方案（BDF2），使用多网格算法解决每个时间步骤产生的隐式系统，以求出现的隐式系统。 然而，大部分关于高阶时间精度的工作都是针对涉及静态网格的案例进行的。 存在涉及相对物体运动的一大类问题，如空气弹性，这需要使用动态变形的计算网格。 本文的目标之一是将为静态网格应用设计的高阶时间准确方法的公式扩展到动态网格问题，使用几何灵活性的非结构化任意拉格朗日-埃勒（ALE）公式。 当使用动态网格时，需要仔细考虑网格速度和其他与几何相关的参数，以便网格变形带来的误差不会降低流动模拟的形式准确性。 离散几何守恒律（DGCL）提供了关于如何评估这些数量的指南。 一阶和二阶时间准确和几何保守方案分别在[11，101中提出和讨论。 Guillard和Farhat已经证明，要至少获得一阶时间精度，必须满足DGCL条件[5]。 对于高阶时间积分方案，必须明确构建一个保留方案设计顺序的DGCL策略。 除了保持时间准确性外，还必须采用高效的解决方案策略，以避免长时间积分问题的计算时间过长。 之前已经研究了非线性和线性多网格方法，用于使用非结构化网格的稳定和非定常流动模拟[12，13，141。 在这项工作中，我们研究了非线性和线性非结构聚集多网格方法对非稳态的时间积分的使用

\section{7.3 控制方程在任意-拉格朗日-尤拉（ALE）形式和基础流动求解器}
D。 J。 毛里普利斯和2。 杨131流动方程，以及用于解决网状变形的方程。 在接下来的章节中，我们首先概述了控制方程和基础流动求解器。 然后，我们讨论了两种时间积分方案的选择和实施：二阶和三阶向后差分方案（BDF2，BDF3），然后是离散几何守恒定律（DGCL）的介绍，这些方案必须尊重这些方案。 然后描述了已经调查的各种网格变形策略，然后是用于流动和网格运动方程的多网格解技术。 然后将这些方案的性能与非黏性三维变形网格问题进行比较。 在第二部分，我们讨论了使用高阶隐式Runge-Kutta（IRK）方法（最高四阶时间精度），基于我们之前对这些方法在静态网格情况下的经验[7]。 给出了一种构建尊重离散守恒律（DGCL）的IRK方案的技术，并在三维非黏性变形网格流问题上演示了四阶IRK方案。 7.3 控制方程在任意-拉格朗日-尤拉（ALE）形式和基础流动求解器中，保守形式的Navier-Stokes 方程可以写成：aU - + V.（ F(U) + G(U)) = 0，其中U表示守恒量（质量、动量和能量）的向量，F（U）表示对流通量，G（U）表示黏性通量。 在a（移动的）控制体 O(t)上积分，我们得到：使用s,,,) udv =的微分恒等式a U(k. ii)dS。 I,,) gdV + JLt) (7.3) 其中k和ii是接口do@的速度和法线，分别，方程(7.2)变成：(F(U) - kU). iidS a + /an(t) G(U). iidS = 0 (7.4)

\section{7.4 高级时间积分和离散几何守恒律}
1 32 动态非结构化网格的驯化积分 考虑U作为单元平均量，这些方程在空间中离散化为：a -(VU) + R(U,B(t), ii(t)) + S(U, ii(t)) = 0 at (7.5)，其中R(U, 8, ii) = (F(U) - kU)。iidS表示ALE形式的离散对流通量，S(U, ii)表示离散黏性通量，V表示控制体。 在离散形式中，8(t)和ii(t)现在表示控制体边界面的时间变化速度和表面法线。 Navier-Stokes 方程通过中心差分有限体积方案离散化，在混合网格上具有额外的基于矩阵的人工耗散，其中可能包括两个维度的三角形和四边形元素，或三维的四面体、金字塔、棱镜和六面体。 使用人工耗散运算符的双通结构实现二阶精度，该结构对应于未分割双谐波运算符。 在流动求解器中，所有类型的元素都使用基于边缘的单一统一数据结构。 采用Navier-Stokes 方程的薄层形式，通过非简单1元素的有限差分近似将黏性项离散化为二阶精度。 对于多网格计算，一阶离散化用于粗网格水平P51上的对流项。 7.4 高阶时间积分和离散几何守恒律 对于非定常流动模拟，最常采用完全隐式时间积分策略，使用多步后向差分公式（BDF）或多阶段隐式Runge-Kutta（IRK）方案。 虽然一阶（BDF1）和二阶（BDF2）向后差分方案是A-稳定的，但高阶多步BDF方案（二阶以上）不是A-稳定的。 然而，BDF3方案的不稳定区域非常小，该方案最常成功用于非定常流动模拟。 A稳定和L稳定的高阶多级IRK方案可以轻松构建。 然而，在每个时间步骤中都需要使用IRK方案解决多个非线性问题，这使得这些BDF方案的替代方案更加昂贵[3]。 在这项工作中，我们研究了BDF和IRK方案，最初集中在BDF方案上，因为以前已经展示了在变形网格背景下使用BDF方案。 方程（7.5）可以使用k步退的一般公式重写-

D。 J。 Mavriplis和Z。 Yang 133 ward Difference scheme as [3, 71 1-k (Yl(vu)n+l + Cai(\textasciitilde{}\textasciitilde{})n+i = \textasciitilde{}t\textasciitilde{}((\textasciitilde{}\textasciitilde{})\textasciitilde{}+l,t\textasciitilde{}+l) (7.6) i=O，其中R现在表示全（对流和黏性）残差。 对于BDF2方案，k = 2，系数如下：a1 = 4，（YO = -2，（Y-1 = i，而对于BDF3，k = 3，系数如下：a1 = v，a0 = -3，a-1 = 3，a-2 = -4。 通过定义非线性残差[7] 3 1-k,"+1(U) =,(un+') \_= a1(VU)"+l + Ca,.( VU)n+i i=O - AtR((VU)"+', tn+') (7.7) 方程(7.6)的解可以通过在每个时间步求解非线性问题P+'(U) = 0得到。 两种不同的方法用于求解上述方程：非线性多网格全近似存储（FAS）聚集方法，以及线性聚集多网格（LMG）方法，在近似非线性牛顿迭代策略中的每个阶段用作求解器。 这两种方法的详细信息可以在参考文献[14，71中找到。 为了在动态非结构网格的存在下应用多步BDF方案，应满足几何守恒律（GCL），以避免降低方案的形式准确性，并保持时间积分方案的稳定性。 几何守恒律的原始声明是由Thomas和Lombard[18]为结构网格引入的。 离散几何守恒律要求状态U =常数是方程（7.5）的精确解。 在这种情况下，我们有S(U,ii) = 0，因为黏性通量是基于U的梯度。 此外，我们有：（F（U）-kU）。iidS = R（U，k，fi）= -UR（k，ii）（7.8），因为闭合控制体周围的对流通量F（U）的积分对于常数U，对于任何空间守恒格式，R是指上述边界积分中第二个项的离散化。 因此，GCL可以以半离散形式表示为：在--R（k（t），ii（t））=0（7.9）几何守恒律必须以离散形式满足，或所谓的离散几何守恒律。 鉴于参考[4]中给出的证明，这一点尤为重要，其中DGCL条件被证明是保留非线性稳定性的必要和充分的

\section{7.5 网格运动策略}
1 34 动态非结构化网格的时间整合基线静态网格时间整合方案。 Lesoinne和Farhat [ll]提供了服从离散几何保守定律的一阶时间精确后差分方案（BDF1），而Koobus和Farhat [9]则推论了服从离散守恒律的二阶精确后差分方案（BDF2）。 我们在工作中利用了BDF2 DGCL合规方案的构建。 此外，由于发现高阶时间积分在不稳定流模拟中优于低阶时间积分[3，71]，遵循离散几何守恒律的三阶时间准确后差分方案（BDF3）也是根据[ll，91和下面介绍的方法推导的。 最终结果包括一个公式，用于计算每个时间步骤的网格点速度和控制体面法线的适当值，这既保留了BDF3方案的形式准确性，又尊重GCL（即承认均匀流动作为精确的解决方案）。 最终方案给出如下：zc\textasciitilde{}i(VU)”+\textasciitilde{} = AtR((VU)n+l, tn+',j.(t),6(t)) (7.10) i=l，其中变量6k表示在时间位置n - 2和n + 1之间的不同正交点评估的控制体边界面的法向量，详见附录，而xn-'、xn-'、x"和xn+'值是指各自物理时间步骤的网格点位置，是前面描述的BDF方案系数。 附录中给出了BDF3的GCL的完整推导。 7.5 网格运动策略 对于具有移动或变形边界的非定常流动模拟，通常符合边界的计算网格必须与施加的边界运动一起位移。 网格运动必须使位移的网格配置保持高质量元素的平滑度，同时避免网格元素的交叉，从而导致无效的负体积元素。 计算变形网格配置的最常见技术是通过构建和解一组偏微分方程，这些方程用于控制受强加的边界条件约束的网格变形。

\subsection{7.5.1 张力弹簧类比}
\subsection{7.5.2 线性弹性类比}
D。 J。 毛里普利斯和2。 Yang 135 7.5.1 张力弹簧类比 张力弹簧类比也许是非结构网格变形最古老、最简单的策略（见示例[2]）。 在这种方法中，网眼的每个边缘都由一个弹簧表示，其刚度与边缘的长度有关。 控制方程与一个简单的拉普拉斯方程密切相关，因为每个坐标方向的位移都变得解耦，并受以下方程支配：（7.11）其中Zij = ((xi - xj)' + (yi - yj)2 + (zi - \textasciitilde{}j)\textasciitilde{}) 4和lcij = d。 参数p通常设置为2 [19]。 虽然这些方程相对容易解决，但这种方法往往会产生折叠或负细胞体积的变形网格，特别是对于黏性流动问题中使用的高长宽比网格。 由于拉普拉斯方程遵循最大原理，因此不难看出，即使在存在高弹簧常数的情况下，这种方法也无法再现固体体旋转，例如，在俯仰翼型的情况下，这将需要远离翼型表面的更大位移。 7.5.2 线性弹性类比 由于这种方法的鲁棒性，各种研究人员使用线性弹性方程来模拟网格变形[l，17，61。 计算网格假定服从线性弹性方程，该方程可以写成：（7.12）其中，在三维中，应力aij、应变eij和位移Ui给出为：

136 动态非结构化网格D=a和其余矩阵的调制积分如下：- 1-u u U 0 0 0 1-u u 0 0 U u 1-u 0 0 0 0 0 0 0 0;-u 0 0 0 0 0 i-u 0 2-u 1 - 0 0 0 0 0 其中E a= (1 + v)(l - 2u) (7.13)，其中E表示弹性模量，u是固体材料的泊松比。 通过引入形状函数N和U = NUe，在我们的案例中作为线性形状函数，并应用标准Galerkin方法，我们得到J（AN）\textasciitilde{}D（AN）U“\textasciitilde{}S = - S，NTfdS（7.14）n，可以重写为KU=F（7.15），其中（B）TD（B）dS，F = - NT fdS，B = AN。 =S，在网格变形的情况下，给出边界位移，因此不需要力向量F的外力fj。 相反，在U位移向量上的Dirchlet条件下，求解了齐次问题KU = 0。 线性弹性方法的一个优点是，大E（弹性模量）的区域将作为固体体位移。 因此，E分布的适当处方可用于避免网格关键区域的严重网格变形。 我们采用的E分布与细胞体积[6]或与变形边界的距离成反比，从而将大部分网格变形降级到网格更粗的区域，并且可以承受更大的相对变形。 事实证明，这对于避免高网格拉伸区域的无效网格单元至关重要。

\section{7.6 加速策略}
\section{7.7 网格运动结果}
D。 J。 Mavraplas和Z。 Yang 137 7.6 加速策略 如前所述，聚集多网格方法用于加速解决在非定常流动方程的每个时间步骤中出现的非线性问题。 聚集多网格最初是作为非结构化网格上的欧拉和Navier-Stokes 方程的稳态求解器而开发的[12，141。 多网格的想法是通过迭代计算对更粗网格级别的精细网格问题的校正来加速精细网格上的求解，该网格的迭代成本较低，并且全局误差分量更容易减少。 图7.la-d显示了聚集多网格程序的示例。 图7.la显示了精细网格，而图。 7.lb-d显示第2、第3和第4级粗网格，其中每个粗细胞是几个细细胞的组合。 在每个网格级别上，多阶段显式方案被用作平滑器来驱动多网格算法。 聚集多网格方法也用于加速网格运动问题的解决，并共享与流动求解器相同的粗层网格。 在这种情况下，Jacobi或Gauss-Seidel平滑器用于所有网格级别，以驱动网格运动方程的多网格算法。 非线性和线性聚集多网格算法已被用于解决流动方程的隐式时间积分问题。 多网格算法的非线性变体由标准全近似存储（FAS）多网格算法组成，其中非线性流动方程在每个网格级别上迭代。 对于线性多网格算法方法，在每个隐性时间步长出现的非线性流动问题由牛顿方案解决，该方案需要在每次牛顿迭代时解决一个大型线性系统，该方法使用线性多网格方法执行。 如前所述[14，71，非线性（FAS）多网格方法的优势包括低存储和高效的非线性启动行为，而线性多重网格法受益于更快的整体cpu时间要求，但以增加内存使用为代价。 上述控制方程的网格变形的各种形式也使用聚集多网格算法得到了解决。 由于此处考虑的网格运动方程是线性的，因此这些方程仅使用线性多网格方案。 7.7 网格运动结果 提出了几个大的偏转/变形问题，以比较由不同网格运动策略生成的网格的质量。 NACA0012 翼型被迫围绕四分之一弦点振荡，攻击角度为Q: = amaz sin(ft)，其中amaz设置为60英寸。 图7.2a-c显示了在翼型表面附近具有大拉伸的黏性流动型网格的结果。 除了弹簧类比网眼mo-

动态非结构化网格的时间积分编号0 02 04 06 0.8 1（a）第1级（c）第3级Od 0 02 04 06 06 1（b）第2级0 02 04 06 08 1（d）第4级图7.1：尺寸非结构网格。 粗聚集多网格水平的插图，用于二-

\subsection{网格运动方程的7.7.1 收敛}
D。 J。 Mavraplis和Z。 Yang 139 tion方程，考虑了两种线性弹性网格运动公式：一种是弹性模量E被规定为整个域的常数值，另一种是E被规定为与细胞面积成反比。 图中对恒定E线性弹性法观察到负面积网格单元（翼型表面内部）。 7.2 b. 弹性的可变模量（E）线性弹性方法被证明是最罗布斯特的方法，因为只有这种方法才能在最大俯仰角下产生有效的网格。 图7.2d-e显示了翼型壁附近中弦区域的网格的特写，说明了弹簧类比方法在翼型表面附近发生的相对网格变形，与线性弹性方法产生的网格相比，该方法在这个区域相对没有变形，因为它倾向于在可变E线性弹性方法下作为固体体平移和旋转而位移。 图7.3a说明了关于机翼体配置的三维非结构网格，其中机翼表面规定了大跨度偏转，从而诱发了网格的变形，该变形是在原始水平位置的机翼产生的。 这个网格在飞机表面附近高度伸展，机翼表面的正常间距约为10V6弦，总共包含473,025个顶点。 线性弹性法（E规定与细胞体积成反比）是在这种情况下唯一能够产生有效网格的方法。 图7.3b显示了网格中的最小单元体积与时间t。 时间t = 2.5对应于图中所示的机翼位置。 7.3a。 7.7.1 图中描述了非黏性二维俯仰翼型问题的网格运动方程的收敛。 7.4a-b，分别用于弹簧类比方程和线性弹性方程（变量E）。 块高斯-塞德尔平滑方法要么在精细网格上用作求解器，要么用作多网格序列的每个网格级别的平滑器。 在这些示例中，为了检查解决方案策略的有效性，残差减少了10个数量级，尽管在动态网格流模拟的背景下通常不需要如此严格的收敛标准。 在每个级别使用6个高斯-塞德尔平滑通道的三级多网格方案，为本例的弹簧类比方程实现了超过10阶的剩余减少。 这种快速的多重网格收敛是预期的，因为弹簧类比方程对应于缩放泊松方程，这些方程可以通过标准多网格方法轻松求解。 在这种情况下，多网格方案比单网格方案实现了接近一个数量级的加速。 没有多网格的线性弹性方程的收敛比弹簧类比方程慢得多，即使对于这个相对简单的测试案例也是如此。 然而，应用于这些方程的多网格算法实现了网格运动方程的收敛

14 0 动态非结构化网格的时间积分 0 05 (a) 弹簧 0 05 I (c) 线性弹性 (b) 线性弹性 (c0nstant E) (d) 弹簧(放大) (e) 线性弹性(放大) 图 7.2：二维黏性网格情况的不同网格运动策略的比较。

\section{7.8 使用反向差分方案的非定常流动模拟}
\subsection{7.8.1 多重网格收敛效率}
D。 J。 Mavriplis和Z。 Yang 14 1 图 7.3：三维黏性网格变形问题和最小细胞尺寸作为跨度偏转函数的插图。 在65个多网格周期中，剩余减少10个阶数，这只比弹簧类比方程慢2倍。 7.8 使用向后差分方案的非定常流动模拟 7.8.1 多重网格收敛效率 为经历强制扭转运动的ONERA M6机翼计算了具有动态解形网格的三维非稳态无黏流动模拟，以检查多网格方案的收敛效率。 M6机翼的三维非结构网格由53961个顶点和287962个四面体组成，如图所示。 7.5a。 使用四级聚集多网格算法，其中相同的聚集水平用于流动和动态网格运动方程求解器。 自由流马赫数为M = 0.84，初始入射率为a = 3.06"。 计算分为两部分。 首先，计算静止机翼上的稳态流量。 这个流场，如图所示。 然后，7.5b用于初始化扭转翼计算，在此期间，机翼被迫围绕四分之一弦线扭转，攻角为a = a，sin（ft），其中a，在翼尖设置为2.51英寸，a，从该值线性变化到翼根处为零；降低的频率为f = 0.1628。 扭转翼的计算非稳力升力和阻力系数是

142 动态非结构化网格的时间积分（a）弹簧fti5 - lo" - 10' - 20 40 60 80 lime step lO"1' ' ' \textasciitilde{} " I' " I' ' ' (b) 线性弹性图7.4：二维非黏性网格问题的网格运动策略的收敛历史。如图所示。 7.5\textasciitilde{}。 这些计算使用二阶反向差分方案（BDF2）进行，时间步进为At = 2.0。 在每个物理时间步骤中，流方程的非线性残差使用线性多网格方案或非线性多网格方案收敛，弹簧类比网格运动方程使用雅各布驱动的多网格方案收敛，每个层面有3个雅各布平滑通道。 图7.6a-b和7.7a-b检查了线性多网格方案的收敛效率，作为平滑通道次数和每次非线性更新的线性多网格周期次数的函数。 在图7.6a-b中，使用每个非线性更新的两个线性多网格周期，在每个网格级别上通过3个平滑通道获得cpu时间的非线性流动解决方案的最佳收敛（即使更多的平滑通道在多网格周期的基础上产生更快的收敛速率）。 请注意，由于平滑通道数量不足，系统会发散。 在图中。 7.7a-b 由于线性多网格求解器在非线性流动方程的近似（一阶）雅各布方程上运行（即线性多网格用于驱动非线性流动方程的近似牛顿方案），因此每个多网格循环的非线性收敛随着线性多网格循环数量的增加而渐近到下限。 就cpu时间而言，在每次非线性更新时通过2个线性多网格周期获得最佳的收敛效率。 使用线性多网格方法的最佳参数，比较了给定时间步骤下非线性流动残余的收敛效率，用于优化的线性多重网格法，以及使用相同粗水平的非线性（FAS）多重网格法，以及三阶段多阶段Jacobi预制

D。 J。 Mavriplis和Z。 Yang 143 图7.5：非结构网格和初始位置的计算密度轮廓，以及扭转ONERA M6机翼黏性测试案例的时变升力和阻力系数。

\subsection{7.8.2 时间准确性验证}
144 动态非结构化网格的时间积分 20 40 60 60 100 120 Herations 10' 10' 1-34 g10。 ClO，10' 1 o* 10' 10" 20 40 60 CPU时间图7.6：平滑迭代次数对扭曲ONERA M6翼壳的线性多网格方案整体收敛的影响。在每个网格级别上更平滑。 在图中。 7.8a，根据非线性残差更新的数量，线性多网格方案的速度大约是非线性多网格方案的两倍。 这是意料之中的，因为线性多网格方案每次非线性更新执行的多网格周期是非线性多网格方案的两倍，并且预计这两个方案每个多网格周期会以相同的渐近速率收敛[14]。 然而，就cpu时间而言，由于线性多网格迭代的成本较低，线性多网格方案比非线性方案快得多。 表7.1显示了非线性流解时间步进解所需的CPU时间，使用非线性（FAS）多网格方法、优化的线性多网格方法，以及将网格运动方程收敛到与流方程等效的精度水平（八阶残差还原）所需的CPU时间。 表中显示，线性多网格方法的效率是流动方程的非线性多网格求解器的四倍多，而网格运动方程在给定时间步进中解决非定常流动解问题所需的时间的10\%或更少。 在实践中，不太严格的收敛公差可用于非稳态动态网格流模拟，但这些求解器的相对性能应该保持大致相同。 7.8.2 时间精度验证 检查了符合GCL的二阶（BDF2）和三阶（BDF3）向后差分时间积分方案的时间精度

D。 J。 Mavriplis和Z。 Yang 145 20 40 60 CPU石灰图7.7：线性多网格周期数（每次非线性更新）对扭曲ONERA M6机翼壳的整体收敛的影响。 I,,, 20 40 60 CPU时间 图7.8：非线性多网格和线性多重网格收敛在非线性迭代和扭转ONERA M6翼壳的cpu时间方面的比较。

动态非结构化网格的驯化集成（a）熵等高线（b）提升历史（c）时间收敛图7.9：流动解和测量的时间误差作为二维振荡圆柱体BDF2和BDF3方案的时间步长大小的函数。

\section{7.9 动态网格问题的隐式Runge-Kutta方法}
D。 J。 毛里普利斯和2。 杨14 7维动态网格黏性流动问题。 测试案例由一个单位直径的振荡圆柱体组成，该圆柱体经过y = Asin（ft）给出的垂直位移，其中A = 0.1和f = 0.1\textasciitilde{}。 二维非结构网格由19012个节点和37632个三角形组成。 圆柱体表面的网格根据规定的圆柱体运动位移，而域的外边界保持固定。 网格变形使用弹簧类比方法建模，Gauss-Seidel 多重网格法用于解决每个时间步骤的网格运动方程。 流动的自由流马赫数是M = 0.2，雷诺数是Re = 185。 图7.9a显示了这种情况在对应于t = 65的时间步骤的熵轮廓。 图7.9b显示了作为物理时间函数的升力系数的比较，对于使用不同时间步骤的三阶BDF3方案，显示使用两个最佳时间步骤计算的情况具有良好的一致性。 在图7.9c中，作为物理时间步骤函数的时间误差以对数-对数格式绘制。 时间误差被定义为升力系数在t = 65时，所考虑的解决方案和使用BDF3方案以0.078的较小时间步骤计算的参考解决方案之间的差异。 图中BDF2方案的斜率为2.0，表明设计精度已实现，而BDF3误差曲线的斜率为2.85，接近3的设计精度。 为了实现5 x 10W3的误差，BDF2和BDF3所需的时间步骤分别为0.258和0.512。 考虑到BDF2和BDF3的CPU时间几乎相同，在这种情况下，BDF3方案的效率大约是两倍，较低的时间误差公差会带来更大的好处。 7.9 动态网格问题的隐式Runge-Kutta方法 隐式Runge-Kutta（IRK）方案提供了一种替代方法，以完全A和L稳定性获得高时间精度，这对于高于二阶的BDF方案是不可能的。 原则上，可以构建任何精确度的IRK方案。 然而，IRK方案是多阶段方法，每次需要多个非线性解NMG I LMG 1 MeshMG 159s I 34s I 3.48s表7.1：使用非线性多网格（NMG）方案和线性多网格方案（LMG）方案，在给定时间步中非定常流动残差减少8个数量级所需的cpu时间的比较，以及使用多网格运动方程的等效收敛水平所需的cpu时间。

148动态非结构化网格的时间积分c1=0 C6 = 1 Un+1步骤，使它们在时间步骤上要昂贵得多。 先前对静态网格的研究表明，四阶精确的IRK方案通常在等效精度水平上优于低阶精确的BDF2甚至BDF3方案，尽管IRK方案的每个时间步骤成本更大，因为高阶方案启用的时间步骤更大[3，71。 从方程（7.5）开始，隐式Runge-Kutta方法的一般形式可以写成：S（vqk = (vqn + (at) C akj\textasciitilde{} (\textasciitilde{}j, tj), IC = 1,2,..., (7.16) 3=1 S (VU)n = (VU)n+l + (At)xbjR(U',tJ) (7.17) 其中 s是阶段数，aiJ和b3是方案的屠夫系数，R现在代表对流和黏性项。 继之前的工作[3,7]之后，我们专注于RK方案的ESDIRK类，它代表显式第一阶段，单对角线系数，对角线隐式Runge-Kutta。 六级ESDIRK方案的屠夫表如表7.2所示。j=1 0 0 0 0 0 0 a21 a66 0 0 0 a31 a32 a66 0 0 0 a41 a42 a43 a66 0 0 a51 a52 a53 a54 a66 0 bl b2 b3 b4 b5 a66 bl b2 b2 b3 b4 b5 a66在表7.2中，ck表示时间间隔中的点，[t,t + At] 这些方案的特点是系数表的三角形形式较低，因此在每个单独的阶段都会产生一个隐式求解。 第一阶段是显式的（arc1 = 0），最后一阶段的系数采用akj = bj的形式，因此方程（5）可以简化为un+l = uk=s（7.18）。在这项工作中，我们采用了六阶段、四阶准确的ESDIRK方案，最初由Kennedy和Carpenter提出的[8]表示为IRK64。 对于动态变形的网格，必须确保IRK方案的实现尊重GCL，而不影响原始时间积分方案的高精度。 从表7.2的复杂性来看：具有阶段数的ESDIRK类RK方案的屠夫Tableau，s=6。

D。 J。 Mam'lis和Z。 Yang 149 BDF3方案的GCL条件的推导（见附录），构建符合GCL的IRK方案似乎是一项艰巨的任务。 为此，开发了一种构建符合GCL的时间集成方案的替代方法[16]。 BDF方案的GCL结构（在附录中概述了BDFS），是基于计算整个时间步骤中使用的网格点位置和速度的适当平均值，基于一系列中间网格配置，这恢复了固定网格时间积分方案的设计精度，同时承认均匀流动作为离散解决方案。 对于高阶Runge-Kutta方案，时间残差的评估由c3系数的位置固定。 这不仅包括流量变量的评估，还包括时变的网格坐标和速度的评估，因为这两个量在控制ALE方程中都显示为自变数。 如果这些数量在给定的c3位置进行分析和评估，可以证明生成的IRK方案保留了设计准确性，但未能尊重GCL [16]。 因此，问题变成了如何构建一个Runge-Kutta方案，该方案在指定的正交点评估项x(t)和k(t)，同时遵守GCL。 虽然网格点坐标通常作为时间的函数可用，特别是如果网格变形由一组偏微分方程（泊松方程、线性弹性）的解支配，但网格速度很少作为时间的函数可用，通常通过有限微分网格点坐标函数获得。 在RK正交点评估的网格速度的近似值，为设计一个也服从GCL的IRK方案提供了必要的额外自由度。 在下表中，我们假设网格点坐标是已知的时间函数，所有量都是在cJ正交点上评估的。 Runge-Kutta方案的每个阶段都给出如下：k（VU）k = (VU)”-A\textasciitilde{}CCY\textasciitilde{}\textasciitilde{}R(U\textasciitilde{},X\textasciitilde{},X\textasciitilde{})，k=1,..., s (7.19)，相应的值：x3 = x(t + cJAt)，以及Xj = k(t + c3At). ii(t + c3At)，c3的值由RK方案给出。 虽然网格点坐标可以及时直接在这些位置进行评估，但必须估计面积分速度X。 面积分速度是直接使用的，与网格点速度相反，因为它们明确出现在非稳的ALE残差中，并且是唯一涉及网格点速度的术语。 这些估计值的构建符合GCL。 在方程（7.19）中设置U=常数，我们得到每个RK阶段的GCL条件为：3=1 k Vk - V" = 在C ak3R(X3)，k = 1，...，s（7.20）J=1

150 动态非结构网格的时间积分 R算子对应于离散曲面积分，该积分是作为单个控制体面的总和获得的：面（7.21）由于体积计算也可以表述为控制体面的总和，我们可能需要方程（7.20）在离散积分的每个控制体面上保持：k（Vk-V”，=在C akjX-jE，Ic = 1，...，s（7.22）j=1，其中（Vk-V”）E表示单个控制体边界面扫过的体积，这与基于顶点的有限体积离散的边缘网格相关联， Xi表示jth RK阶段面集成网格速度的未知值。 由于网格坐标（以及控制体积值）是已知的时间函数，方程（7.22）的左侧可以准确评估，因此可以通过求解系统（此处为四级ESDIRK方案）获得Xi未知数：其中显式第一阶段a11 = 0（c.f. 表1）对应于将XL等同于在上一个时间步骤结束时确定的值。 对于ESDIRK方案，akj系数的矩阵是较低的三角形形式，方程（7.23）很容易通过正向替换求解。 为了验证上述IRK方案的动态网格问题的构建，通过强制扭曲ONERA M6翼构建了三维无黏流动测试案例，马赫数为0.755，入射率为0.016度。 在机翼尖端的四分之一弦点规定了周期性俯仰运动，频率降低到0.1628，振幅为2.51度。 翼根固定，在翼尖和根之间规定了俯仰运动的线性变化，从而导致扭转运动。 前面描述的（53961顶点）和图7.5-a所示的全四面体网格用于此情况。 通过求解受固定外部边界条件约束的弹簧类比方程，在每个时间步骤或阶段计算网格变形。 图7.10显示了IRK64方案的计算升力系数与时间的对比度，显示了IRK64方案的良好时间精度，每个周期使用8个时间步骤。 图7.11

\section{7.10 结论}
D。 J。 Mavriplis和Z。 Yang 151 -002 - 10 20 30 40 50 -0.04 ' ' 'I " I' I' ' ' I " I' ' ' I' 在图7.10：升力系数作为扭曲ONERA M6机翼壳的时间函数，由IRK64 GCL兼容方案计算，使用每个周期8,16,32和64个时间步。说明计算的时间误差作为两个方案的时间步的函数，在t=54时，使用高度解析的IRK64解决方案，时间步为0.3375，（对应于每个周期128个时间步）作为参考解决方案。 BDF方案实现的误差曲线斜率为2，而IRK64方案的误差曲线斜率为3.3。 正如静态网格案例[3,7]所观察到的，IRK64方案通常使用比BDF2方案大一个数量级的时间步骤实现等效的精度。 7.10 结论 提出了提高动态变形网格上隐式时间依赖模拟的准确性和效率的技术。 这些包括使用聚集多网格方法在每个隐含时间步骤中解决非线性流动问题，以及加速网格运动方程的收敛。 流方程中纳入了高阶时间积分技术，与低阶方法相比，这使得使用更大的时间步骤来达到等效的精度水平。 然而，这些时间整合程序必须以尊重离散几何守恒律的方式实施，这已被证明对非线性稳定性来说是必要和充分的。 虽然现有技术用于构建符合GCL的BDF方案，但为IRK方案开发了另一种方法。 迄今为止的结果表明，各自的方案在遵守GCL的同时实现了设计精度。 在未来，一个com-

\section{7.11 致谢}
\section{7.12 参考书目}
152动态非结构网格的时间积分图7.11：在t=54时间t=54的三维扭曲ONERA M6翼案的GCL兼容BDF2和IRK64方案的时间步骤大小的函数的时间误差的比较。对于黏性湍流流，必须对动态变形网格问题进行IRK和BDF方案的偏差。 基于以前在静态网格问题方面的经验，我们期望IRK方法能提供更高的效率，特别是当与优化的线性多网格方案相结合，用于求解流动方程和网格运动方程时。 7.11 致谢 这项工作得到了空军研究实验室、鲍尔航空航天公司管理的分包合同、美国宇航局兰利研究中心的赠款和怀俄明州NSF Epscor计划的部分支持。 7.12 参考书目[I] Baker，T。 \& 卡瓦洛，P. A. 变形四面体网格的动态适应。 AIAA 1999-3253，1999。 [2]巴蒂纳，J。 T。 使用非结构化动态网格的非稳态欧拉翼型解决方案。 AIAA杂志，28（8）：1381-1288，1990年。 [31 Bijl, H., Carpenter, M. H.，Vatsa，V. N.，和肯尼迪，C. A. 不稳定不可压缩Navier-Stokes 方程的隐式时间积分方案：层流。 计算物理学杂志，179（1）：313-329，2002年。

D。 J。 Mavriplis和Z。 Yang 153 [4] Farhat, C., Geuzaine, P., \& Crandmont, C. 离散几何守轳定律和ALE方案的非线性稳定性用于解决移动网格上的流动问题。 计算物理学杂志，1741669-694，2001年。 [5] 吉拉德，H。 \& Farhat，C. 关于几何守恒定律对移动网格的流动计算的意义。 应用力学和工程中的计算机方法，190：1467-1482，2000年。 [6] 约翰逊，A. A. \& Tezduyar，T. E。 模拟多个球体落入充满液体的管子中。 应用力学和工程中的计算机方法，134：351-373，1996年。 [7] Jothiprasad，G.，Mavriplis，D. J.，和Caughey，D. A. 非结构网格上非稳定Navier-Stokes 方程的高阶时间积分方案。 计算物理学杂志，191:542-566，2003年。 [8] 肯尼迪，C. A. 和木匠，M. H。 对流-扩散-反应方程的加法Runge-Kutta方案。 应用数值数学，44（1）：139-181，2003年。 [9] Koobus，B。 \& Farhat，C. 在非结构化动态网格上的半离散黏性通量的隐式时间积分上。 国际流体数值方法杂志，29:975-996，1999年。[lo] Koobus，B。 \& Farhat，C. 用于非结构化动态网格上的流计算的二阶时间准确和几何保守的隐式方案。 应用力学和工程中的计算机方法，170\textasciitilde{}103-129，1999年。[ll] Lesoinne，M. \& Farhat，C. 具有移动边界和可变形网格的流动问题的几何守恒定律，以及它们对空气弹性计算的影响。 应用力学和工程中的计算机方法，134：71-90，1996年。 [12] Mavriplis，D。 J。 用于非结构化网格的多网格技术。 在VKI系列讲座中，第1995-02号，VKI-LS。 1995年。 [13] Mavriplis，D。 J。 关于非结构化[14] Mavriplis的收敛加速技术，D。 J。 线性与非线性多网格方法的评估-计算物理学杂志，[15] Mavriplis，D。 J和Venkatakrishnan，V。 混合元素网格上的Navier-Stokes 方程的统一多网格求解器。 计算流体动力学国际杂志，（8）：247-263，1997年。网格。 AIAA 1998-2966，1998年。用于非结构网格求解器的ods。 175:302-325，2002年。

1 54 动态非结构化网格的时间积分 [16] Mavriplis，D. J和Yang，Z。 动态网格上高阶时间准确模拟的离散几何守恒定律的构造。2005年提交给《计算物理学杂志》的论文。 [17]尼尔森，E。 J。 \& 安德森，W. K。 最近在非结构网格上空气动力学设计优化的改进。 AIAA论文2001-0596，1月。 2001年。 [18]托马斯，P。 D。 \& 伦巴第，C. K。 几何守恒律及其在移动网格上的流量计算中的应用。 AIAA杂志，17（10）：1030-1037，1979年。 1191 Venkatakrishnan，V。 和Mavriplis，D。 J。 隐式方法用于非结构化网格上非稳定流动的计算。 计算物理学杂志，127：380-397，1996年。

\subsection{7-A BDF3的几何扎变定律}
附录7-A BDF3的几何收敛定律 DGCL是通过首先写下时间积分方案推导的，在这种情况下是三阶后向差分方案：--3（7.24）k=l，然后将速度设置为恒定值，从而给出：k=l = u，Atn G（n，x）- - u，L;'' x. hdsdt（7.25），其中x和h分别表示网格点速度和面法线，分别表示法线。 由于上述方程的左侧是确切的（通过每个时间步骤的精确体积计算），因此问题包括找到x和A的适当估计值，这些估计值在右侧积分中使用时验证了上述方程。 由于右侧积分被评估为由三角形面组成的控制体面的求和，我们可以检查单个三角形面的积分为：I2 = lril Lx-hdsdt（7.26），其中ds表示面的面积元素。 我们选择将x和x参数化为三角形中三个角顶点值的线性组合：z(t) == \%zZ(t) f (1 - 771 - V2)Z3(t) k(t) zz vlXl(t) f 7725252(t) + (1 - 71 - 772)\&3(t) (7.27) 其中rll E [0,11, vz E [0,1 - 7711, t E [tn,tn+l] 我们进一步将这些角点值作为固定时间步骤的已知值之间的插值，如：Zi(t) = tn+l(t)zl+' + Jn(t)zy + Jn-l-l -(l-Jn-Jn(t)-Jn+l(t))zGn-', i=1,2,3

156 动态非结构化网格X的时间积分：“图7.12：在单个时间步骤内由三角形面扫描的体积的插图，一致性要求以下条件：tn+l（tn+l）= 1，tn（tn+l）= 0，Jn-l（tn+l）= 0，En+l（tn）=0 tn（t”）= 1 tn-l（t”） = 0，其中Ax13 = x3 - x1和Ax23 = 23 - 22，与和= [En+lAxy\$' + tnA2l;; + tn-I AX:;' + (1 - En-1 - En - tn+i)Ax:y2] x [En+lAz;\$' + <n3\$3 + tn-lAx;cl + (1 - （这对应于将这些表达式代入积分中，我们得到：

D。 J。 Mavraplis和Z。 Yang 157 BDF3方案中使用的四个时间级别之间的两点正交）。 使用六点积分规则来近似I；xwkfk 'gk k=l - Atn I2 M I2 = -其中上标k或pk = 1,2,..., 6表示相应正交点的值。 回到方程的左侧（7.25）：J；= Cyn+1vn+1 + anvn + Qn-lVn-l + Qn-ZVn--2（7.29），我们用体积增量重写：这相当于积分：-an-2 fn-' x。 Adsdt p-2

158 动态非结构化网格的时间积分 t" tn-l 根据[ll]，可以准确评估hY+' 1、x.iidsdt、hn-' 1、x.iidsdt和\&n--5 lA x.fidsdt，导致：JA n -Qn+l --（AX？ +AX；+AX：）。 18 (Ax::' x Ax;:' + ZAxc;", 1 AX;\textasciicircum{}, x Axz3 + -AS;;' 1 Ax:;' x Ax;;' + -AxY3 1 (Ax:;' x Ax;;' + 1 Ax;",-' x Ax;;' + -AX:;' 1 x Ax;::'+ 1 AxY3 x AxF3 + ZAxYZ1 x AxT3+ an+ (Ax;-' + Ax:-')。 18 x AX;\textasciicircum{}+ x AX:;')+ -(AX:-' an-2 + AX;-' + AX;-')。 18 x Ax;;'+ x AX;;') 其中Axr = xq+' - xr，i = 1,2,3和m = n,n - 1，n - 2。 为了满足GCL，我们设置了（7.31）Fn - jn A- A，从而产生关系：1 1 1 1 61 = -(1+ -) 61 = -(1- -) 2d3 2v5和

D。 J。 Mavraplis和Z。 因此，Yang 159和正交点的x和x值以特定时间步长位置的x值给出：xP5 = xP6 = “cj n-l - T2 3 3 在三角形面的三个角顶点的j = 1,2,3。 尊重GCL并保持三阶精度的6和x的平均值

160 BDF3方案动态非结构化网格的时间积分最终给出为：-18 9，使用相应的xp值计算。= \%l a，= Tl anPl = c，an-2 = 2，以及5，法线是

\clearpage
\chapter{第8章 电磁学的显式时域有限元方法 Kenneth Morgan、Mohamed El hachemi、Oubay Hassan和Nigel Weatherill}
\section{8.1 介绍}
第8章 电磁学的显式时域有限元方法 Kenneth Morgan'、Mohamed El hachemi'、Oubay Hassan'和Nigel Weatherill' 8.1 介绍 电磁波传播问题的数值模拟涉及现实几何和波频率，构成了重大的计算挑战。 积分方程技术通常是解决波散射问题的最佳方法[5，61，但基于体积的方法对更广泛的应用有吸引力[12]。 考虑到这一点，我们首先描述了一种基于时域体积的3D电磁问题的方法。 采用的方法是麦克斯韦卷曲方程的有限元解算法，使用显式泰勒-加勒金TG2 时间推进和“H”低阶元素[16，17，151。 该配方通过完美传导障碍物来解决散射问题。 四面体元素用于网格与散射器相邻的区域，为了减少存储和CPU时间要求，常规六面体元素在其他地方使用。 这种方法需要添加金字塔元素，以确保一致的混合网格。 截断的远场非反射条件由英国威尔士斯旺西大学斯旺西SA2 8PP的计算土木和计算工程中心处理。

\section{8.2 电磁散射}
\subsection{8.2.1 控制方程}
162 人工完美匹配层（PML）的计算电磁域[2，31。 具有标准空间分辨率的数值示例表明，该方法能够准确地模拟涉及数十波长距离的波传播的问题。 目前感兴趣的许多问题需要模拟更远距离的波传播。 在这种情况下，无论是在时间和空间上，高阶演算法都是首选，因为它们承诺降低所需的分散化水准。 在这种情况下，通过具有高阶Runge-Kutta 时间推进的基本显式任意顺序不连续Galerkin方法求解Maxwell卷曲方程已被证明对二维和三维应用特别有效[26，11，121。 由于我们的方法是连续的Galerkin程序，一个有趣的替代方法是明确的Fekete点光谱元素方法，该方法已被应用于二维弹性波方程的解[24，131。 我们对显式方案的开发进行了初步调查，将这些元素与高阶Taylor-Galerkin 时间推进程序结合使用。 包括一维波传播和二维电磁波散射的数值结果。 8.2 电磁散射 8.2.1 控制方程 考虑模拟入射单频平面电磁波的散射。 假设这个入射波是由位于远场的源产生的。 虽然这种方法更笼统，但就本文而言，我们将假设散射体是一个被自由空间包围的完美电导体（PEC）。 在自由空间中，Maxwell的卷曲方程可以使用求和惯例以无量值成分形式（8.1）dH* aE,. dE* 3- dH,. - Ejke- 3- - Ejke- at ax k ax k axk中，其中下标j、Ic、C可以取值1,2,3，交替张量由Ejke和E* = (JT;,E\textasciitilde{},E:)\textasciitilde{}和H* = (H7,HZ,Hl)T分别表示电场和磁场强度向量表示。 为了计算方便起见，这些字段根据E* = EZ\textasciitilde{}C + E H* = H\textasciitilde{}\textasciitilde{}\textasciitilde{} + H（8.2）分为入射和散射成分，入射场由问题定义指定，Einc和HanC。 通过将方程（8.2）代入方程（8.1）来实现散射场E和H的公式。 由此产生的方程可能

\subsection{8.2.2 边界条件}
\subsection{完美的导体表面}
\subsection{远场边界和完美匹配的层}
K。 Morgan et a1 163组合生成单向量方程，其中，现在，u=[z]这个线性方程也可以以Fk = AkU的形式书写，每个矩阵Ak中的条目k = 1,2,3，是常数。 8.2.2 边界条件 完美电导体表面 在完美电导体（PEC）的表面，边界条件是电场E*的切向分量应该消失。 这一要求可以表示为n A E = -n A EinC (8.7)，其中A表示向量乘积，n是表面的单位法向量。 远场边界和完美匹配的层 在尝试数值模拟之前，无限的物理域被截断。 应用于截断的外部计算边界的条件是要求散射场应仅由出波组成。 通过在截断域的外部添加一个完美匹配的层（PML）[2]来实现这种条件的建模。 把截断的外边界作为正六面体的表面很方便。 采用的公式遵循了Bonnet和Poupaud 131的工作，因此，通过适当重新定义数量U和Ak，并添加一个源项，方程形式（8.5）在PML中仍然有效。

\section{8.3 网格生成}
\section{8.4 数值解算法}
\subsection{8.4.1 时间离散化}
\subsection{8.4.2 空间中的离散化}
计算电磁学 8.3 网格生成 模拟的起点是构建一个有效的网格。 曲面网格生成器用于将计算域的曲面离散化为三角形元素，元素大小变化受用户定义的网格控制函数的要求控制[20]。 然后，可以使用自动非结构网格生成器来离散计算域fl，使用四面体元素[28，10，211。 虽然四面体元素提供了很大的几何灵活性，但六面体元素更适合填充空间。 一个适当的折合是采用混合网格，其中无结构化的四面体元素在散射器附近使用，而域的外部部分，包括PML，用结构化六面体网格表示。 金字塔元素的过渡层用于以一致的方式连接这两个网格区域[9]。 8.4 数值解算法 Taylor-Galerkin TG2算法[16, 17, 151用于在四面体、金字塔和六面体的一致混合网格上获得方程（8.5）的近似解。 这是Lax-Wendroff方法的Galerkin公式，采用一步有限元实现[14，81。 8.4.1 时间离散化 Lax-Wendroff方法基于截断的泰勒级数展开dU At2 a2U /imi AU=At- +-- at 2 at2，其中AU = U\{mfl) - UIrn），上标\{m\}表示时间t = t的评估，tm+l = t，+ At。 时间导数被空间导数取代，使用方程（8.5），产生时间推进算法。在域的空间离散化之后，该方程的解是从近似变异公式中获得的。 8.4.2 空间中的离散化 采用最低阶连续元素，节点位于元素顶点。 在网格中的每个元素E上，解决方案drn）是inter-

\subsection{8.4.3 计算细节}
K。 Morgan et a1 165在节点值之间，如（8.10）在这里，NJ表示与节点J相关的有限元形状函数，求和延伸到元素E的所有节点J。 时间级tm+l的解是根据方程（8.9）的Galerkin弱变分公式[18]获得的。 在一般内部节点I，结果方程的形式为（8.11）C MI JAU J = 7E\$m' JEI，其中方程（8.11）的总和延伸到节点I和那些通过网格中的元素边缘连接到节点I的节点J，MIJ表示一致的有限元质量矩阵中的条目。 所有积分都被精确地评估。 在位于PEC表面上的节点上，通过本方程中出现的附加边界积分项，通过采用垂直于边界方向的特征分解来应用边界条件[15，221。 可以定义一个对角线块状质量矩阵，第I行的条目MF = CMIJ（8.13）。 由于一致的质量矩阵是对角线主导的，然后方程（8.11）可以通过合并质量矩阵或通过显式迭代[7]（\textasciitilde{}u\textasciitilde{}er+l - AU? \textasciitilde{}。） = - C \textasciitilde{}\textasciitilde{}\textasciitilde{}\textasciitilde{}u.Jter (8.15) JEI 其中 iter = 0,1,2,... 和 A@ = 0. 8.4.3 计算细节 稳定性 该算法是稳定的，前提是Courant 数，C是[18]（8.16）

166 计算电磁学 在这里，hE是元素E的代表长度，最小长度覆盖所有元素和网格。 当使用块状质量近似时，可以使用值/3 = 1，而值/3 = l/\&适用于具有显式迭代的一致质量形式。 数值性能 为了证明该算法的数值性能，请考虑将连续正弦波传播到最初静止的一维半无限区域的问题。 对于X = 15，规定了1.00的目标元素长度h\textasciitilde{}，每个元素的实际长度在0.50 5 hE 5 1.25范围内随机生成。 这是为了复制自动非结构网格发电机可能在三维中产生的网格间距分布类型。 波传播，直到其前缘移动了25倍的距离。 此时的数值解，位于距离波的起始点20倍的点附近，与图中的确切分布进行比较。 8.1. 图中显示了块状质量形式的使用，接近其稳定极限，C = 0.95。 8.l（a）。 数值解的振幅和相位精度很差。 当使用具有显式迭代的一致质量形式接近其稳定性极限时，数值性能会显著提高，C = 0.55。 这如图所示。 8.l（b）。 图8.l（c）和（d）表明，在实际限制范围内减少时间步进大小，可以小幅改善使用质量算法生成的结果，但进一步提高一致质量的结果准确性。 这些结果证实了理论分析[8]，该分析显示了均匀网格上一致质量有限元方案的更高精度。 他们还表明，在这种离散化水平下，该方案可以预期在涉及数十波长距离上传播的模拟中令人满意。 雷达截面的计算 在实际的散射模拟中，解决方案会及时推进，直到达到稳定的周期性条件。 构建了一个完全包围散射器的闭合表面，并计算了一个进一步的周期，在此期间，监测位于该表面的节点的溶液变化。 记录了这些节点上散射电场和磁场成分的振径和相位，通过采用近场到远场变换[l]，这些值用于雷达横截面（RCS）的计算[16]。

\section{8.5 数字示例}
K。 Morgan et a1 167 图8.1：使用TG2算法和线性元素的非均匀网格传播连续正弦波，比较距离波起始点20倍的点附近的精确（-）和数值（0）解：（a）块状质量，C = 0.95；（b）一致质量，C = 0.55；（c）块状质量，C = 0.15；（d）一致质量，C = 0.15。 平行化离子 对于这个解算法，所需的网格大小随着散射器的电气长度的增加而迅速增长。 由此产生的模拟需要使用大量的计算资源，因此，实现已经完全并行化[29]。 8.5 数值示例包括两个示例，以证明该方法的计算性能。 在这两种情况下，提供给自动网格发生器的网格控制功能要求在散射器表面对应于每波长15个元素的网格间距，线性变化为PML内表面每个波长10个元素的网格间距。 PML区有一个波长厚，使用10层元素离散化。 PML位于距离散射器一个波长的最小距离。

\subsection{8.5.1 PEC球体}
\subsection{8.5.2 PEC杏仁}
168 计算电磁学图8.2：直径为15X的PEC球的散射：四面体和六面体网格区域之间的界面。 入射波向x1方向传播，计算是在带有R14000处理器的多处理器ONYX 3800上进行的。 通常发现，仅使用由四面体组成的等效网格的模拟需要大约两倍的CPU时间和大约75\%的存储空间。 8.5.1 PEC球体 第一个例子涉及直径为15倍的完美导电球体对平面单频入射波进行散射。 这用于验证实现，因为这个问题有一个精确的系列解决方案。 使用的网格有192万个四面体、614万个六面体和15万个金字塔。 四面体和六面体网格区域之间的界面视图如图所示。 8.2. 网格中的节点总数为662万个。 该解决方案是通过20个入射波周期推进的，这需要在16个处理器上每周期20分钟的CPU时间。 图中显示了所有视角下计算的和精确的RCS分布之间的出色一致性。 8.3. 8.5.2 PEC杏仁本示例演示了在预测模式下使用该算法，并涉及PEC杏仁配置的散射。 杏仁的主轴位于x1方向，波浪直接进入杏仁的尖端。 波频率是杏仁的电长为21X。 使用的网格由380万个六面体、135万个四面体和0.053万个金字塔组成。 表面三角测量的视图和

\section{8.6 处理更大的电动散射器}
K。 Morgan et a1 cam 16d Ran T - -10 I 图8.3：直径15X的PEC球散射：精确和计算的RCS分布的比较。网格的切割如图所示。 8.4. 节点总数为420万个，模拟通过60个入射波周期进行。 E2在表面计算的等高线分布如图所示。 8.5. 图中显示了使用混合网格方法生成的计算RCS分布与仅使用等效四面体网格获得的RCS分布之间的比较。 8.6. 使用16个处理器，每个周期需要18个CPU分钟的模拟。 8.6 处理电气上更大的散射器 从这些三维示例中可以看出，当散射器的电气长度增加时，所需的网格将很快变得非常大。 当使用相同的网格分辨率尝试在更多周期内传播波时，污染误差的影响加剧了这一困难。 这如图所示。 8.7，它比较了图中考虑的一维连续正弦波传播示例的exact和计算分布。 8.1. 在这种情况下，计算以与以前相同的网格分辨率进行。 然而，当波的前缘移动了165倍的距离时，解现在被绘制在距离波开始点150倍的点附近。 众所周知，使用高阶方案可以缓解这些困难，但为非结构化网格构建显式高阶解算法并非易事。 对于实际的波传播问题，

I 70 计算电磁学图8.4电长21X的PEC杏仁的散射：表面离散的细节和尖端附近的网格切割。使用不连续的Galerkin方法，在空间中具有高多项式顺序和高阶Runge-Kutta 时间推进，已被证明特别有效[26，11，121对于一阶方程组（8.3）。 在二维中，还提出了二阶波方程的连续高阶显式加勒金方案，基于将经典光谱元素方法扩展到三角形，使用为精确近似而不是积分而设计的采样点[24，131。 在试图改进已提出的解决方案算法时，我们采用这种连续Galerkin方法的想法来产生高阶显式泰勒-图8.5：电长20的PEC杏仁的散射：表面Ez的计算轮廓。

\subsection{8.6.1 更高阶的Taylor-Galerkin 时间推进方案}
K。 Morgan et a1 171 5-10 8 -20 -30 0 30 60 90 120 150 180 210 240 270 300 330 -50 I " " " I 图 8.6：电长21X的PEC杏仁散射：仅使用混合网格和四面体的非结构网格计算的RCS分布的比较。 一阶方程组的Galerkin算法（8.5）。 8.6.1 高阶泰勒-加勒金时间推进方案 高阶泰勒-加勒金时间推进方案，如TG3和TG4，可以通过采用泰勒级数扩展到所需顺序来开发。 例如，三阶方案TG3使用泰勒展开（8.17）At3 a3U 6 at3 +--当时间导数被空间导数取代时，使用原始微分方程，Galerkin有限元近似解需要求解形式（8.18）的方程。该方程的构造对PML区域非常相关，其中源项存在于控制方程中。 出于这个原因，多步显式高阶泰勒-加勒金方法，如TG3-2S和TG4-2S，是首选的。 三阶算法TG3-2S使用

172 计算电磁学图8.7：使用TG2算法和非均匀的线性元素网格传播连续正弦波，比较位于距离波起始点150倍的点附近的精确（-）和数值（0）解：（a）聚集质量，C = 0.95；（b）一致质量，C = 0.55；（c）聚集质量，C = 0.15；（d）一致质量，C = 0.15。两步程序d2U tm）+aAt2 - 3 at at2 dU I（m\} At2 a2U nu =at- +-- dt 2 at2 at2 =-（8.19），其中参数a只影响时间膨胀中四阶项的系数。 当时间导数被空间导数取代时，使用原始微分方程，Galerkin有限元近似

\subsection{8.6.2 高阶空间离散化}
K。 Morgan et a1 i r3解需要JEI（8.20）形式的方程解，右侧向量7Zl和722很容易定义，标准mas矩阵出现在两个方程的左侧。 如果将最低阶连续元素用于空间离散化，则可以通过显式迭代再次获得解决方案。 图中显示了具有线性元素和cy = 1/9的TG3-2S算法的数值性能。 8.8，它显示了图中考虑的一维连续正弦波传播示例的精确和计算的分布。 8.7。 由于这些模拟是在同一网格上进行的，比较图。 8.7与图。 8.8表明，仅仅提高时间准确性不足以消除原始方法的感知缺陷。 8.6.2 高阶空间离散化 一维元素 首先考虑基于使用TG2 时间推进的解算法的构建，在一维中具有p阶的有限元。 像光谱元素方法[19]一样，可以通过在高斯-洛巴托-传奇（GLL）正交点的每个元素上定位p+ 1节点来构建p阶元素。 这个正交规则对于2p-1阶的多项式是精确的。 然后，我们可以使用一种方法，即TG2方程（8.11）中的所有项都由足够准确的高斯正交规则进行评估，并通过GLL正交法评估质量矩阵中的条目来定义对角线聚合质量矩阵。 TG2方程组可以大致使用这个块状质量矩阵或使用方程的显式迭代方案（8.15）来求解。 图中显示了具有一维有序立方元素的TG2算法的数值性能。 8.9。 这显示了涉及连续正弦波传播的例子的精确和计算解，之前在图中考虑过。 8.7和8.8。 在这种情况下，对于X = 15，采用3.00的目标元素长度h\textasciitilde{}，每个元素的实际长度在1.50 5 h\textasciitilde{}5 3.75范围内随机生成。 方程（8.16）的定义再次被用于Courant 数 C的计算。 图的比较。 8.7和图。 8.9表明，通过增加元素顺序，在本例中，节点数量大致相同，可以实现解决方案准确性的显著提高。

1 74 计算电磁学（4（4 图8.8：使用TG3-2S算法和线性元素的非均匀网格进行连续正弦波的传播，比较距离波起始点150倍的点附近的精确（-）和数值（0）解：（a）块状质量，C = 0.80；（b）一致质量，C = 0.80；（c）块状质量，C = 0.15；（d）一致质量，C = 0.15。 三角形元素 经典的光谱元素方法没有自然扩展到三角形元素，因为一般来说，不知道三角形上是否存在GLL点。 为了将上述方法扩展到两个维度，我们可以使用p阶的三角形元素，（p + 1）（p + 2）/2节点位于Fekete点[30，24，131。 这些点是根据其插值和近似属性来选择的，而不是为了正交。 存在用于确定Fekete点的算法[25]，它们的位置已被表格化[23，271。 此外，对于p的实际值，可以根据这些采样点构建一个正交规则，该规则具有正权重，并将任何p阶多项式集成到主三角元素上[25]。 然后，方法是通过一个足够的高斯正交规则来评估TG2方程（8.11）中的所有项，并通过使用Fekete点采样的正交规则评估质量矩阵中的条目来定义一个对角线块质量矩阵。 同样，TG2方程组系统可以使用这个聚合质量矩阵或使用显式

K。 Morgan et a1 1 75（a）（b）图8.9：使用TG2算法和非均匀的三次元素网格传播连续正弦波，比较距离波开始点150X的点附近的精确（-）和数值（0）解：（a）块状质量，C = 0.045；（b）一致质量，C = 0.045。 图8.10：使用TG2算法传播连续平面正弦波：p=3阶三角形元素网格的细节。方程的迭代方案（8.15）。 三角形边缘的点正是一维GLL点[4]。 这意味着这些三角形元素与经典的光谱四边形元素自然一致，混合网格实现应该再次成为可能。 为了研究TG2算法在p = 3三角形元素的网格上的性能，请再次考虑连续正弦波传播到最初静止的半无限区域的问题。 这本质上是一个一维问题，但在这里使用三角形元素的网格来解决。 对于X = 15，规定了目标元素长度h\textasciitilde{}为3。 图中显示了所用网格的细节。 8.10. 波传播，直到其前缘移动了15倍的距离。 此时的数值解，用C = 0.045计算，在距离波开始点11X的点附近，与图中的确切分布进行比较。 8.11。 很明显，对于二维模拟，算法的一致质量形式更优越。

\section{8.7 结论}
1 76 计算电磁学图8.11：使用TG2算法和p = 3阶三角元素的平面波传播，比较位于距离波起始点1lX的点附近的精确（-）和数值（0）解：（a）块状质量，C = 0.045；（b）一致质量，C = 0.045。 数字示例包括两个示例，以演示基于TG2 时间推进和p = 3阶三角形元素的算法的使用。 如前例所示，PML区域的厚度为1波长，距离散射器最小1波长。 在这两种情况下，自动三角网格发生器被要求在散射器表面产生每个波长5个元素的网格，在PML的内表面线性变化到每个波长3个元素。 第一个例子是直径为1OOX的完美导电气缸对平面单频入射波的散射。 这个问题可用的精确系列解决方案可以验证实现。 图中显示了所有视角下计算的和精确的RCS分布之间的出色一致性。 8.12。 第二个例子涉及由一个完美导电的U形腔体散射，旋转90度。 腔体的长度是8倍。 Ez和H3等高线的计算分布详情如图所示。 8.13。 图中显示了计算的RCS分布与使用基于TG2 时间推进的算法和等效的线性三角形元素网格计算的分布之间的比较。 8.14. 8.7 结论 提出了一个能够模拟计算电磁学中三维问题的数值程序。 计算机-

K。 摩根等人图8.12：使用TG2算法和p = 3阶三角形元素网格的直径lOOX的PEC圆柱体进行散射：精确（- -）和计算（- -）RCS分布之间的比较。使用自动生成的混合非结构网格离散域，节点位于元素顶点，并使用Taylor-Galerkin TG2 时间推进方案。 解决方案程序是平行化的，并展示了该方法在电磁散射领域的应用。 还对制定更高级解决方案程序的可行性进行了初步调查。 使用TG2 时间推进和p = 3阶的三角形元素解决了二维散射问题。 这些元素的节点位于Fekete点，近似质量矩阵的构建可以通过显式迭代实现解决方案。 在时间和空间中使用高阶离散化应该会进一步扩展能力。 例如，图中说明了具有一维立方和五元元素的TG3-2S算法的数值性能。 8.15。 这显示了涉及连续正弦波传播的例子的精确和计算解，之前在图中考虑过。 8.7、8.8和8.9。 在这种情况下，对于五元元素和X = 15，采用5.00的目标元素长度h\&，每个元素的实际长度在2.50 5 h\& 5 6.25范围内随机生成。 与图的比较。 8.7和图。 8.9表明，在本例中，节点数量大致相同，可以通过这种方式进一步提高解决方案性能。 使用高阶离散化对所需节点数量和可能的三维扩展的影响仍有待研究。

\section{8.8 参考书目}
1 78 计算电磁学（a）（b）图8.13：使用TG2算法和p = 3阶的三角形元素网格通过电长为8X的PEC腔进行散射：（a）E2的轮廓；（b）H3的轮廓。 8.8 参考书目[l] Balanis，C. A.，高级工程电磁学，威利，纽约，1989年。 [2] 贝雷nger，J. P.，有限差分计算机代码中自由空间模拟的完美匹配层，计算物理学杂志114，1994年，pp. 185-200。 [3] 邦内特，F。 \& Poupaud，F.，三角网格的时间有限体积方案的Berenger吸收边界条件，应用数值数学25，1997年，pp. 333-354。 [4]博斯，L.，泰勒，M. A. 和温盖特，B。 A.，张量乘积高斯-洛巴托点是立方体的费克特点，计算数学70，2000年，pp. 1543-1547。 [5] Chew，W.C.，计算电磁学：平滑场与振荡场的物理学，英国皇家学会哲学交易：数学、物理和工程科学，362，2004年，pp. 579-602。 [6] Darve，E。 \& Had，P.，Maxwell方程在所有Fkequencies上稳定的快速多极方法，英国皇家学会哲学交易：数学、物理和工程科学362，2004年，pp. 603-628。 [7] Donka，J. \& Giuliani，S.，基于线性有限元的显式方案生成高阶求解对流算子的简单方法，国际流体数值方法杂志1，1981年，pp. 63-79。

K。 Morgan et a1 -20 \textasciitilde{}30 I,I, -200 -150 -100 -50 0 50 100 150 m 图8.14：使用TG2算法通过电长8X的PEC腔散射：比较使用p = 3（-）阶三角形元素和p = 1（-+-）的三角形元素计算的RCS分布。 [8] Don\&，J. \& Huerta，A.，流动问题的有限元方法，Wiley，Chichester。 2003年。 [9] El hachemi，M.，Hassan，O.，Morgan，K. \& Weatherill，N。 P.，使用H1 有限元法和混合非结构化网格的3D时间主要计算电磁学”，计算流体动力学杂志13，2004年，pp. 391-402. [lo] 乔治，P. L.，自动网格生成。 有限元方法的应用，Wiley，Chichester，1991年。[ll] Hesthaven，J. S。 \& Warburton，T.，非结构化网格上的高阶节点方法。 我。 麦克斯韦方程的时间域解，计算物理学杂志181，2002年，pp. 1-34。 [12] Hesthaven，J. S。 \& Warburton，T.，Maxwell 特征值问题的高阶节点不连续Galerkin方法，皇家学会的哲学舞蹈：数学、物理和工程科学362，2004年，pp. 493-524。 [13] Komatitsch, D., Martin, R., Tromp, J., Taylor, M. A. 和温盖特，B。 A.，使用三角形和四边形的光谱元素方法在二维弹性介质中的波传播，计算声学杂志9，2001年，pp. 703-718。

180 计算电磁学图8.15：使用TG3-2S算法传播连续正弦波，比较位于距离波起始点150倍的点附近的精确（-）和数值（0）解：（a）一致质量和p = 3，C = 0.80阶元素的非均匀网格；（b）一致质量和p = 5阶元素的非均匀网格，C = 0.15的元素网格，顺序为p = 5，C = 0.15的一致质量和非均匀网格。 [14] Lohner，R.，Morgan，K. \& Zienkiewicz，0。 C.，可压缩高速流动的自适应有限元素程序，应用力学和工程中的计算机方法52，1985年，pp. 441-465。 [15] 摩根，K.，哈桑，O.，佩格，N。 E。 \& Weatherill，N。 P.，使用非结构化网格在分段均匀介质中模拟电磁散射，计算力学25，2000年，pp. 438-447。 [16] 摩根，K.，哈桑，0。 \& Peraire，J.，解时域麦克斯韦方程的非结构化网格演算法，国际流体数值方法杂志19，1994年，pp. 849-863。 [I71 摩根，K.，哈桑，0。 \& Peraire，J.，模拟分段均匀介质中电磁散射的时域非结构化网格方法，应用力学和工程中的计算机方法234，1996年，pp. 17-36。 [18]摩根，K。 \& Peraire，J.，流体力学的非结构化网格有限元方法，物理学进展报告62，1998年，pp. 569-638。 [I91 Patera，A.，流体力学的光谱元素方法：通道扩展中的层流，计算物理学杂志54，1984年，pp. 468-488。 [20] Peir6，J.，表面网格生成，在J中。 F。 Thompson，B。 K。 索尼和N。 P。 Weatherill（编辑），网格生成手册，CRC出版社，博卡拉顿，1998年，pp. 19.1-19.20。

K。 Morgan et a1 181 [21] Peraire, J., Peir6, J. \& Morgan，K.，前进前线网格生成，在J。 F。 Thompson，B。 K。 索尼和N。 P。 Weatherill（编辑），网格生成手册，CRC出版社，博卡拉顿，17.1-17.22，1998年。 [22] Shankar，V.，Hall，W. F.，Mohammadian，A. \& Rowell，C.，使用CFD技术的时域电磁学理论和应用，课程笔记，加州大学戴维斯分校，1993年。 [23]泰勒，M。 A.，私人通信，2004年。 [24]泰勒，M。 A. 和温盖特，B。 A.，非四边元的广义对角质量矩阵光谱元素方法，应用数值数学33，2000年，pp. 259-265。 [25]泰勒，M. A.，温盖特，B. A. 和文森特，R. E.，在三角形中计算Fekete点的算法，SIAM数值分析杂志38，2000年，pp. 1701-1720。 [26] Warburton，T.，将不连续Galerkin方法应用于使用非结构化多态hpFinite元素的Maxwell方程，在B中。 科克本，G。 E。 卡尼亚达基斯和C。 W。 Shu（编辑），不连续的Galerkin方法：理论计算和应用。 计算科学与工程讲义11，Springer Verlag，纽约，2000年，pp. 451-458。 [27] Warburton, T., Private Communication, 2004. [a81 Weatherill, N. P。 \& Hassan，O.，具有自动点创建和强加边界约束的高效三维Delaunay三角测量，国际工程数值方法杂志37，1994年，pp. 2005-2040。 [29] Weatherill，N。 P.，Hassan，O.，Morgan，K.，Jones，J. W。 \& Larwood，B.，在非结构网格上实现完全并行的航空航天模拟，工程计算18，2001年，pp. 347-375。 [30]温盖特，B。 A. \& 泰勒，M. A，三角光谱元素的自然功能空间，技术报告98-1 711，新墨西哥州洛斯阿拉莫斯，1998年。

此页面故意留空

\clearpage
\chapter{第9章 估计CFD解决方案T中的网格诱导误差。 I-P。 什叶}
\section{9.1 介绍}
第9章 估计CFD解决方案T中的网格诱导误差。 I-P。 Shih' 9.1 导言 为了仅基于计算流体动力学（CFD）的热流体系统和设备的设计和分析，必须为每个CFD解决方案提供误差界限。 只有对每个CFD解决方案都有信心，设计决策才能建立在坚实的基础之上。 CFD解决方案有四个主要错误来源。 第一个是物理建模不足，没有通过第一原理（如湍流和燃烧）来解决。 第二个是关于边界条件的信息不足（例如，流入和流出边界的信息）。 第三种是非物理效应，如数值扩散、色散和其他虚假模式，当控制偏微分方程（PDE）通过有限微分有限体积或有限元素方法在离散域上离散化为代数方程时。 第四个错误来源来自网格。 对于CFD的大多数用户，特别是那些无法访问源代码的用户，网格和时间步骤大小是用户可以完全控制的解决方案程序中唯一部分。 网格的重要性怎么强调都不为过。 网格必须以足够的细节表示几何形状，并使控制PDE的代数模拟能够解决相关的Aow物理。 对于复杂的稳定和不稳定的三维问题，网格中可以使用的网格点或单元格的数量受到可用的计算资源或需要实际周转的限制'爱荷华州立大学航空航天工程系，爱荷华州艾姆斯，爱荷华州5001 1

\section{9.2 方法分类}
184 计算解决方案时估计网格诱导误差时间。 由于网格点或单元格的数量受到限制，准确性要求将网格点放置在最需要它们解决几何和流动物理的区域（例如，通过r-或h-细化）。 不幸的是，这种不均匀的分布可能会产生所谓的劣质细胞，这可能会在计算解决方案中诱发相当大的误差。 对于大多数工程感兴趣的问题，情况变得更糟，因为生成独立于电网的解决方案通常不可行，因此由于电池质量差和分辨率不足而出现错误。 因此，问题是（1）在计算解决方案中，人们可以相信什么？ （2）来自质量差或不够精细的网格的解决方案的误差限制是什么？ （3）如何改进网格或网格以提高准确性？ 本文描述了最近的工作，该工作试图通过提供估计网格诱导误差的方法来回答这些问题。 9.2 方法分类 一些研究人员已经开发并评估了量化由质量差或不够精细的网格网格引起的PDE解决方案中的错误的方法。 Roache [26]将所有这些方法分为两类：基于多个网格的方法和基于单个网格的方法。 基于多个网格的方法[6, 7, 331需要在一系列越来越精细的网格（至少3个）上生成解决方案。 这些方法可以对错误做出明确的陈述，但在要求更精细的网格上提供解决方案时成本更高，这对现实的工程问题来说可能是禁止的。 如果理查德森的外推被用作判断网格收敛的标准，那么网格必须足够精细，截断的泰勒级数才能有界，然后才能产生有意义的结果。 单网格误差估计方法可以分为代数或PDE。 代数方法假设网格点或单元格的误差是相关网格和解的函数。 在代数方法上做了很多工作，主要与解自适应网格加密有关（例如，见Carey [5]和[2，18-20，28，30，361）。 这些代数误差估计器，也称为网格质量度量，通常只考虑标量场的梯度。 Shih等人[27]和Gu等人[16，171提出了网格质量措施，这些措施还考虑了流场的向量和张量性质，并将解与网格中每个单元格的几何形状和大小联系起来。 单网格PDE误差估计器由Babuska等人首次提出[3，41]，认识到一旦产生的误差可以通过对流和扩散传输到流场的其他部分。 因此，需要一个传输方程来理解误差的产生和演变。 Babuska等人提出了一种推导误差传输方程的方法[3，41用于有限元方法，Ferziger [lo]、Van Straalen等人[29]和Zhang等人[34，351适用于有限差和有限体积方法。 为了说明他们的方法，考虑一个在因变量U上运行的微分算子L [21]：

T。 I-P。 Shih 185 L(U) = f (9.1)，其中f是非同质项。 如果U是一个近似解，那么将其替换为等式。 （9.1）将产生一个剩余的R，因为它不能满足Eq。 （9.1）；即，L（U，）-f = R（9.2）如果等式。 （9.1）和（9.2）是线性的或线性化的，然后减去等式。 （9.2）来自Eq。 （9.1）产生以下误差的传输方程：L（e）=-R或Le=-R（9.3），其中e = U-U，（9.4）是解误差。 由于解中的误差是根据PDE的精确解定义的，因此Eqs给出的误差传输方程。 （9.3）和（9.4）可以解释来自网格的误差和来自数值方案的误差。 实现等式的困难。 （9.3）和（9.4）在对残差进行建模，R。 如果R建模正确，那么Eqs。 （9.3）和（9.4）可以在网格中的每个网格点或单元格上提供误差。 Qin和Shih [21]指出Eqs。 （9.3）和（9.4）仅适用于有限元和光谱等有限扩展方法，但不适用于有限差（FD）和有限体积（FV）等搭配型方法。 这是因为对于FD和FV方法，Eq中的微分算子L。 （9.2）被一个离散运算符取代，这样减去方程。 （9.2）来自Eq。 （9.1）不会产生等式。 （9.3）和（9.4）。 Roache[26]也注意到了这种不一致。 为了纠正这种不一致，Roache [26]建议将网格点或单元格上的所有FD和FV解决方案变成连续函数。 然而，蟑螂[26]认识到，由此产生的连续函数不是唯一的，将取决于使用的插值函数。 由于FD、FV和有限元和光谱等有限展开方法可以通过适当的加权函数[111]在加权残差法下以积分形式统一，因此最通用的误差传输方程应以积分形式统一。 方程给出的微分形式。 （9.3）和（9.4）只有当加权函数是连续的（例如，有限元或光谱方法，但不是FD或FV方法）并且生成的解是真实的（例如，没有弱解）时，才能从积分形式中恢复。 Giles等人[12-151通过伴随变量公式对有限元和FV方法采用积分方法来估计升力和阻力等积分量的误差，并确定将误差最小化的最佳网格点分布

\section{9.3 离散误差传输方程概述}
186 估计这些整数中的网格诱导误差。 Qin \& Shih [21-251采取了不同的方法，该方法也可用于FD、FV和有限元方法，尽管他们只将其应用于FD和FV方法。 它们直接从FDFV方程中推导出离散或不连续的误差传输方程（DETE），而不考虑原始的PDE。 通过忽略原始PDE，他们的方法只提供网格诱导的误差，但他们的方法确实为FV方法提供了每个单元格和FD方法的每个网格点的误差估计。 在本文的其余部分，描述并回顾了Qin \& Shih [21-251）在FD和FV方法中估计网格诱导误差的方法。 首先，引入了DETE概念。 接下来，推导出用FV方法求解的Euler 方程的DETE。 随后是数值实验，这些实验显示了DETE在估计尤拉和Navier-Stokes 方程的FV解中网格诱导误差的有用性。 9.3 离散误差传输方程概述 Qin \& Shih提出的DETE概念[21-251进行如下。 假设FDFV方法，方程中的微分算子L。 （9.1）被离散域上的离散运算元LD所取代，因此FDFV方程对应于等式。 （9.1）是Lo（Ui，h.k）=0（9.5）在上述方程中，Ui是网格点或单元格i的解；h表示网格间距或单元格大小；k表示时间步骤大小。 当h和k被细化为无条件一致的FDFV方程时，解Ui接近网格独立解Ug,i，根据Lax等价定理，如果PDE定位良好且线性，并且FDFV解在数值上是稳定的，则该解也是FDFV方程近似的PDE的精确解。 然而，由于CFD中感兴趣的PDE不是线性的，我们使用限制性较小的术语网格独立解。 与Eq。 （9.3，网格独立解由h给出，其中h是网格间距或单元格大小，k是获得网格独立解所需的时间步大小。 来自Eqs。 （9.5）和（9.6），可以推导出两个不同的DETE。 一个是基于粗网格间距h，在那里得到解。 另一个是基于网格独立网格间距h，。 Qin \& Shih [21]表明，只有基于粗网格间距的那个才有意义，因为将U,i插值到h上是良好的（即所有插值都给出相似的结果），而将粗网格解Ui插值到h上是错位的（即不同的插值可以产生非常不同的结果）。 因此，Qin \& Shih [21]将网格独立解Ug,i插入到Eq中。 （9.5）与粗网（h，k），这给了

T。 I-P。 Shih 187 R,i = LD(ug,i, h, k) (9.7) 如果LD是线性的或不里化的，那么减法方程。 （9.5）来自Eq。 （9.7）给出了离散误差传输方程（DETE），即LD（q，h，k）= Rg，i或LD（h，k）ei = R，i（9.8）e.-U I - g，l.-U。 I（9.9）由于Ug，i在严格意义上不是唯一的（因为它取决于在解决方案被视为网格独立之前不需要更改的有效数字数量），并且可能与确切解决方案不同，因为FDFV方法允许的伪模式，Eqs。 （9.8）和（9.9）只能解释网格诱导的误差，但不能解释来自FDFV 离散化的误差。 这是当前方法的缺陷。 来自Eqs。 （9.7）到（9.9），可以进行以下观察：第一个是，残差可以通过Fq精确计算。 （9.7）虽然需要一个独立于电网的解决方案。 了解确切的残差有助于指导残差的建模及其验证。 第二个是，一旦知道残差（例如，通过建模），那么Eq. （9.8）和（9.9）可用于估计每个单元格或网格点的网格诱导误差。 第三是由Eq定义的残差R,i。 （9.7）是一个局部量（至少对于显式方法），而解误差q作为Eq的解给出。 （9.8）是一个全局函数。 当残差非零时，就会产生错误，然后将其传输。 这表明网格自适应应该在残差高的地方进行，而不是在误差高的地方进行。 Qin \& Shih [21-231表明，如果Eq.定义的“实际”残差。 使用（9.7）；即R,i = LD(ug,i, h, k)或R,i = LD(\textasciitilde{}, k) U,i (9.10)，如前所述，需要计算独立于网格的解，然后由Eq给出的DETE。 （9.8）可以通过以下一维模型方程精确预测线性、非线性、稳定或非稳态的FDFV方程的“不完美”网格上的网格诱导误差：对流扩散方程、线性波方程和非黏性汉堡方程。 他们还发现，在估计弱解中网格诱导的误差时，用于推导非线性FD/FV方程的DETE的线性化过程不需要保守，因为误差传播速度和跳跃条件被保守形式的FD/FV方程捕获。 对于大多数工程感兴趣的问题，计算能力通常禁止生成独立于网格的解决方案，因此需要估计网格诱导的误差。 但是，没有独立于网格的解也意味着Eq中的残差。 （9.8）必须建模。 这是主要的挑战

\section{9.4 Euler 方程的FV解决方案的DETE}
188 在实现两个连续等式时估计网格诱导的误差。 （9.3）和（9.4）以及离散等式。 （9.8）和（9.9）。 误差传输方程的有用性取决于残差的建模程度。 由于Warming \& Hyett [32]的修改方程是FD/FV方程真正表示的PDE，因此它被用于指导残差的建模[21-25，34，351。 Zhang等人[34，351和Qin \& Shih [21-251）表明，如果网格间距或单元格大小足够小，使用修改方程中截断误差的领先项对残差进行建模可以产生出色的结果。 然而，当网格间距或单元格大小过大时，估计误差预测可能相当差[21，221。 最近，Celik等人[8]对Eq中的残差进行了建模。 （9.8）通过扩展关于细胞中心的FDFV方程中的项，并只保留前导项。 这种建模方法本质上与使用修改方程相同，只是时间导数不会被空间导数取代。 由于修改方程和泰勒级数扩展中的截断误差的领先项对网格粗时的残差建模没有用，秦和石[23]提出了一种基于数据挖掘的替代方法。 他们的方法涉及两个步骤。 首先是通过Eq评估“实际”残差来研究残差的行为。 （9.7）由各种质量差的网格以系统的方式创建。 第二步是根据通过统计分析和曲线拟合获得的对残差行为的理解，对残差进行建模。 9.4 Euler 方程的FV解决方案的DETE在本节中，推导出可用于估计以下Euler 方程的FV解中网格诱导误差的DETE：aQ aF aG at ax ay -+-+-=o (9.11a) y-1 2 E = - +-p (u2 + v2) (9.1 lc)，其中p是密度；u和v分别是速度的x和y分量；E是机械能和热能；P是压力；T是温度；y是比热比值的比率。 上面给出的Euler 方程适用于

\subsection{9.4.1 解决方案的有限体积法}
T。 I-P。 Shih 189 二维，可压缩，热量和热完美的气体的无黏流动。 9.4.1 解决方案的有限体积法 用于获取等式解决方案的FV方法。 （9.11）如下[31]。 考虑一个由有限数量的连续和非重叠的单元格离散化的域，其中每个单元格可以有一个任意的形状，如三角形、矩形或其他多边形。 整合等式。 （9.1 la）在体积为Vi的任意第i个单元格上产生（9.12），其中1和7分别是x和y方向的单位向量。 如果细胞不会随着时间的推移而变形，那么调用高斯定理；即等式。 （9.12）变成Qi=- QdV Vi Vi JVi（9.14）如果围绕细胞体积Vi的边界aVi有Ji面，那么Eq。 （9.14）可以写成（9.15a），其中Si,j表示j”面的长度。 等式中每个面的平均通量。 （9.15b）由以下一阶Lax-Friedrichs公式近似，即TVD保留，在上述方程中，下标L和R指的是第i个单元格的第j面的左右单元；v,,,, = (v,,,\textasciitilde{} + v,,,)/2是平均面正常速度；c, = (cL + cR) I2是平均音速。

190 估计网格诱导误差 方程中的时间导数。 （9.15a）由以下两阶段Runge-Kutta方案近似：Q（“=G;+AtY（G”）（9.17a）（9.17b），其中假定第n个时间水平（t”）的解（6英寸）是已知的，并寻求（n+l）'h时间水平（t"）的解（3英寸）。 如果只对稳态解决方案感兴趣，那么可以实施本地时间推进来加速收敛速率。 每个单元格（Ati）的时间步骤大小受以下CFL条件的限制：（9.18）因为F和G是Q到一级的齐次函数。 等式中的F和G。 （9.1磅）也可以写成F = AQ，G = BQ（9.19a），其中0 1 0 0（3-Y）U（1-Y）V Y-1 eu2 +y-lv2 2 2 - uv V U 0（y-I）（u2+v2）u-- YE -（3u 1-Y 2 +v2）+- B（1-y）uv？ /U P2 P 0 0 1 0 - uv V U 0 (1 - Y)U (3 -? Ov Y-1 - Y-3 +y-1,,2 2 2 (y-l)(u2+v2)v-- YE (1-y)uv X(3? +L？） +- m ly P 2 P -=-+-(u c2 Y2 +v2) P Y-1 2 (9.19b) (9.19\textasciitilde{}) (9.19d) 所以方程中的f(Q)。 （9.16）变成

\subsection{9.4.2 FV方法的DETE}
T。 1-P。 Shih 191 f(Q)=AQn, +BQny (9.20) 方程(9.19)和(9.20)是需要的，因为DETE需要线性化。 9.4.2 FV方法的DETE 方程给出的FV方程的DETE的推导。 （9.15）到（9.20）涉及三个步骤。 第一步是将FV方程线性化。 正如Qin \& Shih [21]所指出的，通过使用在寻求误差估计的“粗”网格上计算的变量值进行线性化很重要。 因此，线性化的FV方程的形式如下：（9.21）其中时间导数由等式近似。 （9.17）。 在Eq中。 （9.21），所有带有下标c的变量都使用粗网格上的信息进行评估。 需要注意的是，虽然Eq. （9.21）据说是线性化的，它与非线性化的FV方程相同，因为方程给出的恒等式。 （9.19）。 因此，当在粗网中获得的溶液Qc插入Eq中时。 （9.21），他们将满足Eq。 （9.21）完美无残留。 方程（9.21）只有在将网格独立解插入到其中时才被认为是线性化，因为在这种情况下，系数A和B不会通过使用用于生成网格独立解的网格上的信息来评估。 因此，当网格独立解Qg插入到Eq中时。 （9.21），将产生一个剩余的R，i。 第二步是减去Eq。 （9.21）从Eq中插入Qc。 （9.21）插入Qg。 由此产生的方程构成了方程给出的FV方程的DETE系统。 （9.15）到（9.20），即（9.22a），其中ei由

\section{9.5 DETE的有用性}
\subsection{9.5.1 测试问题1：无黏流动对翼型}
192 估计网格诱导误差（9.22b）是第i个单元格诱导的网格诱导误差。 如前所述，即使有弱解，DETE也不需要以守恒法形式。 这是因为激波的位置和跨越它的跳跃条件是FV解决方案的一部分，守恒格式正确捕获了FV解决方案。 完成Eq给出的DETE的推导的第三步。 （9.22）是对剩余Rg,i进行建模。 如前所述，对残差进行建模的一种方法是使用修改后的方程中的截断项[21、22、34、351。 另一种方法是从数据挖掘中进行曲线拟合[23]。 9.5 DETE的有用性 在本节中，Equ.给出的DETE的有用性。 （9.22）对两个测试问题进行了演示。 这两个问题的目标是说明，如果对残差进行准确的建模，网格诱导的误差可以估计得有多好。 第一个问题是不黏性，可压缩流动在二维翼型（商务喷气式飞机翼型，GLC305 [261）上。 第二个问题是黏性，可压缩流动在相同的翼型上加积冰（944冰形状[26]）。 对于这两个问题，自由流马赫数和压力分别为0.12和20.5磅/平方英寸。 基于自由流条件的雷诺数和翼型弦长度为3.5 x lo6。 翼型相对于迎面而来的流量的攻角为4度。 使用的边界条件是流入和流出边界的特征条件。 对于翼型，无黏流动是防滑条件，黏性流动是防滑条件。 无论是无黏性还是黏性流动，翼型表面都是隔热的。 9.5.1 测试问题1：无黏流动在翼型上。对于这个测试问题，Eq给出的Euler 方程。 （9.11）是有效的，Eqs给出的FV方法。 （9.15）到（9.17）用于生成解决方案。 生成了以下解决方案：（1）一个独立于网格的解决方案Qg和（2）三个粗网格解决方案Qc。 网格独立解决方案是在具有769 x 129网格点的网格上生成的，而较粗网格解决方案是在具有385 x 129网格点、193 x 65网格点和97 x 33网格点的网格上生成的。 较粗的网格是通过从前面的更精细的网格中删除所有其他网格线来获得的。 为了估计在较粗的网格上获得的解的误差，涉及两个步骤：首先，通过将网格无关的解Qg代入Eq的线性化版本来计算确切的残差。 （9.21）在粗网上；即，

\subsection{9.5.2 测试问题2：黏性流动在冰翼型上}
T。 I-P。 Shih 193 第二，使用Eq。 （9.22）与Rg，i由Eq给出。 （9.23）估计粗网格溶液中每个单元格的网格诱导误差。 由于DETE总是线性的，这些方程可以非常有效地求解。 通过比较方程的预测，可以评估估计的网格诱导误差的准确性。 （9.22）和（9.23）与实际网格诱导的误差直接从网格独立解和粗网格解通过Eq计算。 （9.22b）。 对于所有三个粗网格解决方案，发现如果使用实际残差，则由Eqs给出的DETE。 （9.22）可以完美地预测每个单元格的网格诱导误差。 由于残差Rg,i通常未知，使用DETE概念的主要挑战是开发残差模型。 9.5.2 测试问题2：黏性流动在冰翼型上。对于这个测试问题，控制方程是Navier-Stokes方程，而不是Eq给出的Euler 方程。 （9.11）。 Navier-Stokes 方程和Euler 方程之间的唯一区别在于扩散项。 由于当流量的雷诺数高时，对流传输预计将主导扩散传输，因此有兴趣看看基于Euler 方程的DETE是否可以以足够的准确性预测网格诱导的误差。 使用Eq给出的Euler 方程的DETE，估计Navier-Stokes（N-S）解中的网格诱导误差。 （9.22），请注意，将网格独立N-S解U!和粗网格N-S解Uzf替换为Eq。 （9.22）在粗网上都会产生残差；即，读者参考[24-251了解详情。 （9.24）（9.25）其中L：表示Eq的运算符。 （9.22）在粗网上。 减去等式。 （9.25）来自Eq. （9.24）产量

194 估算网格诱导误差方程（9.26）可以提供粗网格N-S解中网格诱导误差的估计。 在实践中，RFf可以通过使用Eq轻松计算。 （9.25），但RFS仍然需要建模。 图1显示了用于生成网格独立和粗网格N-S解决方案的精细和粗网格。 图2显示了使用Eq预测的x-momentum（pu）中粗网格N-S解的网格诱导误差。 （9.26）以及通过直接比较网格独立和粗网格解决方案获得的实际网格诱导误差。 来自图。 2，可以看出Eq的预测。 （9.26）只要R::是准确的，就相当不错。 有关进一步讨论，请参阅[9,25]。 图。 1. 左：网格独立网格。 右：粗网格网格。 图。 2. 左：pu中预测的网格诱导误差。 右：pu中的实际错误。

\section{9.6 最后讲话}
T。 I-P。 Shih 195 9.6 最终评论 在本文开始时，有人指出，如果单独使用CFD会影响设计，则必须知道每个CFD解决方案的误差界限。 CFD解决方案的一个主要错误来源可以归因于网格分辨率不足或差。 这里描述的DETE提供了一种估计网格诱导误差的方法。 这种方法也可用于指导网格的细化。 对此，有人指出，网格应该细化，在残差高的地方，而不是误差高的地方。 然而，本文中描述的方法的有用性取决于残差模型的准确性，这仍然需要更多的研究，特别是对于在非常粗的网格上获得的CFD解决方案。

\section{9.7 参考书目}
196 估计网格诱导的误差 9.7 参考书目 [31 PI [61 [71 Addy, H.E., “调制解调器机翼的结冰和结冰效应”, Babuska, I. \& Rheinboldt，W.，“自适应有限元计算的错误估计”，SIAM数值分析杂志，卷。 15，Babuska, I., Strouboulis, T., \& Upadhyay, C.S., “线性椭圆问题后误差估计器质量的模型研究：三角形的斑驳均匀网格内部的误差估计”，应用力学和工程中的计算机方法，卷。 114，1994年，Babuska，I.，Strouboulis，T.，Gangaraj，S.K.和Upadhyay，C.S.，“有限元法的h版本中的污染误差和回收衍生物的局部质量”，应用力学和工程中的计算机方法，卷。 140，1997年，pp. 1-37。 凯里，G.F.，计算网格：生成、适应和解决方案策略，泰勒和弗朗西斯，华盛顿特区，1997年。 Celik, I., Chen, C.J., Roache, P.J., \& Scheuer, G., 编辑，计算流体动力学不确定性量化研讨会，FED-Vol. 158，ASME流体工程部，夏季会议，华盛顿特区，1993年。 Celik，我。 \& Zhang，W.-M.，“使用Richardson外推计算数值不确定性：一些简单的湍流计算的应用”，ASME流体工程杂志，卷。 117，9月Celik，I.，Hu，G.和Badeau，A.，“单网格误差估计技术的进一步完善和基准标记”，AIAA论文2003-0626，1月。 2003年。 Chi, X., Zhu, B., Shih, T.1-P., Addy, H.E., \& Choo, Y.K., “CFD 前沿冰吸积的公务机翼型空气动力学分析”，AIAA Paper 2004-0560，2004年1月。 Ferziger，J.H.，「数值误差的估计和减少」，在计算流体动力学中不确定性量化研讨会，FED-Vol. 158，ASME流体工程部，夏季会议，华盛顿特区，1993年，pp. 1-8。 芬莱森，B。 A.，加权残差法和变异原理，爱思唯尔科技书籍，1972年。 Giles, M.B., Larson, M.G., Levenstam, J.M., \& Suli, E., “黏性流动中升力和阻力系数有限元近似的自适应误差控制”，报告NA-97/06，牛津大学计算实验室，1997年。 NASAlTP 2000-21003 1，2000年4月。 不。 4，1978年，pp. 736-754. pp. 307-378。 1995年，pp. 439-445。

T。 I-P。 Shih 197 Giles，M.B. \& Pierce，N.A.，“关于Adjoint Euler 方程解的性质”，第六届ICFD流体力学数值方法会议，英国牛津，1998年。 Giles，M.B.，“关于误差分析和最优网格自适应的伴随方程”，在计算流体动力学的前沿，D.A. Caughey和M.M. 哈菲兹编辑，世界科学，1998年，pp. 155-170。 吉尔斯，M.B. \& Pierce，N.A.，“使用辅助Euler 方程改进的升力和阻力估计值”，AIAA论文99-3293，1999年。 Gu, X., Schock, H.J., Shih, T.1-P., Hernandez, E.C., Chu, D., Keller, P.S., \& Sun, R.L., “结构化和非结构化网格的网格质量措施”，AIAA Paper 2001-0652，1月。 2001年。 Gu，X。 \& Shih，T.1-P.，“解决方案自适应网格加密中错误源和错误位置之间的差异化”，AIAA论文2001-2660，CFD会议，2001年6月。 Mackenzie, J., Sonar, T., \& Suli, E., “双曲问题的自适应有限体积方法”，有限元素和应用的数学，由J.R.编辑。 怀特曼，约翰·威利和儿子，纽约，1994年，第19章。 Oden，J.T.，“计算流体动力学中的错误估计和控制”，有限元数和应用的数学，由J.R.编辑。 怀特曼，约翰·威利和儿子，纽约，1994年，第1章。 Peraire, J., Vahdati, M., Morgan, K., \& Zienkiewicz, O.C., “可压缩流动计算的自适应重合”，J. 计算物理学，卷。 22，1976年，第131-149页。 秦，Y。 \& Shih，T.1-P.，“CFD中误差估计的离散运输方程”，AIAA论文2002-0906，2002年1月。 秦，Y。 \& Shih，T.1-P.，“在有限差和有限体积方法中估计网格诱导误差的方法”，AIAA论文2003-0845，1月。 2003年。 秦，Y。 8z Shih，T.1-P.，“离散误差传输方程中残差的分析和建模”，AIAA论文2003-3850，2003年6月。 Qin, Y., Keller, P.S., Sun, R.L., Hernandez, E.C., Perng, C.Y., Trigui, N., Han, Z., Shen, F.Z., \& Shih, T.1-P., “通过离散误差-传输方程估计CFD中的网格诱导误差”，AIAA论文2004-0656，2004年1月。 秦，Y.，Chi，X.，和Shih，T. I-P.，“用欧拉离散错误传输方程估计Navier-Stokes解决方案的网格诱导误差”，AIAA论文2005-0567，2005年1月。 Roache，P.，计算科学和工程中的验证和验证，Hermosa出版社，新墨西哥州阿尔伯克基，1998年。 Shih, T.1-P., Gu, X., \& Chu, D., “网格质量测量和误差估计”，计算场模拟中的数值网格生成，B.K. 索尼，J.F. Thompson，J. Hauser和P.R. 艾斯曼，编辑，密西西比州立大学，ISSG，2000年，pp. 799-808。

198 估计网格诱导误差 Sonar，T.，“基于可压缩流动计算的有限元剩余的强和弱规范细化指标”，科学和工程计算的影响，卷。 5，1993年，pp. 1 1 1 - 127. van Straalen, B.P., Simpson, R.B., \& Stubley, G.D., “流体流动运输有限体积模拟的事后误差估计”，加拿大CFD学会第三届年会论文集，卷。 1，巴尼夫，阿尔伯塔省，1995年6月。 文迪蒂，D.A. \& Darmofal，D.L.，“函数输出的伴随误差估计和网格自适应：对准一维流动的应用”，《计算物理学杂志》，卷。 164,2000，pp. 204-227。 王，Z.J.，「非结构化网格守恒定律的光谱（有限）体积方法：基本制定」，J. 计算物理学，卷。 变暖，R。 F。 \& Hyett，B. J.，“有限差分方法的稳定性和准确性分析的修改方程方法”，J. 补偿。 物理学，卷。 14，1974年，pp. 159-179。 威尔逊，R.V. \& Stem，F.，“海军水面战斗人员RANS模拟的验证和验证”，AIAA论文2002-0904，1月。 2002年。 Zhang, X.D., Trkpanier, J.-Y., \& Camarero, R., “双曲守恒定律有限体积解的事后误差估计”，应用力学和工程计算方法，卷。 185, Zhang, X.D., Pelletier, D., Trkpanier, J.-Y., \& Camarero, R., “Euler 方程误差估计器的数值评估”，AIM杂志，卷。 Zienkiewicz，O.C. \& Zhu，J.Z.，“实用工程分析的简单误差估算器和自适应程序”，国际工程数值方法杂志，卷。 1987年24日，pp. 337-357。 178，pp. 210-251,2002。 2000年，pp. 1-19。 39，不。 9,2001，pp. 1706-1715。

\clearpage
\chapter{第10章使用涡量约束 John Steinhoff和Nicholas Lynn处理流}
\section{10.1 摘要}
第10章 使用涡量约束 John Steinhoff' \& Nicholas Lynn' 10.1 摘要处理涡流 描述了一种新的计算方法，该方法旨在捕获高雷诺数不可压缩流中的薄涡流区域，包括可能分离的边界层和涡流片和可能长距离对流或合并和重新连接的丝。 新方法的主要目标--涡量约束（VC）--是捕捉这些小规模涡流结构的基本特征，并直接在欧拉计算网格上用非常高效的微分方法对其进行建模。 该方法允许对流结构在2个网格单元上建模，在长距离对流时没有数值扩散，不需要特殊逻辑来合并或重新连接。 在本文中，对基本VC方法的描述比以前更全面。 VC和众所周知的冲击捕捉方法之间有非常接近的类比。 首先讨论这些来解释基本动机，因为VC背后的想法与传统的CFD方法有些不同。 VC为更有效地计算复杂流量提供了一些可能性，并提出了结果，包括一些新的探索性研究。 “田纳西大学太空研究所，田纳西州塔拉霍马。 2田纳西州塔拉霍马田纳西大学太空研究所研究生研究助理。

\section{10.2 介绍}
10.2 导言涡量约束 我们描述了一种新的计算方法，用于有效处理高雷诺数不可压缩流中的薄涡流结构。 例子包括附着和分离的边界层、对流涡流片和长丝，以及钝体尾迹的湍流中的小尺度。 新方法-涡量约束（VC）的主要目标是捕捉这些小规模涡流结构的基本特征。 通过基本特征，我们指的是可能分离的边界层和在长距离对流中没有显著扩散的涡流丝，并且可以改变拓扑结构并合并或重新连接。 这是直接在（固定欧拉式）计算网格上用非常高效的差分方法完成的。 这比首先制定模型偏微分方程，然后进行离散近似更有效，因为它允许对薄涡流区域进行建模，分布在1-3个网格单元中。 然后，流的非旋转部分，以及任何更大规模的涡流结构，都可以用传统的离散偏微分方程（Euler 方程）以足够的分辨率求解。 对于这些地区，该方法可以简化为传统的计算流体动力学（CFD）。 本质上，相同的（VC）方法用于对上述所有区域（附着和分离边界层和对流涡流）的小规模结构进行建模。 这分别导致了一种新型的边界层模型、薄对流涡流模型和湍流大涡流模拟（LES）的“子网格”模型。 除了允许使用简单的粗网格外，该方法还可以有效地用于基本方程的低阶时间和空间离散化，因为它消除了由于小涡流尺度的数值扩散而导致的人工传播。 这大大简化了边界条件，并进一步缩短了计算时间。 本文的主要目的是比以前的论文更详细地解释VC背后的基本思想，并探索它为更高效地计算复杂流提供的可能性。 风险投资背后的基本理念和众所周知的冲击捕捉方法之间存在非常密切的类比。 简要讨论了这些，以解释基本动机。 然后对涡量约束方法进行了描述。 作者小组和其他人使用VC提出了许多结果。 在本文中，只提出了作者小组的少数结果，包括一些新的探索性研究。 也提到了其他群体的一些结果。 正在准备对结果的总体审查。

\subsection{10.2.1 基本概念}
Steinhoff \& Lynn 201 10.2.1 基本概念 由于本文涉及一种相对较新的方法，因此将首先简要描述基本动机。 基本目标接近于类似问题，也涉及薄结构--冲击捕获。 因此，在描述新方法之前，将简要描述冲击捕捉方法的类似相关特征，因为它们已经广泛使用了一段时间，对CFD社区非常熟悉。 然后，将审查新方法（VC）的基本概念，并介绍一些最近的新结果。 虽然上述类比对某些人来说可能是显而易见的，但可能更广泛的受众并不明显。 当然，冲击捕捉方法在CFD社区得到了极大的关注，并被证明是极其重要的。 这些方法通常只在冲击区域使用中等大小的非黏性计算网格。 这是可能的，因为只保留了冲击的基本物理学（就解决流动问题而言）。 我们所说的“基本物理学”是指那些影响冲击内部外部流动的特征。 这些特征包括计算的冲击厚度，它不一定和物理厚度一样小，但和物理厚度一样，与问题的主要长度尺度相比，必须很小。 它们还包括通过冲击整合的保护法被保留的要求。 通过这种方式，对于许多不依赖于冲击内部结构细节的问题，通过专门开发的数值“冲击捕获”算法获得了准确的流动解决方案。 在这些方法中，避免了内部冲击结构的偏微分方程（pde）的详细准确解（例如，Navier-Stokes 方程）。 这很重要，因为它避免了结构中非常精细的计算网格的要求，以及那里非常耗时的黏性计算。 这些想法可以追溯到Von Neumann和Richtmyer [41]、Lax [21]等人，涉及pde的“弱解”的概念，在非黏性极限中，可以处理不连续特征。 在离散化之后，这些想法允许在1-2个网格单元格中近似冲击。 这个问题自然会出现，即是否有可能对其他薄特征进行类似的高效“捕获”处理，这应该会产生类似的好处。 然而，一个区别是，与涡流区域不同，冲击特征向内倾斜，在计算过程中自然倾向于陡峭。 因此，建模冲击比建模薄涡流结构和其他类似接触的不连续性更简单，由于数值离散化误差和稳定数值扩散的需要，这些不连续性自然倾向于扩散。 这导致要求添加“陡峭”或“禁止”一词，以防止人为传播。 最早的禁闭类型之一

202 涡量约束接触不连续方案由Harten开发[I71用于一维可压缩流动。 本文中描述的方法-涡量约束（VC）-已经制定，以类似于冲击捕获的方式有效处理难以计算的集中涡流区域。 虽然独立开发，但该方法可以被认为是Harten、Van Leer和其他人的一维可压缩流动方案到多维不可压缩流动的旋转不变扩展[17,401。 最近的一些论文[29,30,32,34-36,42,43]描述了VC用于不可压缩流动。 最近开发了可压缩流动的进一步扩展[4-6,181，但这里不会描述。 与冲击捕获一样，据了解，涉及薄涡流区域内部结构细节的现象将不会被准确处理，除非为它们开发特殊模型。 当然，仍然存在一些物理内部结构的细节很重要的问题。 然而，主要用途是针对大量问题，仅冲击捕获就非常有效。 与有冲击的流动一样，存在许多高雷诺数（Re）流动问题，只要某些基本方面得到准确处理，薄浓缩涡流区域（“covons”）的内部结构细节与重要的流动特征无关。 此外，我们认为用VC捕获的covons可以作为起点，或“零阶”解决方案，可以为内部涡流结构添加额外的物理模型--但只需要\$necessary。 单独VC有效的例子包括对流涡流环[35]。 这些可以在不扩散的情况下对流，但可以在不需要特殊逻辑的情况下合并。 此外，薄薄的翼尖漩涡可以在相对较远的距离（一些翼跨度）上进行计算，并表现出乌鸦不稳定性，包括合并[44]。 上述两种现象的计算与实验密切一致。 在这种情况下，对于非常远的距离（许多公里），该方法可以作为零阶方法，因为湍流最终诱导了非常缓慢的传播。 然后，这种效果可以简单地建模，同样无需使用非常精细的网格或高阶方法。 [44]描述了一项涉及以这种方式将VC用于尾随涡流的初步研究。 例子还包括分离边界层：单独VC对尖锐边缘的分离有效[6, 8,12,45]，但对于与光滑表面的分离，在任何高Re情况下，准确的计算需要涉及边界层 湍流的附加建模术语。 我们目前正在为这些效果开发模型，也没有使用非常精细的“黏性”网格或高阶pde 离散化方法[3,8,9,11,12,27,45]。 应该重申的一个重要点是，在高Re，大多数涡流区域将是湍流。 因此，任何计算方法都必须明确或隐含地涉及内部结构的数值模型，因为直接求解此结构的Navier-Stokes 方程是不可行的。 足够的

Steinhoff \& Lynn 203 在捕获covon时用VC获得的结构，因此可以被认为是这样的模型，但计算效率很高。 此外，该模型本质上是离散的，仅在几个网格单元格上定义，并不是模型pde的准确解。 原因是，即使使用高阶方法，如果只分布在少量的网格单元（2-4）上，也很难解决长距离的薄涡流结构的pde。 这是因为众所周知的事实，即方法的准确性或顺序只是对误差行为的渐近估计，适用于大N，即整个涡流区域的网格单元的数量。 N=2-3可能不足以应用这样的估计值。 由于VC旨在捕捉该功能，而不是准确解决模型pde，因此它绕过了这个问题。 目前还有其他努力直接在网格上捕获小规模，用于LES 湍流模拟。 这些大多使用一维运算符的组合，如原始不连续捕获方案[2,17,40]--它们被适当地称为隐式大涡流模拟，或ILES（151。 与VC相反，那里的重点通常不是涡量。 此外，他们的主要目标是尽可能取消数值扩散，而我们的目标是直接用受控的“模型”结构创建薄而稳定的涡流结构。 因此，VC甚至在某些长度尺度上涉及负的，尽管不发散的，但总扩散。 将VC与另一种用于薄涡流区域-涡流晶格的方法进行比较是有指导意义的[19]。 我们描述了VC如何消除这些方法的主要缺点，但保留了优势。 我们还解释了从涡流晶格方法到VC的过渡如何与从激波拟合到冲击捕获方法的过渡非常相似。 涡流晶格方法有效地涉及附着边界层（“绑定涡量”）内部结构的简单模型和分离涡流片（作为涡流丝的集合）。 然后，这些方案使用基本流体动力学方程来确定绑定涡量的强度和对流片的位置，并用标记数组定义它们。 与VC一样，由于离散化错误，这些方法没有任何人为扩散，与VC一样，它们可用于薄区域内部结构细节不重要的情况。 然而，与VC不同，涡流晶格方案通常仅限于拓扑没有变化的流，例如对流涡流片的合并。 为了允许拓扑结构的变化，如边界层分离或涡流合并，这些方案通常需要特殊的逻辑，并且通常需要创建新的标记数组。 相比之下，VC允许在不增加复杂性的情况下任意更改拓扑结构。 重要的是要强调，VC背后的目标是用另一种更灵活的涡流晶格结构取代这些显式建模的涡流晶格结构，该结构可以“捕获”涡流区域。 这与治疗冲击的过渡非常相似：“激波拟合”方法最初是为定位标记集而开发的，使用冲击跳转条件来定义

\section{10.3 说明性的一维示例}
204 涡量约束 单个冲击表面。 然而，当尝试通过改变冲击拓扑或多次冲击来进行更一般的流动时，事实证明，固定欧拉网格上的冲击捕获更简单、更高效（[39]，第365页）。 首先，将介绍基本的VC方法。 然后，应用上述现象的一些例子。 第一个是捕捉薄而孤立的对流涡流片和长丝，第二个是钝体 湍流唤醒的模型。 第三，分离边界层的处理，就像第二种一样，最近得到了研究，一些方面仍在开发中，尽管已经有许多有希望的结果。看到应用最多。 10.3 说明性一维示例 如上一节所述，VC的基本概念是，薄涡流像冲击一样，仅在几个网格单元上被“捕获”。 当然，这意味着离散方程并不代表内部结构的简单pde的准确解，尽管它们可以被视为奇异（弱）解的近似值。 在这些情况下（对于小规模）的目标是只准确处理某些整数，并保持定性结构（即涡流保持薄）。 这个概念在一个维度中有一个简单的例子，涉及被动标量或集中在一个小区域的“脉冲”的对流。 因此，一个目标可能是保持这个标量的总振幅，让重心以正确的对流速度移动，并确保它保持紧凑--基本上分布在少数细胞上。 表示额外结构的脉冲的额外矩也可以使用额外的场进行传输，但没有尝试对流内部结构的准确、点向表示。 然后，我们考虑以速度c进行标量（@）对流，在连续情况下，它将符合pde a,@ = -a xc@符号化，然后我们定义一个离散方程：= - AtC@y + EY，（10.1）（10.2）其中C是传统的保守离散对流算子，j表示空间网格指数，n是时间步（t = nAt），ET是一个非线性项，旨在保持脉冲。 如果总脉冲振幅和速度要保持保守，与脉冲振幅无关，安永有要求：

Steinhoff \& Lynn 205 1. EY必须是（至少）第一个差，以保持总振幅（假定它从重心迅速消失）。 E必须（至少）是二阶导数，因此，每个时间步，重心位置不会被El改变，因此，重心的速度是正确的，正如原始pde给出的那样。 然后，其中6：是一个离散的二次差分算子，Ff是4及其差的函数，在远场中消失。 F,n必须与方程中的其他项一样，在4中具有1度的均匀性，因此对4的振幅或尺度没有依赖性。 6，？ 如果脉冲分布在太大的区域，F，”必须表示负扩散，使其收缩并放松成固定形状。 6：F：出于同样的原因，如果脉冲太薄，必须变成正扩散。 属性（4）要求F,n是非线性的。 如果它是线性的，负扩散将导致某些模式发散，因为它们将是解耦和独立的。 EY = 6:F," (10.3) 满足上述要求的简单公式是：6'c C=" h (10.4) F" = p@" -a'' (10.5) 其中6]'是中心差，h是网格单元格大小，p和E是常数，是谐波平均值（10.6）F还有许多其他公式也可以使用。 F中的第一项，"作用于稳定中心差算子C，也满足条件（1-3）和（5），第二项满足条件（1-4）和（6）。 由此产生的差分方程可以写成（10.7）

206 涡量约束，其中6；是第二个差分算子，使用y =.2和E = 0进行无限制的计算结果。在100个时间步（1/10英寸穿过网格）后，如图所示。 1. （当然，更高阶的常规CFD方法可能会导致比此处显示的更少的扩散。 然而，与禁闭相比，这些都需要脉冲内更多的网格单元，并且由于累积的数值误差，最终将其传播到更多的单元格上。） 结果也显示在图1中，通过y =.2和E =.5的（periodic 256网格单元格）网格进行1和100次限制。 对于所有这些情况，v = \&/6。 任何给定时间步的确切脉冲高度取决于网格单元格内的重心位置，几乎就像固定的脉冲形状通过网格和每个时间步骤采样的值“移动”一样。 这就是两个受限脉冲图像看起来有些不同的原因。 然而，脉冲仍然被无限限制，除了由于计算机的有限精度而产生非常小的影响（在-lo6时间步后）。 这意味着Eq有一个近似的、平滑的解。 10.7，@(x,t)，就相似性变量而言，z=x-ct。 对于具有薄脉冲形式的初始条件范围，@将放松到这种特定的“孤立波”脉冲解决方案。 这是以数字方式显示的。 然而，孤波的稳定性和被它吸引的初始条件集还没有从数学上推导出来。 已经表明[27]，对于Ll、E和V的值范围，@的符号不能改变。 因此，对于初始非负脉冲，最大值不能超过守恒的总和。 这可以防止解决方案发散。 V，=c，At/h。 （10.8）在小At极限中，对于常数c，@满足C5;(p\$;-\&;)=O。 (1 0.9) 解决方案是A和xo取决于初始振幅和重心，以及4,” = Asech ( y (x - xo - ct)) ( 10.10) (10.1 1) 强调脉冲的粒子样属性的有趣一点是，可以从Eq中推导出“Ehrenfest”型关系。 10.7. (X>”+'-(x)” =(c)”在（10.12），时间步骤n处的重心和对流速度，（10.13）

\section{10.4 涡量约束}
Steinhoff \& Lynn 207 (c)“ = c c,4; /c 4; 7 (10.14) x, = jh. （10.15）其中我们只有方程的@(z)的数值解。 10.7一般At。 然而，如果我们有常数c，并使用：以及<p：的不同形式，那么可以证明有一个精确的闭合形式解，由传播在两个细胞上的三角波组成。 这个结果之前在参考[8，321中显示过。 应该提到Eqs。 10.6和10.7，在j'<j上相加时，表示类似传播步进函数的不连续性，而不是脉冲。 Harten [17]、Van Leer [40]和其他人使用类似的方程来捕捉可压缩流动中的一维接触不连续性和冲击，@作为“限制器”（见第10.4.2节）。p=v/2（10.16）@ = min(v@,“, (1-\textasciitilde{})4,“-\textasciitilde{}); (O<v<l) (1 0.17) 10.4 涡量约束 已经开发了两种具有相似属性的VC公式：第一个“VCI”涉及速度的一阶导数[31, 341，而第二个“VC2”涉及二次导数[27]。 下面将介绍这两个版本。 VC的基本原理，如上面描述的一维对流标量示例，是一个具有稳定结构的解，可以无限期地传播。 两个项控制了这个结构：一个作用于收缩它，一个作用于扩展它，这样结构在传播时保持接近平衡解。 此解决方案可以稳定地受到其他术语（如对流）的离散化误差引起的扰动。 虽然VC方程可以写成pde的离散化（如下所述），但小尺度（结构内）的结果解并不是原始pde的准确甚至近似解。 这是因为VC旨在仅在几个网格单元格上捕获或建模小规模特征，因此离散化“错误”是0（1）。 然而，正如第10.2节所解释的那样，对这个问题至关重要的特征仍然保留下来。 因此，捕获的特征实际上是一个“生活”在网格晶格上的非线性孤波。 目前，有大量的工作正在进行内在离散-r差，而不是jhite差，方程（例如，见参考[41]）。 另一方面，在平滑区域（或大尺度）中，VC可以自动恢复到传统的CFD，然后精确有效地近似pde。

208 涡量约束 对于细涡流区域，我们使用与上述一维示例基本相同的方法：“禁锭”项被添加到传统的离散动量方程中。 虽然我们使用了一个原始变量，而不是一个涡量的公式，但我们会看到，如果我们查看生成的涡量传输方程（对于VC2版本，将在下面描述），它具有相同的“限制”术语，作为第10.3节中描述的一维标量传输方程的多维扩展。 通过这种方式，涡量以紧凑的方式运输，没有数字扩散。 对于一般的非稳定不可压缩流动，控制方程与涡量约束是连续性和动量方程的离散化，并添加了项：v.i=o（10.18）-vp 1 + [pv'q -a P（10.19），其中G是速度向量，p是压力，p是密度，p是扩散系数，包括数字效应，例如，由于第一个右侧项（对流项）的离散化。 （我们假设雷诺数很大，物理扩散比附加项小得多）。 对于最后一个项，\&，E是一个数值系数，与p一起控制对流涡流区域或涡流边界层的大小和时间尺度，s'定义如下。 出于这个原因，我们将括号中的两个术语称为“限制条款”。 向量3对于两个VC公式是不同的，定义如下。 方程10.19涉及常数，Ll和E，这足以解决许多问题。 如果这些不是恒定的，例如，当网格间距不是恒定的或有多个涡流刻度时，那么这些量可以在相应的项中取在微分算子内，以保持动量守恒（在VC2公式中）。 如一维示例所示，表示传播或正扩散和“收缩”或负扩散的一对禁闭术语共同构成了禁闭结构。 当两个项大致平衡时，就会产生稳定的解决方案。 通过这种方式，每个时间步骤都会进行校正，以补偿非恒定外速对流、对流算子中的离散化误差或压力校正引起的涡流结构的任何扰动。 参数，U和\&决定了生成的涡流结构的厚度和该状态的松弛率。 一般来说，对于边界层和对流涡流丝，涡流区域外部的计算流场对内部结构不敏感，因此对参数E和p不敏感，在广泛的值范围内。 例如，一般稀薄、集中的涡流在物理上往往会演变成一个

Steinhoff \& Lynn 209轴对称配置[22]。 此外，即使是快速旋转的非对称配置在短时间内平均时也会大致呈轴对称[23]。 然后，众所周知，轴对称二维涡核外的流动与涡流分布无关，因此只要核心薄（并且灯丝曲率大，因此在垂直于灯丝的平面上，流动大约是二维的），就不会依赖于\&和，U）。 因此，设置这些参数所涉及的问题将与其他标准计算流体动力学方案中设置数值参数所涉及的问题相似，例如许多传统激波捕捉方案中的人工耗散，正如所解释的那样，这些方案非常相似。 此外，对于湍流尾流，初步研究表明，E可用于参数化有限的雷诺数效应，因为它控制了最小解析的涡流尺度的强度（这是当前研究的主题[l 11）。 涡量约束方法的一个重要特点是，对于不可压缩流动，限制项仅在涡流区域中非零，因为扩散项和“收缩”项在这些区域之外都消失了。 因此，即使存在与对流算子相关的二阶各向同性数扩散，并且扩散算子只是二阶，在涡流区域之外，这些项的结果精度可以是三阶或四阶，因为这种扩散只是涡量的负卷曲。 最后一点涉及VC校正在质量、涡量和动量中引起的总变化，在对流涡流的横截面上集成。 可以证明[20]，由于求解器中的压力投影步骤，质量是守恒的，涡量是显性守恒的，因为校正在涡流区域外的消失，而且，（在VC2公式中）动量也是完全守恒的[27]。 当然，动量守恒是前面有一个空间导数运算符的附加项的结果。 在一维示例中，这允许我们为脉冲重心运动写出我们所称的“Ehrenfest”关系。 我们尚未为封闭对流涡流证明这一点，但我们也相信在这种情况下是真的。 然后，涡流重心将以“背景”流动速度的加权平均值移动，不会因自诱导流动而产生影响（至少对于具有恒定背景速度的二维情况）。 这已被数字证明[12，281。 （由于VC1公式中缺乏动量守恒而导致的错误在大多数情况下都很小，下文将讨论。） 许多基本的数值方法可用于空间和时间离散化。 我们在时间上使用简单的一阶欧拉积分，在空间上使用二阶积分。 在传统的CFD方案中，通常必须使用高阶方法，通常是为了减少数值扩散，从而为薄涡流区域扩散。 涡量约束在许多情况下消除了这个问题，并避免了高阶方法的边界条件复杂性和计算成本。 （然而，应该提到，第二个限制或收缩项涉及比其他项更大的差异模版）。

\subsection{10.4.1 基本配方}
\subsection{10.4.1.1 VC1配方}
210 涡量约束 另一个数字问题涉及网格的规律性。 重要的是要认识到，由于对流涡流或分离的边界层直接捕获在网格上，在几个网格单元上，因此不应使用大网格长宽比或快速变化的单元格大小。 如果避免这些，VC将产生接近旋转和平移不变的动力学。 当然，这些问题也发生在冲击捕捉中。 然而，如果长宽比不是太大，可以进行一些更正，以适应非均匀网格。 10.4.1 基本公式 两种不同的公式，VC1和VC2，具有一些不同的动力学，因为它们在收缩项中导数的顺序上有所不同。 最初开发的（VC1）已经在许多出版物中进行了描述，这里将只介绍一些细节；10.4.1.1 VC1 公式化 此公式涉及“收缩项”的表达式，s'没有明确保存动量：;=liX6（10.20）（10.21）（10.22）该项本质上将涡量在沿着其自身梯度2的细涡流区域内对流，从边缘或幅度较低的区域向中心或规模较大的区域。 随着结构收缩和梯度增加，线性扩散的“膨胀”项会增加，直到达到平衡。 （这是对流扩散现象的一个众所周知的属性。） 由于对流集中涡流的快速旋转，任何非保守动量误差几乎完全被取消，事实证明，该方法对许多问题都足够准确。 使用这种方法的人经常忽视应用这种方法的技术性：在早期的论文中对此进行了描述[31]。 由于涡量沿h对流，因此应在收缩项中使用w的上风（以h为单位）值，以避免产生具有相反符号的涡量的“下风”值。 这很容易通过每个节点的加权因子来实现，这些因子依赖于ri和相邻网格节点的单位向量。 文献中提出的大多数VC结果都使用VC1配方。 然而，它们不涉及非常缓慢的背景流动，也不涉及下面讨论的动量守恒问题。 然而，一个重要的一点是

\subsection{10.4.1.2 VC2配方}
Steinhoff \& Lynn 211 在某些情况下，精确动量守恒可能不如其他特征重要（例如，在我们的案例中，确保对流涡流保持薄），不应被视为绝对要求（例如，见基本CFD教科书-Tannehill、Anderson和Pletcher [39]，pg。 60）。 10.4.1.2 VC2配方 仅用于非常精确地确定在非常缓慢的背景速度场中对流的涡流的长期轨迹，才发现保持动量的VC2配方是必要的（对于不可压缩流动）。 这确保了自诱导速度对涡流运动的贡献被完全抵消。 VC2配方的不可压缩流体动力学方程是V的离散化。 +O (10.23)和\$+' =\$"'At?.( <\textasciitilde{})+Vzp\$ +\$x(EG")或类似形式：;"+' =4" -\&.( G<)-Vx(p\&" -\&") (10.24) 其中6" =axsin (10.25)和5; =I\&; I+\& (10.27) 方程10.24比它上面的形式有一些数值优势，因为相同的差分算子作用于和G。 此外，第二个限制项（10.26）是模板的总和，该模板由中央节点（\$计算）及其相邻（-1）节点组成，6是一个小的正常数（-10'），以防止有限精度引起的问题。 当两个相反方向的涡流彼此靠近时，中间可能存在网格单元，其中6没有明确定义，这可能会导致振荡。 为了防止这种情况，如果带有中心节点的模板中任何其他涡量向量的标量积为负数，\$将设置为零。

\subsection{10.4.1.3 边界条件}
212 涡量约束 为了了解VC2的效果，我们取Eq的卷曲。 10.24。 然后我们得到了涡量的运输方程。 例如，在二维中，这个方程，包括禁键项，正是第10.3节中显示的有效一维标量对流方程的多维、旋转不变的推广。 当然，该解决方案仍将反映网格的四倍对称性。 然而，这种效应从涡流区域迅速消失。 此外，围绕涡核的旋转流实际上允许更简单的离散化等式。 10.28，与轴对称对流被动标量分布相比。 这在参考中进行了解释。 [27]。 方程10.26是一个谐波平均值。 它被选择更重地权衡模版中的小值。 众所周知，当其参数的任何值消失时，该术语就会消失，阻止创建相反符号的值（用于一系列参数）。 使用VC，涡流周围区域的总涡量是守恒的，因为它是一个局部术语。 这意味着涡量不能因为这个项而发散，因为当所有值都有相同的符号时，最大绝对值不能大于和的绝对值。 （这也是一维标量对流示例的属性。） 有大量的替代形式可以和谐波平均值一样有效。 然而，我们认为该术语应该有一个平滑的代数形式，以给出平滑的结果。 这应该比涉及逻辑函数的形式（如“minmod”）更适合多维应用，后者在一维应用中提供了良好的效果。 如上所述，10.26等术语以前被用作限制器，但在一维可压缩流动中，据作者所知，不用于多维涡量。 （VC方法是独立开发的，作为一种多维、旋转不变的“禁镌”方法，专门用于薄涡流区域）。a,m = -v。 (M;) + v2 [pm-m(W)] (10.28) 10.4.1.3 边界条件 首先，我们描述了涉及使用VC执行边界条件和创建模型边界层（BL）的一般特征。 然后，我们描述了两种方法：使用浸入式BL，以及使用表面符合网格。 两者都涉及无流动条件。 它们还涉及防滑条件。 后者对分离问题很重要，因为它们确保由此产生的分离BL具有明确定义的涡量。 此外，两者都涉及粗糙的、不黏性大小的网格。 就像将VC用于薄的对流涡流区域一样，它们产生了一个简单、非常经济的BL模型。 该模型不涉及确定详细的时间平均速度曲线，这将需要一个非常精细的、适合身体的网格：相反，它仅在几个粗网格单元格上对轮廓进行建模。 因此，它对钝体流量很有用

Steinhoff \& Lynn 213 大规模分离，与分离BL的位置和强度相比，这种详细的轮廓以及皮肤摩擦是次要的。 对于没有VC的传统CFD溶液，由于粗略的和可能不符合要求的网格导致边界处的巨大数值误差，该涡量将迅速对流并扩散到表面区域，破坏外部溶液的准确性。 然而，VC的使用将涡量限制在附着时沿表面的1-2个网格单元。 就在这个层外，速度是平滑的，接近相邻表面的切线。 然而，这种简单的边界层仍然可以分离，特别是在边缘和不利压力梯度强的区域。 尽管我们专注于VC2版本，但这里应该提到VC1版本的使用，因为它对附加边界层的解释非常简单。 在这种情况下，方程中的向量2。 10.21被定义为局部垂直于边界层区域的表面。 然后，VC1只是正扩散（使涡量远离表面）和涡量向表面的对流的组合。 事实证明，这是一种非常强大和高效的边界层建模方法，结合了外部速度的切向平滑和涡量的正态“压缩”。 已经提出了一些结果，证明了这种方法的有效性[9,11,12,27]。a. 浸入式边界层模型在浸入式表面上强制执行防滑边界条件，首先，表面由在每个网格点定义的平滑「水准集」功能「F」隐含表示。 这只是从每个网格点到物体表面最近点的（有符号）距离--外部为正，内部为负。 然后，在溶液过程中的每个时间步骤中，内部的速度都设置为零。 在使用VC的计算中，这导致沿表面形成一个集中的涡流区域。 重要的一点是，与许多传统方案不同，“切割”单元格不需要特殊的逻辑：与网格的其余部分一样，只应用相同的VC方程。 此外，与传统的浸入表面方案不同，由于细胞大小限制，这些方案是无黏性的，实际上有一个防滑边界条件，这导致边界层具有明确定义的总涡量，并且由于VC，它仍然很薄。 这种方法对与尖锐的分离的复杂配置有效。 此外，即使没有特殊的边界层模型术语，它也可以大致处理与光滑表面的分离，如第10.5节所示。b。 符合网格，湍流边界层模型，我们正在开始在VC框架内为湍流边界层开发详细模型。 这涉及对禁闭参数的演变进行建模，例如，即使在平滑上也能准确预测分离

\subsection{10.4.2 VC2配方与传统不连续性陡峭方案的比较}
214 涡量约束表面。 下一节将介绍一些新的探索性研究的结果。 对这些模型仍有待研究，但基本方法的能力似乎很有希望。 开始这项工作的一个简单方法是使用符合表面的网格。 然而，一个重要的一点是，这种基于VC的方法与传统的RANS方案有根本的不同，后者通常使用涡黏性（EV）类型的项，并在非常精细的网格上离散化（修改的）Navier-Stokes类型的方程，以便对时间平均速度进行建模。 与EV类型方案相比，VC的一个非常重要的特点是，它大大扩展了我们的建模能力：通常，后者的方案只能容纳EV的正值。 如果EV在空间和时间的重要区域是负的，它们往往会因数值不稳定而发散[141。 这意味着建模的BL只能直接膨胀或扩散，而不是收缩。 （当然，可以获得较慢的扩展率和更小的BL厚度，但随后需要更精细的网格和更小的EV值。） 另一方面，VC可以直接模拟收缩；与负涡黏性的传统方案不同，VC不会发散。 这非常有用，例如，在低Re：物理上与光滑表面分离时，分离层往往会过渡并迅速重新连接。 像VC这样的收缩项很容易模拟这种效应[7]。 另一点是，对于符合网格，必须将缩放系数应用于，u和E，这取决于（不同的）单元格大小[33]。 这不是问题，因为使用的是没有大宽高比的非黏性网格。 最后一点涉及VC的使用，即在不利压力梯度区域的延迟分离。 例如，众所周知，湍流 BL倾向于比层压梯度更晚（在不利压力梯度中）分离。 VC可以很容易地用于模拟这个，同样不需要非常精细的网格。 [121 10.4.2 VC2配方与传统不连续性陡峭方案的比较 在本节中，我们首先重新制定第10.3节的一维标量脉冲方程，作为速度接触不连续性的不连续性陡峭方法。 其结果类似于多年来开发的形式，以保持梯度，如接触不连续性，陡峭，并克服对流项产生的平滑。 如上所述，这些方案通常涉及一维可压缩流，并利用许多不同的“限制器”。 如果我们考虑第10.3节一维脉冲的积分（并更改符号），我们有一个仍然陡峭的传播步进函数（见图2）。 然后我们有：Eq。 （10.6）然后变成@Jn = SjVJ？ = VJ？-vJ：，（10.29）

\subsection{10.4.3 VC2配方的计算细节}
SteinhofS\& Lynn 215 a; = 0; (\{JjVY\}) (10.30) 部分求和等式。 （10.7）overj，然后我们有（对于常数c）V v,n+' = v; - - ( v;+l - v;-, ) + pdj'Vi」 - €6, a; (10.31) 以这种形式，限制类似于传统的一维陡峭方案。 然而，真正的流量从来没有一个具有陡峭梯度的单点，并且在其他地方完全恒定的属性。 一般来说，有O（1）（平滑）梯度远离不连续性。 这些平滑的梯度也由陡峭器作用，导致错误，除非使用特殊逻辑来切断陡峭器。 沿着每个坐标轴使用这种一维方案来保持对流涡核紧凑会导致同样的问题，因为速度与远离核心的半径成反比（见图3）。 然而，我们不这样做！ 重要的一点是，我们不应该使用确切的一维限制项，而应该只保留它们的基本结构--它们是速度一导数的函数。 在开发多维公式时，我们应该只使用旋转不变量。 例如，在2-D 不可压缩流动中，唯一作为速度初导数的量是：W=VXijl，（10.32）D = V.q'（10.33），其中k表示平面外方向。 但D = 0对于不可压缩流动，所以我们只有一个选择：这消除了从w + 0开始远离核心的梯度的任何问题，即使a,v和a,u在那里都是0（1）。 此外，它产生了一个更简单的配方。 最后，对于三维涡流丝，我们认为将其限制在垂直于涡旋的二维平面上就足够了，如图4所示。 通过这种方式，我们得出了动量守恒VC公式。a=a((W)。 (10.34) 10.4.3 VC2配方的计算细节 离散化的细节可能并不重要，过去不同的形式用于不同的应用。 下面只描述了一种特定的形式。 对于每个时间步骤（n），执行以下计算：步骤a：对流

216 涡量约束 与传统的不可压缩“分裂速度”方法1201一样，使用类似对流的计算来处理动量方程。 这是（10.35）+n+n +;'=;”-AtV.(q q )的离散保守版本，涉及二阶中心差。 例如，对于x-组件，（10.36）中间涡量，6'，然后通过取G'的卷曲来计算。 步骤b：限制涡量约束用于计算速度增量，使溶液放松成薄的涡流结构（在少量初始时间步骤后）：在等式中。 10.26，基于中心节点和26个最近邻居的谐波平均值用于i6': (10.37) +” q =\&At? X（pc\$'-Ei？'）。 在许多情况下，只有8个最接近的节点的更简单的公式提供了非常接近的结果；需要进一步的研究来确定所需的最小节点数。 然后用二阶差分方案将方程10.37离散化。 步骤c：压力计算 压力的计算使时间步骤IZ + 1的速度是无发散的：v.;n+l =o。 （10.39）这涉及求解泊松方程v”=-v.;”（10.40）

\section{10.5 结果}
\subsection{10.5.1 翼尖的Vortics}
Steinhoff \& Lynn 217在传统的“分裂速度”程序中。 带有交错网格的传统“盒子”方案用于计算右侧。 然后，压力p由（10.41）给出。对于在浸入网格中的近表面上计算压力的情况，使用三线性外推方法来计算实际表面上指定点的压力。 步骤d：速度更新 然后再次使用盒子方案或交错网格计算下一个时间步骤的速度：;“+I = T+Qj（10.42）这与动量方程一致（在. 步骤e：边界条件 边界条件的执行按照第10.4.1.3节所述进行处理。 对于不符合的网格，即浸没的表面，每个时间步骤在身体内的每个网格单元格上，速度都设置为零。 使用符合格栅，在表面应用一个简单的防滑条件，并具有无流动通过条件。 （对于VC1公式，向量i？ 局部定义为与表面垂直。） 10.5 结果 10.5.1 翼尖 在亚琛的Fuer Luft和Raumfahrt研究所的牵引罐中对带襟翼的机翼进行了测量，如图5所示，雷诺数为55,000。 使用激光测速法在60个完整跨度的距离上确定了脱落涡流的轨迹和强度（这在[161]中进行了描述）。 使用VC对该流量进行了初步的三维计算。 为了节省计算机内存，只计算了（对称）流量的一半，并通过将其分解为1.35跨度的横向块来半抛物线化问题。 每个块的上游边界条件取自上一个块的下游块。 与完整的三维计算相比，这只引入了可以忽略不计的误差，因为在这个距离上涡流曲率的上游影响非常小。 流量计算是使用机翼后面的测量速度启动的，仅来自半跨度的流量在对称条件下进行处理。

\subsection{10.5.2 气缸唤醒}
图6显示了218 涡量约束 涡量的震级，比较了四个下游位置的实验和计算。 对称平面在右侧，可以看到（较弱的）襟翼涡旋围绕主翼尖涡螺旋。 对禁闭参数进行了敏感测试。 两个广泛分离的涡流对这些值不敏感。 例如，在另一个尖端（图像）涡流的影响下，主翼尖涡的下降率是正确的。 由于襟翼和翼尖涡与它们的直径相比很近，预计对每个瓣的内部结构会有一些（小）影响。 这被最小化了，因为VC只允许在几个网格单元上捕获涡流。 此外，边界的接近会引起影响。 这只反映在限制参数对翼尖涡周围襟翼涡旋转角度的微小影响上。 详情见参考。 [16]。 10.5.2 气缸唤醒 本研究涉及VC的两种不同特性。 在一个方面，VC用于处理非符合性、规则笛卡尔网格中的浸没表面。 在另一个方面，VC被用作一种新型的LES模型。 记录了浸入式表面的使用，并验证了许多案例的结果。 这两种用途都是当前的研究领域。 显示基于VC2配方的三维圆形圆柱体的结果。 参考中显示了VC1配方的结果。 [12]。 （在该研究中还处理了方形圆柱体。） 两种VC配方在实验中都进行了很好的比较。 之前使用VC1公式以及这里介绍的计算，都是针对Re = 3900的计算。 使用了粗略、均匀的笛卡尔网格18 1 x 12 1 x 61，圆柱体的直径只有15个单元格的浸入边界。 参考[12]显示了使用VC1公式的涡量的计算等曲面，以及尾流速度的平均值和rms平均值。 还显示了与实验数据的比较。 在这项工作中，和这里一样，只调整了一个参数；用于模拟雷诺数效应的禁闭强度。 这在整个领域都是恒定的。 VC2配方的结果如下：图中显示了涡量大小的等表面。 7.a，其中等表面大小的值为最大值的VI。 图中显示了与沿圆柱体后线计算的平均溪流速度相对应的图。 7.b和rms溪流速度波动如图所示。 7.c。 测量的线如图所示。 7.d. 与实验的一致性被认为是好的。 如图所示，气缸表面的压力分布也与实验数据相比非常好。 7.e。 这里重要的一点是，只有通过调整一个参数E，该参数在整个场中是恒定的，计算结果与实验密切相关

\subsection{10.5.3 动态摊位}
Steinhoff和Lynn 219用于图中绘制的所有六条曲线。 7.b和7.c。 需要与不同雷诺数的实验进行额外的比较，以校准此参数的雷诺数依赖性。 必须强调的是，当E增加时产生的跨度不稳定性和混乱行为只是来自三维效应，如物理湍流，而不是由于数值不稳定性：对更广泛的E值范围进行了广泛的研究，比这里对预计没有不稳定性的流动进行了广泛的研究，这些流动只显示了稳定的流动。 这些涉及二维圆柱体上的涡流配对和孤立的、三维的翼尖涡流。 10.5.3 动态失速 该项目的主要目标是为动态失速创建一种计算方法，该计算方法在计算机资源中非常高效，以便它可以用作常规工程工具。其中一个约束是将VC制定为在现有的美国宇航局可压缩代码“TURNS”中实现的修改。 因此，与第10.4节中的描述不同，使用了可压缩的配方。 这可以在参考中找到。 [9，101。 流量是亚音速的，结果并不显著取决于这种差异。 为了简单，我们在本文中只描述了不可压缩的表述。 在Re = 1.9 x106的俯仰VR12 翼型上模拟了流动，经历了动态失速，边界层突然分离。 这导致升力迅速丧失。 这项研究的细节描述了参考[9]和[lo]。 这里的重点是，这种情况与许多钝体流程相似，分离位置和时间是边界层（BL）的重要属性。 其他数量，如皮肤摩擦，对这些流动来说是次要的。 我们的方法是定义一个基于VC的动态模型，该模型直接导致这些数量的正确行为，并且仅在少量的不黏性网格单元格上有效定义BL。 如第10.4节所述，没有像传统的湍流 BL Reynolds平均Navier-Stokes（RANS）方程方案那样，尝试对物理速度曲线的详细时间平均版本进行建模。 这项工作将需要靠近地表的非常精细的网格单元和更多的计算能力。 此外，薄的对流涡流片通常是分离的结果。 这也通过VC方法有效地捕获。 参考[9，101的初步结果显示了音高周期内几次的速度向量，如图8所示。 还显示了计算的C、值以及与风洞测量的比较。 本研究中使用的模型需要VC参数，这些参数在周期中为每个阶段的每个阶段单独规定，但在整个网格中是恒定的（除了网格缩放，除了效率之外，

\section{10.6 其他研究}
\subsection{10.6.1 导弹基地阻力计算}
\subsection{10.6.2 刀片涡旋相互作用（BVI）}
220 涡量约束因子）。 正在开发一个更先进的模型，可以动态计算这些参数。 一个重要的一点是，在标准（1.7 GHz）PC上，二维动态失速计算每个周期只需要大约10分钟。 传统的RANS方案显然需要更多的计算资源[6]。 这些快速运行对开发基于VC的模型有很大帮助。 10.6 其他研究 上述三个案例展示了使用VC来模拟选定的现象。 在过去几年获得的大量结果中，将提到另外三个结果，因为它们是最近的，并展示了在尖锐彗星（导弹基流）下对边界层分离的VC模拟、涡流传播和与翼型（叶片涡相互作用）的相互作用以及湍流的模拟和可视化（用于特效）的额外用途。 正如导言中提到的，正在准备对VC的一般应用进行全面审查。 最近的一些结果可以在http://www.flowanalysis.com上看到。 10.6.1 导弹基地阻力计算边界层导弹基地的分离是使用可压缩版本的VC计算的[6，26，371。 （参考标题。 1261使用inviscid一词仅表示需要一个inviscid大小的网格，计算时间相应很短，即使计算是黏性的。） 在这些论文中，证明VC可以捕捉分离边界层的基本特征，使其保持薄并保持其循环，同时允许在防滑条件下使用经济的粗网格，如第10.4.1节所述。 从论文中可以看出，结果与实验一致。 唯一准确的替代方案是更昂贵的RANS方法，因为传统的、经济的隐性方法会导致不准确的边界层和膨胀扇。 10.6.2 叶片涡流相互作用（BVI）VC经济地模拟BVI集中涡旋传播的能力使最近对二维BVI案例的参数研究成为可能[24，251。 这项研究还使用了涡量约束的可压缩版本。 在这些论文中，还表明计算和实验之间存在极好的一致性。

\subsection{10.6.3 湍流 特效模拟}
\section{10.7 结论}
Steinhoff \& Lynn 221 10.6.3 湍流 特效模拟 对于特效，湍流模拟的一个重要方面当然是它看起来像湍流，这意味着它包括可见的小规模效果。 当然，这本身不足以用于工程目的，但可以被认为是预处理，特别是如果小规模现象在问题中很重要。 VC1被发现比大多数其他方案更有效地、更经济地模拟小规模现象。 考虑到这一点，Ron Fedkiw进行了出色的计算和可视化[13]。 10.7 结论 涡量约束（VC）方法比以前更全面地介绍了。 虽然基本理念与传统的CFD有些不同，但与一些众所周知的计算方法（如激波捕捉）有一些共同之处。 广泛使用这些方法的类比来解释基本动机。 VC的主要目标是有效地计算复杂的高雷诺数不可压缩流，包括具有广泛分离的钝体和长距离对流的脱落涡流丝。 这些流动中几乎所有的涡流区域都是湍流。 这意味着，对于任何可行的计算，它们都必须进行建模。 流动的其余部分是非旋转的，一旦涡流分布是定义的。 此外，这些涡流区域通常非常薄。 出于这些原因，VC的基本方法是有效地对这些区域进行建模。 最有效的方法似乎是直接在计算网格上开发模型方程，而不是首先开发模型偏微分方程（pde），然后尝试在这些非常薄的区域中准确地离散化它们。 一些目标很容易实现，特别是当主流的基本特征对薄涡流区域的内部结构不敏感时。 然后，VC可以很容易地用于仅在几个网格单元上捕获这些区域并传播它们，本质上是作为在计算晶格上“存在”的非线性孤波。 具有这些特征的流动，可以在VC的当前状态下进行处理，包括与边缘和其他明确定义的位置分离的钝体。 这些配置包括复杂的几何形状，可以使用VC轻松“浸入”在均匀的笛卡尔网格中。 这些流还包括涡流丝，涡流丝可以对流，没有数值扩散，即使在任意的长时间内，并且可以在不需要特殊逻辑的情况下自动合并。

\section{10.8 致谢}
222 涡量 Conjhernent 更困难的目标涉及与光滑表面的分离，这取决于边界层的湍流状态。 这些显然需要更详细的建模，包括参数校准。 多年来，大量工人为传统的CFD计划（如RANS和LES）开发基于湍流 pde的模型付出了大量努力。 我们目前正在开始开发相应的基于VC的模型。 这些研究的初步结果，其中一些已经提出，表明可以实现非常大的计算机节约。 10.8 致谢 感谢来自多个来源的资金：主要是陆军研究办公室和陆军航空飞行动力学局，它们从一开始就支持了涡量约束的开发。 亚琛技术大学Luft und Raumfahrt研究所和田纳西大学太空研究所提供了额外的帮助。 此外，与莫菲特球场的弗兰克·卡拉多纳以及与加州大学洛杉矶分校的斯坦利·奥舍尔和其他人的多次讨论得到了认可。

\section{10.9 书目}
Steinhoff \& Lynn 10.9 参考书目 Bohr, T., Jensen, M., Paladin, G. and Vulpiani, A., Dynamical Systems Approach to 湍流, Cambridge, 1998. Boris，J.和Book，D.，“通量校正传输：I SHASTA，一种有效的流体传输算法”，计算物理学杂志，卷。 1973年11月。 布劳恩，C.，“在高攻角下对6:1椭圆体流的流动应用”硕士论文，Institut fur Luft- und Raumfahrt，RWTH Aachen，德国，2000年。 Costes，M.和Kowani，G.，“基于涡量约束的自动反扩散方法”，航空航天科学和技术，卷。 7,2003年。 Dadone，A.，Hu，G.，Grossman，B.，“更好地理解可压缩流动中的涡量约束方法”，AIAA-2001-2639。 AIAA阿纳海姆会议，2001年6月。 Dietz，W.，“涡量约束对可压缩流动的应用”，AIAA-2004-0718。 AIAA里诺会议，2004年1月。 Dietz，W.，“低雷诺数转子和螺旋桨的分析、设计和测试”，SBIR最终报告，2004年9月。 Dietz, W., Fan, M., Steinhoff, J., Wenren, Y., “涡量约束在复杂体流动预测中的应用”，AIAA- 2001-2642。 AIAA阿纳海姆会议，2001年6月。 Dietz，W.，Wang，L.，Wernen，Y.，Caradonna，F.和Steinhoff，J.，“基于CFD的动态摊位模型的发展”，AHS第60届年度论坛，马里兰州巴尔的摩，2004年6月。 Dietz, W., Wenren, Y., Wang, L., Chen, X., and Steinhoff, J., “可扩充套件空气动力学和耦合综合模块，用于预测转子机动载荷载”，SBIR最终报告，2004年5月。 Fan，M.和Steinhoff，J.，“通过涡量约束计算钝体唤醒流，”AIAA-2004-0592。 AIAA里诺会议，2004年1月。

224 涡量约束 r 161 Fan, M., Wenren, Y., Dietz, W., Xiao, M., Steinhoff, J., “使用涡量约束在粗网格上计算钝体流量”，流体工程杂志，卷。 124，不。 4，pp. 876-885，2002年12月。 Fedkiw，R.，Stam，J.和Jensen，H. W.，“烟雾的可视化”，SIGGRAPH 2001年会议记录，洛杉矶，2001年。 Ferziger，J.，“大型涡流模拟”，湍流流的模拟和建模，Ed. 通过T。 B. 加茨基，M.Y. Hussaini和J.L. 牛津的Lumley。 1996年。 Grinstein，F.和Fureby，C.，“具有通量限制方案的高Re流量的隐性大涡流模拟”，AIAA-2003-4100。 AIAA奥兰多会议，2003年6月。 Haas，S.，“使用涡量约束计算长距离尾随涡流”，硕士论文，Institut fur Luft- und Raumfahrt，德国亚琛，2003年。 Harten，A.，“计算冲击和接触不连续的人工压缩方法111，自调整混合方案”，计算数学，卷。 32，不。 142，1978年4月。 Hu, G., Grossman, B., and Steinhoff, J., “可压缩流动中涡流禁闭的数值方法”，AZAA-Journal，卷。 40，2002年10月。 Katz，J.和Plotkin，A.，低速空气动力学，剑桥大学出版社，纽约，2001年。 Kim，J.和Moin，P.，“分数步法对不可压缩Navier-Stokes 方程的应用”，《计算物理学杂志》，卷。 59，不。 2，1985年。 拉克斯，P。 D.，“守恒定律的交比系统11，”Comm。 纯Appl。 数学，卷。 10，1957年。 梅兰德，M。 V.，McWilliams，J. C.和Zabusky，N.J.，“通过丝化分离的二维涡流的轴对称化和涡量梯度增强”，流体力学杂志，卷。 178，1987年。 Misra，A.，Pullin，D.，“大涡模拟的基于涡流的应力模型”，流体物理学，卷。 9，不。 1997年8月8日。

Steinhoff \& Lynn 225 Morvant, R., Badcock, K., Barakos, G. and Richards, B., “使用可压缩涡量约束方法的气翼-涡流相互作用模拟”，第29届欧洲旋转工艺论坛，德国腓特烈港，2003年9月。 Morvant，R.，“使用计算流体动力学的刀片-涡旋相互作用噪声的调查”，博士论文，英国格拉斯哥大学，2004年2月。 Robinson，M.，“涡量约束在隐形导弹力和力矩预测中的应用”，AJAA-2004-07 17。 AIAA里诺会议，2004年1月。 Steinhoff，J.，Dietz，W.，Haas，S.，Xiao，M.，Lynn，N.和Fan，M.，「模拟流体动力学和声学中的小尺度特征是非线性孤波」，AIAA-2003-0078。 AIAA里诺会议，2003年1月。 Steinhoff，J.，Fan，M.和Wang，L.，“涡量约束-最近的结果：湍流 Wake模拟和新的保守配方”，不可压缩流动的数值模拟，Ed。 由M。 M。 哈菲兹，世界科学，2003年。 Steinhoff, J., Fan, M., Wang, L. and Dietz, W., “集中流和被动标量作为孤波的对流”，SIAM Scientific@c计算杂志，卷。 2003年12月19日。 Steinhoff，J.，Mersch，T.，Underhill，D.，Wenren，Y.和Wang，C.，“计算涡量约束：涡流主导流的非扩散欧拉方法”，UTSI预印本，1992年。 Steinhoff，J.，Mersch，T.，Wenren，Y.，“计算涡量约束：二维不可压缩流动”，理论和应用力学的发展，第十六届东南理论和应用力学会议记录，1992年。 Steinhoff，J.；Puskas，E.；Babu，S.；Wenren，Y.；Underhill，D.，“使用孤波在长距离上计算薄特征”，AIAA会议记录，第13届计算流体动力学会议，pp. 743-759，1997年。

226 涡量 Con\$nement Steinhoff，J.和Raviprakash，G.，“使用涡量约束计算刀片-涡流相互作用的Navier-Stokes”，AIAA-95-016 1。 AIAA Reno会议，1995年1月。 Steinhoff，J.，Senge，H.和Wenren，Y.，“计算涡流捕获”，UTSI预印本，1990年。 Steinhoff，J.和Underhill，D.，“T032用于“涡量约束”的修改-相互作用涡旋环计算的应用”，流体物理学，卷。 1994年8月6日。 Steinhoff，J.，Wenren，Y.和Wang，L.，“使用涡量约束分离高雷诺数不可压缩流的高效计算”，AIAA-99-33 16-CP 1999。 Suttles, T., Landrum, D., Greiner, B., and Robinson, M., “导弹空气动力学涡量约束技术的校准：第一部分-表面限制”，AIAA-2004-0719。 AIAA里诺会议，2004年1月。 Szydlowski，J.和Costes，M.，“使用RANS和DES模拟NACA0015 翼型的静态和动态摊位配置的流动”，第四届十年一次的航空力学专家会议，2004年1月。 Tannehill，J.，Anderson，D.和Pletcher，R.，计算流体动力学和传热，Taylor \& Francis，1997年。 Van Leer，B.，“走向终极保守差异计划。 11. 单调性和守恒性结合在二阶方案中，”J. 补偿。 物理，卷。 14，1974年。 Von Neumann，J.和Richtmyer，R.D.，“流体动力冲击的数值计算方法”，J. 应用程序。 物理，卷。 21，不。 3，1950年。 Wang, C., Steinhoff, J. and Wenren, Y., “涡流-固体体相互作用问题的数值涡量约束”，AIAA杂志，卷。 33，1995年8月。 Wenren, Y., Fan, M., Wang, L., and Steinhoff, J., “涡量约束在复杂体流动预测中的应用”，AIAA杂志，卷。 41，2003年5月。

Steinhoff \& Lynn 227 [44] Wenren, Y., Steinhoff, J., and Robins, R., “飞机尾随的计算”，SBIR最终报告，1995年6月。 [45] Wenren, Y., Fan, M., Dietz, W., Hu, G., Braun, C., Steinhoff, J., Grossman, B., “使用涡量约束对现实旋转船流的高效尤拉计算：近期结果的调查”，AIAA-200 1-0996。 AIAA里诺会议，2001年1月。

涡量约束 ttttt ttt t 禁闭c。 1次通过（900个时间步骤）与限制100次通过1110次通过（100个时间步骤）图1。 有和没有限制的脉冲的一维标量对流。 图2。 I-D不连续捕获方案和脉冲对流的比较。

SteinhofS \& Lynn 229 \textbackslash\{\}\textbackslash\{\}\textbackslash\{\}\textbackslash\{\}\textbackslash\{\}\textbackslash\{\}\textbackslash\{\}\textbackslash\{\}\textbackslash\{\}, 5, 'lillIfJ// 'l\textasciitilde{}lli\textasciitilde{}//\textasciitilde{} llL\textbackslash\{\}\textbackslash\{\}\textbackslash\{\}\textbackslash\{\} tVtI\textbackslash\{\}l\textbackslash\{\}! tL\textbackslash\{\}\textbackslash\{\}t\textbackslash\{\} /！ /Lflll I1 i，我\textbackslash\{\}\textbackslash\{\}！ \textbackslash\{\}\textbackslash\{\}\textbackslash\{\} /////lllll ////llllllt \%\textbackslash\{\}\textbackslash\{\}\$\textbackslash\{\}\textbackslash\{\}, ////////ll\$l. \textbackslash\{\}\textbackslash\{\},\textbackslash\{\}\textbackslash\{\}\textbackslash\{\}, /////1/111,11. \textbackslash\{\}\textbackslash\{\}\textbackslash\{\}.,\textbackslash\{\}\textbackslash\{\}I。 C/,\#/fl/l,,i \textbackslash\{\}\textbackslash\{\}\textbackslash\{\}\textbackslash\{\}\textbackslash\{\}\textbackslash\{\}\textbackslash\{\} 'r,,,/lll,,l' \textbackslash\{\}\textbackslash\{\}\textbackslash\{\}\textbackslash\{\}\textbackslash\{\}\textbackslash\{\}, d,1///1/11J1..*\textasciitilde{}/r/,ll,,/l.. 我\textbackslash\{\}.----- --. *.,,,.,..... - -cccIc, CF-r.... C，，，，，图3。 涡流周围的速度向量场显示了选择旋转不变的陡峭器的基础。 正常平面图4。 三维禁闭操作员

230 涡量约束图5。 拖水罐实验的示意图。 两个涡流从机翼的两侧脱落，一个从翼尖脱落，一个从襟翼内侧脱落。

Steinhoff \& Lynn U omegs 0 0175 -0 w l""l" '1""II Y 图6.a 4.05翼展后的涡量分布。 左：实验数据。 右：计算结果。 计算中的涡流倾向于比实验中稍微快一点向左弯曲。ompga 0.0175 -0 w Y图6.b 5.4翼展后的涡量分布。 左：实验数据。 右：计算结果。

涡量约束 Y图6.c 9.45翼展后的涡量分布。 涡流的旋转角度仍然与实验数据一致。 左：实验数据。 右：计算结果。 N 46 23 43 ec Y J I I 图 6.d 涡量 10.8翼展后的分布。 左：实验数据。 右：计算结果。

Steinhoff \& Lynn 233 图7.a 涡量 涡量约束的气缸流的等表面（，U = 0.15，\textasciitilde{} = 0.325）

234 涡量约束图7.b平均流向速度曲线。 符号是实验数据（，U = 0.15，\textasciitilde{} = 0.325）0.3 0.1 02i 0 -\_\_-图7.c流式雷诺应力。 符号是实验数据（，U = 0.15，\textasciitilde{} = 0.325）

Steinhoff \& Lynn 235 a xlD= 1.06 1.54 2.02 图7.d 圆形圆柱的测量位置 0 20 40 60 80 100 120 140 160 180 -1.5 e 图 7.e 圆柱表面的时间平均 压力系数 分布。 圆圈表示Norberg的实验数据（[141，1987年）

涡量约束 ' i ' 'Ok ' is' ' 0;s, ' '; 图 8.a 0" 攻角的动态停摊位结果。 左：压力分布。 右：流动场向量。 I/5 图8.b 10英寸攻角的动态摊位结果。

Steinhoff \& Lynn 237 - CPOATA。 I* a9，n.m I wo：- XIC图8.d 20'攻角的动态停摊位结果。 左：压力分布。 右：流动场向量。 图8.e 17英寸攻角的动态摊位结果。 \textasciitilde{} GPVC f\_ CPDATA 0 026 06 076 1 图8.f 10英寸攻角的动态摊位结果。

此页面故意留空

\clearpage
\part{三。 流量稳定性和控制}
第111部分 流量稳定和控制

此页面故意留空

\clearpage
\chapter{第11章 合成和脉冲喷射器R的流量控制应用。 阿加瓦尔，J。 瓦迪洛，Y。 谭，J。 崔，D。 郭，H。 耆那教 W。 卡里和W。 W。 树荫处}
\section{11.1 摘要}
第11章 合成和脉冲喷射器R的流量控制应用。 Agarwal，'J。 瓦迪洛，'Y。 谭，'J。 崔，'D。 郭，'H。 耆那教 W。 卡里，2和西。 W。 Bower2 11.1 摘要 近年来，使用合成和脉冲喷射执行器控制壁界流和自由剪切流的有前途的方法备受关注。 许多研究人员通过实验取得了各种令人印象深刻的流量控制结果，包括传统推进喷射器的矢量化、悬崖体的空气动力学特性的修改、翼型的升力和阻力控制、平板边界层的皮肤摩擦的减少、圆形喷射器的增强混合，以及外部和内部流动分离和腔体振荡的控制。 最近，有人试图用数值模拟其中一些流场。 上述几个流场已使用湍流模型的雷诺平均Navier-Stokes（RANS）方程进行数值模拟，少数使用直接数值模拟（DNS）进行有限。 在模拟中，喷射出口处的简化边界条件以及腔体和唇部的细节都包括在内。 在本文中，我们描述了五个不同流场的模拟结果，涉及虚拟空气塑造，推力-华盛顿大学机械和航空航天工程系，圣。 路易斯，'波音公司，圣。 路易斯，密苏里州63166。 密苏里州63130。

\section{11.2 介绍}
242 合成和脉冲喷射矢量的流量控制，合成喷射与湍流边界层的相互作用以及分离和腔体振荡的控制。 这些模拟是使用RANS方程与单方程或双方程湍流模型一起进行的。 11.2 导言 近年来，人们非常重视先进的空气动力学技术的发展，该技术基于空气动力学流体力的流体修改，无需传统的控制面，如襟翼、扰流板和可变机翼扫荡，即可覆盖多种飞行条件。 流体改性（或流量控制）可以通过使用由智能控制系统动态操作的微表面效应器和流体装置乳化器来实现。 这些新的主动流量控制（AFT）技术有可能大幅提高飞机性能并减轻重量。 许多研究人员使用脉冲或合成（振荡）喷射执行器通过实验取得了各种令人印象深刻的流量控制结果，包括传统推进喷射器的矢量化、虚张声势体的空气动力学特性的修改、机翼的升力和阻力控制、减少平板边界层的皮肤摩擦、增强圆形喷射器的混合以及外部和内部流动分离和腔体振荡的控制。 合成射流是通过在腔内使用振荡表面形成的。 它完全由被控制的流体产生。 它以周期性方式用压电隔膜产生。 流动周期性地通过孔进入和离开腔体。 合成喷射器的独特之处在于不需要流体管道。 合成射流已被证明在许多应用中具有重要的控制权，并具有紧凑且净质量通量为零的额外好处。 Glezer和Amitay最近对合成喷气式飞机及其一些应用进行了出色的审查[5]。 另一方面，脉冲喷射具有稳定和振荡速度分量，净质量流量和动量注入流量。 本文介绍了使用合成或脉冲喷射器对五个不同流场进行流量控制的数值模拟结果，涉及虚拟空气形成、推进喷射的推力矢量、合成喷射与湍流边界层的相互作用、向后台阶后面的再循环区域的控制以及腔体振荡的控制。 足够的

\section{11.3 CFD 采用的求解器}
R。 Agarwal等人。 尽可能将243个数值模拟与实验数据进行比较。 11.3 CFD 采用的求流器 在本文第3节报告的五个流场模拟中使用了两种不同的求解器。 两者都是众所周知的RANS求解器，称为INS2D [ 111和WIND [2]。 “INS2D'是美国宇航局艾姆斯研究中心开发的不可压缩雷诺平均纳维尔斯托克斯（RANS）求解器。 它用于第3节（b）、（c）和（d）部分中报告的模拟，分别对应于使用合成喷射的推进喷射的矢量化，使用合成喷射器向后一步的分离控制，以及合成喷射与湍流边界层的相互作用。 INS2D使用伪可压缩性方法求解广义坐标中稳态和时变流动的二维不可压缩流动的连续性方程和RANS方程。 对流项使用基于通量差分分裂的上风三阶精确微分方案计算。 差分方程使用广义最小残差法（GMRES）求解，该方法稳定，能够以非常大的伪时间步进运行，导致每个物理时间步快收敛。 该代码有两个用于计算涡黏性的湍流模型--Spalart-Allmaras（SA）一方程模型[141和门特剪切应力传输（SST）双方程模型[9]。 SA模型用于使用r”NS2D执行的所有计算，在第3节（b）、（c）部分中报告，第3节（a）部分报告的计算，使用合成喷射对翼型进行虚拟空气动力学形状修改，以及（e）部分使用脉冲吹动控制亚音速腔剪切层，采用多区结构网格可压缩RANS求解器WIND。 WIND [2]已经开发出来，现在得到了NPARC联盟的支持，该联盟是美国宇航局格伦研究中心（GRC）、波音、普拉特和惠特尼和美国空军阿诺德工程开发中心（AEDC）之间的伙伴关系，该中心是为了提供面向应用的国家流动模拟代码而成立的。 WIND还有几个湍流模型--SA模型、SST模型和大型涡流模拟（LES）模型。 最近，Bush和Mani [3]修改了该代码，包括Spalart[15]的独立涡流模拟（DES）公式，用于SA模型，并由Strelet[16]将其扩展到k-E模型。 此扩展允许在流域中以一致的方式组合RANS和LES模型。 虽然WIND是一个3D代码，但第3节（a）和（e）部分中报告的模拟是2D的。 （4.

\section{11.4 结果和讨论}
244 合成和脉冲喷射的流量控制 在除第3（b）和第3（d）节外的所有模拟中，执行器的边界条件应用于喷射器的出口处；不包括执行器腔的影响。 在模拟合成喷射与第3（d）节中描述的湍流边界层的相互作用中，包括了执行器腔的影响。 以下边界条件用于腔体底部：（11.1）（11.2）v（x，y = const，t）= A sin0 t u（x，y = const，t）= 0（11.3）在Eq.（ll.l）中，A=UJ b/W，其中UJ是合成射流的速度振幅，W和b分别是腔体和射流孔的宽度。 11.4 结果和讨论 在本部分，我们展示了合成和脉冲喷射对五个不同流场的影响的典型选择结果。 对于每个案例，详细的结果都可以在引用的参考资料中找到。 （a）使用合成喷射执行器fXJ对翼型进行虚拟空气动力学形状的修改 最近，实验表明，通过修改翼型的表观空气动力学形状，翼型在低攻角下的压力阻力可以显著降低，而升力变化最小。 这种虚拟的空气动力学形状修改可以通过在翼型表面上部（前缘下游）附近创建一个小域来实现，该域足以取代局部流线，以改变局部压力分布。 在最近的一次实验[I]Glezer和他的同事们故意在24\%厚的Clark-Y 翼型的上表面附近创造了这样一个相互作用域，使用一个合成喷射执行器，紧挨着表面安装的小尺寸的被动障碍物的下游。 我们对实验观察到的翼型压力分布的流体修饰进行了数值模拟，导致减压阻力[181。 为了执行计算，使用了CFD求解器WIND。 对亚声速流动进行了计算，经过24\%厚的Clark-Y 翼型，上表面有和没有合成喷射的三角形凸起。 计算和实验数据之间取得了合理的一致性。 图1显示了翼型周围的双区网格，带有凸起和合成喷射器。

R。 Agarwal等人。 245 图2显示了计算压力分布与基线Clark-Y 翼型（无颠簸）在0.1马赫、弦雷诺数 Re、=381,000和攻角a=3"的实验数据的比较。 图3显示了计算压力分布与Clark-Y 翼型的实验数据的典型比较，其凸起和合成喷射频率为f=850Hz，动量系数Cp =1.2e-3；得到了合理的一致性。 参考[181提供了为实现最大阻力降低和最小升力变化而进行的一系列计算的广泛细节。图4和图5分别显示了升力和阻力系数与振荡喷射频率的变化。 这些数字表明，在更高的射流频率下，可以显著降低阻力，而升力变化很小。 I I 图1：围绕24\%厚的Clark-Y 翼型的双区网格，带有凸起图2：Clark-Y 翼型的计算和实验压力分布；M=0.085，Re=381,000，a=3.0”。

5.00E02 4.50E-02 4.00802 3.50E-02 WOE-02 3 2.50802 2.00c02 - 1.50802 - 5. QQE-03 - o.ooE+oo > 使用合成和脉冲喷射器进行流量控制................................................ -使用Burrp和SJ计算基线热量的空气箔计算Orag系数- -。 A- -。 仅使用BUW的翼型计算hag系数，，，I I，，I -0.2 0.5 1.5 J图3：Clark-Y 翼型的计算和实验压力分布，在x/c=.22处安装一个凸起和合成喷射；M=0.085，Re=3.8e5，a=3.0”，e850 Hz，CP=1.2e-3。 7.00801 5.00801 4. WE-01 2. WE-01,,................................................. ---c--ConWLedUftCosfflcientforAirtalw)hBurrpandSI ---. I---“CD\textasciitilde{}UtBdLdtCDefficienlforAirfolwi\textasciitilde{}Bumpanly comwted LiftCosffElentforeasellne CaSB o.ooE+oo -,,,,,,,, 0.0 200.0 4W.O W.0 800.0 1WO.O 1200.0 1400.0 1600.0 频率（W）图5：计算压力阻力系数作为带有凸起和合成喷射的Clark-Y 翼型的执行频率的函数；M=0.085，Re=38 1,000，a=3.0°，Cp=l.2e-3。

R。 Agarwal等人。 24 7（b）用合成喷射器对初级喷射器进行矢量控制 最近实验表明[131，从喷嘴出口的亚音速初级喷射器的推力矢量控制可以通过单个或多个合成喷射器执行器完成。 我们进行了广泛的数值模拟[6]，以量化合成射流的各种参数，即宽度、速度振幅、速度、频率、相对于主射流的位置、角度和数字，对主射流的矢量角度的影响。 这些模拟提供了导致这些参数的最佳值的信息，以实现最大的矢量效益。 这些模拟也有助于澄清负责矢量的物理机制。 数值分析的结果被用于确定控制F-18飞机推力矢量的合成喷气装置的要求（这些结果未发表，不能包括在这里）。 在这里，我们展示了一个典型的模拟结果；参考[6]中给出了额外的模拟和其他细节。 图6显示了数值研究中使用的典型配置（包括初级喷射和合成喷射）的示意图。 它还展示了Smith和Glezer的实验中获得的非强制主射流（无合成射流）和强制主射流（带合成射流）的Schlieren照片[131。 对于给定的一组初级喷射器和合成喷射的流动参数，在实验中观察到初级喷射的30英寸矢量，如图6所示。 我们的图6：主射流（PJ）与合成射流（SJ）的矢量化；Uave=5.16 ds，ReuO=383，el120 Hz，dpull/h=3.556

248 使用合成和脉冲Jc进行PJ f5rced PJ（带腔）强迫PJ（无腔）图7：用合成喷射（SJ）的初级喷射（PJ）矢量化计算；Uave=5.16 m/s，ReuO=383，el120 Hz，dpull/h=3.556 dpuII = PJ和SJ h之间的距离= SJ H的宽度= PJ U的宽度，，= SJ U的速度振幅，，= PJ Re\% = SJ f的雷诺数的速度 = SJ计算的频率非常详细地再现了Smith和Glezer的实验结果。 图7显示了数值模拟的典型结果；应该注意的是，合成喷射的腔体的包含对于从数值模拟中获得“正确”的结果非常重要。 图8-10分别显示了合成射流的三个参数，即dpull=合成射流和初级射流之间的距离，f=频率和Uamp=速度振幅，对初级射流的矢量角度的影响。 保持其中两个参数固定，剩余参数的最佳值可以提供最大的向量效益。 这些计算是通过修改不可压缩的雷诺平均纳维尔斯托克斯求解器INS2D来执行的。

R。 Aganual等人。 249 图8：主射流矢量角与dpun的变化；Reu,,=383, U,=7.0 ds, el120 Hz 图9：主射流矢量角与f的变化；Reuo=383，Uav，=7.0 d\textasciitilde{}，dp，l+h=3.556

250 用合成和脉冲喷射器控制流量 图10：主喷射器矢量角的变化与U,,,; U,,=7.0 m/s, (c) 再循环控制！ 向后的流动区域；步进el120 Hz,d,ll/h=3.556使用合成喷射器最近，Rediniotis等人a1 [lo]的实验表明，通过在步进上使用合成喷射，可以显著减少通道中反向步进后面的再循环流区域。 我们进行了一项数值研究[S]，以调查单个和多个合成喷射器对台阶后面分离的流动区域的影响。 数值模拟使我们能够进行广泛的参数研究，其结果记录在参考[S]中。 图11显示了Rediniotis等人1 [lo]在实验中使用的实验配置和流动参数的典型计算模型的结果。 计算出的速度曲线与实验曲线相当好。 更重要的是，如表1所示，计算和实验的重新连接长度非常一致。 这些计算是使用双区网格与Reynolds-Averaged Navier-Stokes求解器INS2D进行的。 （d）合成射流与平板湍流边界Laver r41的相互作用 最近，Honohan [7]通过实验研究了合成射流与湍流边界层的相互作用。 他考虑了两组不同的

R。 Agarwal等人a）没有S\textasciitilde{}IC kt x（=4 c）的Stcadv流的Vomiritv Contoras，没有SJ的峰值吸力行程的Vodcity轮廓d）Stuady Flow wtthout SJ x tfm的速度Vacton c）SJ x（毫米）Pal吹风的速度向量9个Pel吸力行程的速度向量图11：涡量轮廓和速度曲线，在台阶的垂直面上使用合成喷射（SJ）；平均流入速度=lO cds，速度振幅为SJ=40 cds，频率为SJ=20赫兹，SJ宽度=2 111111，合成喷射距离底壁6至8毫米

252 使用合成和脉冲喷射器的流量控制 流量条件 没有合成喷射的稳定流量 表1的峰值吹气冲程：图11各种流量条件的重新连接长度 实验性重新连接长度 8小时 7.95小时 1.6小时 1.6小时 合成喷射器 合成喷射器的峰值吸力行程0.6小时 0.6小时 湍流边界层和合成喷射器参数（如表2所示），因此在第一种情况下，合成喷射器的离散涡流与整个湍流边界层相称，在第二种情况下，它们与接近的内部黏性子层缩放 湍流边界层。 在本文中，我们用案例1的实验数据介绍了数值模拟的一些结果。 本案例和案例2结果的更详细结果见参考[4]。 U，案例（ds）1 8.0 2 10.0 表2：Honohan16 6 0 UJ f LO（mm）（mm）（mh）（Hz）（mm）12.7 1.4 31.4 300 33 31.0 3.2 39.3 1130 11 Case Sr Re C，h/b W/b H/b（mm），057 1214.23 1.91 4.9 21 38,046 406.026 0.51 13 6.3 94 '\# =（pJbuO2）'（pcomco2）Sr =2 A f bl/2 / UJ = b/ Lo 首先，计算基线案例（无喷射）。 图12显示了从实验、当前计算和七次幂定律中获得的合成射流位置的顺流速度曲线。 在模拟中，这个速度曲线被强加在左边界上的边界条件。 计算的速度曲线与靠近墙壁的实验非常一致。 然而，区域0.2 < yl6 < 0.5有细微差异。 原因

R。 Agarwal等人。 253 彻底调查了这一差异。 它与网格的细度、数值方案的顺序或湍流模型无关。 在合成喷射位置和合成喷射下游的计算和实验中，边界层厚度6是相同的。 总体而言，如图13所示，平板湍流边界层的计算和实验速度向量是一致的。 图15显示了在合成喷射的一个时间段（=2n）内每6O0（=xI3弧度）的速度向量和涡量场的模拟及其与实验结果的比较。 总体而言，模拟和实验结果非常一致，但观察到了显著差异，特别是在靠近合成喷射的下游区域的速度向量中。 这些差异在@=O、n/3和5x13时更加明显；即在“吹”周期的早期和“吸”周期的后期。 然而，除了@=O和x/3之外，与实验数据相比，模拟中很好地捕捉到了合成喷射的离散涡流的产生、它们在“吹”期间的演变和下游对流。 与计算相比，合成喷射下游涡流的强度在实验中要强得多。 在吸力循环期间，腔内形成一对离散的涡流。 同样，在@=5n/3和2n（与+O相同）时，在合成喷射下游实验中观察到的涡流的强度比计算中强得多。 基于这些结果，可以得出结论，模拟在大多数情况下很好地捕捉到了合成喷射与湍流边界层相互作用的物理，然而，与实验的显著定量差异仍然存在，特别是在“吹”周期的早期和“吸”周期的后期，在喷射的下游附近。 解决这些差异需要对实验数据和模拟进行更仔细的检查和评估。 正如参考资料[4]所报告的那样，案例2也获得了类似的结果。 尽管如此，尽管模拟和实验结果之间存在这些差异，但模拟和实验都很好地捕捉到了合成喷射对交叉流湍流边界层的影响的性质。 合成喷射改变了壁附近的压力梯度，导致喷射下游的涡量厚度大幅减少（与非强迫流动相比），并由于相互作用域的位移效应导致的有利压力梯度的存在而减少了湍流边界层的厚度。

254 使用合成和脉冲喷射的流量控制（e）使用脉冲吹塑控制亚音速腔剪切层 我们使用多区流动求解器 WIND在二维腔体上对亚音速层和湍流进行了时间准确的数值模拟[2]。 为了验证这些类型计算的代码，使用了Rowley等人[12]论文中报告的测试用例。 图16显示了一个周期腔体振荡的四个不同时间阶段的涡量轮廓。 这些对涡量的计算以及时间平均的溪流剖面（此处未显示）与Rowley等人的计算非常一致。 为了打破腔体上的剪切层结构，我们在腔体的前缘采用了垂直于溪流方向的脉冲吹风。 如图17所示，脉冲吹塑的数值结果表明，在打破剪切层结构方面取得了显著成功。 计算中的脉冲射流马赫数是Mj= 0.2 + O.lsin（0.902t），其中02 =780Hz是第二个罗西特模式的频率。 图18显示了有和没有脉冲吹的腔体下游边缘的冲击压力。 很明显，脉冲吹风显著降低了压力振荡。 参考0.8 0.6 - 0.4 0.2 - 图12：案例1中合成喷射位置的流向速度曲线给出了此案例以及（2x1）和（15x1）腔体的详细流场解决方案

R。 Agarwal等人。 模拟速度向量实验涡量轮廓和速度模拟涡量轮廓图13：案例1中没有SJ的涡量轮廓和速度向量图14：案例1的整个腔中模拟涡量轮廓在@=O

256 使用合成和脉冲喷射器的流量控制（a. 1）模拟速度向量（a.2）实验涡量轮廓和速度-1 0 1 a 3 4（a.3）模拟涡量轮廓（a> o=o图15：案例1中的锁相涡量轮廓和速度向量

R。 Aganual等人。 25 7（b. 1）模拟速度向量（b.2）实验涡量轮廓和速度-1 0 1 2 3 4（b.3）模拟涡量轮廓（b）@=d3（c. 1）模拟速度向量（c.2）实验涡量轮廓和速度

用合成和脉冲喷射器控制流量 -1 0 1 2 3 4 (c.3) 模拟涡量轮廓 (c) @=2r1/3 (d. 1）模拟速度向量（d.2）实验涡量轮廓和速度-1 0 1 2 3 4（d.3）模拟涡量轮廓（4 @=n

R。 Agarwal等人。 259（e. 1）模拟Velocitv矢量（e.2）实验性涡量轮廓和速度-1 0 1 2 3 4（e.3）模拟涡量轮廓（e）\$=4z/3（f. 1）模拟速度向量（f.2）实验涡量轮廓和速度

\section{11.5 结论}
\section{11.6 致谢}
260 合成和脉冲射流的流量控制-1 0 1 2 3 4（f.3）模拟涡量轮廓图1 S（续）：案例1中的锁相涡量轮廓和速度向量（x和y由6归一化，涡量刻度如图14所示）（f）+=5n/3 11.5结论在本文中，使用合成和脉冲射流对五个不同的流场进行流量控制的一些结果，涉及翼型的虚拟空气形成、推进喷射的推力矢量、合成喷射与平板湍流边界层的相互作用、向后台阶后的再循环区域的控制和 已经提出了腔体振荡的控制。 这些应用表明，合成和脉冲喷射可以用作有效的流量控制设备。 模拟和实验在捕捉流场的整体特征方面有合理的一致性。 本文中提出的所有五个案例都是低速流。 需要额外的实验数据和数值模拟来涵盖实际应用的现实流动配置。 11.6 致谢 波音-St.提供的财政支持。 路易斯到J。 瓦迪洛，Y。 Tan和D。 Guo得到了感激的认可。 作者还想感谢教授。 Glezer，博士 Amitay和博士 佐治亚理工学院的Honohan，提供了关于病例3（a）、3（b）和3（d）的实验信息，以及许多有用的讨论。

R。 Aganval等人。 261 没有脉冲吹---。 *.- -2\% C) 'ikmiesR. contoun a, 1.4 I..............................................'-"Ix* I). 图16：非定常流动在（4x1）腔上的涡量轮廓；M，=0.6，Ree=58 3。 脉冲吹---，prrG:"-- -v*。 Z.-*图17：涡量 非定常流动在脉冲喷射（4x1）腔上的轮廓；M，=0.6，Re，5 8.8。

\section{11.7 参考书目}
262 使用合成和脉冲喷射器的流量控制图18：无脉冲吹气（11.7参考书目Amitay，M.，Horvath，M.，M.，Michaux，和Glezer，A，“使用合成喷射执行器在低攻击角度的虚拟空气动力学形状修改”，AIAA论文01-2975，2001。 Bush，R.H.，“NPARC联盟的流动求解器生产”，AIAA Paper 88-0935，1988年。 布什，R. H.和Mani，M.，“高亚网格剪切的双方程大涡流应力模型”，AIAA论文2001-2561，2001。 Cui，J.，Agarwal，R.，Cary，A. W.，“合成喷气式飞机与湍流边界层相互作用的数值模拟”，AIAA论文2003-3458，第33届AIAA流体动力学会议和展览，佛罗里达州奥兰多，2003年6月23日至26日。 Glezer，A.和Amitay，M.，“合成喷气式飞机”，Annu。 牧师。 流体机械，卷。 34，Guo，D.，Cary，A.W.和Agarwal，R.K.，“用合成喷射器对初级喷射的矢量控制的数值模拟”，AIAA论文2002-3284，第一AIAA控制会议，St. 2002年6月24日至27日，密苏里州路易斯。 Honohan，A.M.，“合成喷射与交叉流动的相互作用和空气动力学表面的修改”，博士。 论文，乔治亚理工学院，2003年5月。 Jain，H.，Agarwal，R.K.和Cary，A. W.，“2002年合成喷射对反向再循环流动影响的数值模拟，第503-529页。

R。 Agarwal等人。 263步，”AIAA论文2003-1125，AIAA第4届航空航天科学会议和展览，内华达州里诺，2003年1月6日至9日。 9. Menter，F.R.，“空气动力学流动的Zional Two-Equation k-o 湍流模型”，AIAA Paper 93-2906，1993年。 10. Rediniotis，O.K.，KO，J.和Yue，X.，“合成喷气式飞机，其简化顺序建模和流量控制应用”，AIAA论文99-1000，37”航空航天科学会议和展览，内华达州里诺，1月。 1999年。 11. Rogers，S.E.和Kwak，D.，“时间精确不可压缩的Navier-Stokes方程的迎风差分方案”，AIAA杂志，卷。 28，12。 Rowley，C.W.，Colonius，T.和Basu，A.J.，“关于矩形腔流的自持续振荡”，J. 流体机械，卷。 455，2002年，pp. 315-346。 13. Smith，B.L.和Glezer，A.，“使用合成喷射执行器在自由剪切流中影响的矢量和小规模运动”，AIAA论文97-0213，35”AIAA航空航天科学会议，内华达州里诺，1997年1月。 14. 斯帕拉特，P。 R.和Allmaras，S. R.，“空气动力学流动的单方程湍流模型”，La Recherche Aerospatiale，卷。 1，1994年，pp. 5-21。 15. 斯帕拉特，P。 R.，Jou，W. H.，Strelets，M.，Allmaras，S. R.，「关于LES翼的可行性和混合RANSILES方法的评论」，lst AFOSR Int. 1997年8月4日至8日，洛杉矶拉斯顿，DNSLES会议。 16. Strelets，M.，“大规模分离流动的独立涡流模拟”，AIAA论文2001-0879，39\textasciitilde{}AIAA航空航天科学会议和展览，内华达州里诺，8日至1月1日。 200 1. 17. Tan，Y.，Agarwal，R. K.，鲍尔，W. W.和Cary，A. W.，“使用脉冲喷射和空气光学分析对空腔体的剪切层的流量控制”，AIAA论文2004-0928，AIAA 42”d航空航天科学会议，内华达州里诺，2004年1月5日至8日。 18. Vadillo，J.L.和Agarwal，R.K.，“使用合成喷射执行器对翼型的虚拟空气动力学形状修改的数值研究”，AIAA论文2003-4158，331d AIAA流体力学会议，佛罗里达州奥兰多，2003年6月23日至26日。 不。 2，1990年，pp. 253-262。

此页面故意留空

\clearpage
\chapter{第12章 使用电磁场在圆形圆柱体上控制流动分离：数值模拟 Brian H. Dennis和George S. 杜利克拉维奇}
\section{12.1 命名法}
第12章 使用电磁场在圆形圆柱体上控制流动分离：数值模拟 Brian H. Dennis'和George S. Dulikravich2 12.1 命名磁通密度，kg A-' s-'恒压下的比热，K-l m2 s-\textasciitilde{}平均变形张量率，s-l材料导数，s-l电位移向量，A s rn-'电场强度，kg m s-\textasciitilde{} A-'电动势强度，kg rn sP3 A-l磁场强度，A rn-l电流密度，A rn-2 - E - k=\&+gxB - H - I单位张量- J = - J，+凝胶'德克萨斯大学阿灵顿分校机械和航空航天工程，德克萨斯州阿灵顿，76019 '机械和材料工程，佛罗里达国际大学，10555 W。 美国佛罗里达州迈阿密弗拉格勒街 33174

\section{12.2 介绍}
266 流动分离 电磁场控制电传导电流，每单位体积A m-'总磁化，A m-l磁力强度，A m-l压力，每单位体积kg m-1 sp2总极化，每单位体积A s mP2总电荷，A s mP3传导热通量，每单位体积kg s-\textasciitilde{}热源，kg m-l s-\textasciitilde{}时间，s柯西应力张量，kg m-l sp2绝对温度，每单位质量K内能，m2 s-'流体速度，m s-'电介电常数，kg-' mP3 s4 A'真空介电常数，kg-' m-3 s4 A2极化电常量，kg-' mP3 s4 A2相对 电介电常数流体黏度，kgm-ls-l电势，Vm改性静压，kgm-sP2流体密度，kg mP3应力张力的偏转部分，kg m-l sP2磁渗透性，kg m A-' s-\textasciitilde{}真空磁导率，kg mA-2s-2相对磁渗透性磁磁渗透性，kg mA-2s-2基于B电易感性涡量，s-l 12.2导言近年来，人们对耦合物理或多学科现象的模拟越来越感兴趣。 计算机处理器技术的进步最近使研究人员能够考虑代表复杂耦合问题的大型微分方程组。 多学科分析的一个例子是在外部施加的电磁场影响下模拟流体流动。 在外部施加的电场和可忽略的磁场的影响下，含有电荷的流体流的研究被称为电流体力学或EHD。 研究没有电荷且仅受外部影响的流体流动

B. H。 丹尼斯和G。 S。 Dulikruvich 267磁场被称为磁流体动力学或MHD 1181。 有许多关于EHD和MHD模型的出版物[19，131，它们的数值模拟和应用[6，4，3，5，7，17，9，8，11 虽然相当复杂，但现有的EHD和MHD数学模型往往代表了实际物理学的简化和不一致的模型[ll]。 在外部施加和内部产生的电场和磁场的综合影响下对流体流动的研究通常被称为电磁流体动力学（EMFD）[12，lo]。 然而，这种组合电磁场与流体流动相互作用的数学模型非常复杂，需要大量在公开文献中找不到的流体的新物理属性。 因此，在外部施加的电场和磁场的综合影响下，流体流动的实际数值模拟中应该使用一个稍微简化的数学模型。 在不可压缩流体的情况下，KO和Dulikravich[15]推导出了这种称为二阶电磁流体力学（EMHD）的非线性模型。 这是一个二阶理论，与完整的EMFD模型的所有基本假设完全一致[12，lo]。 基本假设是，电场和磁场、应变率和温度梯度相对较小。 此外，与传统的牛顿流体一样，平均变形张量中二阶和更高的项被忽视了。 仅保留d、E、B、VT中二阶的术语。 由于直到最近还没有完整的EMHD模型，而且由于EMHD模型的更简单版本相当复杂，仍然很难找到涉及电场和磁场以及流体流动的综合影响的出版物。 本文的目的是提出受电场和磁场组合影响的圆形圆柱体流动的数值结果。 这里提出的结果表明，电磁场可用于消除稳定流动中的流动分离。 此外，结果还表明，在非定常流动的情况下，电磁场可用于消除周期性涡流脱落。 这些模拟是使用简化的EMHD模型进行的，用于电传导不可压缩流体的二维平面流的情况。 方程用最小二乘有限元法（LSFEM）离散化，并在单个处理器工作站上解决。 将呈现同构牛顿流体的稳态和非稳态层流的数值结果。 数值方法的准确性也根据磁流体动力学的简单分析解决方案进行了验证。 应该指出的是，德国的一个研究小组对圆柱体周围流场的类似影响进行了数值预测和实验验证[20，211。 然而，它们对磁铁和电极的排列与本页中呈现的排列完全不同，这表明磁铁和电极有多种配置

\section{12.3 EMHD的二阶分析模型}
268 流动分离 电磁场控制，能够产生相同的流场变化。 12.3 EMHD的二阶分析模型 本节总结了在电磁力[15]的综合影响下控制不可压缩流动的完整偏微分方程系统，使用通过二阶理论得出的构成方程。 具体来说，极化和磁化向量被定义为（12.1），它表示具有纯瞬时响应的介质[16]。 应力张量的偏差部分被定义为电流传导向量被定义为，而热传导通量被定义为然后，麦克斯韦方程变为（12.5）（12.6）（12.7）

\section{12.4 最小二乘有限元法}
B. H。 丹尼斯和G。 S。 Dulikravich，而Navier-Stokes 方程变成（12.9）+ \&.---.- D（E\& B DB Dt pm Dt'请注意，在这个EMHD模型中，不可压缩流体的物理特性，Xei XB，Pv，01，02，04，05，07，61，62，64，65，67，K8，Q，可以是恒定值或仅温度函数。 12.4 最小二乘有限元法 本节12.3中描述的偏微分方程组使用最小二乘有限元法（LSFEM）离散化。 我们首先查看LSFEM的一般线性一阶系统[14, 21 - - Lu=f ( 12.12)，其中（12.13）- =ax =ay =用于二维问题。 系统的残差由H表示。 R(a) = \& - f (12.14) d d L = A1 - + A2- + A3 -- - 我们现在定义了域R 1(a) = H(aIT上的以下最小二乘函数I。 R（u）dx dy（12.15）

\subsection{12.4.1 简化EMHD的非维一阶订单}
270 流动分离 电磁场控制 然后通过取I的变异并将结果设置为零来获得弱声明。 （1 2.16）使用等阶形状函数，对于所有未知数，向量写成（12.17）T u = \textasciitilde{}\&\{Ul,U2,U3:..., i=l，其中\{u1,u2,ug,...,\textasciitilde{},\}\textasciitilde{}是有限元第i节点的节点值。 将上述g的近似引入弱语句会导致一个线性代数方程组m=F-（12.18），其中K是刚度矩阵，u是未知数的向量，向量是力-12.4.1简化EMHD的非维一阶形式描述EMHD流的完整偏微分方程组包含许多参数，这些参数指的是目前未知的物理属性。 我们选择只使用已知材料特性的项，而不是用这些参数的猜测值进行完整的数值模拟。 在这种情况下，我们通过只保留包含ril和01的源项来简化方程，因为这些值适用于各种流体。 由于对近似函数施加了更高的连续性限制，将LSFEM用于包含高阶导数的方程组通常很困难。 出于这个原因，在应用LSFEM之前，将系统转换为等效的一阶表格会更方便。 在电磁流体力学的情况下，二阶导数通过引入涡量，g作为额外的未知物进行转换。 此外，我们假设流动是不稳定的，但等温的，没有带电粒子。 在这种情况下，不需要能量和电荷传输方程，与电荷相关的源项被丢弃。 我们还只考虑静电场和静磁场。

B. H。 丹尼斯和G。 S。 Dulikravich 2 71 v ' g* = 0 (12.19) a?' 在+g*。vg* + \&v x g* +vp* - gg* x B* x B* -M\&* x B* = 0 (12.20) g* - v x g* = 0 (12.21) V.B* = 0 (12.22) V x B* = Rm< x B* -+ B\&* (12.23) V. E* = 0 (12.24) v x E* = 0 (12.25) vq5* = - E* (12.26) 其中g* = gv0 1, g* = g)Ov,', B* = BB,', E* = EL0 A\&',4* = q5Aq5ClI p* = y* = y Li'。 在这里，LO是参考长度，210是参考速度，Aq50是参考电势差]，Bo是参考磁通密度。 为了方便起见，在论文的其余部分将删除*上标。 IC* = 5 L，非维数由以下方式给出：应该注意的是，由于对电极应用物理上有意义的边界条件的便利性，电势被引入为一个额外的变量。 对于电场方程]麦克斯韦方程的一阶形式不包括电势。 由于静态电场最常见的边界条件是以电位为条件给出的，因此有必要为电势添加方程（12.26）] 4。 我们现在以一阶系统的一般形式编写上述系统（12.12）。 虽然用（12.19）-（12.26）编写的整个系统可以用LSFEM处理，但发现以迭代方式分别求解流体和电磁场方程更经济。 在这里，为流体系统（12.19）-（12.21）编写了一个一阶的一阶系统，并用上标流体表示。 在这里，流体方程中的时间导数项使用后退的Euler方案进行近似。 （12.28）一阶系统也以一般形式写成电磁场方程（12.22）-（12.26），并用上标em表示。 此外，流体方程中的非线性对流项与牛顿的线性化

272 流动分离电磁场控制方法导致一个适合使用LSFEM治疗的系统。 对于二维情况，我们指定磁场的z分量，并且x和y分量为零。 对于许多工程应用，Rm和B2的大小通常很小，因此我们预计x-y中的电流诱导磁场与外部施加的磁场的大小相比可以忽略不计。 速度2和电场E的x和y分量保持未知，而其z分量假定为零。 为了简单起见，我们只考虑不包含自由带电粒子的流动。-蕨类植物=\{ ii.y.-i L\} (12.30) 可以通过一个简单的迭代过程找到满足上述所有方程组的解。 首先，（12.30）中给出的系统是电场的求解。 (12.29)中的系统是用上一个时间步骤的电场和速度求解的。 在这里，从上一个迭代中取的数量用下标0指定。 如果问题非常非线性，这些方程可以在每个时间步骤中迭代。 在这种情况下，每个时间步骤的迭代都会重复，直到达到指定的收敛公差。 通常不到5次迭代就能实现两个系统的残差范数减少3.5个数量级。

\subsection{12.4.2 准确性验证}
B. H。 丹尼斯和G。 S。 Dulikravich 273 Ht 12.4.2 准确性验证 10.0 很难验证EMHD代码的准确性。 这是因为这种方程没有分析解。 然而，MHD流的分析解决方案确实存在。 在这里，我们将使用这样的分析解决方案来验证代码的MHD部分。 LSFEM对MHD的准确性是根据Poiseuille-Hartmann流的分析解决方案测试的[13]。 普瓦伊苏伊-哈特曼流是两个固定板之间的导电和黏性流体的一维流动，其均匀的外部磁场正交施加在板上。 假设墙壁在y = fL，并且墙壁上的流体速度为零，并且流体在恒定压力梯度的影响下向x方向移动，那么速度曲线由Rm 210（rn s-1）Lo（m）（5，-1 S -1）Bo（TI p（Hm-l）c（0-lm-l）dpldx（Pa m-l）cosh（Ht）-cosh（y）sinh（Ht）u（y）= -- aBy”dx 6 x 10W7 1.0 0.6 0.01 1.0 1 x 10-6 0.6 1.0（12.31）流体的运动诱导x方向的磁场，并由sinh（y）- f sinh（Ht）cosh（Ht）-1（12.32）测试用表12.1给出的参数运行一个测试案例，网格由2718组成 抛物轮三角形元素。 图12.1显示了速度曲线的计算和分析结果。 图12.2显示了感应磁场的计算和分析结果。 对于这两种情况，人们可以看到分析解决方案和LSFEM解决方案之间的一致性非常好。 表12.1：用于Poisuille-Hartmann流量测试问题的参数

\section{12.5 数值结果}
274 流动分离 电磁场控制 7.00E-01 6.00E-01 5.00E-01 4.00E-01 - - 3.00E-01 2.00E-01 1.00E-01 O.OOE+OO 0.5 1 1.5 -1.00E-01 t 图12.1：速度的计算和精确值 12.5 数值 数值 EMHD的LSFEM配方现在将用于圆形圆柱体上流动的电磁控制。 图12.3说明了极子和磁铁的配置。 在这种配置中，气缸分为两个电极，一个在顶部和底部，磁铁稍微放在尾迹区域的下游。 这种配置是现实中的3D，但可以通过指定磁场的z分量并计算电场和流场的x-y分量来近似2D。 已知的磁场被假定是均匀的，并被施加到x-y平面中。 流体被认为是导电的，流过一个单位直径的圆圈。 我们计算了稳态和非定常流动情况，以观察电磁场对流动模式的影响。 表12.2显示了两个测试案例的相关非维度参数。 图12.4所示的混合三角形和四边形网格用于所有计算。 网格由818个抛物面元素和2458个节点组成。 在气缸表面没有应用滑边界条件，而在入口和外部域的顶部和底部应用了自由流条件。 出口边界上的条件是自由的。 电位在电极表面指定，从而产生电位差

B. H。 丹尼斯和G。 S。 Dulikravich 275 0 0.5 1 1.5 2 Y+L（m）图12.2：诱导磁场的计算和精确值，该磁场迫使电流流过导电流体。 图12.5显示了电势的计算分布以及静电场的电场线。 EMHD系统中的源项直接涉及电场强度E，因此我们预计场线的形状将对流动模式产生强烈影响。 在本例中，每个电极都使用恒定电位，因此场线在域中平稳分布。 在第一种情况下，Re = 37的稳定流是在没有电场或磁场的情况下计算的。 图12.6显示了这种经典流动的流线和压力分布。 当电场和磁场应用于流动时，速度和压力分布发生了巨大变化，如图12.7所示。 关于这个结果的一个有趣的说明是，虽然气缸后面的分离已经消除，但1.0 x 10- 1.0 x 10-背面的压力表12.2：用于测试案例的非尺寸编号

276 流动分离电磁场控制均匀流动图12.3：磁铁和电极的测试配置0 20 10图12.4：混合非结构网格使用的气缸明显减少。 这种低压是由于当电磁场引起的身体力存在时，气缸后面的流速增加。 此外，如果电极的极性颠倒，则会观察到相反的效果。 图12.8显示，反向极性电磁场实际上增强了分离强度。 然而，尽管分离更强，但气缸后壁的压力会增加。 再一次，这种压力的增加是由于能量通过外部电磁场插入到流动中。 在第二种情况下，首先在没有电场或磁场的情况下计算Re = 100的非定常流动。 流量是从最初均匀的流动场开始计算的，直到达到周期性状态。 在这种情况下，使用了时间步数At = 0.1，并在900个时间步数后达到周期性流量。 图12.9显示了粒子在大约1400个时间步数的一瞬间的痕迹。 此时可以清楚地看到特征性的涡流脱落模式。 这个案子已经办完了

\section{12.6 结论}
B. H。 丹尼斯和G。 S。 Dulikravich 277 图12.5：再次计算电场线和电势轮廓，但在1200个时间步后，静态电磁场被激活。 通过1400个时间步骤，图12.10中的粒子痕迹清楚地显示了电磁场对跟踪结构的影响。 到2000个时间步进标记，流量达到稳定的、与时间无关的状态。 如图12.11所示，由于周期性涡流脱落的消除，流量变得更加稳定。 这种配置中结合的电场和磁场具有很强的阻尼作用，完全抑制了通常在Re = 100时看到的涡流脱落。 由此产生的流场是稳定的，如图12.12所示。 12.6 结论 同时应用和相互作用的电场和磁场以及不可压缩的同构黏性流体流动的统一理论模型已被表示为一阶偏微分方程系统的耦合序列。 这些系统使用最小二乘有限元法在二维中离散化，并集成在非结构化计算网格上。 数值结果与同时施加均匀垂直电场和均匀水平磁场的无限平行板之间稳定层流的测试案例非常一致。 该方法用于模拟有和没有外部施加的电场和磁场的圆形圆柱体上的流动。 结果表明，某种安排

\section{12.7 致谢}
\section{12.8 参考书目}
2 78 流动分离 电磁场控制图12.6：无电场和无磁场的稳定流动的压力场和流线型电极和磁铁可用于消除稳定流动中的流动分离，并抑制非稳定流动中的涡流脱落。 12.7 致谢 这项工作是在主要作者担任东京大学客座副教授期间完成的。 主要作者感谢东京大学前沿科学研究生院的支持。 第二位作者感谢通过计算数学计划管理的NSF赠款DMS-0073698提供的部分支持。 12.8 参考书目[l] Colaco，M。 J.，杜利克拉维奇，G. S.，和Martin，T. J。 墙电极的优化，用于凝固过程中自然对流效应的电流体动力学控制。 材料和制造工艺，19（4），2004年。 [2]丹尼斯，B。 H。 电磁流体流的模拟和优化。 博士论文，宾夕法尼亚州立大学，宾夕法尼亚州大学公园，12月。 2000年。

B. H。 丹尼斯和G。 S。 Dulikravich 279 图12.7：电场和磁场稳定流动的压力场和流线[3] Dennis，B。 H。 \& Dulikravich，G。 S。 晶体生长中熔融流动的磁场抑制。 国际热流体杂志，23（3）：pp。 269-277，2002年。 [4]丹尼斯，B。 H.，埃戈罗夫，I. N.，Han，Z.-X.，Dulikravich，G. S。 \& 波洛尼，C. 涡轮机械级联的多目标优化，以实现最小损耗、最大负载和最大间隙与弦比。 国际涡轮和喷气发动机杂志，18（3）：201-210，2001年。 [5] 杜利克拉维奇，G. S。 电磁流体力学和凝固。 在D。 A. 签名人，D。 De Kee和R。 P。 Chhabra，编辑，非牛顿流体流动和流变学的进步，B部分，流变学系列第8卷，第9章，第677-716页。 爱思唯尔出版社，1999年。 [6] Dulilcravich，G。 S.，Ahuja，V.，和Lee，S. 用磁场和降低重力对三维凝固进行建模。 国际热与质量杂志Tbansfer，37（5）：pp. 837-853，1994年。 [7]杜利克拉维奇，G。 S.，Choi，K.-Y.，\& Lee，S. 稳定不可压缩层流中涡量的磁场控制。 在D。 A. 西格纳，J。 H。 金，S。 A. 警长和H。 W。 Colleman，编辑，电流和测量不确定性发展研讨会，ASME WAM'94，1994年。 伊利诺伊州芝加哥，11月 6-11，1994年，ASME FED-Vol. 205/AMD-Vol。 190，（1994），pp. 125-142。

280 流动分离 电磁场控制图12.8：具有反向电场和磁场的稳定流动的压力场和流线[8] Dulikravich，G。 S.，科拉科，M. J.，丹尼斯，B. H.，T. J。 马丁，T。 J.，Egorov-Yegorov，I。 N.，\& Lee，S.4。 优化强度，以及在凝固过程中控制熔体流动的磁铁的方向。 材料和制造工艺，19（4）：695-718，2004年。 [9]杜利克拉维奇，G。 S.，科拉科，M. J.，Martin，T. J.和Lee，S. 凝固复合材料中的磁化纤维取向和浓度控制。 复合材料的J.，47（15）：pp。 1351-1366，2003年。[lo] Dulikravich，G。 S。 \& Lynn，S. R。 统一电磁流体动能（emfd）：介绍性概念。 国际非线性力学杂志，32（5）：913-922，1997年。[ll] Dulikravich，G。 S。 \& Lynn，S. R。 统一电磁流体力学（emfd）：数学模型调查。 国际非线性力学杂志，32（5）：923-932，1997年。 [12] Eringen，A. C. \& Maugin，G。 A. 连续九的电动力学；流体和复杂介质。 Springer-Verlag，纽约，1990年。 [13] 休斯，W。 F。 \& Young，F.-J. 流体的电磁动力学。 John Wiley and Sons，纽约，1966年。 [14] Jiang，B.-N. 最小二乘有限元法。 Spring-Verlag，柏林，1998年。

B. H。 丹尼斯和G。 S。 Dulikravich 281:t 3 时间=l40.3 sec I I 0 10 图12.9：无电场和无磁场的非定常流动的粒子痕迹[15] KO, H.-J \& Dulikravich, G. S。 电磁流体力学的完全非线性模型。 国际非线性力学杂志，35（4）：pp。 709-719，2000年。 [16] 拉赫塔基亚，A。 \& Weiglhofer，W. S。 关于磁电介质的因果构成关系。 1995年，IEEE电磁兼容性国际研讨会。 佐治亚州亚特兰大，8月14日至18日，（1995年），pp. 611-613。 [17] Meir，A. J。 \& 施密特，P. G。 Mhd流量的分析和有限元模拟，应用于海水阻力降低。 减少海水阻力国际研讨会论文集，纽波特，7月22日至23日，第401-406页，1998年。 [18] Stuetzer，0。 M。 磁流体动力学和电流体力学。 [19]萨顿，G。 W。 和谢尔曼，A。 工程磁流体动力学。 Mc- [20] Weier, T., Gerbeth, G, Mutschke, G, Lielausis, O., \& Platacis, E. 通过定位在气缸表面的电磁力在电解质溶液中气缸唤醒稳定的经验。 实验热和流体科学，10:84-91，1998年。 [21] Weier, T., Gerbeth, G., Posdziech, O., Lielausis, O., \& Platacis, E. 关于人体周围流动的电磁控制的一些结果。 流体物理学院刊，5（5）：534-544，1962年。 格劳山，纽约，1965年。

0.3: 0.2: -0.1 -0.1: -0.2 7 -0.3: -0.4 流动分离 电磁场控制 EM 应用. 应用. -0.5 I,,,, I,,,, I,, 8,,,, I, 75 100 125 150 175 时间（见）图12.10：海水阻力减少国际研讨会后速度的v组分的时间变化，纽波特，7月22日至23，第395-400页，1998年。

B. H。 丹尼斯和G。 S。 Dulikravich 283 时间=140.3 秒 \_*....... 10 10 图12.11：在120秒时打开电场和磁场的非定常流动的粒子痕迹图12.12：电场和磁场诱导的Re=100稳定流动的压力场和流线

此页面故意留空

\clearpage
\chapter{第13章 跨声速流动在扁平翼型 Alexander G上的分叉。 库兹敏}
\section{13.1 导言}
第13章 跨声速流动在扁平翼型 Alexander G上的分叉。 Kuz'min' 13.1 介绍 在高亚音速自由流速度下，翼型的上表面和/或下表面附近存在局部超音速带。 如果翼型的曲率在中弦区域很小，那么流动可能会在每个表面附近表现出几个超音速区。 当自由流马赫数 M增加时，这些区域通常合并，当M减少时，它们分裂成更小的区域。 最近，Kuz'min和Ivanova [8, 5, 121研究了通道中凸起上的黏性跨声速流动，并展示了与超音速带分裂/合并相关的不稳定性。 不稳定性意味着在边界条件的小扰动下，稳定流动结构的突然变化。 结构不稳定性的概念有助于理解前几年揭示的跨声速流动的非独特性（Jameson [7]，Bafez和Guo [3，41，Caughey [l]）。 对于对称的机翼，在[lo，6，141中分析了非独特性和不稳定性之间的联系。 在[ll, 91.中讨论了一些不对称的机翼。 在本文中，我们在对称翼型上进一步分析了非黏性跨声速流动。 升力系数 C对自由流的依赖-'数学和力学研究所，St. 圣彼得堡州立大学，圣彼得罗德沃罗茨大学大道28号。 彼得堡198504，俄罗斯。

\section{13.2 问题陈述和数值方法}
\section{13.3 升力系数作为M函数的分析}
286 讨论了跨声速流动位的分叉、对称解的稳定性以及与攻角变化有关的迟滞。 13.2 问题陈述和数值方法 我们考虑翼型 y(z) = *0.06.\textbackslash\{\}/1 - (2x - 1)4 (1 - d2 1 ', O<Xll, (13.1)上的二维不黏性可压缩流动，其曲率在x = 0.482时达到最小值0.0126。 计算域的C型远场边界位于距离翼型的15到18弦长度。 马赫数毫米和攻角a在远场边界上给出，而经典的滑移条件则在翼型上规定。 为了获得问题的数值解，我们采用了NSCKE有限体积求解器，其中Euler 方程使用Roe方案[13]在非结构网格上的空间离散化。 使用Van Albada类型限制器进行MUSCL重建可以获得二阶精度。 流入和流出边界条件的数值建模基于Steger-Warming通量向量分割技术[15]。 时间积分是用明确的四级Runge-Kutta方案执行的。 使用本地时间推进策略计算了稳态解决方案。 初始数据要么是由自由流条件定义的统一状态，要么是之前为Mm和a的其他值获得的稳定流场。 计算主要在733 x 215网格点的三角形网格上进行，该网格聚集在翼型附近和激波附近。 在489 x 143的较粗网格上进行测试计算表明，激波位置的误差大大增加。 另一方面，对主网格的改进在计算的流场中只产生了微不足道的修正，同时导致获得稳态解所需的CPU时间不成比例地增加。 该方法通过在通道中计算跨声速流动并与用方案EN02获得的Euler 方程的解决方案进行比较来验证[8，51。 在类似网格上获得的结果非常一致。 此外，通过计算Q = 1.25度的NACA 0012 翼型基准问题，证实了该方法的准确性[2]。 13.3 图1显示了升力系数作为Mm的函数的升力系数与Mm的图，用上述方法在Q = 0时计算。 得到了地块的左分支

A. Kuz'rnin 287 cL c I I I I I I, I I I 'M, 0.820 0.824 0.828 0.832 图 13.1: 升力系数 Ch作为M的函数，对于翼型 (13.1) at (I: = 0. 通过计算M, = 0.818的流量，并逐步将马赫数增加到0.8275。 在每个步骤中，之前的稳态被用作初始数据。 计算出的流场相对于z轴对称，涉及两对超音速带（图2a）。 图中所示的图的右分支。 1是通过计算Mm = 0.836的流量，然后逐步将马赫数减少到0.829获得的。 相应的流量是对称的，在翼型的每个表面上都显示出单个超音速区。 奇异马赫数 M，M 0.8283被确定为在不重组流量的情况下可以延长左右分支的极限。 对于图中描述的对称流给出的初始数据。 2a，扰动a = 0.07度导致上表面超音速带突然合并，并放松到本质上不对称的稳态。 如果从0.07返回到零，则获得的结构1+2（翼型的上表面有一个超音速区，下表面有两个超音速区）将被保留（图2b）。 当自由流Much数在区间Mmin < Mm < Mmax（13.2）中以Mmin M 0.8228和M,,, x 0.8315变化时，这种结构会持续存在（见图中的上分支。 1）。 同样，（I：）的负扰动意味着过渡到图中所示的下分支。 1. 图3显示了间隔（13.2）对网格加密的依赖性。 它表明，M,in和M,的值的绝对误差可以估计为0.0003。 因此，获得值的最后几位是

288 跨声速流动的分叉 0 0.4 0.8 X 0 0.4 0.8 X 图13.2：跨声速流动中的马赫数等高线在翼型（13.1）在M，= 0.824，a = 0：b）不对称稳态对应于图中的A1点。 1. a）对称稳态对应于图中的A2点。 1，相当值得怀疑。 然而，我们不四舍五入值以避免额外的错误。 图的比较。 1，在（13.1）[6，141中用（1-do）而不是（1-d2）的类似翼型获得的结果表明，在手头的情况下，间隔的长度（13.2）和C的最大值都大于（6，141。 这归因于翼型在中弦区域的曲率较小。 与此同时，间隔长度（13.2）和CL的最大值都比翼型的长度小，其中平中间部分在[3]中考虑。 图4显示升力系数作为M的函数，在a = 0.1度。 在这种情况下，区间（13.2）扩大，因为Mmin转移到0.8206，而M，，扩大到0.8333。

\section{13.4 关于变异的稳定性分析}
A. Kuz'min 0'823：Mmin,, NO. of.grid 0.829 0.827 0.825 图13.3：网格加密的区间（13.2）端点的偏移。 13.4 如果采用图中所示的不对称流量，则分析ctl变化的稳定性。 初始数据为2b，并将Q从0度减少到-0.04度，保持M，= 0.824固定，然后计算展示了上表面超音速带的分裂，并过渡到结构2+2（翼型的每个表面都有两个超音速带相邻）。 这对应于从图5b所示的点A1跳到图中间分支的左端点，如箭头所示。 之后，返回Q = 0沿着中间分支移动到点A2，而Q进一步降低到-0.06度会导致与下表面相邻的超音速区域突然合并，并过渡到结构2+1。 后者意味著跳到图5b中图的下部分支。 然后，恢复a = 0会沿着下分支移动到点A3。 作为M，从0.824增加到0.8255，地块CL（Q）的中间分支缩小并旋转到几乎垂直的位置（图5c）。 因此，对于对称流，相对于Q变化的稳定性范围变得非常窄。 在M，=M，中间分支消失，而与从上分支过渡到下分支相关的迟滞，反之亦然，达到其最大值（图5d）。 事实上，对称流不仅在M, = M时不稳定，而且在包含M,的自由流马赫数的一定间隔中是不稳定的。 例如，在区间0.8248 f Moo 5 0.8305（由图中的虚线段表示。 1）将Q = 0替换为Q = 50.04度足以触发突然

\section{13.5 结果摘要}
\section{13.6 结论}
290 跨声速流动 cLl的分叉 0.1 '.-P -0.2 -O-l" 0.820 0.824 0.828 0.832 0.836 Moo 图13.4：升力系数 CL与M，对于翼型（13.1）在a = 0.1度。过渡到不对称状态1+2或2+1。 因此，在实践中，在翼型上不可能存在稳定对称流动的马赫数范围似乎取决于自由流湍流水平。 在M，= 0.831 > Ms，CL与a的图与图5c所示相似。 对于马赫数字0.818 < M，< Mmin和Mmax < M，< 0.836，在a轴的两个独立区间观察到迟滞（图Sa，e，f）。 13.5 结果摘要 升力系数 CL作为两个变量a和M的函数，可以用空间中的曲面（a、M、CL）来说明。 图6显示了位于上半空间CL > 0的表面片段。 表面有一个轻微的接缝，从表面上部的El点到Ez点，以及右下角的E3点到E4点。 接缝对应于流动模式，其中斜冲击与冲击的脚部相遇，终止了局部超音速带。在原点附近。 图8显示了通过将上述表面的边缘投射到平面（a、M、）获得的分叉曲线。 图7显示了表面CL（CY，M，）13.6的正负部分结论对于在中弦区域扁平的翼型（13.1），数值分析显示Euler 方程在某些范围内的多个解

\section{13.7 参考书目}
A. Kuz'min 291自由流参数M和a。 事实证明，对称流动对于接近奇异值M的自由流马赫数的小扰动是不稳定的。 这项工作得到了俄罗斯基础研究基金会的支持，赠款为03-01-00799。 作者感谢B。 Mohammadi为跨声速流动模拟提供NSC2KE求解器。 13.7 参考书目[l] Caughey，D.A.，非稳态跨声速流动和空气弹性研究，Proc。 IUTAM Symp。 Transsonicum IV，H。 Sobieczky（编辑），Kluwer，2003年，pp. 41-46。 [2] Delanae，M.，Geuzaine，Ph. \& Essers，J. A.，非结构化自适应网格上可压缩流的二次重建方案的开发和应用，AIAA论文97-2120，1997年，pp. 250-260。 [3]哈菲兹，M。 M。 和郭，W。 H.，跨声速流动的非唯一性，Acta Mechanica 138，1999年，pp. 177-184。 [4]哈菲兹，M。 M。 和郭，W。 H.，激波数值模拟的一些异常，pt I：非黏性流动、计算机和流体28，编号。 415，1999年，pp. 701-719。 [5] 伊万诺娃，A. V.，通道中Inviscid 跨声速流动的结构不稳定性，J. 工程物理学和热物理学76，编号。 2003年6月，pp. 58-60。 [6] 伊万诺娃，A. V。 \& Kuz'min，A。 G.，跨声速流动在机翼上的非独特性，Izvestiya Akademii nauk，Mekhanika zhidkosti i gaza（翻译：Proc. 俄罗斯科学院，气体和液体力学），编号。 4，2004年，pp. 152 - 159。 [7]詹姆森，A.，翼翼承认尤拉方程的非唯一解，AIAA论文91-1625，1991年。 [8] 库兹敏，A。 G.，激波与Sonic Line的相互作用，Proc。 IUTAM Symposium Transsonicum IV，H。 Sobieczky（编辑），Kluwer，2003年，pp. 13-18。 [9] Kuz'min，A. G.，跨声速流动在单数自由流马赫数字下翼翼上的不稳定性，Proc. 第四届欧洲应用科学和工程计算方法大会\textasciitilde{}ECCOMAS 2004，P. Neittaanmaki，T。 罗西，K。 Majava，和0。 皮罗诺（编辑），2004年。

292 跨声速流动的分叉[lo] Kuz'min，A。 G。 \& 伊万诺娃，A. V.，Inviscid 跨声速流动的结构不稳定性，预印本01-04，数学和力学研究所，St. 彼得堡国立大学，2004年。[ll] Kuz'min，A。 G。 \& 伊万诺娃，A. V.，跨声速流动在翼型上的结构不稳定性，J. 工程物理和热物理学77，编号。 5，2004年，pp. 144 - 148。 [12] Kuz'min，A. G。 \& 伊万诺娃，A. V.，与超音速域的合并/分裂相关的跨声速流的结构不稳定性，理论和计算。 流体动力学17，斯普林格，2004年（印刷中）。 [13] Mohammadi，B.，NSC2KE的流体动力学计算：用户指南，1.0版，INRIA Tech。 报告RT-0164，1994年。 [14] 塞米奥诺夫，D。 S.，跨声速流动在翼型上的不稳定制度，数学建模16，莫斯科，2004年（印刷中）。 [15] Steger，J. \& 温暖，R。 F.，非黏性气体动力学方程的通量向量分割，应用于有限差分法，J. 补偿。 物理。 40，不。 2，1983年，pp. 263-293。

A. Kuz'min 0.2 a) - 0- h -0.4 -0.2 0 0.2 0.4 a, deg 图13.5: 升力系数 C, 对比翼型的攻击角a (13.1): a) M, = 0.820, b) M, = 0.824, c) M, = 0.8255, d) Mm = 0.8283, e) M, = 0.834, f) M, = 0.836。

跨声速流动的Bafurcataon图13.6：位于半空间CL>0中的表面CL（M\textasciitilde{}，\textasciitilde{}）的碎片：a），b）-从两个不同视角的视图。

A. Kuz'min 295 图13.7：原点附近的表面CL（M\textasciitilde{}，a）的片段。 0.1 -0.1 - 0.2 0.820 0.824 0.828 0.832 0.836 图13.8：分叉曲线显示了发生流动重组的自由流条件。 虚线段表示相对于扰动a = f0.04 deg的对称流动不稳定范围。

此页面故意留空

\clearpage
\chapter{第14章 Euler计算对超细圆锥体的涡流对稳定性的研究 Jinsheng Cai、Her-Mann Tsai、Shijun Luo和Feng Liu}
\section{14.1 摘要}
第14章 尤拉计算研究细圆锥体上涡流对的稳定性 Jinsheng Cai'、Her-Mann Tsai'、Shijun Luo2和Feng Liu2 14.1 摘要 Euler计算研究了圆锥和平板三角翼的细锥体组合中静止对称和不对称涡流对的形成和稳定性，在高攻角无侧滑。 欧拉求解器基于并行、多块、多网格、有限体积法，用于覆盖网格上的三维、可求解、稳定和非稳态Euler 方程。 静止涡流配置首先通过在稳态或时间准确模式下运行尤拉代码来捕获。 在获得静止的涡流配置后，将短暂的不对称扰动引入到机翼的左侧和右侧的短暂吸力和短暂的吹风，然后以时间准确模式下运行尤拉代码，以确定流动是否会恢复到原来的不受干扰的状态，还是演变成“新加坡国立大学Ternasek实验室，新加坡肯特里奇新月，”，加州大学欧文分校机械和航空航天工程系，119260。 加利福尼亚州 92697-3975

\section{14.2 导言}
298 涡流对的稳定性 不同的稳态或非稳态解。 检查了涡核的细节，以评估欧拉计算在解决涡流结构中的有用性。 计算结果与蔡、刘和罗开发的稳定性理论非常吻合（流体机械学J.，卷。 480，2003年，pp. 65-94）和可用的实验数据。 该协议证实了这样的结论，即绝对类型的流体动力学不稳定性可以成为涡流不稳定的机制，并证明了欧拉方法对涡流主导的高角攻击力在锋利的物体上流动的稳定性研究和模拟初级涡流的有用性。 14.2 导言 最初对称的涡对在纤细的尖体上变得不对称，因为攻角增加超过一定值，即使在零侧滑时也会导致大的不对称侧力和力矩。 这种涡流破坏对称性的机制尚不清楚。 尽管在涡旋不对称性上花费了大量实验、理论和计算努力，但在对漩涡不对称的理解、预测和控制方面存在许多分歧。 亨特[22]、爱立信和雷丁[18]以及尚皮尼[lo]审查了该主题。 许多实验观察[15，23，471和数值研究[12，13，26，201）发现，靠近鼻尖的微不对称扰动在高攻角下会产生强烈的流动不对称。 似乎毫无疑问，涡流不对称是在顶端区域触发、形成和发展的，而前体的后部和后圆柱体（如果有的话）对顶点区域的不对称性影响不大。 高点扰动的演变在决定整个身体的流动模式方面起着重要作用。 由于任何细长尖体的尖端部分几乎都是圆锥体，因此对圆锥体的高攻击角流进行了分析研究。 使用Bryson[2]的分离涡流模型，Dyer、Fiddes和Smith[16]发现，除了静止对称涡流解外，当攻角大于半顶点角的约两倍时，圆锥体上还存在静止的不对称涡流流解，即使分离线假设在对称位置。 皮德和史密斯后来对这些静止涡流的稳定性进行了分析研究。 [35]他们在稳定性分析中处理的干扰是空间性的，而不是时间性的。 原始的站位对称/不对称锥形涡流的位置略有变化，被引入到体尖附近的流动中。 在细长体理论的假设下，计算了下游方向这种扰动的初始变化率。 如果所有扰动都衰减，溶液是稳定的，而如果任何扰动增加，则溶液是不稳定的。 这种不稳定机制通常被称为

J。 蔡，H-M。 蔡英文，S。 罗tY F。 刘299对流不稳定。 Dagani[l2, 131和Levy, Hesselink和Degani[26]使用时间准的Navier-Stokes方法进行数值计算，研究了3.5口径切向ogive-圆柱体在低速旋转的分离涡流。 他们发现，为了在数值计算中获得并保持不对称的涡流模式，有必要在体顶点附近保持固定的小几何扰动。 一旦鼻尖上人工引入的小“不完色”被去除，涡流就会恢复对称。 数值计算中的这一发现，加上对德加尼和托巴克[14]的实验观察，使他们相信，在概念上与皮德等人研究的对流不稳定机制相似。[35]是导致在具有尖头的旋转的细长体上，否则可能是对称的涡流的不对称的开始。 Degani和Tobak的实验表明，奥格圆柱体上的涡流模式连续和可逆地取决于所有攻角的受控尖端扰动。在一项单独的数值研究中，Hartwich、Hall和Hemsch [20]报告了不可压缩的三维湍流 Navier-Stokes 方程的不对称涡流场解，用于攻角为40英寸的3.5口径切向圆柱体，在计算中没有固定的几何不对称。 有人声称，不对称解决方案是由计算中的计算机四舍五入错误触发的。 因此，不对称可以由瞬态扰动引起。 这种不对称的途径被称为绝对不稳定。 有实验观察支持绝对不稳定性的概念，例如，当在50-60"[15，471和迟滞效应[34，11）的入射角度变化时，存在双稳定配置。 绝对不稳定性机制也在理论上进行了研究。 使用Legendre[24]的简化分离涡流模型，Huang和Chow [21]成功地从分析上表明，在零侧滑时，纤细平板三角翼上的涡对可以静止，并且在小锥形扰动下是稳定的。 使用相同的流动模型，蔡、刘和罗[5]在圆锥流动和经典细体理论的假设下，开发了静止圆锥涡对在一般细圆锥体上的稳定性理论。 他们在稳定性分析中处理的干扰是时间性的或短暂的，而不是空间性的。 小位移被引入静止涡流位置，然后被移除。 假设位移的涡流仍然是圆锥形的。 干扰是全球性的，而不是局部性的。 Cai, Luo, and Liu[6, 8, 71 扩展了参考[5]中描述的方法，以研究静止不对称涡流对在细长锥体和有和没有侧滑的翼体组合上的稳定性。 在他们的研究中，在初始时间引入了扰动。 如果涡对的两个涡流在30-60英寸的初始动作后恢复到原来的静止位置。

900 涡流对扰动的稳定性，涡对是稳定的；如果该对的任何一个涡流离开其原始位置，涡对是不稳定的；如果涡流周期性地围绕静止点移动或保持在受干扰的位置，涡对是中性稳定的。 这与上述讨论的绝对不稳定类型有关。 爱立信[l7]提出了不同的观点，他声称没有直接的流体力学不稳定性证据。 由于实验总是以几何微不对称为主，因此永远无法实现完全对称的流动状态，这将是显示扰动效应的预处理。 他认为，流动不对称的机制是不对称的流动分离和/或不对称流动重新连接。 然而，有一些观察结果很难与这种观点相协调。 例如，当攻角低时，对称涡流就存在。 只有当攻角超过一定值时，才会出现流动不对称。 此外，正如Cai等人1所示，Shanks[40]观察到的具有尖锐前缘的三角翼上的不对称涡流，这是由不对称的重新连接引起的，这更可能是由安装在模型的lee-side入射平面上的短垂直鳍存在流体动力学不稳定引起的。[ 5]。 从静态（固定）攻击角度来看，对流不稳定性概念和绝对不稳定性概念截然不同。 它们可能会导致同一涡流场的不同稳定性结论。 对于Pidd等人a1.[35]研究的圆锥的情况，他们表明，静止对称圆锥涡流对在入射参数的窄带中对流稳定，而静止不对称圆锥涡对对流稳定，例外情况微不足道。 然而，Cai等人[7，81证明，无论假设对称还是不对称分离线，对称和不对称的静止锥形涡流对都是绝对不稳定的。 因此，在某些条件下，在对流稳定性方面，稳定的圆锥对称和不对称涡流对可能存在于圆锥上，而在绝对稳定性方面，圆锥上的稳定涡流对要么不存在，要么必须是非圆锥形配置。 在逻辑上，满足对流和绝对类型的稳定性条件对于任何圆锥对称或不对称漩涡模式的配置在流动中持久是必要的。 最近，Cummings、Forsythe、Morton和Squires[ll]对高攻击角度不对称性给出了另一种解释。 对流不稳定性假说指出，随着攻角增加到高范围，任何程度的不对称性（包括对称情况）都是可能的。 绝对不稳定性假说指出，随着攻击角度增加到一定水平，将发生分叉，这将产生两个镜像不对称解之一，中间不对称解和对称都是不可能的。j从动态系统的行为角度来看

J。 蔡，H-M。 蔡英文，S。 罗和F。 刘301随着攻角的增加，对流不稳定性可以称为不稳定分叉，而绝对不稳定性是分叉。 目前还没有完全理解身体运动和前体流场之间的相互作用。 导致涡流不对称的基本机制仍然是一个争论的话题。 值得注意的是，预测分离涡流的行为在计算中也非常困难，因为计算也不是完美的模拟。 虽然一些研究人员满足于允许他们的数值算法提供触发涡流不稳定20、30、411所需的物理扰动，但使用一种不增加任何未知扰动水平的算法更可取。 将扰动明确地添加为几何或流场扰动会更好。 Levy、Hesselink和Degani[25]表明，某种算法中包含的固有偏差在流程中会产生异常不对称。 为了对涡流稳定性进行数值研究，数值算法需要对称，以便在分配适当的初始和边界条件时获得基本的静止对称或不对称涡流解。 高攻角的流动对人工黏度或数值耗散特别敏感。 Hartwich[lS]表明，出口边界条件中一阶精确差分方案导致的过度数值耗散抑制了对称性破坏。 在他的Navier-Stokes计算中，Thomas[43]在攻角为20°的5英寸半角锥体上对超音速黏性流动的计算中，发现体尖端附近的网格分辨率不足会产生虚假的不对称。 由于剪切层分离区域和涡流核心区域以及身体表面附近的区域都存在高梯度，在这些区域中具有足够网格密度的适当网格拓扑对计算研究至关重要。 对于尖刃几何，分离点固定在前缘，独立于雷诺数。 Euler 方程的数值方法引入的耗散应该模仿物理黏性并导致分离。 就像分离点对真实黏度不敏感一样，它也应该对人工黏度不敏感。 一旦分离和产生的涡量建立，涡流运动的动力学，即它与相邻物体表面的相互作用，本质上是不可见的，因此Euler 方程可以充分描述。 虽然尤拉计算中没有模拟由机翼背侧黏性效应引起的二次涡流，但它们对初级涡流的影响很小。 因此，使用欧拉代码进行涡流在细长锋利的物体上的稳定性研究是使用Navier-Stokes代码的有吸引力的替代方案，Navier-Stokes代码需要更大的计算资源，并受到湍流建模的经验主义的影响。 本文中考虑的计算模型是圆锥形的。 既不会研究后缘效应，也不会研究涡流分解。 在以下章节中，数值方法和流动模型是

\section{14.3 欧拉求解器和流动模型}
502 描述的涡流对的稳定性。 验证了低速锥形流动假设，并在验证的基础上实现了圆锥形网格来计算静止对称和不对称涡流及其稳定性。 研究了计算的初级涡核作为流动中的一个重要参与者。 计算结果与已知的理论分析和可用的实验数据进行了比较。 最后得出了结论。 14.3 欧拉求解器和流动模型 众所周知，欧拉求解器可以自动捕获与扫过的锋利前缘分离的剪切层及其螺旋卷起到流场的涡核中。 虽然欧拉解中没有涡流的次要特征，但再现了流场的总体主导特征，即初级涡流配置及其与体表面的相互作用。 目前的尤拉求解器基于三维可压缩尤拉/Navier-斯托克斯方程的平行、多块、多网格、有限体积代码。 仅使用代码的欧拉选项。 基本数值方法使用中心差分方案，混合了二阶和四阶人工耗散和显式Runge-Kutta型时间行进。 生成的代码保持了对称性。 非稳定时间准确计算是通过使用具有双时间步进的二阶时间准确隐式方案来实现的。 求解器已针对一些稳定和不稳定的案例进行了验证[27、28、29、391。 实施了一种新开发的过量网格技术[3]，以促进网格加密在高涡量的域中。 图14.1说明了由圆形锥体和平板三角翼组成的纤细锥翼体模型。 还显示了两个分离涡流，体直线性坐标系（2，y，z）和涡旋圆柱坐标系（a，r，O）。 连接涡流和机翼前缘的剪切层（图中未显示）在参考理论中被忽略。[5]但在本文的计算中被考虑。 在计算中，自由流马赫数 M设置为0.1，以近似于不可压缩流动。 参考[4]报告了平板三角翼和翼体组合的一些计算。 我们将在本文中的讨论限制在一种配置上，在零侧滑（p = 0）的流动条件下，delta-wing的体半径与机翼半跨度比y = b/s = 0.7和半顶角6 = 8"的配置！ = 28英寸。 纤细体上的高攻角流动的特点是Sychev相似性参数[42] K = tana/tanE，而不是单独a或E。 对于当前模型，K = 3.7833。

\section{14.4 计算网格和边界条件}
J。 蔡，H-M。 蔡英文，S。 罗t3 F。 刘303图14.1：纤细的锥形翼体组合和分离涡流。 14.4 计算网格和边界条件 众所周知，圆锥体上的亚声速流动不能是严格的圆锥形。 然而，如果圆锥体是纤细的，那么流动几乎是圆锥形的。 如参考[45]所示，1961年Wed6日，在水隧道中观察到E = 15英寸的三角形薄翼，a = 20英寸。 Navier-Stokes的计算也证明了这一点，例如 Thomas、Kirst和Anderson [44]以及我们之前在Ref中的研究。 [4]。 原则上，圆锥流可以用适当的修改方程在二维平面上求解。 然而，本研究坚持在三维网格上使用三维代码，以允许计算非完全锥形流量。 然而，对于本文中研究的近圆锥流，可以使用圆锥网格。 过量锥形网格旨在匹配局部流量梯度，并促进计算的并行处理。 在图中。 14.2，（a）部分给出了入射平面上的网格，（b）部分显示了右半出口平面上的网格，尽管在以下计算中使用了完整的交叉平面网格。 远场侧边界是一个圆锥形表面，它与圆锥体共享相同的顶点，距离体轴25秒，其中s是局部机翼半跨度。 网格在身体的尖端被束成一个点。 由于使用了以细胞为中心的有限体积法，因此在顶点上不会遇到数值困难。 圆锥流计算的纵向只需要几条网格线。 然而，径向和圆周的网格非常精细

涡流对的稳定性图14.2：平板三角翼和圆形锥体的翼体组合的锥形网格，E = 8英寸，y = 0.7，（a）在入射平面上，（b）在右半部出口平面上。 为了清楚起见，只有每4条线在径向和圆周方向上显示。为了稳定性研究的目的，必须使用横向平面中的方向来解析涡流场。 退出平面上全网格的特写视图如图所示。 14.3. 网格由三层组成：iqner层在8个块中有5 x 177 x 581个网格点；中间层在2个块中有5 x 49 x 385个网格点；外层也有5 x 49 x 257个网格点。 内层有两个子层，每个子层有4个块。 内层子层的上半部分有5 x 81 x 387网格点，外层子层的上半部分有5 x 97 x 387网格点。 网格点总数为671,475个。 在由AMD Athlon XP1600+ CPU组成的8处理器并行集群计算机上，双倍（64位）精度的一次迭代的计算时间约为一秒。 零法向速度边界条件适用于机翼和机身表面。 翼尖前缘的Kutta条件在尤拉代码中自动得到满足。 基于特征的条件用于电网的上游边界。 在下游边界上，所有流量变量都被推断。 计算从均匀的自由流开始，直到连续性方程的最大残差减少11个数量级以上。 正如Siclari和Marconi[41]所指出的，需要如此严格的收敛标准，特别是对于高攻角流动的稳定性研究。

J。 蔡，H-M。 蔡英文，S。 罗和F。 刘305图14.3：出口平面上平板三角翼和圆形锥体的翼体组合的圆锥形网格的特写视图，为了清楚起见，图中只有每4条线绘制，E = 8英寸，y = 0.7。 表14.1：三个不同密度的网格。 中间 5 x 65 x 385 5 x 49 x 385 5 x 49 x 321 网格 内层 外层 精细 5 x 49 x 257 基线 5 x 49 x 257 粗 5 x 49 x 225 5 x 209 x 641 5 x 177 x 581 5 x 89 x 387 用于计算的网格密度表示根据网格加密研究和可用计算资源的平衡确定的网格密度。 上面给出的网格被视为基线网格。 计算也在我们可用计算资源范围内的精细网格和表14.1所列的粗网格上进行。 流模型的静止对称解由欧拉求解器在稳态模式下获得，自由流作为初始解，使用三个不同的网格。 表14.2比较了涡流核心中心的计算位置x/s和y f s，以及计算的总速度U/Um和压力系数 C在涡流核心中心。 粗网格上的解决方案产生了静止涡流中心的显著位置偏移和流量参数的相当大的变化，而精细网格上的解决方案产生了微不足道的位置偏移和微不足道的流动参数变化

\section{14.5 静止对称和不对称解决方案及其稳定性}
\subsection{14.5.1 时间不对称扰动}
306 涡流对的稳定性 表14.2：计算的涡核位置和流动属性。 网格x/s y/s U/Um C，精细0.5734 0.9178 2.7742 -9.4063 基线 0.5741 0.9174 2.7288 -9.0903 粗 0.6066 0.9156 2.4357 -7.1355到基线网格上的解决方案。 因此，基线网格用于以下计算。 14.5 静止对称和不对称解决方案及其稳定性 静止涡流配置首先通过在稳态模式下运行欧拉代码来捕获。 获得静止涡流配置后，将由机翼左侧和右侧短时间的吸力和吹气组成的小瞬态不对称扰动引入到流动中，然后以时间准确模式运行尤拉代码，以确定流动是否会恢复到原来的无干扰状态，还是演变成不同的稳定或不稳定的解决方案。 前一种情况表明，原始的静止涡流配置是稳定的，而后一种情况证明它不稳定或中性稳定。 该流动模型的静止对称和不对称涡流配置及其稳定性在参考[7]中通过参考中提出的圆锥形细长体理论进行了分析。 [5]。 14.5.1 时间不对称扰动 不对称扰动包括通过机翼上表面的两个狭窄的锥形槽的吸力和吹气，因为它们是理论上最不稳定的运动模式[5]。 吸气和吹气槽对称地位于涡流芯的下方。 扰动在时间准确欧拉计算的初始时间周期0 < t < 1中被激活，其中t是非维时间。 在机翼的右侧，向下游方向看，吹气速度Vj的定义如下。 左侧槽的吹气速度以相同的方式定义，但方向相反。 两种相同类型的不对称配置

\subsection{14.5.2 静止对称涡流}
J。 蔡，H-M。 蔡英文，S。 罗和F。 刘307图14.4：平板三角翼和圆锥体的翼体组合上对称溶液的交叉流平面上的压力轮廓，E = 8"，y = 0.7，（Y = 28"，/3 = 0”。使用扰动。 扰动（A）有y1 = 0.73s，yz = 0.78s，Vo = 2.0Um，扰动（B）有y1 = 0.90s，y2 = 0.95s，Vo = -1.33Um。 瞬时最大吹动量通量发生在t = 112。 两种配置基于翼面积的瞬时最大吹动量通量系数分别为c/，= 0.1和c/，= 0.04。 相应的吹气力大约比在高攻角下作用在身体上的法力少一个数量级。 与扰动（B）相比，扰动（A）的吸气和吹气槽更靠近翼根。 此外，吹气速度有相反的方向，即吸力和吹气之间存在交换。 图中的箭头标记了扰动（A）的吸力和吹气速度的位置和方向。 14.6. 14.5.2 静止对称涡流 为了研究对称和不对称涡流对的静止位置及其稳定性，必须考虑包括入射平面两侧在内的整个流动空间。 首先通过在流动的整个空间中以稳定模式运行欧拉求解器来搜索静止涡流。 以均匀的自由流为初始解决方案，在三层基线网格上，计算以双倍精度执行，直到最大残差减少11个数量级。 发现稳态解确实与入射平面对称。 图14.4给出了

\subsection{14.5.3 静止对称涡流的稳定性}
308 涡流对的稳定性图14.5：平面三角翼和圆形锥体的翼体组合上对称解的四个交叉流平面中纵向速度分量的轮廓，c = 8"，y = 0.7，a = 28O，p = 0。横流平面上的计算压力轮廓。 两个涡流核心的中心可以清楚地看到对称地位于z/s = 0.5741，yls = +0.9174。 图14.5显示了沿机翼-机身组合的四个横流平面中纵向速度分量的计算轮廓。 在图中，连接机翼前缘和涡核的剪切层出现了。 显然，这个解决方案代表了一个静止的对称涡流，在以下讨论中对其进行稳定性检查。 14.5.3 静止对称涡流的稳定性 为了研究静止涡流的稳定性，将上面得到的对称解用作t = 0的新初始条件，扰动（A）在0 < t < 1时激活，如图中的两个箭头所示。 14.6. 时间准确的欧拉代码用于模拟扰动流动的演变。 在时间准确的尤拉计算中，每单位增量t都要执行50个实时步骤。 对于每个实时步骤，都执行伪稳流计算，直到最大残差减少四个或更高。 图14.6给出了计算的不对称解的交叉流平面中的压力轮廓。 可以看出，即使在最初的扰动消失很久之后，受干扰的流动也不会恢复到最初的静止构态。 相反，它越走越远，直到它到达一个新的稳态解，其中带圆圈的两条线标记了两个涡流中心的轨迹，轮廓是

J。 蔡，H-M。 蔡英文，S。 罗和F。 Liu 909 图14.6：扰动（A）后不对称溶液横流平面中的压力轮廓应用于平板三角翼和圆锥体的翼体组合上的对称静止涡流溶液，E = 8英寸，y = 0.7，a = 28英寸，p = 0；以及扰动涡核的轨迹。 '0 8 16 24 32 40 48 56 时间（t）图14.7：涡核位置与扰动（A）后的时间应用于平板三角翼和圆锥体的翼体组合上的对称静止涡流溶液，E = 8"，y = 0.7，a = 28"，/3 = 0。

\subsection{14.5.4 静止不对称涡流的稳定性}
31 0 涡流对的稳定性图14.8：在平板三角翼和圆形锥体的机翼体组合上，E = 8"，y = 0.7，a = 28"，B = 0。新获得的稳态解在时间准确计算结束时，新获得的稳态解的恒定压力线。 图14.7显示了涡核位置X/S和y/s与无量值时间t。 新款稳态解高度不对称。 与原来的对称解相比，左涡旋向机翼表面向下移动了一小段距离，而右涡旋则明显向上游离机翼表面和内侧。 左侧较低的涡核在x/s = 0.3448和y/s = -0.9127。 右上方涡核在x/s = 1.6320和y/s = 0.5992。 图14.8显示了沿机翼-机身组合的四个交叉流动平面中纵向速度分量的轮廓。 请注意，扰动只在短时间内施加0 < t < 1，而计算继续进行，没有任何外部施加的扰动或不对称，从t = 1到t = 64（见图的abscissa。 14.7.） 这只能通过以下事实来解释：初始对称涡流解（在非常严格的收敛标准下获得）不是一个稳定的配置，此外，新的不对称解是另一个静止的涡流配置。 14.5.4 静止不对称涡流的稳定性 上述计算表明，除了这种高攻角条件下的对称涡流配置外，还存在静止不对称涡流配置。 这种不对称的涡对在小扰动下是否稳定还有待观察。

\subsection{14.5.5 不对称涡流的镜像}
J。 蔡，H-M。 蔡英文，S。 罗和F。 刘'... 图14.9：在平板三角翼和圆锥体的翼体组合上应用扰动（A）后，对不对称溶液的横流平面上的压力轮廓，E = 8英寸，y = 0.7，a = 28英寸，p = 0；以及扰动涡核的轨迹。 为此，时间重置为零，以便时间准确的尤拉计算以新的静止不对称解作为初始条件继续进行。 扰动（A）再次施加在初始单位时间段（0 < t < 1）。 只是为了多样性，这次交换了吸力和吹风方向。 新方向由图中的箭头显示。 14.9. 计算的涡核位置与时间t如图所示。 14.10. 可以看出，解回到了最初的不对称解。 图中涡流中心周围的实线显示了流动演变过程中涡流的轨迹。 14.9. 尽管涡流的偏差持续了相当长的时间，特别是对于上涡流，但两个涡流都回到了原来的位置，因此静止的不对称涡对在小扰动下是稳定的。 14.5.5 不对称涡流的镜像 另一个问题是上面得到的静止不对称解是否取决于特定的初始扰动。 为了回答这个问题，将不同的扰动扰动（B）应用于图中的收敛对称解。 14.5. 扰动（B）的功能形式与扰动（A）相同，只是它的吸力和吹气位置互换，并靠近翼尖，其最大速度V更小。 足够的

312 涡流对的稳定性 Y g >!!!!! '0 5 10 15 20 25 30 35 40 45 50 时间（t）图14.10：涡核位置与施加扰动（A）后的时间，在平板三角翼和圆形锥体的翼体组合上反向吸/吹到不对称溶液，E = 8"，y = 0.7，a = 28O，，8 = 0...- i., \textbackslash\{\}. 'i. /图14.11：在平板三角翼和圆锥体的翼体组合上对对称溶液施加扰动（B）后，不对称溶液的横流平面上的压力轮廓，E = 8英寸，y = 0.7，Q = 28英寸，8 = 0； 以及扰动涡核的轨迹。

\subsection{14.5.6 当前欧拉求解器的对称性质}
J。 蔡，H-M。 蔡英文，S。 罗和F。 刘313 I I I I I I -。 UI f H 8 n s 0.8 0.4 '0 8 16 24 32 40 48 56 64 时间（t）图14.12：涡核位置与时间应用扰动（B）到平板三角翼和圆锥体组合的翼体组合的对称解决方案后的时间，E = 8"，y = 0.7，Q = 28"，p = 0。和以前一样进行相同的时间准确计算。 同样，初始收敛解在初始（短时间）不对称扰动下与对称性分离，最终收敛到不对称涡流配置，如图中的压力轮廓所示。 14.11. 涡流核心从初始对称位置到静止不对称位置的轨迹由图中的圈线标记。 14.11. 图中显示了两个涡流的无维坐标的时间历史。 14.12. 这一次，由于吸力和吹气的方向颠倒，左涡旋向上向右移动，而右涡旋向下移动，最终形成了图中所示解决方案的完美镜像。 使用扰动（A）获得的14.6和14.7。 由于扰动（B）与扰动（A）不仅在吸力和吹气的方向上不同，而且在强度和位置上也不同，因此计算表明，除了精确的镜像的可能性外，流动不对称与初始扰动无关。 显然，镜像静止不对称涡对也必须稳定。 因此，流动是双稳定的。 14.5.6 上述计算还证明了算法的对称性和本研究中使用的三维时间准确欧拉求解器的计算机代码。 这种对称性对于流动不稳定性的数值研究是非常可取的。 否则，图中获得的静止对称涡流。 14.5将难以捉摸，因为当前欧拉求解器的数值对称性质中的不对称性

\subsection{14.5.7 与稳定性理论预测的比较}
\subsection{14.5.8 与稳定性实验数据的比较}
314 Voytex对的稳定性计算可以提供必要的扰动来触发物理不稳定，从而推动流动走向不对称解。 14.5.7 与稳定性理论预测的比较 在参考[7]中，通过蔡等人的分析方法研究了一些机翼体组合。[ 5]。 对于计算模型的情况，即平板三角翼和圆形锥体的机翼体组合，€ = 8"，y = 0.7，Q = 28"，B = 0和K = 3.7833，这里总结了分析结果。 1. 对称和非对称的静止涡流对都存在。 相对于对称涡流，不对称对的上涡旋位于一个小舷外，下涡位于一个小内。 2. 静止对称涡对在对称扰动下是稳定的，而在反对称扰动下是不稳定的。 3. 静止的不对称涡对在对称和反对称扰动下都是稳定的。 在反对称扰动下，上涡旋比下涡更不稳定。 目前对翼-体组合上的涡流的欧拉计算完全符合关于形成静止对称和不对称涡流对及其在小扰动下稳定性的分析预测，甚至包括上涡旋比下涡移动到其未受干扰位置的时间更长和更大的区域，如图所示。 14.10，这表明理论预测的上涡不如下涡稳定。 14.5.8 目前的计算结果与Murri和Rao[32]制作的前体横向/机身配置的低速风洞实验数据进行了比较。 测试的前体是圆锥体和平板三角翼的圆锥组合，y = 0.77，E = 12'。 后体是一个圆形圆柱体。 基于机身直径的雷诺数是Re = 1.9 x 10'。 测试表明，在K = 3.7833或Q = 38.54"时，存在稳定的非零偏航矩。 它表明稳定的静止不对称涡对在流场上占主导地位。 值得注意的是，测试模型的y值比计算模型的y值大10\%。 然而，它足够接近，可以用来确认静止稳定不对称涡对存在的当前计算结果。 与稳定性实验数据的比较

\section{14.6 涡核的结构}
\subsection{14.6.1 计算结果}
J。 蔡，H-M。 蔡英文，S。 罗和F。 Liu 315 两个镜像双稳定不对称涡流配置在高攻角下旋转的尖鼻细长体上，就像图所示的当前计算模型一样。 在一些风洞实验中观察到了14.9和14.11。 例如，Zilliac等人[47]在低速和低湍流级风洞中使用基于气缸直径的雷诺数的六组分应变计平衡仪测量了3.5口径Ogive气缸体的总侧力。 斜鼻的半顶角为16.26英寸。 滚动角度的总体侧力系数是50" < a < 60"或等效4 < K < 6的方波曲线。 在这种情况下，不对称只有两种稳定状态或双稳定状态。 在整个滚动角度范围内，侧力系数从恒定正值突然切换到相同大小的恒定负值，反之亦然，没有发现中间侧力系数。 结果表明，测试模型尖端几何形状的微观变化对流场的下游发展有很大影响，双稳定的李流场的存在是不稳定的结果。 Ng和Malcolm[34]在流量可视化水隧道中进行了实验。 测试模型是6\%比例的F/A-18正体。 模型鼻子几乎是轴对称的，半顶角约为30英寸。 在约60英寸或K = 3的攻角以上，流动可视化显示，流动变得双稳定。 偏航矩可以通过位于背风侧和鼻尖附近的喷嘴吹来的瞬态和小量质量流量喷射器在两个基本稳定的值之间切换。 切换后保持吹气只会使偏航时刻增加相对较小。 在a = 65"和67.5"或K N 4处观察到类似的结果。 值得注意的是，目前的计算模型是y = 0.7的纤细翼体组合，而上述两个测试模型仅为细长体，即y = 1。 然而，对测试模型的双稳定流场的观察在质量上证实了计算和理论结果。 14.6 涡核的结构 上述计算的流场在横流平面上显示身体背向侧的封闭和几乎同心的圆形图。 这些代表由与机翼前缘分离的剪切层的螺旋卷起形成的初级涡流的核心。 本节检查了共蟟涡芯的结构，并与实验观察进行了比较。 14.6.1 计算结果 上一节中介绍的静止对称解的右涡旋被选为一个典型的例子。 核心的中心位于

31 6 涡流对的稳定性 0.4 08 09 1 11 图14.13：平板三角翼和圆形锥体的翼体组合的横流平面上涡流芯区域圆锥网格的局部和详细视图，E = 8"，y = 0.7。x/s = 0.5741，y/s = 0.9174。 图14.13显示了用于解析涡核区域的本地但详细的网格配置。 纵向速度分量u,/U,、纵向涡量分量wzs/U和总压力损失系数C,t的轮廓分别如图14.14、14.15和14.16所示，其中CPt = (pt - pt,)/(p,U\&/2)，pt是局部总压力，pt是自由流总压力。 三个流动参数的所有轮廓几乎都是围绕涡流中心的同心圆。 所有三个流量参数的值都向中心增加。 为了进一步检查涡核的结构，两条正交网格线穿过涡核。 两条网格线的交点是离涡流中心最近的网格点。 沿着径向和圆形网格线绘制各种流动参数的分布与从两条网格线的交点的距离T。 图14.17给出了uz/Uw和w,s/U的分布，与沿径向（从中心体）和圆形（围绕体）网格线的r/s的分布。 图14.18显示了Cpt和静态压力系数 C的分布，与沿两条网格线的r/s的分布。 图14.19给出了速度分量ua/U和ue/U与r/s沿两条网格线的分布。 在这里，速度沿着圆柱坐标（a、T、O）的方向分解成分量，其中轴a与涡芯轴重合。 可以看出，沿着两条垂直网格线的每个流动参数的分布几乎一致，它们几乎对称

J。 蔡，H-M。 蔡英文，S。 罗tY F。 刘图14.14：在平板三角翼和圆锥体的翼体组合上的对称溶液的交叉流平面上涡流核心区域纵向速度分量的轮廓，E = 8英寸，y = 0.7，= 28英寸，B = 0，A（uZ/Um）= 0.2。 图14.15：在平板三角翼和圆锥体的翼体组合上对称溶液的横流平面上，涡流核心区域的纵向涡量分量，E = 8英寸，y = 0.7，a = 28英寸，I3 = 0，A（wZs/Um）= 5。

318 涡流对的稳定性图14.16：涡流核心区域总压力损失系数的轮廓，在平板三角翼和圆形锥体的翼体组合的对称溶液的交叉平面上，E = a"，y = 0.7，QI = 28"，8 = 0，AC，= 0.2。 6.0 I I I I I 图14.17：速度和涡流的纵向分量与从涡流芯中心距离的分布，沿着径向和圆形网格线穿过涡流芯中心，在平板三角翼和圆形锥体的翼体组合的对称解决方案的横流平面上，E = a"，y = 0.7，a = 28"，，8 = 0。

J。 蔡，H-M。 蔡英文，S。 罗tY F。 Liu 319 - r - - -I - - - - -'0'00.4 -0.3 -0.2 -0.1 0 0.1 0.2 0.3 0.4 图14.18：总压力损失系数和静态压力系数与涡流芯中心沿径向和圆形网格线穿过涡流芯中心的距离分布，在平板三角翼和圆形锥体的翼体组合上的对称溶液的交叉流平面上，E = 8"，y = 0.7，CY = 28'，p = 0.到原点r/s = 0，表明涡核中的流动几乎是轴对称和圆锥形的。 涡量向核心中心急剧增加，并在中心达到最大值。 图。 14.17表明，高度旋转区域的边缘位于r，= 0.2s，即旋转或涡核的半径。 在旋转区域内，黏性扩散完全平滑了速度分布的梯度，不能再检测到剪切层。 在涡流或旋转核心内，静态和总压力向涡流中心下降，并在涡流中心达到最小值。 轴向速度分量ua向涡流轴增加，并在涡流轴上达到最大值，而周长速度分量UQ首先向涡流轴增加，在达到涡流轴附近的最大值后，在涡流轴上急剧减少到零。 最大UQ的位置定义了子核心的边缘，在这个子核心内，速度和压力的巨大梯度占主导地位，数值黏性力占主导地位。 这个子芯是一个黏性子芯，其中会导致人工的总压力损失。 从图14.19和14.18中，子核的半径，r,, = 0.08s。 计算出的径向速度分量比其他两个速度分量小一个数量级，因此没有显示出来。 在涡核中，径向速度分量指向核心轴。 它首先增加

\subsection{14.6.2 与实验数据的比较}
320 涡流对的稳定性 I I I I,,,,l,,,,l,,,,,,L,,,,I,I,,,,,,,,,, -2.0 - - - - - - - - - - - - - - - - - -1- - - - - -0.4 -0.3 -0.2 -0.1 0.1 0.1 0 0.2 0.4 图 14.19：速度的轴向和圆周分量与涡流芯中心距离的分布，沿着平板三角翼和圆形圆锥体的对称溶液的翼体组合的对称溶液的横向和圆形网格线的涡流平面上穿过涡流核心中心，6 = 8"，y = 0.7，Q = 28"，p = 0。向涡流轴，然后在核心轴急剧减少到零。 14.6.2 与实验数据的比较 可以想象，涡核在流动体相互作用中起着重要作用。 本小节将计算的涡核与已知的实验数据进行比较。 到目前为止，目前作者没有已知的翼体组合的实验数据。 仅使用一对关于尖锐前缘三角翼的可用实验数据。 两个测试模型的几何形状几乎相同，设置在大约相同的攻角上，但雷诺数明显不同。 测试1：Carcaillet、Manie、Pagan和Solignac[S]使用三维激光测速仪和1米研究低速风洞压力探头的五孔压力探头测量了锋利前缘平板三角翼上的涡核。 测试的机翼在a：= 20英寸处有E = 15英寸，基于根弦的雷诺数，CO，是0.7 x lo6。 这里使用了60\%\textasciitilde{}的测量横流平面的数据。 参考图11给出了静压、总压力损失、流向速度分量和沿主涡核中心的横向速度分量的测量分布。 [9]。 它们与图14.18和14.19中报告的相应计算分布相似。

J。 蔡，H-M。 蔡英文，S。 罗和F。 Lau 321 表 14.3：计算的涡流核心参数与测试数据的比较。 流动变量计算测试1测试2（uz）max/Um 2.7 3.0 3.0（ug）maz / urn 1.3 1.4 1.3（Wz）max slum 53 152 224（Cp）min -9.2 -11 -13（Cpt）min -2.9 -2.8 -4.8 rc/s 0.2 0.2 rscls 0.08 0.08 0.03 测试2：Verhaagen和Kruisbrink [46]在低速和湍流水平约为0.05\%的风洞中使用五孔压力探头测量了前缘涡旋的圆锥部分的流动特性。 该模型是一个锋利的前缘平板三角翼，E = 14'，a = 20.4'。 雷诺数是3.8 x lo6，基于模型根弦长度，CO。 测量横流平面为50\%\textasciitilde{}。 参考图1、10和11显示了轴向涡量、轴向速度分量、周长速度分量、静压和沿着穿过主涡核中心的横向的总压力损失的测量分布。 [46]。 它们也与图14.17、14.19和14.18中报告的相应计算分布相似。 可以看出，计算分布的趋势与测试结果非常一致。 表14.3列出了各种特征流量参数的大小，并进行了比较。 尤拉计算很好地预测了大多数重要的流动参数。 特别是，预测的总压力损失水平是非常现实的。 考虑到这是一个纯粹的虚假数字人工制品，计算和实验之间的这种一致性可能是假的。 然而，在使用欧拉方法对锋利边缘三角翼的其他自由涡流模拟中也发现了类似的观察，例如 Murman和Rizzi[31]，Rizzi [38]，以及Powell，Murman，Perez和Baron [36]。 他们进行了系统研究，其中改变了各种计算参数。 特别是网格空间和人工阻尼系数的变化幅度很大。 他们发现，尽管涡流区域在较粗的网格和/或具有高阻尼常数上更分散，但总压力损失的大小对所有计算参数不敏感。 Rizzi[38]声称，总压力损失与网格大小的不变性似乎是由解中的奇点引起的。 此外，Rizzetta和Shang [37]报告说，Euler解的总压力轮廓与Navier-Stokes计算的总压力轮廓几乎相同，除了二次流动区没有重新-

\section{14.7 总结和结论}
322 在黏性结果中产生的涡流对的稳定性。 就像尖锐前缘的分离对黏度不敏感一样，子芯的总压力损失对黏度不敏感。 翼的锋利边缘和涡核的中心都是欧拉微分方程的奇点。 数值耗散平滑了奇点。 围绕尖锐边缘的涡量的生成和涡核中心的总压力损失都对数值耗散的实际大小不敏感，只要有一些。 值得注意的是，计算的最大涡量低于实验数据。 事实上，正如Nelson和Visser [33]所指出的，不同研究人员的最大轴向涡量的实验结果差异很大。 从对每次调查实验测量中使用的网格分辨率的检查中，他们发现最高的推导涡量值对应于最好的网格分辨率，反之亦然。 计算的最大轴向涡量的较低值可能是由于计算网格分辨率不足。 根据图。 14.13，沿着径向和周长方向分别有大约100 x 40的网格点，位于涡核的横截面，在子核心中大约有50 x 20的网格点。 网格似乎仍然不够好，无法解决子核心中的流量。 14.7 总结和结论 我们简要回顾了关于高攻角下细体上对称和不对称分离涡流的稳定性的不同理论。 文献中提出了对流和绝对稳定性机制。 此前，作者为细圆锥体开发了绝对稳定性标准。 在本文中，我们提出了一个使用重叠网格的时间准确三维欧拉代码来研究涡流稳定性问题。 Euler 流动求解器计算了在理论分析确定的典型流动条件下的完整三维流动，这些流动条件从稳定到不稳定的制度。 固定涡流配置首先通过在稳态或时间准确模式下运行尤拉代码来捕获。 获得静止涡流配置后，将由机翼左侧和右侧短时间的吸气和吹气组成的瞬态不对称扰动引入流动，并以时间准确模式运行欧拉代码，以确定流动是否会恢复到原始不受干扰的状态，还是演变成不同的稳定或不稳态解决方案。 还检查了尤拉计算解析的涡核结构的特征。 从计算中得出以下结论。 1. 欧拉方法可用于有效研究分离涡旋在翼体组合上的稳定性。 计算的欧拉解决方案自动满足尖锐前导的库塔条件

\section{14.8 参考书目}
J。 蔡，H-M。 蔡英文，S。 罗9欧元F。 纤细三角翼的刘323边缘，捕捉从前缘脱落的自由剪切层，并将其发展成机翼侧的紧凑和连贯的旋转核心。 必须使用精细网格、64位计算的高精度和严格的收敛标准来解决紧密的涡流。 数值算法也应该对称，以保持计算的静止对称涡流配置的对称性。 2. 有限但明智地选择的计算案例与作者之前的理论分析和一些实验观察非常一致，因此支持这样的结论，即绝对类型的流体动力学稳定性是导致许多典型的涡流在高攻角下在纤细的圆锥体上破坏对称性。 3. 欧拉计算再现了涡核的基本特征。 可以根据涡流中切向速度的变化识别出显著旋转流的清晰核心和子核心。 Euler方法很好地模拟了涡核的基本特征。 特别是，尽管流动模型具有隐性，但子核心的预测总压力损失是现实的。 看来，Euler 方程溶液中的数值黏度在奇点附近适应，如机翼的尖锐前缘或非黏性涡旋的中心，从而产生合理接近实黏性流动的解。 14.8 参考书目[l] Bernhardt，J. E。 和Williams D. R。 不对称前体涡流的比例控制。 AIAA杂志，36（11）：2087-2093，1998年11月。 [2] Bryson，A. E。 圆形圆柱体和圆锥体上的对称涡流分离。 J。 应用程序。 机械。 （ASME），26：643-648，1957年。 [3] Cai, J., Tsai, H-M., \& Liu, F. 用于多网格和并行计算的黏性计算的重叠网格求解器。 AIAA论文2003-4232，2003年6月。 [4]蔡，J.，蔡，H-M.，罗，S. 和刘，F。 涡流对在细锥体上的稳定性-理论和数值计算。 AIAA论文2004-1072，2004年1月。 [5] Cai, J., Liu, F., \& Luo, S. 二维和三维细锥体对称涡流的稳定性。 J。 Fluid Mech.，480：65-94，2003年4月。

324 涡流对的稳定性 [6] Cai, J., Luo, S., \& Liu, F. 对称和不对称涡流对在纤细的锥形翅膀和体上的稳定性。 AIAA论文2003-1101，2003年1月。 [7] Cai, J., Luo, S., \& Liu, F. 对称和不对称涡旋对在纤细的锥形翼体组合上的稳定性。 AIAA论文2003-3598，2003年6月。 [8] Cai, J., Luo, S., \& Liu, F. 对称和非对称涡流对在纤细的锥形机翼和体上的稳定性。 流体物理学，16（2）：424-432，2004年2月。 [9] Carcaillet, R., Manie, F., Pagan, D., \& Solignac, J. L。 前缘涡流在75度扫过三角翼-实验和计算结果。 ICRS 86-1.5.1，1986年9月。[lo] Champigny，P. 高攻击角度的侧力。为什么、何时、如何？ La Recherche Aerospatiale，（4）：269-282，1994年。[ll] Cummings，R。 M.，Forsythe，J. R.，莫顿，S. A.，\& Squires，K. D。 高角度攻击流预测中的竞争挑战。 航空航天科学进步，39（5）：369-384，2003年5月。 [12]德加尼，D。 几何扰动对涡旋不对称的影响。 AIAA杂志，29（4）：560-566，1991年4月。 [13] 德加尼，D。 大量发生时流过身体的不稳定。 AIAA [14] Degani，D。 \& Tobak，M。 受控尖端扰动对流动不对称的影响的实验研究。 流体物理学A，4（12）：2825-2832，1992年。 [15]德克斯特，P。 C. \& 和亨特，B. L。 在高攻角下，滚动角度对细长体转流的影响。 AIAA Paper Journal，30（1）：94-100，1992年1月。 1981-0358，1981年。 [IS] 戴尔，D。 E.，Fiddes，S. P。 \& 史密斯，J. H。 B. 入射时锥体的不对称涡流形成\textasciitilde{}一个简单的非黏性模型。 航空季刊，33，第4部分：293-312，1982年11月。 [17]爱立信，L. E。 零侧滑处高α涡旋不对称的来源。 飞机杂志，29（6）：1086-1090，1992年11月至12月。 [18]爱立信，L. E。 \& 雷丁，J。 P。 旋转体上的不对称流动分离和涡流脱落。 在M.J. Hemsh，编辑，《战术导弹空气动力学：一般主题，航天和航空学进展》，第141卷，第391-452页，纽约，1992年。 美国航空航天学会。

J。 蔡，H-M。 蔡英文，S。 罗和F。 刘325 [19] 哈特维奇，P. M。 圆锥体上涡流的对称性断裂-理论和数值实验。 AIAA期刊，32（5）：1013-1020，1994年5月。 [20] 哈特维奇，P. M.，霍尔，R. M.，和Hemsch，M. J。 由小表面缺陷控制的涡流不对称的Navier-Stokes计算。 航天器和火箭杂志，28（2）：258-264，3月-4月。 1991年。 [21]黄，M。 K。 \& 周，C. Y。 前缘涡对在纤细的三角翼上的稳定性。 AIAA杂志，34（6）：1182-1187，1996年6月。 [22] 亨特，B。 L。 不对称的漩涡在纤细的物体上力和唤醒。 AIAA论文1982-1336，1982年。 [23]拉蒙特，P. J。 带有层压、过渡或湍流分离的倾斜倾斜圆柱体周围的压力。 A IAA杂志，20（11）：1492-1499，1982年11月。 [24] Legendre，R. Ecoulement au voisinage de la pointe aavant d'une aile a forte flBche aux incidences moyennes。 La Recherche Aeronautique, Bul- letin Bimestriel, De L'Ofice National D'Etudes Et De Recherches Aero- nautiques, Jan-Feb 1953. [25] Levy, Y., Hesselink, L., \& Degani, D. 未能保存对称性的算法产生的流动异常不对称。 AIAA杂志，33（6）：999-1007，1995年6月。 [26] Levy, Y., Hesselink, L., \& Degani, D. 对几何扰动和流动不对称之间的相关性进行系统研究。 A IAA杂志，34（4）：772-777，1996年4月。 [27]刘，F。 \& 詹姆森，A. 三维级联的多网格导航-斯托克斯计算。 AIAA杂志，31（10）：1785-1791，1993年10月。 [28]刘，F。 \& 郑，X。 一种强耦合时间行车方法，用于用多网格求解navier-stokes和k-w 湍流模型方程。 计算物理学J.，128：289-300，1996年。 [29]刘，F。 \& Ji，S。 使用多网格navier-stokes方法进行非定常流动计算。 AIAA杂志，34（10）：2047-2053，1996年10月。 [30]马可尼，F。 在超音速流中围绕尖锐锥体的不对称分离流动。 在第11届实习生的程序中。 流体力学数值方法会议，第395-402页，1988年7月。 [31] 默尔曼，E。 M。 \& Ftizzi，A。 Euler 方程应用于带有前缘涡流的锋利边缘三角翼。 1986年4月7日至10日，法国普罗旺斯省议会帕拉伊斯，计算流体动力学在航空应用的AGARD会议记录中，第15-1-15-13页。

326 涡流对的稳定性[32] Murri，D. G。 和Rao，D。 M。 在高攻角下用于偏航控制的可激活前体带的探索性研究。 AIAA论文87-2557，1987年9月。 [33] Nelson，R. C. 和维瑟，K。 D。 分解三角翼漩涡--涡量在分解过程中的作用。 AGARD CP-494，pp. 1990年1月21日至21日。 涡流空气动力学。 [34] Ng，T。 T。 \& 马尔科姆，G. N。 使用前身吹气和吸力进行空气动力学控制。 AIAA文件91-0619，1991年1月。 [35] Pidd，M. \& 史密斯，J. H。 B. 不对称的涡流流在圆形圆锥体上。 在涡流空气动力学中，AGARD CP-494，第18-1-11页，1991年7月。 [36]鲍威尔，K。 G.，默尔曼，E. M.，佩雷斯，E. S.，\&男爵，J. R。 圆锥形Euler 方程的涡流溶液中的总压力损失。 AIAA杂志，25（3）：360-368，1987年3月。 [37] Rizzetta，D。 P。 \& 尚，J. S。 前沿涡流的数值模拟。 AIAA杂志，24（2）：237-245，1986年2月。 [38] Rizzi，A。 具有100万个网格点的Euler 方程的三维解决方案。 AIAA杂志，23（12）：1986-1987，1985年12月。 [39] Sadeghi, M., Yang, S., \& Liu, F. AIAA论文2003-1347的平行计算，1月耦合navier-stokes/csd方法。 2003年。 [40] 香克斯，R. E。 前缘扫角为70英寸至84英寸的六个平板机翼的静态和动态稳定性导数的低亚音速测量。 美国宇航局TN D-1822，1963年。 [41] Siclari，M。 J。 \& 马可尼，F. 计算表现出不对称涡流的Navier-Stokes解。 AIAA杂志，29（1）：32-42，1991年1月。 [42] Sychev，V。 V。 三维高超音速气体以高攻角流过纤细的物体。 数学与机械杂志。 （苏联），24：296-306，1960年。 （431 Thomas，J. L。 雷诺数对锥体上空超音速不对称流动的影响。 航空杂志，30（4）：488-495，Ju1.-Aug. 1993年。 [44]托马斯，J。 L.，Kirst，S. T.，和Anderson，W. K。 Navier-stokes在低长宽比机翼上涡流的乘法。 AIAA杂志，28（2）：205-212，1990年2月。 [45] 范戴克，M. Fluid Motion的专辑。 抛物面出版社，1982年。照片90，第54页。

J。 蔡，H-M。 蔡英文，S。 罗63 F。 刘327 [46] Verhaagen，N. G。 \& Kruisbrink，A。 C. H。 前缘涡旋的牵引效应。 AIAA杂志，25（8）：1025-1032，1987年8月。 [47] Zilliac，G。 G.，Degani，D.，和Tobak，M. 革命细长体上的不对称涡流。 AIAA杂志，29（5）：667-675，1991年5月。

此页面故意留空

\clearpage
\chapter{第15章 上游条件对全球分离Oleg S时湍流边界层速度赤字曲线的影响。 里佐夫}
\section{15.1 介绍}
第15章 上游条件对全球分离Oleg S时湍流边界层速度赤字曲线的影响。 Ryzhov' 15.1 介绍 不可压缩流动的分离是流体力学中最复杂的p--3nomena之一。 自Prandtl[13]提出边界层概念以来的几十年来，它无法接受理性的分析处理。 Stewartson[19]、Neiland [ll]和Messister [lo]以层流运动的三层流动模式的形式提出的Prandtl想法的延伸，为理解围绕小型局部障碍物再循环的不可压缩流体气泡的起源提供了线索。 Sychev [20]和Smith [17]是第一个连续解决层边界层与圆柱体等光滑体大量分离问题的人。 湍流边界层提出了一个严峻的挑战，因为有必要引入闭包假说，以便使数学“加州大学戴维斯分校机械与航空工程系，加利福尼亚州戴维斯95616

\section{15.2 奇异的不黏性压力梯度}
330 湍流边界层 分离说明完成。 不管这种假设如何，Yajnik [22]和Mellor [8]为平均湍流运动开发了渐近展开。 该理论的一个新颖特点是，它涉及一个不确定的方程组。 尽管如此，他们的结果涵盖了速度赤字定律和von Kkmh的对数定律。独立地，Sychev和Sychev [21]努力解决从光滑表面大规模分离湍流的问题，而没有诉诸特定的闭合假说。 这项工作后来受到Melnik [9]的批评，他分析了边缘湍流分离，导致Goldstein类型的奇点。 最近，Barenblatt、Chorin和Prostokishin[3]在Barenblatt中提出的不完全相似性的背景下讨论了非零压力梯度边界层的渐近特性。湍流、非零压力梯度边界层的实验数据在Schubauer \& Klebanoff [16]、Coles [6]、Samuel \& Joubert [15]、MaruSiC \& Perry [7]和Castillo \& Walker [4]中可用，仅提及有关该主题的基本论文中，仅提及该主题的基本论文。 Perry等人[la]在风洞测试的基础上建议的模型对我们的目的尤为重要。 然而，在上述文献中没有报告对湍流边界层速度和皮肤摩擦曲线在大规模分离阈值处的直接测量。 PI。 15.2 奇异的无黏性压力梯度 自由流线从身体表面分离是大规模分离的典型特征。 这种分离在压力梯度中产生了一个奇点，该奇点支配着速度场。 根据Sobey [18]中可用的结果，压力梯度（15.1）沿着分离点上游的长度2 < 0，无论是层流还是湍流。 然而，层流和湍流流之间存在根本区别。 Brillouin-Villat条件对层压黏性压力分布施加了严重的约束b+ =O（15.2）。 因此，理论上将与前体在角度坐标8 M 55"（Sobey [18]）分离。 这个估计与实验数据（Achenbach [l]）无关。 Sy-chev [20]假设层隔离时边界层外的流动服从（15.1），但bi取决于雷诺数 R通过bl = bl（R）= by'R-A + 0作为R + m。 （15.3）25 2

\section{15.3 控制方程}
0. S。 Ryzhov 331 Smith[17]获得的全三层解决方案表明hio）M 0.22。 因此，hi -+ 0作为R + DC）与inviscid分析一致。 Brillouin-Villat条件（15.2）不适用于湍流。 实验数据提供了强有力的证据，表明在完全湍流环境中，hi = O（1）（15.4），分离点在角坐标8 M 120"（Achenbach [l]）下远移。 这大致对应于无黏性极限解，其中自由流线变得与迎面而来的溪流方向平行于无穷大。 15.3 控制方程 让x和y分别表示沿着物体表面和垂直的局部坐标，u和v代表相应的平均速度分量，u'，v'意味着脉动。 湍流边界层中的二维运动方程可以写成au av -+- = 0，ax ay（1 5.5a）（15.5b）au ax ay u-+v- =，其中v像往常一样表示运动黏度，沿着物体的压力来自（15.1）。 运动方程右侧的黏性项离近壁区域很小，下面会被忽略。 这个方程组没有闭合，因为额外的项T\textasciitilde{}\textasciitilde{} = - (u'v')进入（15.5b）。 我们的目标是制定一种独立于湍流模型上任何具体假设的渐近方法，将雷诺应力与平均运动量联系起来。 Yajnik [22]和Mellor [8]首先将类似的想法应用于附着的边界层和通道流。 如果湍流脉动被认为是小的并丢弃的，（15.5b）简化为Cole \& Aroesty [5]引入的与吹硬问题相关的薄非黏性边界层的方程。 因此，过渡和湍流分离的数学基础与这两种现象是共同的。 连续性方程（15.5a）表明引入了流函数，这样

\section{15.4 隐形子层1}
332 湍流边界层 分离axc' 'u = -- 因此，x-动量方程（15.5b）减少到+- a\$ a2+a+az\$ ap aTzy ayaxay axay2 ax ay - - - - - - - - - (15.6a) ( 1 5.6b) (15.7)，省略了v@。 正如多结构边界层理论中经常发生的情况一样，速度场被分为两个域，每个域都由根据湍流脉动所起的作用具有不同属性的子层组成。 在预分离运动的第一个域中，压力分布来自（15.1），而在发生分离的第二个域中，压力来自特定的非黏性/非黏性相互作用，事先未知。 非黏性/非黏性相互作用的性质类似于在早期非线性过渡阶段控制索子形成的性质（Ryzhov \& Bogdanova- Ryzhova [14]）。 以下分析的预分离流域可能会受到上游条件和分离时奇异压力梯度的影响。 15.4 Inviscid子层1（15.3）和（15.4）之间的区别至关重要。 层预分离流可以通过使用基于线性化的方法来处理，b+作为小参数。 然后，以分离点为中心的域的缩放自然出现，并提供三甲板理论的典型参考长度和规范函数（Sychev [20]；Smith [17]）。 在压力梯度中处理由强奇点（15.4）控制的湍流分离时，不允许线性化。 势预分离运动的不黏性子层的速度场来自伯努利积分U；（x）-4b；（-x）；-2bl（-X）+... = 1（15.8）因此，在x中达到三阶精度，我们有一个扩展U1（x）=1 + 2b；（-x）i +（bl - 2bi）（-x）+... （15.9）显示导数\%的奇异行为。 因此，雷诺应力会影响二阶速度场。

\section{15.5 外层湍流子层2}
\section{15.6 外部 湍流 子层 3}
0. S。 Ryzhov 333 15.5 外湍流子层2 这里是渐近展开的形式 \$ = 20 (Y) + (-44 \$21 (y) + (-1 \$22 (Y) +... （15.104 7 = 720（y）+... （15.10b）由进入速度分布的奇点决定（15.9）。 将（15.10a，15.10b）替换为（15.9）表明，来自边界层上游历史的零阶函数\$20（y）仍然是任意的，而\$21来自d+20 d\$2l d2*20 dY dY -\textasciitilde{}7@21 = 2b+。 因此，速度场的一阶近似值完全不依赖于湍流雷诺应力。 然而，下一个近似值是，因此，雷诺应力会影响二阶速度。 这是外部湍流子层2与非黏性电位子层1的主要区别，其中湍流脉动在二阶速度场中不发挥作用。 渐近膨胀（15.9）和（15.10a，15.10b）分别保持在黏性子层1和外部湍流子层2中的渐近膨胀（15.9）和（15.10a，15.10b）的匹配很容易实现，因为15.6外湍流子层3位于下方的这个子层中的解是由极限y + 0中速度场的行为决定的。 让20美元--t aoyQfl +...，作为y ---f 0。 （15.12）然后从（15.10a）和（15.11）中得出

\section{15.7 压力主导的流动模式}
334 湍流边界层 分离 2bi \$ -+ aoya+l+... + [(a + 1) (1: 2a)ao (-z)' +... + (a + 1) aoAy" +... （15.13）到一阶精度。 因此，a的值决定了湍流子层3中的流动模式。 奇异压力梯度的主导影响。 在这种情况下，方括号中的第一个项在大小上与（15.13）右侧的前序项相当，前提是y - (-z) zz。 子层3的自相似变量由11 a>- (=--- Y (-X)" 2 (15.14) 上游条件的主导影响引入。 如果方括号中的第二个项在大小上与（15.13）右侧的前序项相当，则该制度就会生效。 然后y - (-z)i和自相似变量11 a<- (=- Y (-X)\% 2 (15.15)完全不依赖于a。当a取值使(15.13)中的渐近无效时，（15.14)和(15.15)给出的定义是一致的。 相反，我们有压力梯度与上游条件保持平衡。 两者都是（15.16）20美元（y）的极限，因为y -+ 0是建立位于湍流子层2下方的扰动场的决定性因素。 15.7 压力主导的流动模式 我们以U+l \$ = (-z) 4a \$30 (<) (15.17) 的形式寻求零阶解决方案，将（15.17）替换为（15.7），其中（定义为（15.14）。 右侧省略的雷诺应力术语给我们留下了（1 5.18）

0. S。 Ryzhov 335 这个方程可以通过传递到30美元作为一个新的自变量，并作为所需的函数进行积分。 因此，（15.18）在（15.19）的强度上减少到（15.19）（15.20），C是一个任意常数，通过匹配外部湍流子层2的一阶解来确定。 在6 -+ 00的极限中，我们从（15.13）中推导出两个项渐近点4（1-a）b\$ 1 - 2a，导致一个唯一值c =（a + 1）Z uo*进入（15.20）。 值得将上述解决方案与Sychev和Sychev [21]在von Mises变量中得出的解决方案进行比较。 在这里采用的符号中，这个解决方案与指数a --+ 00 in（15.12）一起读。 因此，流函数随着y + 00呈指数增长，并且无法通过极限条件（15.10a）匹配。

\section{15.8 与实验的比较}
\section{15.9 结论}
336 湍流边界层 分离 20.5 21.0 21.5 22.0 23.0 - - - Xin (xS - x)in 92.4 86.4 80.4 74.4 68.4 62.4 56.4 50.4 44.4 38.4 U - UCC 0.495 0.483 0.450 0.442 0.404 0.409 0.396 0.368 0.357 0.314 0.276 (\&)* x 102 6.00 5.44 4.10 3.82 2.66 2.80 2.46 1.83 1.62 0.97 0.58 表15.1：来自G.B. 舒博和P.S. Klebanoff，NACA Rept。 2133，1950年。 15.8 与实验的比较 基于物理论证和实验数据，Perry等人[12]提出了一个在壁层中有三个区域的模型。 在测试拟议的模型时，他们在平均速度分布中发现了广泛的半功率区域，有些区域延伸到自由流。 这与（15.16）相关，前提是a = i。 从理论上讲，a = i不是（15.12）中指数的唯一值，因为我们看到其他值，包括a > i和a < i，也是可能的。 然而，具有a = \$的制度是值得注意的，因为强压分度在确定湍流边界层中的扰动模式时与上游条件保持平衡。 与Schubauer和Klebanoff[16]的研究进行验证上述理论的再一次比较，其与预分离流动相关的数据总结在表15.1中，并显示在图15.1中。 在他们的工作中，X显然是一个位置x\textasciitilde{}f，= 25.7英尺= 308.4in的完全分离，U表示迎面而来的溪流速度。 在上述符号中，“最佳拟合线（k）* x 10”= -2.417 + 0.086（x，-x）与这个理论预测一致，表现出一些散点。 15.9 结论适用于预分离湍流的渐近方法是基于薄非黏性边界层方程组的领先顺序。 足够的

\section{15.10 参考书目}
0. S。 日州337（XS-X），”图15.1：基于表1的依赖性，由自由流线势理论确定的强奇异压力梯度创造了驱动机制。 与无黏性/无黏性相互作用引起的压力梯度的类似方法导致了一个包含soliton的模型（Ryzhov \& Bogdanova-Ryzhova [14]）。 因此，非线性过渡阶段和完全湍流流动之间存在深刻的平行关系。 但是，分离上游的扰动模式基本上也取决于迎面而来的溪流的历史。 与实验数据的初步比较证实了理论预测，无论湍流模型如何。 15.10 参考书目[l] Achenbach，E. 1968 在横流中，局部压力和皮肤摩擦在圆形圆柱体周围的分布，最高可达R = 5 x lo6。 J。 流体机械。 34，625-639。 [2]巴伦布拉特，G.I. 2003年缩放。 剑桥大学出版社。 [3]巴伦布拉特，G.I.，Chorin，A.J. \& Prostokishin，V.M. 2002年具有非零压力梯度的湍流边界层模型。 程序。 美国国家。 学术。 科学。 99，5772-5776。 [4] 卡斯蒂略，L. \& 沃克，D.J. 2002 上游条件对湍流边界层外流的影响。 AIAA J. 40，1292-1299。

338 湍流边界层分离[5] Cole，J.D. \& Aroesty，J. 1968年吹风问题-表面注入的不黏性流动。 英特。 J。 热质量传递11，1167-1183。 [6] 科尔斯，D. 1956年湍流边界层的守夜法。 J。 流体机械。 1，191-226。 [7] MaruSiC，I。 \& 佩里，A.E. 1995年边界层湍流结构的墙醒模型。 第2部分。 进一步的实验支持。 J。 流体机械。 298，389-407。 （Http://www.mame.mu.oz.au/ivan）。 [8]梅勒，G.L. 1972年大雷诺数，动荡边界层的渐近理论。 英特。 J。 工程师。 科学。 10，851-873。 [9] 梅尔尼克，R.E. 1989年湍流分离的渐近理论。 Competers \& Fluids，17，165-184。[lo] Messiter，A.F. 1970 平面后缘附近的边界层流动。 暹西姆J。 应用程序。 数学。 18，241-257。[ll] Neiland，V。 是的。 1969年对超音速流中层流边界层分离理论的贡献。 Izv。 阿卡德。 Nauk SSSR，Mekh。 Zhidk. i Gaza (4), 53-57（俄语；英文翻译：Fluid Dyn。 （4），33-35，1972）。 [12] Perry，A.E.，Bell，J.B. \& Joubert，P.N. 1966年逆压梯度湍流边界层的速度和温度曲线。 J。 流体机械。 25，299-320。 [13] Prandtl，L. 1905年Uber Flussigkeitsbewegung bei sehr kleinen Reibung。 在：Verhandl。 I11实习生。 数学。 孔。 海德堡，pp. 484-491，莱比锡，Teubner。 [14] Ryzhov，O.S. \& Bogdanova-Ryzhova，E.V. 1997年，与不稳定边界层相关的孤波的强制生成。 高级。 应用程序。 机械。 34，317-417。 1151 Samuel，A.E. \& Joubert，P.N. 1974年，边界层在日益不利的压力梯度下发展。 J。 流体机械。 66，481-505。 [16] Schubauer，G.B. \& Klebanoff，P.S. 1950年对湍流边界层分离的调查。 NACA TN编号 2133. [17] 史密斯，F.T. 1977 流过光滑表面的不可压缩流体的层层分离。 程序。 R。 SOC。 隆德。 A356，443-463。 [18] 索贝，我。 J。 2000 交互式边界层理论简介。 牛津大学出版社。

0. S。 Ryzhov 339 [19] Stewartson，K. 1969年在平板后缘附近的流动中11。 数学16，106-121。 [20] Sychev，V.V. 1972年关于层压分离。lzv。 阿卡德。 Nauk SSSR，Mekh。 Zhidk. i Gaza (3)，47-59（俄语；英文翻译：Fluid Dyn。 （3），407-417，1974）。 [21] Sychev，V.V. 和Sychev，Vik.V. 1980年在湍流分离。 Zh.vychisl。 马特。我数学。 菲兹。 20，1500-1512（俄语；英文翻译：USSR Comput。 数学。 数学。 物理。 20，133-143，1980）。 [22] Yajnik，K.S. 1970年湍流剪切流的渐近理论。 J。 流体机械。 42，411-427。

此页面故意留空

\clearpage
\chapter{第16章 高超音速磁铁o-流体-动态相互作用离子J。 S。 尚}
\section{16.1 摘要}
第16章 高超音速磁铁o-流体-动态相互作用离子J。 S。 Shang' 16.1 摘要 流动控制最有效的等离子体-流体-动态相互作用来自对剪切层生长速度的电磁扰动，并由随之而来的强黏性-无黏性相互作用放大。 采用漂移扩散和简单现象学等离子体模型的计算效果表明，使用电流体动态相互作用作为超人格流动控制机制的有效性。 实验观察充分证实了数值结果。 磁流体动力学相互作用在洛伦兹力中引入了一种额外的机制，作为流量控制机制。 然而，由于计算模拟的霍尔效应，这种方法也带来了额外的挑战。 分离流量抑制的电磁力已被证明能够在非常有限的范围内为黏性相互作用区域的延迟剪切层注入能量。 基于简单现象学等离子体模型的数值模拟表明了预期目的的可行性。 然而，这种磁流体-动态相互作用模式的有效性需要进一步验证。 俄亥俄州代顿市莱特州立大学机械与材料工程系研究教授

\section{16.2 命名法}
\section{16.3 导言}
16.2 B E J n S 2, Y, z U Pm U E ff 16.3 高超音速磁流体动力学命名法磁通密度 电通密度 电场强度 电流密度 带电粒子数密度 斯图尔特数，aB2L/pu 速度向量 笛卡尔坐标 保守变量 磁渗透性 电介电率 电导率 介绍 最近的磁流体动力学（MFD）研究在跨学科科学的基本理解方面取得了令人印象深刻的进展 [ll, 27, 26, 1, 22, 2, 2, 171. 这种增加的物理维度来自导电流动介质的行为。 流体介质的运输特性急剧改变源于通过电离激发空气分子的更高内部自由度。 Sears等人1 [25]在五十年代末和几年前，俄罗斯的阿贾克斯计划提倡使用电磁力来提高空气动力学效能[7，121。 提出了创新和有吸引力的想法，但受到我们分析这些极其复杂的物理学的能力的限制。 对这些革命概念的评估充其量是没有定论的。 然而，对MFD的恢复兴趣为使用电磁力控制流量提供了一些可实现的机会。 在这些努力中，流动介质必须是导电的，弱电离气体经常在高超声速流动中遇到。 然而，电离分数仍然相对较低，对于再入条件，电导率被限制在100 mho/m左右[25]。 因此，带电粒子的浓度必须富集。 产生具有足够带电粒子数密度以进行强MFD相互作用的体积等离子体所需的能量是巨大的[34,21,15,24,191。 即使使用更简单的电子碰撞机制来产生等离子体，电离电位总是低估能量需求，因为非平衡能量级联到振动兴奋、重组和附着过程[15，241。 过去，强大的MFD相互作用被外部施加的磁场进一步增强。 电磁力的相对大小和相互作用的流体惯性-

J。 S。 Shang 34 3 ing现象由一个相互作用参数，Stuart数，S = uB'L/\textasciitilde{}u [34，211]来描述。 尽管磁场的存在会通过霍尔效应和部分电离气体的离子滑移带来额外的并发症，但当磁场强度高，流体惯性低时，MFD相互作用占主导地位。 由于磁场强度在磁极处是最大的，而在剪切层的内侧区域的惯性是最小的，因此最强烈的MFD相互作用总是发生在固体表面附近。 弱电离气体的描述是实验和计算研究工作最具挑战性的问题[20，31，161。 在大多数实验室实验中，等离子体是由电子与控制表面上嵌入的电极碰撞产生的[26，1，221。 这种机制经常在特定位置采用，以丰富应用中的带电粒子数密度。 产生的等离子体由处于高度激发状态的电子组成，但重离子保留了周围环境的热力学条件[15，241。 最重要的是，电子碰撞产生的表面等离子体远非稳定的热力学平衡状态。 对于使用电磁力的流量控制，等离子体的电离程度、电导率、电流密度和电场强度至关重要。 不幸的是，为高超声速流动开发的高温热力学和化学动力学的所有众所周知的方法都不适用于描述电子碰撞过程产生的弱电离气体。 Surzhikov和Shang的漂移扩散等离子体模型成功地模拟了磁场中的发光放电[32，331。 数值结果证实了经典的观察，即带电粒子数密度与电极相邻，因此在等离子体护套区域上方的电导率要大得多[15，241。 在5马赫等离子体通道的直流放电测量中也观察到这种行为[20，31，161。 这种高电导率和强施加磁场相结合的条件大大确保了强大的电磁-流体-动态相互作用。 事实上，这种使用表面发光放电的相互作用证明了高超声速流动控制的适用性。 然而，这些模拟受到极高海拔环境的限制，压力只有几毫米汞柱。 为了将等离子体流量控制应用于广泛的应用，必须缓解等离子体生成对低压或低密度条件的限制。 有许多表面等离子体生成过程可以在高压环境中可靠地发挥作用，以增强弱电离气体的低电导率[3，61。 黏性-无黏性相互作用是超声速流动的独特特征之一，因为边界层在尖锐的前缘附近的存在不再像其他流动制度那样被忽略不见。 剪切层的位移厚度产生向外流动偏转，导致压缩波的形成，并最终凝聚成激波。 诱导的激波反过来修改了边界层结构，以关闭

\section{16.4 控制方程}
344高超音速磁流体动态相互作用回路。 经典的高超声速流动理论将Hayes和Probstein在尖锐前缘上的无黏性相互作用描述为压力相互作用[13]。 感应压力的大小由单个相互作用参数X = M3(C/Rey)lI2为特征。 压力相互作用参数对自由流马赫数立方体的依赖性大大放大了其在高超声速流动中的重要性。 一些实验和计算研究一直专注于通过引入流动控制的电磁效应来利用这种可能性[20，31，161。 这种控制机制可以在微秒内瞬间激活，以偏转固定控制表面上的流动，模拟偏转的襟翼。 实验和计算努力的初步结果表明，事件链，包括磁流体动力学和无黏性相互作用，使用简单的表面发光放电构成了非常有效的高超声速流动控制。 为了追求这一有前途的创新，目前的努力试图描述一个可能更新流体动力学研究前沿的基本概念。 在这个过程中，概述了为高超声速流动控制开发有效等离子体执行器的一些成就以及可能的研究机会。 16.4 控制方程 经典磁流体动力学（MHD）的控制方程由时间依赖性可压缩Navier-Stokes 方程和时域中的麦克斯韦方程组成。 等离子体是一种气体介质，其中远距离库仑力主导了带电粒子的全局行为。 等离子体的临界长度尺度是德拜屏蔽长度。 在这个距离内，周围的带电粒子有效地保护了任何两个粒子相互相互作用。 因此，等离子体是全球中性的，导电电流aD/\&在微波频率之前可以忽略不计[34]。 调用这些传统的MHD近似值，剩余的控制方程退化为法拉第归纳法、简化安培电路定律和两个高斯发散定律的电和磁通密度，以及可压缩Navier-Stokes 方程的完全补。 经过一些重新排列后，MHD 控制方程变成：aP - + v ' (pu) = 0 at dB - + V。 (uB - Bu) = -V x [(V x B/pe)/O] at (16.1) (16.2 T)] = (V x B/pe)'/a (16.4)

J。 S。 Shang 34 5 上述非线性偏微分方程由八个因变量组成：五个流体动力学和三个电磁变量（B、E和J）。 在非黏性极限中，方程组进一步简化为所谓的理想MHD方程，它构成了一个非严格双曲微分方程系统；此外，这些方程是非凸的[4]。 共有四个不同的波在MHD场中传播，Alfven、慢速和快等离子体波是横波，以及纵向声波，使波结构更加复杂[4，231。 波主导运动方程的公式也变得非常复杂，因为洛伦兹力垂直于磁通密度B和电流密度J。 因此，与正常磁通密度相关的特征值值为空。 MHD方程的其余七个特征值也可以局部退化，相互吻合，这取决于磁场的相对大小和极性[4，23，281。 对于大多数航空航天应用，磁雷诺（伯尔尼=p，ouL）远小于1，这意味着与外部施加的场相比，诱导磁通密度可以忽略不计[34，211。 在这种情况下，法拉第的麦克斯韦方程的归纳法可以从其余方程中解耦。 这种近似值现在将重点从研究电磁波运动转移到磁-空气动力学相互作用。 在这个配方中，洛伦兹力和焦耳加热在修改后的Navier-Stokes 方程中显示为源项。 由此产生的控制方程、初始值和边界条件都大大简化了。 使用这个治理方程组获得了大量数值模拟，并证明了它们能够准确预测广泛的磁-空气动力学相互作用[31、32、33、8、14、18、9、10、291. aP - + v。 （Pu）=0在* +V。 （Puu+？ If-?) - J x B = 0 at (16.5) (16.6) (16.7) -+V. ape [peu+Q+u.( \$-?)] -E.J=O at 由于基本偏微分系统主要由可压缩的Navier-Stokes 方程组成，其源术语中的电磁变量不会修改特征向量，因此Navier-Stokes 方程的所有传统通量分裂公式和求解程序都可以直接使用。 因此，方程组可以很容易地转换为通量向量形式，并通过CFD社区开发的数值程序求解[9,10,29,35]。 （16.8）其中U = U（p、pu、pv、pw、pe）

\section{16.5 等离子模型}
34 6 高超音速磁流体动力学空气动力学变量的边界条件是直截了当的；所有速度分量都施加了防滑条件，恒定壁温度或隔热壁条件描述了固体表面的条件，消失的压力梯度条件提供了密度的值。 在远场，流动需要回到冲击包络之外的无扰动状态。 在高超声速流动中，传统的无变化条件适用于下游远场。 对于电磁变量，一般来说，电场强度的切向分量和磁通密度的法向分量在介质界面上是连续的。 16.5 等离子体模型 远距离库仑力主导了带电粒子的集体行为；部分电离气体的一个基本属性是其电中性倾向。 这种内在特征，以及电子和离子之间的巨大质量差异，直接影响了等离子体的动力学。 带电粒子运动的两个基本机制是漂移速度和普通扩散[24]。 这种行为与气体排放的产生方式无关。 热碰撞产生的部分电离气体在高温环境中得到了广泛的研究，等离子体的建模建立在解离气体的Lighthill模型和电离的Saha方程上[34，211。 高温气体的解离或电离的平衡程度可以通过其各自的特征温度和相关的分配函数来计算。 电子碰撞电离被广泛用于流量控制，但物理模型并不为人所知。 部分原因是物理现象非常复杂，涉及气体和固体原子结构水平的相互作用。 然而，足够强度的电场通过中性气体的电子冲击电离产生电子离子对。 由于电导率，电流流经供应电极的外部电路[15，241。 在电极之间的放电中，电流由传导和位移电流分量组成。 在直流电场中，只有导电电流流动，由电子和离子成分组成。 电子分量是电极二次发射产生的电子数密度雪崩增长的结果。 在交流电场中，位移电流随着频率的增加而增加，而二次发射的重要性则减少。 最广泛采用的介电屏障放电（DBD）通常在大气压下工作，电极间隙中的放电表现像飘带，由于介电层上的局部电荷积累，随机过渡灯丝被电流限制淬灭，并由交流场恢复[3，61。 从现象学的角度来看，运动的驱动力

J。 S。 Shang 34 7 带电粒子是漂移速度和扩散，其中电子和离子之间的不同扩散速度抑制了电子扩散。 扩散的这一成分被称为双极扩散。 同样，在自维持等离子体中，控制体中带电数密度的变化率主要通过电离产生和通过重组消耗来平衡。 漂移扩散理论给出了双组分等离子体物种浓度的连续性方程[32，331。 （16.9）ane at \textasciitilde{} + V. re = a(E, P)lr,l - pn,ni (16.10) 低磁雷诺数极限的电场强度必须满足电荷守恒方程和全局中性条件，才能与漂移扩散公式或调用的广义欧姆定律兼容。 V. [(pi + pL\$)nE + (D,* - De)Vn] = 0 (16.11) 其中re = -DeVn, + n,peE和l? I = -DiVni + nipiE分别是电子和离子通量密度。 在这个公式中，a(E,p)和,B是第一个汤森电离系数和重组系数，p和pi是电子和离子流动性，D和Di轴分别是电子和离子扩散系数。 在大多数电动力学公式中，电场强度被电势函数取代：E = -V4（16.12）兼容性条件，Eqs。 （16.11）和（16.12），导致了电势4的众所周知的泊松方程[32，33，8，14，18，9，10，291。 根据定义，电流密度由以下方式给出：J = e（ri - re）（16.13）在存在磁场的情况下，可以通过明确修改离子和电子的流动性，以及漂移扩散理论的扩散和电离系数来包括霍尔效应[28，81. pe = pe/（l+H，2）（16.14）\textasciitilde{}2 = +pi）/ [（I + H，”）pi +pel（16.15）De = De/（l+Hz）（16.16.16）D，* =（0，peD+）/（Fe + pi）（16.17）a（E，P）= a（E，P）/（l+ H，”）（16.18）

34 8 高超音速磁流体动力学，其中He = wem，C和Hi = wimiC是电子和离子的霍尔参数，we和wi分别是电子和离子的Cyclotron或Larmor频率[21，15，241。 在这里，C是光速。 对于漂移扩散模型，阴极的适当边界条件需要强制执行零电位，电子数密度与次级电子发射系数成正比[32，331。 假定阳极反射所有离子；离子数密度消失，电势由电极上的差异规定。 在介电表面上，带电数密度可以忽略不计，电位的向外法向梯度为零。 对于所有电子碰撞电离，基本放电结构由电场维持。 为了数值模拟，等离子体可以通过电场强度和电导率来完全描述。 这些基本数据通常由大多数实验观察收集[20，161。 Gaitonde采用一个简单的表面放电现象学模型来研究熵层的稳定性、冲击边界层相互作用以及三维scramjet的性能[9、10、29、351。 为了获得最大的灵活性，采用了修改后的高斯分配来描述电极之间的电导率。 在简单的现象学等离子体模型中，霍尔效应和离子滑移可以包含在广义欧姆方程中[21]。 E=g。 [E = u x B - ahe(J x B) + ai,(J x B x B)] (16.19)，其中ffhe和ais分别是与霍尔效应和离子滑移相关的系数。 鉴于流体动态运动的声速和电动力学的光速之间的特征速度差异很大，治理方程是松散耦合的。 在这种方法中，用磁空气动力学方程迭代求解时间尺度要短得多的电动力学方程。 电动力学方程由带电粒子种类的连续性方程和电荷数密度方程的守恒组成，在大多数情况下，电势的泊松方程组成。 采用SOR（Successive Over Relaxation）方案来求解五对角矩阵系统中的控制方程[31，32，331。 磁空气动力学和电动力学方程的所有数值结果都是以名义上的二阶空间和时间分辨率获得的。 马赫数 5.15流中楔形模型中嵌入式电极上计算的带电粒子数密度如图所示。 16.1. 1.2千伏的电场在78.4 Pa（0.59 Torr）的环境压力下维持发光放电。 在研究的配置中，阴极被放置在阳极的上流中，距离前缘为2.22厘米，电极之间的距离为3.81厘米。 相同电极的整体尺寸为宽0.64厘米，长3.18厘米。 计算的离子

\section{16.6 电子流体动态相互作用}
J。 S。 Shang 34 9 图16.1：电极上的离子数密度分布，阳极上方的离子数密度分布与实验观察非常一致[20，161。 然而，阴极正上方的离子数密度曲线的计算结果和实验结果之间存在很大差异。 然而，这两个结果之间的差异仅限于一个数量级，类似于微波吸收技术和Lang-muir探测获得的数据之间的测量差异。 差异的原因仍在研究中。 16.6 Elect ro-Fluid-Dynamic相互作用 在图的Schlieren图片中清楚地显示了尖锐前缘楔形附近的黏性-黏性相互作用的行为。 16.2. 该图像是在5.15的自由流马赫数下录制的，密度为5 x 1O3kg/m3，静态温度为43 K。 该模型的总尺寸为（3.81 x 6.67）厘米。 在这些条件下，基于楔形模型长度的雷诺数为1.08 x lo5，楔形模型平坦上表面的边界层预计为层[20，161。 边界层的位移厚度的增长显然使流向外偏转，并在模型的上表面和下表面都产生压力。 无黏性相互作用甚至在楔形模型的平板表面产生斜冲击。 在模型的后缘，Chapman-Rubesin常数为0.65的压力相互作用参数为0.65。 在这张照片中，边界层离开楔形模型表面，作为自由剪切层继续下游。

350 高超音速磁流-D动力学图16.2：超音速流中的楔形楔形Schlieren 细长体的尖锐前缘附近的黏性-不黏性相互作用表明，剪切层位移厚度的增长率的扰动可以大大放大。 随后的压力相互作用进一步放大了扰动，为流量控制产生显著的压力高原。 这种非侵入性电磁扰动似乎非常有吸引力，因为带有施加电场的嵌入式电极可以点燃发光放电，施加的横向磁场也会将洛伦兹力添加到MFD相互作用中。 事实上，发光放电对剪切层的结构产生了两个不同的扰动。 首先，电极加热提高了表面温度，其次，焦耳加热提高了发光放电域中表面上方的气体温度。 这两种机制都对剪切层结构产生热扰动。 图中显示了马赫5.15高超音速流中尖锐前缘楔形的直流放电。 16.3。 等离子体是在嵌入平板表面的两个电极之间产生的。 外部电路中施加的1.2千伏电场维持了50毫安的总电流。 等离子体的最大电子数密度为3 x lOI2/cm3，电极温度估计为600 K [24, 201。 在0.59 Torr的环境压力和43 K的静态温度下，空气密度为1.33 x lOI7/cm3；电离程度为2.25 x，电导率为1 mho/m。 从这个意义上说，发光放电在电极上提供了真正的弱电离气体。 电磁场从两个方面改变位移厚度的增长率；改变运动场结构和引发热量

J。 S。 Shang 351 图16.3：壁区域高超声速流动交换中楔形的直流放电。 可以操纵洛伦兹力和焦耳加热来改变边界层的轮廓。 当等离子体生成的电极嵌入模型表面时，大量的局部等离子体加热来自两个不同的来源。 其中之一是体积焦耳加热，另一个释放热源是加热电极的传导。 对于发光放电，电极通常达到接近600 K的表面温度[24，201。 由于加热机制大不相同，传导加热的特征时间尺度以毫秒为单位，焦耳加热的时间尺度要短得多。 等离子体加热会影响边界层的热和速度曲线。 电空气动力学相互作用的计算温度轮廓在复合演示中进行了描述。 图16.4包括简单的高超音速边界层、具有电极加热的相同边界层的计算结果，最后是楔形模型上的发光放电。 清楚地表明，传导和焦耳加热对高超音速剪切层的结构产生了显著的扰动。 焦耳加热的效果比电极加热要明显得多。 使用等离子体模型的两种计算，无论是通过漂移扩散论还是更简单的现象学方法，都与屏蔽停滞温度探针获得的测量结果合理一致[20，161。 图16.5比较了电极上的无量值温度（由43 K的自由流温度归一化）。 计算结果最突出的特点是观察到焦耳加热占主导地位，导致墙壁附近的气体温度比表面更热

高超音速磁流体动力学平板I GlowDischarge 1 图16.4：电极的温度轮廓和阴极和阳极的焦耳加热温度。 在图中。 16.6，焦耳加热的效果将边界层结构的扰动转化为增强的黏性-非黏性相互作用，由楔形表面上方的压力曲线所描述。 通过与经典超人压力相互作用的计算结果的直接对比，观察变得更加容易，其中诱导压力主要集中在楔形的前缘附近。 压缩波迅速凝聚，在前缘附近形成斜冲击。 另一方面，焦耳加热会触发阴极和阳极上方的额外压缩波。 这些波浪融合在一起，最终与来自尖锐前缘的斜激波合并。 由此产生的激波在楔形表面产生更高的压力上升。 这种诱导压力高原聚集在楔形的前缘附近，成为产生高超音速车辆控制俯仰矩的有效手段。 由于压力相互作用现象，板块前缘区域的高超声速流动场以有利的压力条件为主[13]。 电空气动力学相互作用改变了与电极相邻的多个压缩和膨胀域的膨胀流。 边界层结构必须响应流向压力梯度。 表面剪切力使用远下游值作为参考计算，并且随着与前缘距离的增加，熟悉的逆平方根衰变清晰地表现出来。 皮肤摩擦系数Cf的表面剪切应力如图所示。 16.7揭示边界层对等离子流控制机制的结构响应。 表面剪切向上游减少，并且

J。 S。 Shang 353 aza 0.18 a.16 0.14 0.12 0.10 am an6 an4 an2 ana m - Comp [Wthndo) I\_\_\_\_e\_\_\_\_ Comp (Anode) olts(hthodo) D WtalAnoda) 1\% me SJJ ua 711 ED sa Ian iia i2.a 136 i4n 160\%.a 178 i8.a TrdTlnf 图16.5：停滞温度曲线比较图16.6：EFD中的压力曲线

\section{16.7 磁铁o-流体-动态相互作用}
354 高超音速磁流体动力学 1i 5 10.0 B €4 r u 611 0.0 00 0.1 02 oa 0.4 06 06 0.7 O.B 09 in 1.i iz UL 图16.7：根据局部表面压力梯度，EFD相互作用中的表面剪切应力在电极的下游增加。 使用等离子体模型的实验和计算结果证实了表面传导和体积焦耳加热的相对重要性。 在图中。 16.8，使用漂移扩散模型的磁空气动力学方程的解与测量结果进行比较。 类似比较的对应物也如图所示。 16.9，基于现象学等离子体模型。 这两个数值结果都与数据合理一致。 更重要的是，发光放电诱导了真正的电气动力学相互作用，这仅靠电极加热是不可能的。 在60瓦的非常低电平电源下，电空气动力学相互作用会产生1度的等效表面偏转。 根据这一结果，每个电极长度的等离子体执行器所需的功率缩放为每厘米每度18.90瓦。 从本质上讲，结合的计算和实验研究表明，高超音速流中楔形表面的发光放电引起的电空气动力学相互作用可以作为一种有效的流量控制机制。 然而，等离子体生成过程将有效应用范围限制在低环境压力环境中。 16.7 磁铁o-流体-动力学相互作用 磁-空气动力学相互作用的全部影响需要在外部施加的磁场的存在下使用等离子体[25，71。 横向磁场

J。 S。 Shang 355 4.6 - 411 - 3s - E f 25- 211 - 15 - 图16.8：EFD对表面压力的影响（漂移扩散模型）----3-- ElectrDdQHearlng Jou! eHsaling \{ - Dam E " B 2.6 1.6 OR 0.1 02 03 0.4 05 06 0.7 OS 09 I11 1.1 I2 图16.9：EFD对表面压力的影响（现象学模型）

356 高超音速磁流体动力学对等离子体产生深远的影响，特别是改变了等离子体结构的特征，包括电极护套的特性[4，231。 基本现象受磁场中电子流动性降低的支配。 在平行板发光放电中，放电柱根据横向磁场的极性向方向连续漂移[32，331。 对于高超声速流动控制中的发光放电，对电动力学结构知之甚少，但只能从基本碰撞过程进行分析，并尝试从实验观察中获得更好的理解。 目前的数值模拟复制了实验安排，电极嵌入模型楔形表面，与其前缘平行，以产生与气流对齐的导电流向量[20,16]。 然后，横向磁场施加在垂直于电流的等离子体通道上。 通过这种排列，横向磁场产生洛伦兹力，J x B，该力要么将带电粒子排出（J x B > 0），要么约束（J x B < 0）到电极上。 预计，在带电粒子被从板表面排出的情况下，离子和中性粒子之间的动量交换通过非弹性碰撞增强了其次的黏性-非黏性相互作用。 相反，受约束的放电粒子运动抑制了压力相互作用的强度。 初步实验观察证实了这一发现，但在广泛的磁场强度上存在明显的不确定性[20，161。 从本质上讲，磁-空气动力学相互作用的这种净结果揭示了洛伦兹力的收缩表面放电和加速之间的补偿效应。 实验中的表面等离子体已经在异常放电状态下运行；外部施加的磁场会产生额外的放电不稳定性，阻碍了实验观察[15，241。 在测试条件下，B = 1.0 Tesla的电子Larmor频率为1.76 x 10l1弧度/s，电子重粒子碰撞频率估计在3.9\textasciitilde{}10g/s到2.3 x 1O1'/s之间。 在这些条件下，等离子体在高超音速MHD通道中的行为不是碰撞主导的，最大和有效霍尔参数大约在7.6到45.1之间。 为了缓解不确定性，模拟限制在较低的磁场强度，0.2 T < B < 0.2 T，然后霍尔参数减少到单位级。 从实验观察来看，外部施加的磁场对发光放电的主要影响是抑制阴极层上的可见放电[20,16]。 在没有外部施加的磁场的情况下，发光放电在阴极层上最明显，并集中在电极之间的介电表面上方。 等离子体也被楔形模型以外的高超音速流向下游对流。 通过驱逐洛伦兹力，放电在电极上变得更加均匀。 阴极层上方的可见等离子体域通过降低的电离程度被施加的磁场显著抑制。 当

J。 S。 Shang 4s - 4a - 3s -.......... FIR1 Plat0............，... JXBcO。bn2T JXBSD。 BtP.2T。“-mta，B4。 Pl mg842T:: I oa 1.0 211 311 40 \%a 6.0 711 81 9.0 10.0 X（om）图16.10：MFD对表面压力磁极性的影响被逆转，发光放电被限制在模型表面的窄层中。 在数值模拟中，洛伦兹力的这种效应很容易从电极上计算的带电粒子密度轮廓中检测到，如图所示。 16.10. 在没有外部磁场的情况下，带电粒子密度主要集中在阴极和阳极层上。 等离子体模型的最高电离度在电极附近达到8.11 x 1O1'/cc的值，该值与测量值一致。 然而，计算结果预测了阴极层上方的数据。 在外部施加的磁场的存在下，放电模式会大大改变。 喷射洛伦兹力（J x B > 0）确实将带电粒子从电极上推开，并扩大了放电域。 与没有外部磁场的发光放电相比，电极的电离程度也降低了24.1\%。 放电域增加的趋势与磁场的相反极性相反。 带电粒子被强烈收缩到电极上，阴极和阳极层的最大值为1.06 x 101'/cc，放电域明显减少。 在发光放电场中，电场力总是在J x B加速度上占主导地位，等离子体几乎无碰撞，但洛伦兹力的影响被清楚地证明了。 然而，从图中电极上的计算表面压力中很容易检测到洛伦兹力的影响。 16.10. 计算结果通常会高度预测压力测量值，数据也显示出异常大的数据散射带

358 高超音速磁流体动力学由于在应用和均匀磁场的存在下不稳定的发光放电模式。 喷射洛伦兹力（J x B > 0）确实将带电粒子从电极上推开，并扩大了放电体。 不出所料，磁场会产生更强的电磁扰动，这种扰动可以通过黏性-无黏性相互作用来放大。 因此，在相同的电场强度下，感应的表面压力大于电空气动力学相互作用。 放电域增加的趋势与磁场的相反极性（J x B < 0）相反。 通过抑制流线的向外偏转，诱导的表面压力也会相应地降低。 然而，维持计算稳定性也需要增加努力。 使用漂移扩散模型的数值结果捕捉到了相反极性施加的磁场之间的差异，而现象学模型在预测感应压力方面不太准确[as]。 MFD相互作用也应用于分离流量控制[2，17，351。 在这些应用中，电磁力被引入分叉流场，而不是以小扰动的形式引入，而是使延迟剪切流通，以克服不利的压力梯度。 MFD分离的流量控制来自有源洛伦兹力，因此控制机制需要存在强大的外部施加的横向磁场。 在压缩坡道上的流动中遇到了最图解和最经典的黏性-非黏性相互作用[30，51。 在这个流动场中，迎面而来的边界层在压缩坡道的上游分离。 增厚的剪切层诱导了一系列压缩波，最终在斜坡上凝聚成斜激波。 分离的流动区域可能很广泛，导致高压高原和重新连接的热点。 在这方面，MFD分离流量控制抑制了黏性-非黏性相互作用，并提高了退化的空气动力学性能。 然而，黏性-非黏性相互作用不会获得任何杠杆来放大磁-空气动力学扰动[24，19，201。 使用等离子体进行分离流量控制是基于与等离子体执行器不同的机制，依靠电磁扰动，并通过黏性-无黏性相互作用放大。 流动分离是最剧烈的流体动态分叉之一；动态状态的过渡是沿流方向的不利压力梯度的结果。 通过给剪切层通电以克服不利的法向应力或压力梯度来实现分离流动的抑制。 因此，流量控制的能量需求是明确的，通过黏性-非黏性相互作用的优势不再存在。 所有已知的分离流量控制都与冲击边界层相互作用有关[17，351，最近的数值模拟利用一个简单的现象学模型，在马赫数为14.1和雷诺数为103,680的24度压力坡度上进行分离流量抑制

\section{16.8 结束性讲话}
J。 S。 Shang 359基于从前缘到角落的距离[35]。 在数值模拟中，将表面等离子体结构域强加在整个角落区域（0.8 < x/L < 1.2）。 均匀的横向磁场正交施加在二维流动上，斯图尔特数被分配到一个恒定的单位值。 等离子体域的表面温度从297.2的恒定值提高到555.6 K，以模拟等离子体生成过程中的电极加热。 对于现象学等离子体模型，等离子体域外缘的电导率降低到0.001 mho/m。 控制电场强度的参数描述为k = -E,/u,B = 1.5。 由于该模型不一定代表详细的物理，而是代表可行性研究，因此模型和数值模拟都没有优化。 Navier-Stokes（基线）的速度曲线和压缩角（x/L=l.O）的磁空气动力学方程并排呈现在图中。 16.11. 计算出的黏性非黏性相互作用导致上游到压缩角的分离流动区域，事实上，分离流动区域的物理维度从x/L = 0.36到x/L = 1.42，与实验观察和以前的计算一致[35]。 对于磁空气动力学计算，这个等离子体域由0.8 < x/L < 1.2的条带描述，并垂直延伸至y/L = 0.1。 电磁场，主要是洛伦兹力，加速了流向速度分量，完全消除了反向流动区域。 与此同时，表面附近焦耳加热的强加速度和高值也导致压缩角的皮肤摩擦和斯坦顿数分布的激增。 无论如何，分离的流动区域在数值模拟中被完全抑制[35]。 在图中。 16.12，给出了四个压力分布。 它们描绘了非黏性、冲击边界层相互作用和磁空气动力学方程的解。 磁空气动力学方程的两个解是在不同的垂直等离子体域获得的。 前面提到的垂直放大的等离子体域能够消除冲击-边界-层相互作用的分离流动。 数值模拟表明，洛伦兹力加速了剪切层的内部区域，以覆盖下游的高压高原。 然而，计算结果尚未从实验观察中得到任何证实。 16.8 结论性评论 等离子体执行器使用电磁扰动对剪切层结构，并由高音速流中的黏性-无黏性相互作用放大，通过实验观察和两个等离子体模型的数值模拟得到了充分证实。 基本机制是导电电极和体积焦耳加热的综合效应。 的时间尺度

360 高超音速磁流体动力学 0.15 0.14 0.12 0.1 1 WHD 无 Sep。 \textbackslash\{\} 0.05 0.04 0.03,, - \_\_--- 0.01 ' i - \_,\_ - 7 - 1 u 0.5 图16.11：没有MFD控制的速度曲线图16.12：没有MFD控制的表面压力分布

\section{16.9 致知}
\section{16.10 参考书目}
J。 S。 Shung 361焦耳加热在微秒范围内估计，比导电电极加热和随后的对流短数级。 嵌入固定和不可移动板表面的等离子体执行器的电流体-动态相互作用就像板正在执行俯仰运动一样。 实验和计算结果都表明，压力高原诱导的流量控制有效，压力高原可以缩放到每厘米每度18瓦的电极尺寸。 与电流体-动态相互作用相比，磁流体动力学相互作用在施加的外部磁场下表现出更大的放大。 漂移扩散等离子体模型已成功应用于数字模拟，霍尔效应被明确纳入公式中。 从实验和计算来看，霍尔效应显著限制了带电粒子的流动性，并改变了表面放电模式。 除此之外，由此产生的放电不稳定性掩盖了测量和计算精度。 使用简单的现象学等离子体模型进行磁空气动力学相互作用的二维计算表明，在14.1的马赫数下，完全抑制压缩坡道上的分离流动的可能性。 对这一现象的实验验证尚未完成。 16.9 致知 博士的赞助 J。 Schmisseur和博士。 F。 非常感谢空军科学研究办公室的Fahroo。 教授的宝贵贡献。 S.T. 俄罗斯科学院的Surzhikov，博士。 D.V. Gaitonde和R。 空军研究实验室的Kimmel，以及教授。 J。 莱特州立大学的Menart受到诚挚的表彰。 16.10 参考书目[l] Artana, G., D'Adamo, J., Lger, L., Moreau, E., \& Touchard, G., “使用电流体力学执行器的流量控制”，AIAA Journal, Vol. 40，2002年，pp. 1773-1779。 [2] Bityurin，V.，Klimov，A.和Leonov，S.，大气中超音速飞行的先进流量/飞行控制概念的评估，AIAA 99-4820，弗吉尼亚州诺福克，11月。 1999年。 [3] Boeuf，J.P.，等离子体显示面板：物理学，最近的发展和关键问题，J. 物理学D；应用物理学，卷。 36，3006，pp. R53-R79。

362 高超音速磁流体动力学[4] Brio，M. \& Wu，C.C.，理想方程的迎风差分方案磁流体动力学，J. 补偿。 物理学，卷。 75，1988年，第400-422页。 [5] Dolling，D.S.，五十年的激波/边界层相互作用重新搜索：下一步是什么？ AIAA J. 卷。 39，8月 2001年，第1517-1531页。 [6]埃利森，B。 \& Kogelschatz，U.，非平衡体积等离子体化学处理，IEEE Trans。 血浆科学，卷。 19，1991年，第1063-1077页。 [7] F'raishtadt, V.L., Kuranov, A.L., \& Sheikin, E.G., 使用MHD系统 [8] Gaitonde, D.V., 开发高超音速飞机中的3D非理想磁共体求解器，技术物理学，卷。 1998年11月，第 1309.动力学，AIAA 99-3610，1999年6月。 [9] Gaitonde，D.，三维非理想磁气动力学的高阶解程序，AIAA J.，卷。 39，不。 1，2001年，pp. 2111- 2120. [lo] Gaitonde，D.V.，使用MGD能量旁路的三维流通Scramjet模拟，AIAA 2003-0172，2003年1月。 [11] Ganiev, Y., Gordeev, V., Krasilnikov, A., Lagutin, V., Otmennikov, V., \& Panasenko, 通过等离子体和热气注射减少空气动力学阻力，J. 热物理学和热传递，卷。 14，不。 1，2000年，第10-17页。 [12]古里雅诺夫，E.P. \& Harsha，P.T.，阿贾克斯：高超音速技术的新方向，AIAA预印本96-4609，11月。 1996年。 [13] 海斯，W.D. \& Probstein，R.F.，高超声速流动理论，学术出版社，1959年。 [14] Hoffmann，K.A.，H-M Damevin和J-F Dietiker，超音速磁流体动力流的数值模拟，AIAA 2000-259，2000年6月。 [15] Howatson，A.M.，《气体排放导论》，第2版，Perga-mon出版社，牛津，1975年。 [16] Kimmel, R, Hayes, J., Menart, J., \& Shang, J.S., 地表等离子体排放对5马赫边界层的影响，AIAA 2004-0509，内华达州里诺，2004年1月5日至8日。 [17] Leonov, S., Bityurin, V., Savelkin, K., \& Yarantsev, D., 放电对分离过程和冲击的影响 超音速气流位置，AIAA 2002-0355。 内华达州里诺，2002年1月。 [18] MacCormack，R.W.，磁流体动态的保守形式方法，AIAA 2001-0195，内华达州里诺，2001年1月。

J。 S。 Shang 363 [19] Macheret, S.O., Shneider, M.N., \& Miles, R.B., 磁流体动力学 and Electrohydrodynamic Control of Weakionized Plasmas of Hypersonic Flows, AIAA J., Vol. 42，2004年7月，pp. 1378-1387。 [20] Menart, J., Shang, J.S., Kimmel, R., \& Hayes, J., 磁场对5马赫风洞中产生的等离子体的影响，AIAA 2003-4165，佛罗里达州奥兰多，2003年6月。 [21] 米奇纳，M. \& Kruger，C.H.，部分离子气体，John Wiley \& Sons，纽约，1973年。 I221邮报，M.L. \& Corke，T.C.，使用等离子体执行器的分离控制：固定和振荡机翼，AIAA 2004-0841，2004年1月。 [23] Powell, K.G., Roe, P.L., Myong, R.S., Gombosi, T., \& Zeeuw, D.D., AIAA 95-1704-CP, 1995, pp. 661-671。 [24] Raizer，Yu。 P.，气体排放物理学，Springer-Verlag，柏林，1991年。 [25] Resler，E.L.，Sears，W.R.，磁空气动力学的前景，J. 航空。 科学1958，卷。 25，1958年，pp. 235-245和258。 [26] Roth J. R.，Sherman，D.M.和Wilkinson，S.P.，发光放电表面等离子体的电流体动力流控制，AIAA J. 卷。 37，2000年，pp. 1166-1172。 [27] Shang，J. S.，用于高超音速钝体减阻的等离子体注射，AIAA J. 卷。 40，不。 2002年6月，pp. 1178-1186。 [28] Shang，J.S.，计算流体动力学中的共享知识，电磁学和磁空气动力学，航空航天科学进步，卷。 38，2002年，pp. 449-467。 [29] Shang, J.S., Gaitonde, D.V., \& Updike, G.A., 模拟高超声速流动控制的磁空气动力学执行器，AIAA 2004-2657，俄勒冈州波特兰，2004年6月。 [30] Shang，J.S. \& Hankey，W.L.，压缩坡道上的超音速湍流的数值解，AIAA J. 卷。 10月13日 1975年，pp. 1368-1374。 [31] Shang, J.S \& Surzhikov, S.T., 用于超人体流动控制的磁-空气动力学相互作用，AIAA 2004-0508，内华达州里诺，2004年1月5日至8日。 [32] Surzhikov，S.T. \& Shang，J.S.，磁场中的发光放电，AIAA 2003-1054，第41届航空航天科学会议，内华达州里诺，2003年1月6日至9日。

364 高超音速磁流体动力学 [33] Surzhikov，S.T. \& Shang，J.S.，用于磁场二维发光放电的双组分等离子体模型，J. 补偿。 物理学，卷。 199，2004年，pp. 437-464。 [34]萨顿，G.W. \& Sherman，A.，工程磁流体动力学，McGraw-Hill，纽约，1965年。 [35] Updike，G.A.，Shang，J.S.和D. V。 Gaitonde，使用磁空气动力学相互作用进行高超音速分离流动控制，AIAA 2005-0164，内华达州里诺，2005年1月。

\clearpage
\part{四。 多相和反应流}
第四部分多相和反应流

此页面故意留空

\clearpage
\chapter{第17章 使用AUSM+-up方案计算多相流量 Meng-Sing Liou和Chih-Hao Chang}
\section{17.1 摘要}
第17章 使用AUSM+-up方案计算多相流 Meng-Sing Liou'和Chih-Hao Chang' 17.1 摘要 本章介绍了AUSM+-up方案的扩展到所有速度和所有可压缩性状态下的多相流计算。 将详细描述描述多相流动的两种说明，即混合和双流体模型。 第一种方法将多相流体视为由实流体状态方程描述的混合物，并应用热力学平衡来处理液-蒸汽相变。 第二种方法是非平衡模型，它通过运输方程分别解决每个相，相之间存在压力平衡。 提出了分层流动模型，以包括细胞界面通量中的相间效应。 新的AUSM+-up方案在两种方法中都采用。 数值结果表明，新方案在模拟一些涉及液体和蒸汽的相当具有挑战性的数值问题方面是有效的。 计算是稳健和稳定的，与分析、实验和其他计算结果相比，结果是准确的，其中相位界面及其演变被敏锐地捕捉到，而无需对界面进行特殊处理。 美国俄亥俄州克利夫兰市刘易斯菲尔德的美国宇航局格伦研究中心推进系统部门，44135

\section{17.2 导言}
368 多相流的USM+-up方案 17.2 引言 在自然界、生物系统和工业设备的广泛情况下可以找到两相流，它们通常涉及多样化和复杂的机制。 虽然物理模型可能特定于某些情况，但求解控制方程的数学公式和数值处理可以是通用的。 基于连续力学，我们将流体视为由两个相互作用相（或材料）组成的混合物，在任何给定时刻占据空间的同一区域。 因此，我们将根据需要在单个阶段中提供有关每个阶段的信息，以及它们之间的相互作用。 然而，这些相互作用术语带来了额外的数字挑战。 例如，不能保证数学方程在所有条件下都是双曲的。 此外，由于时间和速度尺度的不同差异，流体可压缩性和非线性变得尖锐，使数值程序更加复杂。 在过去的十年里，已经开发了一系列AUSM计划[16、11、25、41。 事实证明，这些方案是准确、简单、稳健的，并且容易扩展到其他类型的守恒定律，因此提供了其他现有方案的有吸引力的替代方案；最近在[la]中给出了摘要。 本文旨在提供适合解决任意流速和任意可压缩性水平的一般多相流问题的扩展。 先前的努力集中在求解「混合」模型方程上[14，151。 本文将专门讨论所谓的双流体模型方程的解。 将推导出离散方程，揭示以前未包含的数值通量。 该方案将逐个术语进行分析，以说明细胞内（控制体）和细胞之间的各种界面相互作用。 最近，AUSM系列已扩展到多相流计算，例如，在Refs中。 [14，15，5，171。 Paillkre等人1.[17]解决了一个包含界面源项的双流体模型系统。 与空气动力学流动的通常方程不同的几个特征大大增加了复杂性。 与单相流相反，处理多相流的一个主要困难是缺乏一组共同的控制方程，其中提出了不同级别的近似和相交换项建模。 拟议的模型基于多年来积累的实验数据，并针对特定的配置和流动制度进行。 因此，从物理学的角度来看，它们是高度现象学的。 此外，由此产生的系统具有严重的数学后果，例如：非保守形式、长度和时间尺度差异导致的刚度、热力学状态变量的复杂表示、非混合性等。 因此，这些直接影响了数值解的各个方面，包括准确性、效率和稳定性。 最重要的是，它施加了严重的限制

\section{17.3 多相流动的控制方程（模型）}
\subsection{17.3.1 热力学平衡模型 [14]}
M.-S. Liou和C.-H. Chang 369 多相流代码的稳健性。 单相流程序的扩展非常直接。 因此，在过去十年中，我们看到积极的研究试图设计一种用于求解多相流方程的稳健、准确和通用的数值方法。 为了克服方程的非混合性质的可能困难，提出了一些额外的模型项。 它们本质上通常是临时的（尽管有一些物理考虑），主要用于数字目的。 多相流动的特点是存在相分离界面，无论是微观还是宏观尺度，都存在不同的（即不连续的）流体和动态特性。 例如，在气液流体系统中，根据界面的拓扑结构分为三种不同的流动制度[9，261：（1）分离流动，如分层流动和环状流动，（2）混合流动，如蛞蝓和气泡环状流动，以及（3）分散流动，如气泡流动和液滴流动。 这些不同的流动制度通常用不同的构成规律来描述。 由于两相流可以从一个制度到另一个制度不断演变，因此在一般情况下，如果不是不可能，在一般情况下对两相流的数学描述是困难的。 一般做法是利用某些平均过程，无论是时间、空间还是统计过程，通过这些过程推断出宏观描述，导致描述精细结构的丧失。 17.3 多相流动的控制方程（模型）双相流动通常采用两种根本不同的公式；它们是混合模型和双流体模型。 前者，包括热力学平衡模型和漂移通量模型，将流动系统视为混合物，控制方程与单相方程相似，具有描述相位分布的附加信息。 后者通过为每个阶段使用控制方程来单独考虑每个阶段，并使用表示相间质量、动量和能量传递的耦合（相互作用）项。 在下面，我们将简要介绍这些方法。 对于模型和数学属性的更深入的讨论，感兴趣的读者可以在一些优秀的参考资料中找到它，例如[23，241。 17.3.1 热力学平衡模型 [14] 该模型属于混合模型，特别假设所有相位都处于平衡状态，因此相间具有相同的速度分量、压力和温度。 这种模式在文献中经常被称为，包括-

370 在我们之前的论文中，多相流 ing的USMf-up方案[14]，均质均衡模型。 认为称其为热力学平衡模型更合适，因为描述两相共存的条件2完全基于热力学平衡原理。 它允许通过热力学平衡从一个阶段过渡到另一个阶段。 因此，相的共存只能从热力学原理来处理，如麦克斯韦原理或饱和假设。 热力学平衡模型的方程组看起来与单相流体流动的方程组完全相同。 为了清楚起见，我们将只呈现一个空间维度的方程，而忽略黏性效应。 在其他维度中包含相应的术语是相似且直接的。 然而，独立于空间维度的多相流特有的项将在需要时保持，特别是在双流体模型的后期。aQ dF -+-=o. ax（17.1）保守变量以Q = [p，pu，pEIT给出。 黏性通量是F = [pu，pu2 + p，puHIT。 变量的符号和定义相当标准，因此在此省略。 然而，必须强调的是，这些变量是混合物的变量。 当流体位于两相区域时，会施加额外的条件，否则由任何相的状态方程来描述，例如液体或蒸汽。 该模型的描述可以在[14]中找到。 描述多相的热力学状态方程可以包含在单相实流体模型中，其中可压缩性系数Z（P，2'）是一个来自经验数据或范德华类型公式的复杂代数函数。 Peng-Robinson方程[18]的典型等温线如图所示。 压力密度图上的17.1。 清楚地表明了蒸汽状态，其中压力几乎随密度线性变化，以及液体状态，需要很大的压力变化才能诱导密度变化。 对于给定的压力和温度，Eq的溶液。 （17.2）返回一个或三个可压缩系数Z或密度值，前者对应于单相区域（液体或蒸汽），后者对应于蒸汽和液体共存的两相区域。 两相区域内压力的相应密度显示为点A-C。 A和C代表饱和的蒸汽和液体状态，而B在物理上毫无意义。 对于特定的温度，“允许的”两相区域受D和E的压力值的约束，D和E是局部极值。 这些21n这个连接的基因座，一个由两个相位组成的系统，据说是异质的。

M.-S. Liou和C.-H. Chang 371 2. OE+07 t.BE+07 1.5E+07 T c T，单相液体1.2E+07 - 上（蒸汽）旋轴p（Nlm"1 1 OE+07 - 单相蒸汽7.5E+06 - 亚稳定蒸汽5。 OE+06 - n蒸汽压力/4(/ - - 亚稳定液体2.5E+06较低（液体）旋节-双pKe区域1i1,,III IILI\textasciitilde{}IIIII\textasciitilde{}III IZJ 5 10 20 25 30 p (Cmo\textasciitilde{}elm\textasciitilde{}) 图17.1：压力与摩尔密度（临界温度以下的等温线）。三重和临界点之间的温度压力值定义了液体和蒸汽旋节曲线，将两相区域分为亚稳定蒸汽、不稳定和亚稳定液体区域。 两相区域不稳定部分的信息没有物理意义，也没有用处。 对于自旋数值（D-B-E）之间的密度，可以证明声学特征值是复杂的，这意味着欧拉系统在时间上不是双曲的，并且积分方程的传统时间行进程序是不正常的。 同样值得注意的是，在低温下，液体自旋压可能是负的，这意味着液体的模拟膨胀可能会在亚稳定区域产生合理的密度，但非物理压力。 在亚稳定分支A-D和C-E中，它们在物理上是可行的，并且具有独特和可能有吸引力的特性。 亚稳定蒸汽分支上的状态也被称为过低冷却，因为蒸汽存在于低于饱和蒸汽温度的温度，对应于A-D的压力。 同样，C-E上的状态被称为过饱和，它们可以在不蒸发的情况下达到低于C-E温度的饱和液体压力的压力。 此外，在液体和蒸汽旋点之间的特定压力下，系统处于平衡状态。 这种压力被称为蒸汽压力pvap(T)，通过克劳修斯-克莱佩龙方程与温度直接相关。 或者这等同于设置蒸汽和液体

a2（m/s）多相流动的USM+-向上方案--\textasciitilde{}彭·罗宾逊----声学特征值----理论结果i I I I I I I I液体I两相结蒸汽两相-50000图17.2：声速a2与摩尔密度（350 K的辛烷值）的平方，取自[15]。fugacities达到相等值。f（G，T，pvap）= f（zi，T，Pvap）。 （17.3）通过迭代，蒸汽压力图作为温度的函数计算。 因此，两相混合物的自变量数量从两个减少到一个，这与吉布斯相位规则一致。[lO]这设置了图中的A-B-C线。 17.1。 给定由时间积分方法确定的网格点密度和温度的更新值，通过引入缩放参数，例如（蒸汽）空隙分数（17.4），可以沿着这条线定位一个特定点，其中上标“I”和“II”指的是饱和液体和蒸汽状态，即分别是C和A点。 然后，可以相应地确定混合物的相关热力学量，例如内能或焓。 [14] 多相流的独特特征之一是，混合物的音速非常小，以每秒米为单位，如图所示。 17.2.

\subsection{17.3.2 双流体模型}
M.-S. Liou和C.-H. Chang 3r3 因此，在相变期间向局部“超音速”流动条件转变是一种明显的可能性。 这种方法的一个例子表明了其描述多相流动的能力，涉及半球圆柱体上的水流。 如果压力差（气穴数，K = 2（pm - pvap）/pmU\&）足够低，流量可以进行空化。 图17.3和17.4对应于半球/圆柱几何上的空化水流。 使用了两种流动模型。 第一个是基于桑切斯-拉科姆状态方程的热动力平衡模型[21]。 第二个使用额外的方程来描述相位属性的传输，如体积或质量分数，从而允许非平衡效应。 在[15]中，使用的方程是spy，dpY，u +-- - 0，在dX（17.5），其中Y是蒸汽的质量分数。 图17.3比较了两个模型的密度轮廓。 如图所示，两者都导致空化气泡界面的锐利捕捉，但它们对腔体在“清醒”区域坍塌的预测上有所不同。 图17.4显示，有限速率模型如果仔细校准，可以产生优于均衡模型的预测。 17.3.2 双流体模型 该模型的基本思想是，每个相由自己的一组质量、动量和能量方程来描述，相之间的相互作用通过源项来处理。 它可以进一步分为单压和双压模型，这取决于假设两个阶段的压力是否相同。 单压假设还意味着某些物理现象，如重力、表面张力、黏度等被忽视。 考虑到流体是可压缩的、不可融合的和相互渗透的，多相流的多流体单压模型由1D系统的六个方程组成：其中向量量是下标“1”和“g”的。 （17.6）pauH（17.7）表示液相和气（蒸汽）相，respec-

374 AUSM+-up方案的多相pow图17.3：密度轮廓：液体水流在半球/圆柱体几何上，取自[15]。 请注意，在单压力模型中，两个阶段的压力都是共同的。 更多的物理学（流动结构）可以通过源术语包括在内，但考虑这个问题超出了本文的范围。 然而，如果源项完全基于动量和能量的相间传递，即sz + s，= 0（17.8）（17.6）和（17.7）包含10个未知数，每个相位（a、p、u、e，p），并且有四个额外的辅助方程，包括两个状态方程，一个约束a1 + ag = I，以及一个压力假设pi = pg = p。 因此，未知数的代数系统是封闭的。 这个方程组，Eqs。 （17.6）和（17.7），提出了两个数值困难：（1）源项是非保守形式，（2）系统一般是非双面的。 如果不仔细处理，源术语可能会引起严重的数值困难，因为这些术语产生的强烈刚度，表现为不同的时间和长度尺度。 数字后果可能是不连续性以错误的速度传播或影响振荡。 由于非多波性，一个初始值问题，指出方程中的六个方程是有指导意义的。

M.-S. Liou和C.-H. Chang 375 0.8 0.6 =0.4 K=0.4（平衡）K=0.3（均衡）Kz0.2（平衡）K=0.4，Rouse和McNown数据K=0.3，Rouse和McNown数据Kz0.2，Rouse和McNown数据A - K=0.4（有限速率）- K=0.3（有限速率）- Kz0.2（有限速率）--- --- \_- 0.2 - - -0.2 --0.4 I -IIIIIIIillillII 1/11 IIIIIII III I 1 2 3 4 5 sld 图17.4：表面压力分布：液态水在半球/圆柱体几何上的流动，取自[15]。在Hadamard [6]意义上变得不好，这意味着解决方案不会持续依赖初始数据。 过去曾多次通过在模型中添加正则化项来消除非双交困难，例如界面压力校正项[l]、双压模型[20]或虚拟质量项[22]。 即便如此，仍然需要隐式运算符或额外的数字耗散来使计算稳定。 因此，溶液中通常会发现过多的涂抹。 基于Roe型近似，已经开发了一个多流体模型的弱公式[24]。 但弱解并不是独一无二的，因为它取决于构建耗散矩阵的路径选择。 此外，多流体模型的特征系统非常复杂，很难找到相关特征矩阵的分析形式。 因此，使用罗伊型或戈迪诺夫方案来计算双流体系统下的数值通量是繁琐或复杂的，特别是当状态方程（EOS）复杂时。

\subsection{17.3.3 多相分层流体模型}
376 AUSM“-多相流 Gas Gas A的up方案\_\_\_---液体液体图17.5：不同阶段之间相间术语的说明。 为了处理多流体模型中的非保守项，之前的所有研究都假设相位之间的动量和能量的转移仅在同一单元内发生。 也就是说，只有图中用A标记的相间项。 17.5被考虑，他们将互相取消。 当虚分数函数是连续的时，这是正确的。 然而，当流动场中存在接触不连续性时，我们发现相邻细胞之间的相间项（在图中用B标记。 17.5）必须考虑在内。 因此，在下文中，我们建议使用分层流体模型来澄清多相流计算中需要单相流动中没有的额外注意的问题。 17.3.3 多相分层流体模型 在[2]中引入了分层流动模型，以描述每个离散细胞，其中每个流体都被限制在分离的控制体内。 类似和不同相之间的界面在控制表面上定义。 因此，可以考虑同一单元内或相邻单元之间不同相之间界面的数值通量。 在我们早期的工作中，分层流动模型已被应用于几个一维案例[2]，包括空水激波管问题、兰森水龙头问题和相分离问题等。 在本文中，我们将扩展分层流模型来求解二维流。 分层模型通过在不同区域中定义不同的流体，非常适合有限体积法的做法，因为保守定律可以分别应用于流体的每个区域。 因此，我们将用它来构建支配每种流体流动运动的离散方程。 此外，尽管Eq的推导。 （17.6）是基于空隙分数函数ai是连续的假设，我们将使其更通用，并允许空隙分数函数分段连续，从而识别细胞边界的相间接口。 我们通过根据分层流动模型定义每个阶段的控制体来开始我们的方法。 图17.6说明了连续体的概念

M.3. Liou和C.-H. Chang 377 Fluid 2 \_---- ----\_\_\_ Fluid 1 图17.6：一维分层流的插图，一维分层流的连续体概念。从二维流中减少。 两个不同的流体，被空隙分数函数ai隔开，被认为是在两个平行板之间流动的。 假设流量仅取决于纵向变量x和时间t。 我们还假设ai是接触不连续情况下的分段连续函数。 从分层流动模型的一般离散版本中，我们可以定义网格单元内气体acgh的控制体，如图所示。 17.7. 图中的虚线表示ai的重建功能。 图中可以找到三种不同类型的接口。 17.7，包括气-气（ab和3）、液-液（cd和g）和气-液界面（bc、cg和fs）。 基于此插图，我们可以定义细胞边界处每种接口的有效长度，j f 1/2如下。\&g-g = min ((@g)L, (ag)R) 6,-i = max (0, -Aag) (17.9) 61-, = max (0, -Aa,) = rnax (0, Aa,) A(.) = (')R - (')L hg-, + 6l-l + 6,-1 + 61-, = 1 其中 ( 1 7.10) (17.11) 并具有互排属性 6g-l.61-, = 0 (17.12) 我们有 6,-1 表示气液界面的有效长度，其左侧为气流体，右侧为液体。 61-的定义是相似的，但左侧是液体，右侧是气体。

378 多相流 j-112 i j+1/2的USM+-up方案图17.7：一维分层流的插图，离散概念。 控制体中每种流体的控制方程可以写成J aipi dV + \$(aipi?).? Ids = 0 \& J aipiv'dV + \$(aipi??).? Ids + \$ rip 6 dS = 0 \&J\textasciitilde{}\textasciitilde{}piE\textasciitilde{}dV+\textasciitilde{}\&JaidV+\$(aipi\textasciitilde{}i\textasciitilde{}).iidS = 0 (1 7.13) 其中积分被接管，是包含两种流体的整个控制体 V及其相关的控制表面S。 然后，Eq。 （17.13）可以离散化并写入一个维度，因为6t是时差运算符，以及\&（.） = (.)n+l - (.)”。 源项，油是相间传递项，例如涉及p*的项，并出于数值目的添加到上述方程中。 细胞边界处的通量分别表示对流和压力通量CF:和F[的总和。 它们写成CFf =（FC）。z 3+1/2 - (Ff)j-1/2 (1 7.15) (17.16) 其中

\subsection{17.3.4 对流通量}
M.3. Liou和C.-H. Chang 379 值得注意的是，(F:)j+l12对应于相邻细胞之间通过细胞边界交换的压力力，而(F:)j代表同一细胞内不同流体之间交换的压力力。 我们将在以下章节中分别讨论这些术语。 17.3.4 对流通量，Ffjhtl，2如图所示。 17.7，气体流体只能通过细胞边界上的气-气体和气-液界面流入细胞j，没有气体或水会流经相界面g。 跨气-气界面（2和8）的对流通量可以通过AUSM+-up方案[13，21作为6g-g [(ag)1/2(M\textasciitilde{})1/2(P\textasciitilde{})L/R -k (D\textasciitilde{})g]（\$g)L/R (17.19)给出，其中\$ = (1, u，以及“L/R”状态由以下规则选择：（17.20）和（17.2 1）\&L'如果[... ]在等式中。 （17.19）> 0，'gL/R = \{ QgR，否则。 细胞接口马赫数 MI/和压力扩散项Dp以及AUSM类型方案特有的其他功能将在附录中给出，以保持完整性。 然而，感兴趣的读者可以在[11，131]中找到有关这些计划的其他细节。 气体在相间界面（bc和fs）上的对流通量是由气体到界面的位置和速度决定的。 当气体位于相间界面左侧与6,-1 \# 0时，如果(ug)\textasciitilde{} > 0，气体将流过界面。 因此，界面上的对流通量为69-1 max(0, (Ug)L) (\&)L (17.22) 同样，当气体在相间界面的右侧时，对流通量为61-9 min (0, (ug)R) (\$g)R (17.23) 请注意，我们假设气体和液体相互渗透。 因此，当我们考虑气流的对流通量时，我们可以忽略液流的影响。 等式的总和。 （17.19）、（17.22）和（17.23）给出了细胞边界上气体的对流通量。 液体的对流通量可以类似地推导。 然后，我们可以得到对流通量Ff的一般形式如下：Ft = 62-2 [(.2)1/2(M2)1/2(Pi)k + (OP,i] (\$i)lC + 62-2' ma@, (4L) (1Cli)L + 621-i dn(0, (ui)\textasciitilde{})(\$2)\textasciitilde{} (17.24)

\subsection{17.3.5 压力通量，}
多相流的380 AUSM+-up方案，如果i = g，则a' = 1，如果i = 1，则i' = g。 17.3.5 在细胞边界上，压力力作用于气体的控制体，透过气体-气体和气体-液体介面。 考虑细胞j和（j+l）之间的细胞边界，气体-气体界面的压力通量（\$）可以通过AUSM+-up方案给出[13，21作为：压力通量，F:j*l,2和Ffj（F:-g)j+i/2 =（hg-g）j+i/2[P\&,）（（Mg）L）PL + PG）（（（Mg）R）PR+（DUlg]。 3+1/2（17.25）注意，在气体的控制表面上总是有一个气体-气体接口。 但只有当相间界面与控制体（dg-l > 0）接触时，才能对相间界面施加压力。 因此，对于单元j，我们将相间界面的压力通量定义为：（17.26）（F：-l）j+l/2 =（bg-l）j+l/2 [P\&）（Mg-l）pL + PG）（Mg-l）PR]，3+1/2（17.27）再次，函数P*和D在附录中给出。 作为第一次尝试，我们只是采取了一种简单的方法，将相同的参数Mg-l应用于拆分函数P+和P-，通过简单平均进行评估。 然后，单元j的压力通量（F，P）j+l/2将是方程的总和。 （17.25）和（17.26）。 请注意，Eq的压力通量。 （17.26）明确交换单元格j中的气体和单元格中的液体（j + 1）之间的压力。 细胞j的压力通量（Ff）j-1/2可以类似地推导，事实上它与j + 1/2的压力通量完全相反。 然后，对于单元格j，我们可以以一般形式将压力通量定义为：以及

\subsection{17.3.6 界面压力修正期限}
M.-S. Liou和C.-H. Chang 381 可以证明，对于具有恒定压力和速度的移动接触不连续性。 施加在每种流体上的压力力将自动平衡，得出C Ff = 0。 因此，Eq中只留下对流通量。 （17.14），确保通过我们的方法准确捕获接触不连续性。 17.3.6 方程的原始多流体模型。 （17.6）（p* = 0）姿势不好。 需要额外的界面压力校正项p*，使其摆好。 在目前的工作中，将采用Bestion [l]提出的界面压力更正条款。 我们有界面压力校正项（17.31），其中Qg和Qz是流体的状态，0是正常数。 如果压力校正项p"包含在0 2 1中，方程（17.6）将得到良好的定位。 参考Eq。 （17.6），界面压力校正项强加在相间接口上。 因此，我们遵循这个想法，并将p*应用于控制表面上的所有相间接口。 我们可以用一般形式oi = (O,OZ, U20i)T (17.32)写Oi，如：其中17.3.7 时间积分我们为时间积分使用四步Runge-Kutta方法。 方程（17.14）的时间离散化可以写成：=（Wi，Mi，\&i）T（17.34），其中上标m和s是当前和下一个时间步骤的索引，w是Runge-Kutta方法每个子步骤的参数常数。 为了计算下一个时间步骤的原始变量，状态方程

382 流体多相流（EOS）的USMf-up方案需要用于上述方程，以便解码变量。 对于本文研究的问题，以下EOS用于表示气相和液相。 事实证明，特殊的液体EOS，称为硬化气体模型，含有理想气体EOS作为特殊情况。 硬化气体模型[8]表示为el = -（yj + P，PI和水[17]的参数为（17.35）p，= 8.5 x lo8 Pa（17.36）J kg。 K' yl = 2.8, (Cp)z M 4186 - 请注意，对于理想气体来说，pa是零，对于液体流体来说是一个非常大的常数。 它使压力和密度场之间的耦合对液体非常弱，密度的变化极其微不足道，即使压力梯度强加在流量上。 另一方面，密度场的小振荡会导致压力的剧烈跳跃。 由于这个EOS，得到了一个压力的二次方程，(p')' + ((-A + B + ("19 - l)pm) p" - AC + - l)Bpm = 0 (17.37)，其中(17.38) c = rzP,+(yz - l)pm 压力p"是上述方程的正根。 然后可以很容易地推导出其他原始变量，9 2 w, \& -1s a; =, a;=1-a;请注意，方程中的系数。 （17.37）非常大，因为p，包括液体。 最大系数和最小系数之间的比率

\section{17.4 计算示例和讨论}
\subsection{17.4.1 Ransom的水龙头问题}
M.3. Liou和C.-H. Chang 383 初始状态 过渡状态 稳定状态 图17.8：水龙头问题的插图。在我们的一些测试案例中，可能大到1016。 即使我们在计算中使用双倍精度，结果中也可能出现巨大的数值误差。 因此，使用额外的牛顿迭代方法来提高ps的准确性。 牛顿迭代方法是标准的，这里不会讨论。 一般来说，两三次迭代就足以驱动lop5下ps的数值误差。 然而，我们仍然可以发现，当ai + 0时，其他变量的数值误差会放大。 因此，在我们的模拟中通常应用a（关于lop6到lo-'的下限。 17.4 计算示例和讨论 17.4.1 赎金的水龙头问题 参考图。 17.8，Ransom [19]引入的水龙头问题由通道内的空气和水射流组成。 起初，被空气包围的水柱以10米/秒的恒定速度移动，并受到引力的施加，使水向下加速（重力对空气的影响可以忽略不计）。 随着水的加速，水柱的宽度越来越窄。 如图所示。 17.8，空隙分数波正向出口移动，当其波前移动出计算域时，流量将变得稳定。 流入边界条件的指定如下。 空气的空隙率ag = 0.2，空气和水的速度ug = u1 = 10 m/s，

384 多相流 0.5的USM+-up方案，，，，1，，，，1，，，，，分析，t 0.5s 0.45 1 -.-iB-.- t = 0.1s I --la-- t=0.2s - -.. B.*- t 0.3s 2 4 6 8 0.150' I' I ' ' ' I ' ' ' I I I I Ransom的水龙头问题，时间= 0.5秒，N=500 \textasciitilde{}=2.0，Dp=0.0，Du=0.0，CFL=0.5，RK-4图17.9：水龙头问题中空隙分数剖面的时间演变。温度Tg = TZ = 300" K，压力从内点外推。 在流出边界中，我们设置压力p = lo5 Pa，所有其他原始变量都从内部流动中推断出来。 图17.9显示了水龙头水问题的空隙分数剖面的时间演变。 我们设置了F = 2.0，IC\textasciitilde{} = K.，= 0，因为流量中没有尖锐的压力梯度。 0.5的Courant 数 CFL和500个单元格用于计算。 分析溶液是通过假设水是不可压缩的，空气和压力变化的影响可能被忽略而得出的。 [17]空隙分数的分析解如下。 计算结果非常接近分析解决方案。 轮廓中没有明显的振荡。 图17.10展示了界面压力校正项p*如何影响多流体模型。 在水龙头上应用了不同的0值

\subsection{17.4.2 空气-水激波管问题}
M.3. Liou和C.-H. Chang 0.55 0.5 '11,\textasciitilde{},',,,,1',,' - - 分析'I - - - --m- o=o II - -.*- oz2.0 -..-.-.- 1 -+--- o=lO.O II II - Irn (T= 5.0 图17.10：基于水龙头问题中不同a的空隙分数剖面。 当a = 0时，多流体模型在数学上是不对的。 在图中的波前可以看到空隙分数的大超冲。 实施界面压力校正项（a>1）可以有效地消除超支。 然而，空隙分数曲线被涂抹为增加，表明界面压力校正项引入了额外的数值耗散。 17.4.2 空气-水激波管问题 第二种情况涉及一个更关键的初始条件，其中基本上纯空气和水分别位于隔膜的两侧，并且对流量施加了非常大的压力差。 我们有（p，ag，ui，Ti）\textasciitilde{} =（10' Pa，1 - E，0 m/s，308.15 KO）（p，as，ui，T\textasciitilde{}）R =（lo5 Pa，E，0 m/s，308.15 KO）（17.41），其中E = 1.0 x lop7。 结果如图所示。 17.11，CFL = 0.5，a = 2和K\textasciitilde{} = K，=1。

\subsection{17.4.3 冲击-气泡相互作用问题}
386 AUSM+-up多相流方案 正如我们在上一节中提到的，当ai -+ 0时，数值误差将被放大。 消失相的Auid状态可以被忽视，因为该相只占据流体的消失部分。 在这种情况下，当ai -+ E时，流体的速度和温度曲线可能会出错误。 然而，它不会影响流动的平均（混合）量，因为它的影响非常微不足道。 这在图17.11中可以看到，其中我们展示了由空隙分数加权的平均轮廓。 我们可以看到粗网格和细网格的结果都非常匹配。 请注意，在这种情况下，过渡区域附近的压力/空虚波会消失，因为当ai -+ 0时，其相应的特征值接近ui。 因此，这些波与接触不连续性相吻合。 接触不连续性和激波被当前方法敏锐地捕捉到，没有可见的振荡。 17.4.3 冲击-气泡相互作用问题 研究了移动的水下冲击与气泡的相互作用。 初始条件与Hankin使用的情况基本相同[7]。 气泡（直径6.0毫米）浸入水中，其中心位于原点。 传入的冲击最初位于z = -4.0mm。 冲击前的流体状态为p = 1.013250 x 105Pa，u = O.Om/s和T = 292.98K1，冲击后的流体状态为p = 1.6 x lOgPa，u = 661.81m/sl和T = 595.14K。 激波的马赫数是M = 1.509。 模拟是在大约149,000个单元格的网格上计算的。 模拟结果的时间演变如图所示。 17.12. 我们观察到，在水激波撞击气泡后，产生了强烈的反射稀释波，并将相对微弱的冲击传递到气泡中（t = 0.6 - 3.0，\textasciitilde{} sec）。 气泡中的冲击强度相对较弱（不超过传入冲击强度的0.1\%），因此很难识别气泡内的压力轮廓。 然而，在马赫数轮廓中可以清楚地看到激波，这里没有显示。 稀疏波产生的水射流不断将气泡推向新月形。 气泡最终破裂，水射流与气泡前方的静止水相撞（t = 3.6\textasciitilde{}秒）。 碰撞产生几个激波，向各个方向放射性地传播。 与此同时，由于施加的极高压力（O（3 x 108Pa））（t = 4.8\textasciitilde{}秒），分离的气泡被压缩成非常小的体积。 之后，激波开始向外分散，其强度降低，由于周围水的压力缓解，导致空气小玩意膨胀。 这个案例清楚地展示了我们的方法在捕获复杂的接口和激波系统方面的能力，而没有在结果中发现严重的涂抹。

M.3. Liou和C.-H. Chang 387 n l.lE+09 1E+09 9E+W 8E+08 7E+08 6E+O8 ' 5E+08 4E+08 3E+O8 PE+08 1E+08 - N=5W -1c+oe\textasciitilde{} ' ' ' ' ' ' ' ' 10 4 X X F 3w " I ' ' ' ' I ' ' ' 0 I 6 8,1,3076' ' ' ' '; ' ' ' ' '; ' ' ' ' ' '; ' ' ' ' 1 X X (a,p,u,T),:(a,p,u,T),=(1-\textasciitilde{},2.0e7,0,308.15):(\textasciitilde{},1.0e7,308.15),\textasciitilde{}=1. OxlO\textasciitilde{}\textasciitilde{} cs = 2.0，D，=1。 O，D，=1。 O，CFL = 0.5，lter = 400，N = 500 图17.11：空气-水激波管问题的状态概况。 初始条件是（p，agr ui，q）\textasciitilde{} =（1.0 x lo9，1 -E，0,308.15），（p，ag，ui，T\textasciitilde{}）R =（1.0 x 105，\textasciitilde{}，0，308.15），E = 1.0 x lo-'和N = 5000。

388 多相流动的USMf-up方案 图17.12：冲击气泡相互作用问题的压力和空隙分数轮廓。

\subsection{17.4.4 冲击水柱相互作用问题}
\section{17.5 结束的讲话}
M.-S. Liou和C.-H. Chang 389 17.4.4 冲击-水柱相互作用问题 在这种情况下，我们在原点有一个水柱（直径6.4毫米），在位置z = -4.0毫米处有一个进入空气激波。 初始条件基于R的论文中的案例。 Nourgaliev等人。 [3]。 冲击前的流体状态为p = 1.0 x 105Pa，u = O.Om/秒，T = 347.0K，冲击后的流体状态为p = 2.3544 x 105Pa，u = 246.24m/秒，T = 451.2K。 模拟中使用了大约18.6万个单元格的网格。 图17.13显示了模拟结果的时间演变。 我们发现，当事件冲击击中水柱时，激波的一部分被传输到水中，其中一部分被反射。 由于水的不可压缩性比空气大得多，水柱内的激波比空气中的激波移动快得多（t = 4.0psec）。 此外，空气中的压力波很容易传入水中，但压力波很难像前例那样从水传播到空气中。 因此，传递到水柱中的压力波基本上被限制在水中，很少被“泄漏”出来。 我们发现海浪在水中快速地来回弹跳，在两个分支进入的激波再次合并之前，气压场基本上不受水柱内发生的事情的影响（t 2 18.5psec）。 17.5 结论 本文描述了在所有速度和所有可压缩性状态下计算多相流的AUSM系列方案的扩展。 我们提出了两种描述多相流的方法。 第一种方法将多相流体视为由实流体状态方程描述的混合物，并应用热力学平衡原理来处理液-气相转变。 这种方法简单而有力，可以自然地描述相变的存在，几个阶段的共存。 可以包括额外的运输方程来解释非平衡效应。 第二种方法基于非平衡模型，该模型通过运输方程分开解决每个相位，相位之间的压力平衡。 提出了分层流动模型来构建数值通量，其中相间效应包含在细胞界面通量中。 新的AUSM+-up方案是详细的，并在两种方法中都采用。 数值结果表明，基于新方案的修改在模拟一些涉及液体和蒸汽的相当具有数值挑战性的问题方面是有效的。 由于密度、速度和音速的巨大差异，所考虑的问题在数值上是僵硬的，因此跨越了小到大的马赫数。 发现计算是稳健和稳定的，与分析、实验和相比，结果是准确的

390 AUSM+-up多相流动方案图17.13：冲击-水柱相互作用问题的压力和空隙分数轮廓。

\section{17.6 致谢}
\section{17.7 书目}
M.-S. Liou和C.-H. Chang 391其他已发表的计算结果。 相界面及其演变被敏锐地捕捉到，没有对界面进行特殊处理。 17.6 致谢 第一作者要感谢：（1）Jack R. 北卡罗来纳州立大学的Edwards进行了几次富有启发性的合作，（2）加州大学圣巴巴拉分校的Nam Dinh对多相流研究的鼓励和启发，以及（3）Richard A. 美国宇航局格伦研究中心的Blech，以支持他在此处介绍的主题上的管理支持。 17.7 参考书目[l] Berger，M. J。 \& 奥利格，J. 卡塔尔代码中的物理闭合定律。 核工程与设计，124：229-245，1990年。 [2]张，C.-H. \& Liou，M.3。 一种使用ausmf方案模拟可压多流体流动的新方法。 AIAA论文2003-4107，2003年。 [3] 丁，N.，努尔加利耶夫，R. \& Theofanous，T。 可压缩多相流的直接数值模拟：激波与分散多材料介质的相互作用。 技术报告，第五届多相流国际会议，2004年。 [4]爱德华兹，J。 R。 用于Navier-Stokes计算的低扩散通量分割方案。 计算机与流体，26：635-659，1997年。 [5]爱德华兹，J。 R.，富兰克林，R. K.，\&Liou，M.3。 所有速度下真实流体流动的低扩散通量分割方法。 AIAA杂志，38：1624-1633，2000年。 [6]哈达马德，J. 线性偏微分方程中的柯西问题的讲座。 多佛出版社，纽约，1952年。 [7] 汉金，R. K。 S。 守恒形式的多相可压缩流动的Euler 方程：冲击气泡相互作用的模拟。 计算物理学杂志，172：808-826，2001年。 [8]哈洛，F。 H。 和阿姆斯登，A。 A. 所有流速的数值流体动力学计算方法。 计算物理学杂志，8:197-213，1971年。 [9]石居，M。 两相流动的热流动力学理论。 Eyrolles，巴黎，1975年。

392 多相流的AUSM+-up方案[lo] Kestin，J。 热力学课程。 布莱斯德尔出版公司，马萨诸塞州沃尔瑟姆，1966年。 [11] Liou，M.-S. AUSM的续集：AUSM+。 计算物理学杂志，129：364-382，1996年。 还有美国宇航局TM 106524，1994年3月。 [12] Liou，M.-S. 在造-ausm-家庭的十年。 AIAA文件2001-2521-CP，第15届AIAA CFD会议，2001年6月11日至14日。 1131 Liou，M.-S. 进一步开发ausm+方案，以实现所有速度的强大和准确的解决方案。 AIAA论文2003-4116-CP，第16届AIAA CFD会议。 2003年6月23日至26日。 [14] Liou，M.-S. \& Edwards，J. R。 低马赫和多相流量的ausm方案和扩展。 VKI系列讲座1999-03，VKI，比利时，1999年。 [15] Liou，M.-S. \& Edwards，J. R。 所有速度的多相流的Ausm系列方案。 在M。 M。 Hafez，编辑，不可压缩流的数值模拟，第517-543页。 世界科学，2003年。 [16] Liou，M.-S. \& Steffen，C. J.，小 一个新的通量分裂计划。 计算物理学杂志，107:23-39，1993年。 还有美国宇航局TM 104404，1991年5月。 [17] Pailkre, H., Core, C., \& Garcia, J. 关于将AUSM+方案扩展到可压缩双流体模型。 计算机和流体，32：891-916，2003年。 [18]彭，D.-Y. 和罗宾逊，D。 C. 一个新的双常数状态方程。 印度。 英语。 化学。 Fundam.，15（1）：59-64，1976年。 [19] 赎金，V。 H。 数值基准测试。 在G。 F。 休伊特，J。 M。 Del-hay，和N。 Zuber，编辑，多相科技，卷。 3. 半球出版社，华盛顿特区，1987年。 [20]赎金，V。 H。 \& Hicks，D。 L。 两相流动的双曲双压模型。 计算物理学杂志，751498-504，1988年。[all Sanchez，I. C. \& 拉科姆，R. H。 经典流体的基本分子理论：纯流体。 物理化学杂志，80（21）：2352-2362，1976年。 [22] Stadtke, H., Franchello, G., \& Worth, B. 基于通量向量分割的多维两相流的数值模拟。 核工程与设计，17：199-213，1997年。 [23] Stewart，H. B. \& Wendroff，B. 两相流程：模型和方法。 计算物理学杂志，56:363-409，1984年。

M.-S. Liou和C.-H. Chang 393 [24] Toumi, I., Kumbaro, A., and Paillere, H. 两相流的近似黎曼求解器和通量向量分割方案。 VKI讲座系列1999-03，第30届计算流体动力学，冯·卡曼研究所，比利时，1999年。 [25] Wada, Y. \& and Liou, M.-S. 用于冲击和接触不连续性的准确而强大的通量分裂方案。 SIAM科学和统计计算杂志，18：633-657，1997年。 [26] Wallis，G。 B. 一维两相流。 麦格劳-希尔，纽约，1969年。

\subsection{17-A数值通量公式}
附录17-A 数值通量公式 在本文中，我们在数值通量公式中使用了以下定义。 让Mfi)W = \$(M * IMI), (17.42) M\$)(M) = \&+(M f 1)2 然后，每个相的界面马赫数被定义为，分裂压力函数被定义为最后，质量通量和压力通量中使用的压力和速度扩散项由(17.46) AM max(1 - M2, 0) (PL - PR)给出。a D, = K,, K, = 1.0 - D, = K,?;)(a) pG)(a) pa('U\textasciitilde{} - 'UR); K, = 1.0 (17.47)，其中AM = MG) (ML) - M;) (ML) - M(4) (MR) + MG)(MR) (17.48)参数u, p和M是L和R状态值的平均值。

\clearpage
\chapter{第18章 用于计算复杂几何中多相流的有限体积界面追踪法 Metin Muradoglu}
\section{18.1 导言}
第18章 复杂几何中多相流计算的有限体积界面追踪法 Metin Muradogul 18.1 介绍 多相流的模拟是notoriouL-Jr困难的，主要是由于存在变形相界。 各种数值方法已经开发出来，并成功应用于广泛的多流体和多相流问题[l9、24、25、29、341。 尽管取得了这一成功，但仍需要取得重大进展，特别是对于涉及与复杂固体边界的强相互作用的多相流的准确计算。 计算多相流量的最流行的方法分为四类：第一类是前捕获方法，如液量（VOF）[10，241和水平集[lg，25，281方法。 在这些方法中，正面通过体积分数分布（VOF）或距离函数的零级集间接捕获。 受限制的“机械工程系，科克大学，Rumelifeneri Yolu，Sariyer 34450伊斯坦布尔，土耳其。

396 Yabe[34]的有限体积/前包装方法插值剖面（CIP）方法也属于这一类别。 第二类是边界拟合网格方法，其中每个阶段都使用单独的边界拟合网格[33]。 第三类是具有移动网格的拉格朗日方法[ll]。 本研究中使用的第四种方法是界面追踪法[29，321。 Unverdi和Tryggvason[32]开发的界面追踪法基于Navier-Stokes方程的单场公式，并将不同相视为具有可变材料特性的单个流体。 在这种方法中，固定的尤拉式网格用于流体流动，介面由单独的拉格朗日网格明确跟踪。 Peskin[2l]开发的浸入边界方法用于平稳地离散材料特性的跳跃，并处理表面张力的影响。 界面追踪法与有限差分流动求解器相结合，已成功应用于广泛的多相流问题，但除了Udaykumar等人1之外，几乎所有的几何形状都相对较简单[29]。[ 31]。 Tryggvason等人最近对该方法进行了审查。 [29]。 在许多工程和科学应用中，如微流体系统[27]、孔隙尺度多相流动过程[l7，181和生物系统[7，221]，能够准确模拟气泡/水滴与弯曲固体边界之间的强相互作用非常重要。 界面追踪法有很多优点，如简单和缺乏数字扩散。 然而，它的主要缺点可能是难以维持拉格朗日标记点和欧拉体拟合曲线或非结构网格之间的通信。 在本研究中，开发了一个有限体积/界面追踪法，以利用身体拟合曲线网格计算复杂几何中的分散多相流。 开发了一种高效且稳健的跟踪算法，用于跟踪身体拟合曲线网格中的前标记点。 跟踪算法使用辅助正则笛卡尔网格，它可以很容易地扩展到非结构化网格。 前跟踪方法扩展到车身贴合的曲线网格，并与新开发的有限体积法相结合，以促进相位和复杂固体边界之间强相互作用的准确高效建模。 有限体积法基于双（或伪）时间步进方法的概念，是为不可压缩层流的时间准确计算而开发的。 该方法用于计算复杂几何中的二维（平面或轴对称）分散多相流[l4，151。 首先测试该方法的振荡跌落，并将结果与分析解决方案进行比较。 Han和Tryggvason[8]早些时候研究的直通道中，由于重力导致水滴的运动也得到了验证。 发现目前的结果与Han和Tryggvason[8]获得的结果非常一致。 然后将该方法应用于计算Hemmat和Borhan[9]实验研究的收缩通道中水滴的浮力驱动运动。 最后，平面二维

\section{18.2 垫子血液配方}
M。 该方法的Muradoglu 397版本用于计算通过绕组通道移动的液滴中的混沌混合[l6]。 在下一节中，简要回顾了控制方程，并将其转换为任意的曲线坐标系。 然后，第3节描述了有限体积/前跟踪算法。 结果在第4节中提出和讨论，在第5节中得出一些结论。 附录介绍了人工可压缩性参数的选择分析。 18.2 Mat 血液配方 本节仅简要描述控制方程的轴对称流，但指出二维平面流的方程基本上可以以相同的形式表达。 轴对称流动的不可压缩流动方程可以用向量形式的圆柱坐标写成af dg af, ag, ag + - + - = - f - + h, +fb, dt dr dz dr dz (18.1) where and In Eqs. (18.1)-(18.3)，r和z是径向和轴向坐标，t是物理时间；p、p和p是流体密度、动态黏度和压力；w和w分别是r和z坐标方向的速度分量。 黏性通量向量中出现的黏性应力由方程中的最后一个项给出。 （18.1）表示浮力和表面张力产生的体力，由fb = -给出。（ Po - p)G - ralcnb(x - xf)ds，（18.5）

398 有限体积/前跟踪方法，其中第一个项表示浮力引起的身体力，po和G分别是环境流体的密度和重力加速度。 等式中的第二个项。 （18.5）表示表面张力引起的身体力，S、xf、n、K、n、S和ds分别表示狄拉克三角洲函数、前部位置、表面张力系数、平均曲率的两倍、界面上的外单位法向量、界面的表面积和界面的表面积元素。 在Eq中。 （18.1），假定流体具有恒定特性，因此流体颗粒的密度和黏度保持恒定，即（18.6），其中实质性导数定义为= \& + u。 V。 如Eq所示。 （18.1），连续性方程与动量方程解耦，因为它没有任何时间导数项。 为了克服这一困难，并能够使用时间行进解算法，人工时间导数项以r是伪时间的形式添加到流动方程中。 解向量w、不完全恒等矩阵I1和预处理矩阵 rW1由W= \{ 7r\textasciitilde{}", \}, I1= [ 0 i 01, r-1= [ f:], (18.8) r'uz 001 0P给出，其中p是要确定的预处理参数，它具有速度的维度。 请注意，关系q = pI1w已在Eq中使用。 （18.7）。 根据附录中给出的分析，预处理参数，O由（18.9）给出，其中'GO是单位顺序的常数，Uref和！ 分别是速度和长度尺度，雷诺数 Ree被定义为Reg = U,,f! /Po。 方程（18.7）可以转换为一般的曲线座标系E = E（r，z），77 = dr，z），（18.10），所得方程的形式为dhw dphw ahF dhG dhF，dhG，r-1- + 11- +-+-=- +- + h（h，+ + fb），（18.11）ar at at av at a7

\section{18.3 数值方法}
M。 Muradoglu 399 Auxilary Uniform urvilinear Grid r CartesianGnd r " " 图18.1：计算中使用的三种类型的网格。 控制方程在固定的欧拉曲线网格上求解，不同相位之间的界面由由连接标记点组成的拉格朗日网格表示。 辅助均匀笛卡尔网格用于保持曲线网格和拉格朗日网格之间的通信。其中h = rp? - rqzC表示变换的雅各布。 向量hF = z,f- r7g, hG = -zcf + rtg, (18.12)，分别表示变换的黏性和黏性通量向量。 18.3 数值方法 一旦确定了材料特性和表面张力力，原则上任何标准时间行进算法都可以用于求解方程。 （18.11）因为这些与通常的连续流动方程的形式相同。 图18.1概述了本方法中使用的三种类型的网格。 固定曲线网格用于求解守恒方程（Eqs. （18.11）），而低维的拉格朗日网格用于跟踪分隔不同相位的界面。 辅助均匀笛卡尔网格用于保持曲线网格和拉格朗日网格之间的通信效率。

\subsection{18.3.1 流动方程的积分}
400 A 有限-体积/界面追踪法 18.3.1 跟随Caughey[5]，一个求解Eq的双参数数值方案家族。 （18.11）可以写成流动方程的积分（18.14），其中（...） P和（...）”分别表示伪和物理时间水平。 参数cp支配物理时间导数的近似值，6决定了该方法在伪时间中的隐含性水平。 请注意，当伪时间达到稳态时，我们有wp + wn+l。 Cp和0的三个组合特别令人感兴趣，对应于物理时间的不同近似值，如表18.1所示。 在本研究中，三点向后隐式方法贯穿始终。 请注意，在当前公式中，术语h和fb在伪时间中被明确处理，尽管可以将h包含在隐式运算符中。 增量AT表示子迭代的时间步，而At表示通常不同的物理时间步。 校正AwP = wP+' - wp是根据求解方程的子迭代计算的。 (18.15)线性化为(hF)p+l = (hF)p + APAwP + O(Ar2), (\textasciitilde{}LG)\textasciitilde{}" = (hG)p + BpAwp + O(Ar2), (hF,)p+l = (hF,)p + AEAw? + O(Ar2), (hG,)p+l = (hG,)p + B[Aw: + O(A?)), (18.16)

M。 Muradoglu 4 01 cp 6' 方案精度顺序1隐式后退欧拉法一阶1/2隐式梯形法二阶1三点后向后向隐式方法二阶表18.1：物理时间积分方案。其中不黏性和黏性雅各布矩阵被定义为dhF、dhG、dhF dhG Ap=\{\%-r' Bp=\{x)y; B:=\{bw,\} (18.17)，WE =然后写成和wV = F。 等式线性近似。 (18.15)可以aE Aw' -Rp, (18.18) 其中r-l s = - +p(2++1, AT 2At (18.19) and P - hh,] + 8hh;l" dh(F - F,) + dh(G - G,) drl dh(F - F,) dh(G - G,) + - h(h, + fb)] n. （18.20）arl在Eq。 （18.19），，\&'在解过程中近似为@'+' = pn+l。 方程（18.18）表示一个线性方程组，可以通过各种方法求解修正Awp，但遵循Briley和McDonald[2]以及Beam和Warming[l]，计算效率分解为（18.21），当空间导数通过三点近似离散化时，可以使用块三对角线求解器分两个步骤求解。

\subsection{18.3.2 前-跟踪 Met hod}
402 A 有限-体积/前-跟踪方法 空间导数使用以细胞为中心的有限体积法进行近似，该推导数相当于均匀笛卡尔网格上的二阶中心差，四阶数值耗散项类似于Caughey[4]的项被添加到Eq的右侧。 （18.21）防止奇偶解耦合。 请注意，数值耗散术语在伪时间中被明确处理。 由于伪时间的准确性不感兴趣，除了预处理方法外，还使用类似于Caughey[5]的多重网格法和本地时间步进方法来进一步加速伪时间的收敛速率。 还实现了AD1方法的对角化版本，类似于Pulliam和Chaussee[23]的方法，其中只有对流项被隐式处理，所有其他项在伪时间中显式处理。 Muradoglu和Gokaltun[l4]描述了对角化方案的平面二维版本。 由于本文中研究的所有案例基本上都在斯托克斯极限，因此发现块对角线版本比对角化版本更有效。 因此，本研究使用了AD1方法的块对角线版本。 18.3.2 Front-Tracking Met hod 在Unverdi和Tryggvason[32]开发的界面追踪法中，不同相位之间的界面由带有连接标记点的拉格朗日网格表示，如图所示。 18.1. 标记点可以被视为以局部流速移动的流体粒子。 为了保持拉格朗日网格和固定曲线网格之间的通信，有必要在每个物理时间步骤中确定曲线网格中标记点的位置。 虽然在均匀的笛卡尔网格中确定标记点的位置是一项简单的任务，但在一般曲线或非结构网格中跟踪它们要困难得多。 为了克服这一困难，并保持跟踪在计算上的可行性，最近开发了一种新的跟踪算法，并发现它非常稳健和计算高效[15]。 跟踪算法使用图中概述的辅助均匀笛卡尔网格。 18.1，并将曲线网格上的粒子跟踪减少到带有查找表的常规笛卡尔网格上的粒子跟踪。 这里强调的是，跟踪算法不限于结构化网格，并且可以很容易地适应非结构化网格。 由于在正笛卡尔网格上跟踪粒子是一项微不足道的任务，新的跟踪算法具有计算效率，使界面追踪法成为使用曲线或非结构化网格在复杂几何中计算分散多相流的可行工具。 此外，由于跟踪算法是通用的，它可用于跟踪其他方法中的拉格朗日点，例如广泛用于求解湍流反应流的概率密度（PDF）模型方程的基于粒子的Mote Carlo方法[l3]。 足够的

M。 Muradoglu 403跟踪方法的详细信息可以在Muradoglu和Kayaalp[l5]中找到。 由于流动方程是在曲线网格上求解的，但表面张力是在前面计算的，因此有必要通过适当的分布函数将表面张力转换为身体力。 这涉及以保守的方式近似曲线网格上的delta函数。 让q5f是每单位表面积的界面量，它应该转换为每单位体积给出的网格值q5g。 为了确保总值在平滑中保持，我们必须有（18.22），其中Av是网格单元格的体积。 遵循Tryggvason等a1.[29]，通过写入（18.23）满足这种一致性条件，其中41是前值的离散近似值4fl 4ij是网格值q5gl的近似值，因为是前元素1的面积，wfj是网格点ij相对于元素1的权重。 为了一致性，权重必须满足CWZ = 1，（18.24），但可以通过不同的方式选择[29]。 在本研究中，从xp =（rp，zp）平滑的网格点ij的权重被定义为张量积内核，形式为Wij（Xp）= K（TT）K（TZ），（18.25），其中T，= Irp - rijl/rmax和T，= IzP - ZijI/zmax。 请注意，rmax和z,,,是网格节点的最大距离，前面数量q5ij分别在T和z方向上分布。 这里使用的K的函数形式是\$ - 8(1 -?) F2 if i 5 0.5 ij;(l - 33 if 0.5 < i 5 1.0 (18.26)，大约是对称的i = 0，并且是分段立方，有连续的一导数和二阶导数。 权重wij被归一化，以满足Eq给出的一致性条件。 （18.24）。 权重wij也用于将网格值（如速度场）从曲线网格插值到前点。 密度和黏度等材料属性是根据10计算的，否则K（F）=（18.27）

\subsection{18.3.3 整体解决方案程序}
4 04 有限-音量/前-跟踪方法，其中下标o和d分别指环境和滴流体。 指示函数4的定义是，它在水滴内为单位，在水滴外为零，根据Tryggvason等人[29]，它是通过求解泊松方程0'4 = vh. vh4, (18.28)得到的，其中Vh是梯度算子的离散版本。 跳跃Vh4使用Peskin分布[29]和Eq分布在相邻的网格单元上。 （18.28）然后在每个水滴附近的均匀网格上求解。 在均匀网格上计算指标函数后，使用双线性插值将其插值到曲线网格上。 事实上，在曲线网格上有效地求解泊松方程是可能的，但上述程序似乎很稳健，并为本工作中研究的问题提供了足够流畅的解决方案。 每个前元素上的表面张力是按照Tryggvason等人描述的程序计算的。[ 29]。 小前元素上的表面张力可以计算为SF，=]A，rulcnds。 （18.29）使用二维线的曲率定义，即Kn = ds/ds，并考虑问题的轴对称性，Eq. （18.29）可以积分生成bF，= r4sz - SL）- Asm，，（18.30），其中e，表示径向的单位向量。 曲线s的切向量是通过拉格朗日多项式拟合通过每个元素的端点和相邻元素的端点直接计算的，方式与Tryggvason et a1所描述的相同。[ 29]。 18.3.3 整体解决方案程序 上述有限体积和前向跟踪方法的组合如下。 在从物理时间水平n（t，=n）推进解中。 At）到n + 1级，首先使用显式欧拉方法预测新时间级n + 1的标记点的位置，即\%;+' = Xp” + Atvp”，（18.31），其中X和V表示前标记点的位置和从相邻的曲线网格点插值到前点的速度

\section{18.4 结果和讨论}
M。 Muradoglu 4 05 X，，分别。 然后使用预测的前位置评估材料特性和表面张力，如-n+l -n+l -n+l,,n+l = p(X, ); pn+l =p(X, 1; \$+l =f\& 1。 （18.32）然后通过求解流动方程（Eq.）计算新物理时间水平n + 1的速度和压力场。 （18.11））通过FV方法进行单个物理时间步，最后将前点的位置校正为At 2 x;+l = x; + -(v; + v;+1)。 （18.33）在此步骤之后，使用更正的前部位置重新评估材料特性和身体力。 该方法在时间和空间上都是二阶准确的。 除Eq.中的fb以外的所有术语。（ 18.11）在物理时间中隐含地处理，因此物理时间完全由精度考虑和稳定性约束决定，主要是由于表面张力。 完美反射边界条件用于前标记点的实体边界，即由于数值误差而越过实体边界的前标记点相对于向内法向量反射回计算域。 如果前标记点靠近边界，则以保守的方式将前属性分布在曲线单元格上，即权重仅为计算域内的单元格定义，并归一化以满足Eq给出的一致性条件。 （18.24）。 网格属性以类似的方式插值到前端点。 拉格朗日网格最初是均匀的，通过删除小元素并以Tryggvason等人描述的方式拆分大元素，在整个计算过程中几乎保持均匀。[ 29]。 初始前元素大小通常设置为0.75A1，其中A1是曲线网格单元格的最小尺寸。 在模拟过程中，在每个物理时间步骤中，小于0.5Al的元素被删除，大于A2的元素被拆分，以保持拉格朗日网格几乎均匀，并防止形成比网格大小小得多的摆动。 18.4 结果和讨论 该方法首先对浸入另一种流体中的水滴振荡的经典问题进行了验证。 振荡是由液滴表面配置的初始扰动引起的，并因黏性消散而阻尼。 计算结果在振荡频率和阻尼率方面与分析解决方案进行比较。 它还验证了Han和Tryggvason[8]早些时候使用有限差分/前跟踪（FD/FT）方法研究的直通道中自由坠落的运动。 然后应用该方法来计算自由上升的水滴的运动，因为

\subsection{18.4.1 振荡跌落}
406 Hem-mat和Borhan实验研究了连续收缩通道中的有限体积/前跟踪方法浮力[9]。 最后，该方法的平面二维版本用于计算通过绕组通道移动的液滴中的混沌混合[161。 18.4.1 振荡跌落 第一个测试案例涉及振荡跌落的经典问题。 振荡是由对液滴表面配置的初始扰动引起的，并因黏性阻尼而阻尼。 考虑一滴扰动半径为（18.34）的液体，其中a是等效的液滴半径，是扰动的大小。 请注意，这种形式的变形保留了半径a[6]的下降体积。 当下降被扰动时] 表面张力尝试将接口拉回球形配置，导致表面振荡。 Lamb[l2]表明，振荡频率由T=a i+-（i+3cos2e）给出，周期为（18.35）（18.36），其中n是振荡模式，n = 2是主模式，0是表面张力系数。 忽略环境流体黏度的影响，黏性阻尼时间常数由（18.37）给出，其中v是滴Auid的运动黏度。 图18.2显示了变形参数y的演变，定义为轴向下降半径与径向下降半径的比率。 初始扰动设置为C = 0.05。 计算域在径向方向为6a，在轴向方向上扩展到12a]，并通过在轴向和径向方向延伸的192 x 384笛卡尔网格解析，以更好地解析水滴附近的流动。 周期性边界条件用于轴向，无穿透、完美滑动边界条件用于径向。 滴流体的密度和运动黏度分别比环境流体的密度和运动黏度大50倍和400倍。 对于最低（n = 2）模式，时间由理论周期非维度化。 理论阻尼率也绘制在

\subsection{18.4.2 浮力驱动的直线通道中的坠落}
M。 Muradoglu 4 07 图18.2：浸入宿主流体中的振荡液滴的变形y与非维时间，扰动系数< = 0.05，pd/po = 50和pd/po = 400。 理论上的衰减率由振荡曲线下方和上方的虚线表示。 图。 18.2. 从这个图中可以看出，用当前的方法很好地预测了振荡频率和阻尼率。 例如，在t/rth = 2.0时，计算振荡频率和理论振荡频率之间的误差小于1\%，这表明了本标准测试案例的当前方法的准确性。 18.4.2 直通道中浮力驱动的坠落下降 第二个测试案例涉及Han和Tryggvason之前研究的直通道中浮力驱动的坠落下降[8]。 物理问题和计算域如图所示。 18.3. 如图所示，环境流体完全填充了刚性圆柱体，由于引力，比环境流体密度更高的液滴向下加速。 这个问题由四个非维参数[8]支配，即Eotvos数Eo（可互换称为Bond数，Bo），Ohne-

有限体积/前跟踪方法Z环境流体Po *bo图18.3：直线通道中浮力驱动的坠落体的物理问题和计算域的示意图。索格数Ohd，密度和黏度比定义为Pd PO I** = - (18.38)，其中Ap = Pd - p，是液滴和环境流体之间的密度差，gz是引力加速度，d是初始液滴直径。 基于环境流体的Ohnesorge数的定义与Oh, = e相似。 下标d和0分别表示液滴和环境流体的属性。 非维时间定义为t t* = - (18.39) 在本节中介绍的所有计算中，计算域在径向方向上为5d，在轴向方向上为15d。 防滑边界条件 IhG。

M。 Muradoglu I 1 EO = 12（本）I I I I I I I I I I I I I I I I I I I I I I I I I I I I I I I I I I I I I I I I I I I I I I I I I I I I I I I I I I I I I I I I I I I I I I I I I I I I I I I I I I I I I I I I I I I I I I I I I I I I I I I I I I I I I I I I I I I I I I I I I I I I I I I I I I I I I I I I I I I I I I I I I I I I I I I I I I I I I I I I I I I I I I I I I I I I I I I I I I I I I I I I I I I I 当前结果在每组的左侧绘制。 每列中两个连续下降之间的间隙代表下降在固定时间间隔内行驶的距离，最后一个界面分别绘制在t* = 39.6和t* = 44.01，Eo = 12和Eo = 24的情况。应用于圆柱体壁上，轴对称条件应用于中心线。 下降重心最初位于（r，，z，）=（0,124。 计算域由128 x 768正笛卡尔网格解析。 网格在径向方向上拉伸，以便有更多的网格点靠近中心线。 Ohnesorge数，密度和黏度比在Ohd = 0.0466（Oh，= 0.05），pd/po = 1.15和pd/p，=1中保持恒定。 在图中。 18.4，Eo = 12和Eo = 24的下降演变与有限差/前跟踪（FD/FT）方法获得的结果一起呈现[8]。 如图所示，目前的结果与FD/FT方法获得的结果在定性上非常一致。 为了量化本方法的准确性，速度由\&Z非维化，滴重心V的非维速度如图所示。 18.5a与FD/FT结果一起。 从这个图中可以清楚地看到，目前的结果也与FD/FT方法的定量一致性非常好，表明

\subsection{18.4.3 持续收缩通道中浮力驱动的上升下降}
410 A 有限体积/前背法图18.5：（a）中心体速度（b）Eo = 12和Eo = 24病例的下降体积与t*的百分比变化。 虚线是当前结果，实线是FD/FT结果。当前方法的准确性。 最后，滴水量的百分比变化如图所示。 18.5b作为Eo = 12和Eo = 24病例的非维时间的函数，并再次与FD/FT方法的结果进行比较。 该图表明，当前方法和FD/FT方法的体积变化是相同的顺序，在当前方法中，模拟结束时的最大体积变化小于2.5\%。 与FD/FT方法相比，当前方法的体积变化相对较大，可能归因于当前结果中的数值误差更大，部分原因是物理时间步骤更大，部分原因是插值和分布算法。 18.4.3 连续收缩通道中的浮力驱动上升下降 之前的测试案例证实了本方法的准确性。 这个测试案例涉及Hemmat和Borhan[S]实验研究的具有周期性波纹的黏性水滴通过垂直毛细管的浮力驱动运动。 计算设置如图所示。 18.6a。 毛毛管由一个26厘米长、周期性收缩的圆柱形管组成，有6个波纹。 管子的平均内半径分别为R = 0.5厘米，波纹的波长和振幅分别为h = 4厘米和A = 0.07厘米。 悬浮液是二乙二醇-甘油混合物，用DEGG12表示。 各种UCON油被用作滴液。 表18.2总结了滴液和悬浮流体的特性，每个系统都使用了与Hemmat和Borhan[9]相同的标签。 一个com-

M。 Muradoglu 411系统DEGGl2对实验设置的详细描述可以在[9]中找到。 图中绘制了包含8 x 416个网格单元格的粗网格的一部分。 18.6b显示模拟中使用的车身网格的整体结构。 计算浮力下降的平均上升速度以及下降的下降形状，并将结果与一系列管理参数的实验数据[9]进行比较，即无维下降尺寸，K，定义为等效球形下降半径与平均毛细管半径的比率，下降与悬浮流体黏度的比率，A，流体密度y的相应比率，以及债券数Bo = ApgZR2/a，表示浮力与界面张力的比率；Ap和u表示下降和悬浮之间的密度差和界面张力 流体分别是，gz是引力加速度。 在本节中介绍的所有结果中，水滴最初位于z = 1.5h的环境流体中，该流体完全填充了圆柱形管，最初处于流体静力学条件下。 沿中心线应用对称边界条件，在圆柱管的顶部、底部和侧面使用防滑边界条件。 水滴最初是静止的，由于浮力而开始上升。 水滴最初是球形的。 结果以上标“*”表示的非维量表示。 无维坐标定义为z* = z/h和r* = r/R。 时间和速度分别用Tref =和Kef =“py2）变得无维。 首先，对水滴形状的定性分析如图所示。 18.6\textasciitilde{}。 在本图中，为DEGG12系统绘制了通过收缩通道的黏性水滴形状演变的一系列图像，非维水滴尺寸为K = 0.54、0.78和0.92。 计算在32 x 1664网格上进行，物理时间步骤为At* = 1.641，在每次子迭代中，残差会减少三个数量级。 如这些图所示，当大水滴（K>0.7）达到收缩时，其前缘沿着毛细管壁轮廓挤压穿过喉部。 一旦前半月板清除喉咙，其上升速度随着进入发散横截面而增加，而下降的后缘仍然被困在喉咙后面，类似于实验观测[9]。 水滴形状、速度场和水滴附近的压力轮廓如图所示。 18.7 对于K = 0.92的DEGG12下降，当它通过喉咙po pd Po Pd CT（mPa.s）（mPa.s）（kg/m3）（kg/m3）（N/m）87 115 1160 966 0.0042表18.2：计算中使用的两相系统。

环境流体P，。 IJ.O h A 有限-体积/前-D-acking方法 图18.6：（a）收缩通道中浮力驱动上升下降的计算设置的示意图。 （B）包含8 x 416个单元格的粗计算网格的一部分。 （c）DEGG12系统的浮力下降的快照，分别从左到右下降大小K = 0.54、0.78和0.92。 每列中两个连续的水滴之间的间隙表示水滴在固定的时间间隔内行驶的距离，最后一个界面分别在t* = 2831.3、3693.0和5416.4时从左到右绘制。就在喉咙之后，以更好地显示解决方案的整体质量。 最后，由波纹波长缩放的垂直下降尖端位置和由参考速度Vref缩放的下降尖端上升速度与图中的非维度时间进行绘制。 18.8适用于各种水滴尺寸的DEGG12系统。 收缩的延迟效应在这些数字中可以清楚地看到大跌落，即K。 > 0.7。 还可以看出，在所有情况下，水滴都会迅速加速并达到周期性运动。

\subsection{18.4.4 通过蜿蜒通道移动的液滴中的混沌混合}
M。 Muradoglu 413 图18.7：DEGG12附近的速度向量（右部分）和压力轮廓（左部分）下降6 = 0.92，同时通过（a）喉咙和（b）膨胀区域。 网格：32 x 1664，At* = 1.641。 18.4.4 液滴在蜿蜒通道中的混沌混合 最后的测试案例涉及平面二维设定中穿过蜿蜒通道的液滴中的混沌混合的计算建模[l6]。 通过混沌对流在液滴内的混合已被用于以高时间分辨率和低消耗的样品进行动力学测量[3，261。 这个测试问题的目的是展示该方法在涉及变形滴相和固壁之间强相互作用的复杂几何中计算分散多相流动的能力。 通道几何形状如图所示。 18.9和问题的完整描述可以在Muradoglu和Stone[l6]中找到。 计算域由包含1024 x 64网格单元的身体拟合网格解析。 体积流量是根据完全发展的通道流量在入口处指定的，压力在出口处是固定的。 使用环境流体属性将流动初始化为单相稳定流动，然后将圆柱形水滴放置在环境流动中，如草图所示。 分子混合被忽略，只考虑混沌对流的混合。 混合模式由被动示踪粒子可视化，这些粒子最初在液滴内随机均匀分布。 最初占据液滴下半部分的粒子被识别为“红色”，而其余的是“蓝色”。 足够的

\section{18.5 结论}
414 A 有限-体积/前-跟踪方法 0 500 1000 1500 2000 2500 3000 3500 4000 t t t（4（b）图18.8：与K的DEGG12系统滴落的非维时间t*相比，滴头的非维垂直位置（左图）和非维上升速度（右图）。 = 0.54,0.78和0.92。 网格：32 x 1664，At* = 1.641。示踪粒子的对流方式与前标记点相同。 控制非维度参数被识别为雷诺数数Re = poUidc/po，毛细管数Ca = p0Ut/u，黏度比X = pd/po，密度比y = pd/po，液滴大小与通道入口高度的比率C = dd/dc，其中po和Pd是环境和下降相流体密度，po和ud分别是环境和流体黏度。 基于入口速度和波纹波长，非维物理时间定义为t* = tUi/L。 图中绘制了在非维时间框架t" = 0,3.3 x lop3、7.3 x lop3、11.3 x 10F3、14.7 x 18.0 x lop3和22.0 x下拍摄的混合模式的快照。 18.10演示在通过模型绕组通道移动时，滴中混合模式的演变。 在这种情况下，非维数设置为Ca = 0.025，Re = 6.6，C = 0.7576，X = 1.0。 混合模式在图的顶部图中放大。 18.10 更好地展示混合过程的细节。 如图所示，混沌对流在通过绕组通道移动时以二维下降的形式发生，混合模式在质量上与实际的三维混合模式相似[20]。 18.5 结论 已经开发了一种有限体积/前跟踪（FV/FT）方法，用于复杂几何中分散的多相流的乘法。 该方法基于流动方程的单场公式，并将不同的相视为具有可变材料特性的单个流体。 流动方程

\section{18.6 参考书目}
M。 Muradoglu 415 图 18.9：计算中使用的通道草图。 体积流量是根据完全发展的通道流量在入口处指定的，压力在出口处是固定的。 使用环境流体特性，将流量初始化为单相稳定流量，然后立即将圆柱形水滴放置在环境流量中。通过FV方法在适合身体的曲线网格上解决，并使用单独的拉格朗日网格来表示不同相位之间的界面。 FV方法基于双时间推进和时间积分的概念，由块对角交替方向隐式（ADI）方案完成。 开发了一种新的跟踪算法，该算法利用辅助均匀笛卡尔网格来跟踪曲线网格上的接口，并被发现是稳健和计算高效的。 该方法用于解决二维（平面或轴对称）分散多相流动，并已成功应用于几个测试案例，包括振动液滴的经典问题，连续收缩通道中的浮力驱动上升液滴，以及通过绕组通道移动的平面二维液滴的混沌混合。 发现本方法是精确建模复杂几何中分散多相流的可行工具。 18.6 参考书目[l] Beam，R。 M。 \& 温暖，R。 F。 可折信Navier-Stokes 方程的隐式因数分解方案，AIAA J。 16，393（1978）。 [2]布里利，W。 R。 \& 麦当劳，H. 用隐式技术求解树维可压缩Navier-Stokes 方程，讲义物理学35，105，纽约出版社（1974）。 [3] Bringer，M. R.，Gerdts，C. J.，Song，H.，Tice，J. D.，和Ismagilov，R. F。 “依靠液滴混沌混合的化学动力学微流体系统”，Phil。 翻译。 R。 SOC。 隆德。 A，362，1087，（2004）。

416 A 有限-体积/fiont-跟踪方法图18.10：在非维时间框架t* = 0,3.3 x lop3、7.3 x lop3、3.3 x 18 x和22 x分别从左到右拍摄的混合模式的快照。 顶部图是通道中显示的相应散点图（底部图）的放大版本。 （Ca = 0.025，Re = 6.6，X = 1.0，C = 0.76，Re = 6.6，K = 0.33，网格：1024 x 64。） [4] Caughey，D。 A. 欧拉方程的对角线隐式多网格算法-[5] Caughey，D. A. 过去不稳流的隐式多网格计算[6] Che，J. H。 多相流动的数值模拟：电流体动态和液滴凝固，Ph。 D。 论文，密歇根大学，安娜堡，（1999）。 [7] 福奇，L. \& Gueron，S. Eds。 生物流体中的计算建模。 AIAA J.，26，841（1988）。方形横截面的圆柱体，ComputersfYFluids 30，940（2001）。动力学，Springer-Verlag，纽约（2001）。 [8] 韩，J. \& Tryggvason，G。 轴对称液滴的二次分解：I。 由恒定的身体力加速，物理。 流体11（12），3650（1999）。 [9] Hemmat，M。 \& Borhan，A。 周期性收缩毛细血管中水滴和气泡的浮力驱动运动，化学。 英语。 社区。 150，363（1996）。[lo] Hirt，C. W。 \& Nichols，B. D。 自由边界动态的流体体积（VOF）方法，J。 计算机。 物理。 39，201（1981）。[ll] Hirt，C. W.，阿姆斯登，A. A.和Cook，J. L。 所有流速的任意拉格朗日-欧拉计算方法，J。 计算机。 物理。 135，203（1997）。 [12] 羔羊，H. 流体力学，道尔出版社，纽约，（1932年）。

M。 Muradoglu 417 [13] Muradoglu, M., Jenny, P., Pope, S. B.，和Caughey，D. A. 湍流反应流的PDF方程的一致混合有限体积/粒子方法，J。 计算机。 物理。 154（2），342，（1999）。 [14] Muradoglu，M. \& Gokaltun，S. 通过正弦收缩的浮力光滴的隐含多网格计算，ASME J. 应用机械。 71，1，（2004）。 [15] Muradoglu，M。 和Kayalp，A。 D。 “复杂几何中分散多相流的有限体积/前跟踪计算的高效追踪算法”，J. 计算机。 物理，待提交，（2004）。 [16] Muradoglu，M。 \& 斯通，H. A. “通过蜿蜒通道移动的插头中的混沌混合：一项计算研究”物理。 流体，待提交，（2004）。 [17] 奥尔布里希特，W. L。 \& Leal，L。 G。 不可混合液滴通过收敛/发散管的蠕动运动，J。 流体机械。 134，329（1983）。 [18] 奥尔布里希特，W. L。 多孔介质中多相流的孔隙规模原型，Ann。 牧师。 流体机械。 28，187（1996）。 [19] Osher，S。 \& Fedkiw，R. P。 级别设置方法：概述，J。 计算机。 物理。 169（2），463（2001）。 [20] 奥蒂诺，J. M。 “混合运动学”，英国剑桥：剑桥大学。 新闻，（1989）。 [21] Peskin，C. 心脏血流的数值分析，J。 计算机。 物理。 25，220（1977）。 [22] Pozrikidis，C.，编辑。 胶囊和生物细胞的建模和模拟，Chapman \& Hall/CRC（2003）。 [23] Pulliam，T. H。 \& Chaussee，D。 S。 隐式近似因式分解算法的对角线形式。 J。 计算机。 物理。 39，347（1981）。 [24]斯卡多夫利，R。 和Zaleski，S。 自由表面和界面流动的直接数值模拟，Ann。 牧师。 流体机械。 31，567（1999）。 [25] 塞西安，J. A. \& Smereka，P. 流体接口的液位设置方法，Ann。 牧师。 流体机械。 35，341（2003）。 [26] Song, H., Tice, J. D.，和Ismagilov，R. F。 “一个用于及时控制反应网络的微流体系统，”Angew。 化学。 英特。 Ed.，42，768，（2003）。

418 有限-体积/前-跟踪方法[27] Stone，H。 A.，斯特罗克，A. D.，和Ajdary，A. 小型设备的工程流程：微流体学走向芯片实验室，Ann。 牧师。 流体机械。 36，381（2004）。 [28] Sussman，M.，Smereka，P. \& Osher，S. 不可压缩两相流的计算解决方案的水平集方法，J。 计算机。 物理。 144，146（1999）。 [29] Tryggvason，G.，Bunner，B.，Esmaeeli，A，Juric，D.，Al-Rawahi，N.，Tauber，W.，Han，J.，Nas，S.，\& Jan，Y.-J. 用于计算多相流的界面追踪法，J。 计算机。 物理。 169（2），708（2001）。 [30] 图尔克尔，E。 预处理解决不可压缩和低速可压缩流动的方法，J。 计算机。 物理。 72，277（1987）。 [31] Udaykumar，H。 S.，Kan，H.-C.，Shyy，W. \& Tran-Son-Tay，R. 固定笛卡尔网格上任意几何形状的多相动力学，J. Gomput。 物理。 137，366（1997）。 [32] Unverdi，S。 O.，\& Tryggvason，G. 用于黏性不可压缩多相流的界面追踪法，J。 Gomput。 物理。 100，25（1992）。 [33] Ryskin，G. \& Leal，L。 G。 流体力学自由边界问题的数值模拟，第2部分。 气泡通过静止液体的浮力驱动运动，J。 流体机械。 148，19（1984）。 [34] Yabe, T., Xiao, F., \& Utsumi, T. 用于多相分析的约束插值曲线（CIP）方法，J. 计算机。 物理。 169（2），556（2001）。

\subsection{斯托克斯极限的18-A最佳人工可压缩性}
附录18-A 斯托克斯极限中的最佳人工可压缩性 人工可压缩性参数/3应以这样的方式指定，使其在伪时间内将最佳渐近收敛定为稳态。 众所周知，O必须与对流主导流中的速度尺度成正比[30]，即，如果Re >> 1。 这里提出了一个分析，以确定斯托克斯极限中/3的最佳值。 在斯托克斯极限中，流动方程变成v.v = 0，vp-pv2v = 0。 在添加伪时间导数项后，Eqs。 （18.40）变成dV Px + vp-pv2v=o。 然后取动量方程的发散得出dT p-p。 V) + v2p - pV(8. V）=0。 将连续性方程代入方程中。 （18.42），我们得到（18.40）（18.41）（18.42）（18.43）（18.44），其中v = p/p是运动黏度。 假设周期性边界条件，方程的空间傅里叶变换。 （18.44）由d2p dfi - + YX2;t; + = 0，dr2（18.45）给出，其中fi是p的傅里叶变换，x是波数向量，x2 = x.x。 方程（18.45）与质量弹簧阻尼器系统的形式相同。 然后以fi = \&ear的形式寻找解决方案，（18.46）

420 A 有限-体积/fiont-跟踪方法，其中ljo仅是空间的函数，并替换Eq。 （18.46）进入等式。 （18.45）得出Q2 + ux2a + (px)2 = 0的以下特征方程，（18.47）可以求解为（18.48）因为（xp）'总是正的，两个根的实部都是负的，即它总是被迫到稳态。 然而，如果平方根中的项小于或等于零，则可以获得最佳阻尼，这满足了1 P2 2 px)2。 （18.49）现在让长度刻度l指定为l = x-1 max (18.50)，其中xmax是最大波数，那么人工可压缩性参数可以指定为12 21 p2 = a (./e) = uref - 4Rez ' (18.51)，其中Uref是参考速度，Ref = U,,p? /p是基于Uref的雷诺数。 将此表达式与高雷诺数情况下p的最佳值相结合，人工可压缩性参数可以指定为（18.52），其中rcp是一个单位的常数。 方程（18.52）给出了整个雷诺数范围内P的几乎最佳值，尽管在某些情况下，它可能需要调整常数rcp以获得最佳收敛。

\clearpage
\chapter{第19章 湍流火焰的计算建模 Stephen B. 教皇}
\section{19.1 引言}
第19章 湍流火焰的计算建模 Stephen B. Pope' 19.1 导言 湍流燃烧对计算建模构成了严峻的挑战。 根据燃料的不同，可能涉及10、50或1000种化学物种；燃料通过一组复杂的非线性化学反应反应形成燃烧产物（和微量的污染物物种）。 这发生在包含大量长度尺度和时间尺度的湍流中，这使得直接数值模拟在未来几十年内难以解决。 然而，在统计建模方法和相关的计算算法方面取得了良好的进展。 在本文中，我们回顾了湍流反应流PDF方法的最新进展，重点是康奈尔大学关于非预混合湍流燃烧的工作。 在概述PDF方法之后，第19.2节描述了两个火焰的最新计算，然后第19.3节讨论了湍流混合建模的重要问题。 湍流反应流的建模方法[22][11]可以根据两个属性进行广泛分类：第一，流和湍流如何表示；第二，湍流-化学相互作用是如何建模的。 流和湍流的主要方法是[27]：雷诺平均Navier-斯托克斯（RANS）湍流建模；大涡模拟（LES）；和直接数值模拟（DNS）。 目前，RANS是主导方法“纽约伊萨卡康奈尔大学西比利机械和航空航天工程学院14853-7501

\section{19.2 湍流火焰的PDF计算}
422 湍流 火焰用于应用，而LES是许多研究的焦点[30][23]。 虽然DNS是一个强大的研究工具[37]，但其适用范围受到（计算机能力）的严重限制，对于被动流来说，比起惰性流，更是如此。 重要的是要认识到，湍流-化学相互作用需要在RANS和LES中进行建模[25]。 大规模的湍流 mO-tions在动量、热量和物种的运输中起着主导作用，因此这些在LES中由解析场很好地代表。 但在反应性流动中，特别是燃烧中，分子混合和反应的基本过程发生在最小的（子网格）尺度上，因此需要在LES中进行统计建模，就像在RANS中一样。 本文涉及PDF方法[24][10][11]，即通过解流体成分（和其他变量）的联合概率密度函数（PDF）的运输方程来建模湍流-化学相互作用的方法。 PDF方法的主要优势是：它们普遍适用（而不是局限于均匀预混合或双流非预混合问题）；所考虑的流体变量的湍流波动完全通过其联合PDF表示；并且任意复杂和非线性化学反应可以在不近似的情况下进行处理。 RANS设置中使用最广泛的两种PDF方法是组合PDF方法和速度-频率-组合方法。 在前者中，使用了RANS 湍流模型（例如，k-E或雷诺应力），PDF方程中的湍流运输项被建模为梯度扩散。 相比之下，速度、成分和湍流频率的联合PDF的建模传输方程提供了完整的闭包[36]：不需要分离平均动量和湍流模型方程；湍流对流传输是封闭形式，因此避免了梯度扩散近似。 在LES设置中，过滤密度函数（FDF）[25]表示构图的分布（在所有尺度上），从概念上讲，它表示以解析的流域[ll]为条件的PDF。 近年来，基于组合FDF[13]和速度组合FDF[31]开发了LES/FDF组合的组合。 在下一节中，回顾了最近两个非预混合湍流火焰的PDF计算，以说明该方法准确表示有限率湍流-化学相互作用的能力。 然后，在第19.3节中，我们讨论了分子混合建模的现状，这是PDF方法的主要建模问题。 19.2 湍流火焰的PDF计算 在非预混合湍流火焰中，有限速率化学效应是否显著取决于Damkohler数，Da，即字符的比率-

\subsection{19.2.1 引射式喷气火焰}
\subsection{19.2.2 在Vitiated Co-Flow中解除喷气火焰}
S.B. 教皇423角色混合和反应时间尺度。 在高Da下，基于平衡化学或稳定层形火焰的简单模型可以成功（例如，[3]）。 但随着Da的减少，偏离平衡和火焰行为变得明显，局部灭绝并最终发生全球灭绝[2]。 最近的几个实验在Da范围内研究了这种现象，通过改变射流速度[2][8]或通过改变温度，从而改变反应时间尺度来实现。 对于模型来说，在Da范围内重现观察到的行为是一个很好的测试和挑战。 19.2.1 驾驶喷气式火焰巴洛和弗兰克[2]火焰当之无愧地受到了很多关注。 这些实验的价值在于：燃烧器的设计；诊断的质量；以及所涵盖的条件范围。 通过改变喷射和先导火焰的速度，产生一系列减少Da的六种火焰（指定为A、B、...、F）。 火焰D是湍流，局部灭绝很少；火焰F表现出显著的局部灭绝，非常接近全球灭绝；而火焰E介于两者之间。 使用几种不同的方法对火焰D进行了许多计算，但对更具挑战性的火焰E和F[1]的计算相对较少。 Lindstedt等人1描述了Barlow和Frank火焰的速度-成分-频率联合PDF计算。[ 19]，徐和波普[39]和唐，徐和波普[35]。 在后一项作品中，甲烷和氮氧化物化学由19种增强还原机制[34]和EMST混合模型[32]来描述。 [39][35]给出了计算和实验数据（包括条件均值和PDF）之间的详细比较。 在这里，我们只展示了“燃烧指数”BI的结果，这是有限率效应的总体衡量标准。 燃烧指数可以基于不同的物种，并且是局部的：值1对应于完全燃烧，值0对应于完全灭绝。 图19.1显示了基于COZ和CO的燃烧指数，作为三个火焰D、E和F的轴向距离的函数。 可以看出，一般来说，计算准确地代表了在30射流半径（由于局部灭绝）观察到的最小BI，然后是下游恢复（由于重新点燃）。 此外，火焰D、E和F之间的局部灭绝的增加被准确计算。 19.2.2 在Vitiated Co-Flow Cabra等人[4]中提升的喷射火焰通过实验研究了由Hz/N2喷射形成的提升火焰的新配置，该火焰在Vitiated Co-Flow中（在1045K左右）。 假设这种火焰的稳定机制与冷共流中升起的火焰的稳定机制有很大不同。 具体来说，在喷射出口附近混合燃料和被烧焦的共流，导致热、稀薄的混合物随后自动点燃，从而锚定火焰。

湍流 火焰 424 1.5 6-1 cd 0.5 1.5 M n ” 0 50 100 0 50 100 x/R, “/Ri I (4 (b) 图 19.1：基于（a) COz 和 (b) CO 相对火焰 D, E 和 F 轴向距离的燃烧指数。 PDF计算（行）和实验数据（符号）的比较。 在悉尼大学使用该燃烧器的进一步实验中，被烧蚀的共流温度在狭窄的范围内变化（1010K-1045K），导致升力高度在45个射流直径到5个射流直径之间（随着温度的升高）。 这些火焰有几次PDF计算[4][20][la][5]。 从后一项工作中，我们在图中展示了。 19.2（a）使用氢燃烧的两种化学机制计算升降高度：穆勒机制[21]；和李机制[18]，其中更新了一些形成率和焓。 李机制的计算似乎与实验数据非常一致。 然而，重要的是要认识到，被烧焦的共流温度的实验不确定性，加上升降高度对这个数量的显著敏感性，使实验误差条大于计算之间的差异。 可以得出的结论是，两种PDF计算都再现了实验升空高度（在误差条内），并且它们对化学机制的细节表现出敏感性。 图19.2（b）显示了使用穆勒力学和三个最广泛使用的混合模型计算的升空高度，即：与平均值交换相互作用（IEM）模型[38][9]；修改后的Curl（MC）模型[7][14]；和欧几里得最小生成树（EMST）模型[32]。 可以看出，对于这种流程，混合模型的选择没有很大的敏感性。 第19.3节进一步讨论了这个问题。

\section{19.3 湍流混合的建模}
S.B. 波普混合机制（EMST CW.0）dmerenl mixlng mDdels的比较图19.2：硫化共流中氢喷射火焰的升降高度与共流温度（a）Li和Mueller化学机制的比较（b）MC、IEM和EMST混合模型的比较。 19.3 湍流混合建模 在湍流对惰性流动的研究中，有许多标量混合研究，其中主要重点是单个惰性标量（例如，温度过剩）的平均值、方差和导数统计。 湍流燃烧中的混合问题涉及得多。 通常，有20阶的成分；分子混合对PDF形状的影响很重要（不仅仅是方差和协方差的衰变率）；混合率的波动在影响局部灭绝方面很重要；并且（特别是在预混合燃烧中）反应和混合可以强耦合。 三种最常用的混合模型，IEM、MC和EMST，存在已知的理论缺陷，但同时，在某些情况下，它们可以产生定量准确的结果。 目前该领域的研究目标是划定这些不同模型的适用范围。 上面描述的升升火焰的计算说明了三种混合模型都产生相似结果的情况。 相比之下，Ren和Pope [29]研究了一种部分搅拌反应器（PaSR），其中观察到截然不同的行为。 举例来说，图。 19.3显示了氢燃烧案例的三个模型给出的温度与混合物分数的散点图。 除了化学之外，反应堆的特点是停留时间T\textasciitilde{},,,和湍流混合时间T\textasciitilde{}T,\textasciitilde{}\textasciitilde{}。 对于T\textasciitilde{}\textasciitilde{}\textasciitilde{}/T\textasciitilde{}\textasciitilde{}\textasciitilde{}的小值（例如，T\textasciitilde{}\textasciitilde{}\textasciitilde{}T\textasciitilde{}\textasciitilde{}\textasciitilde{} < 1/20），PaSR近似PSR，三个混合模型产生基本相同的结果。 但对于T,\textasciitilde{}\textasciitilde{}/T,,的较大值，标量方差变得显著，如图所示。 19.3，形状

426 湍流 火焰m... I b）I图19.3：使用不同混合模型（a）IEM（b）MC（c）EMST）PaSR中温度与混合分数的散射图。 图19.4：使用不同混合模型的PaSR的化学计量与停留时间的平均温度条件。模型预测的联合PDF可能有很大不同。 对于相同的PaSR测试案例，图。 19.4显示了条件平均温度（化学计量）作为固定rmix/rre的停留时间的函数，= 0.35。 与PSR一样，爆炸发生在rres的临界值下，如图中的星号所示。 19.4. 可以看出，这三种模型的临界停留时间有很大不同，EMST的弹性最大，而IEM的弹性最小。 Subramaniam和Pope [33]在一个截然不同的测试案例中得出了类似的结论。 在湍流混合的背景下，一个有趣的发展是多重映射条件（MMC）方法的开发[17]。 这可以被认为是条件矩闭包[16]和振幅映射闭包[6][26]之间的婚姻。 对于涉及n、物种的湍流反应流，假设（在组成空间中）所有组成都位于维度n，<n的流形上。 （这个假设也被探索了

\section{19.4 致知}
\section{19.5 参考书目}
S.B. 教皇427教皇[28]。） 流形由应用映射闭包的nr个引用变量参数化。 作为对MMC映射闭包方面的基本测试，从对称三三角函数初始条件获得在各向同性湍流中演变的两个标量的联合PDF的分析解决方案。 在联合PDF的演变过程中采用的形状与Juneja和Pope[15]从DNS获得的形状非常一致。 一段时间内，混合模型应该仍然是一个活跃的研究领域，因为它们是RANS和LES方法中PDF方法的关键元素，而且目前的模型有几个被重视的缺点。 还有一些关于现有模型性能的问题需要回答。 在非预混湍流燃烧理论中，熄灭和点火事件分别与斯卡消散的大小值有关。 如何使用没有明确表示标量耗散分布的现有混合模型来准确计算这些现象？ 19.4 致词 很高兴将这篇论文献给David A. Caughey在他六十岁生日之际。 这项工作得到了空军科学研究办公室的资助。 F-49620-00-1-0171，能源部，赠款编号DE-FG02-90ER14128。 1 9.5 参考书目[11 Anon.，湍流非预混合火焰测量和计算国际研讨会。http://wuw.ca.sandia.gov/TNF/abstract.html，2004年。 [2]巴洛，R。 S。 \& 弗兰克，J. H.，Proc。 燃烧。 Inst，27:1087-1095，1998年。 [3] 巴洛，R. S.，史密斯，N. S。 A., Chen, J.-Y., \& Bilger, R. W.，燃烧。 火焰，117:4-31，1999年。 [4] Cabra, R., Myhrvold, T., Chen, J.-Y., Dibble, R. W.，Karpetis，A. N.，\& Barlow，R. S.，Proc。 燃烧。 Inst，29：1881-1888，2002年。 [5] Cao, R., Pope, S. B.，和Masri，A. R.，（准备中），2004年。 [6] Chen, H., Chen, S., \& Kraichnan, R. H.，物理。 牧师。 Lett.，63:2657-2660，[7] Curl，R. L.，AIChE J.，9:175-181，1963年。 1989年。

428 湍流 火焰 [8] Dally，B。 B.，弗莱克彻，D. F.，和Masri，A. R.，燃烧。 理论建模，21193-219。 1998年。 [9]多帕佐，C。 \& O'Brien，E。 E., Acta Astronaut., 1:1239-1266, 1974. [lo] Dopazo, C., In P. A. Libby和F。 A. 威廉姆斯，编辑，湍流反应流，第7章，第375-474页。 学术出版社，伦敦，1994年。[ll] Fox，R. O.，湍流反应流的计算模型。 剑桥大学出版社，纽约，2003年。[la] Gordon，R.，Masri，A. R.，教皇，S. B.，和Goldin，G. M.，Proc。 燃烧。 Inst，30：（已提交），2004年。 [13] Jaberi，F. A.，Colucci，P. J., James, S., Givi, P., \& Pope, S. B.，J. Fluid Mech.，401：85-121，1999年。 [14] Janicka, J., Kolbe, W., \& Kollmann, W., J. 非平衡。 Thermodyn，4\textasciitilde{}47-66，1977年。 [15] Juneja，A. \& 波普，S. B.，物理。 流体，8:2161-2184，1996年。 [16] Klimenko，A. Y。 和比尔格，R。 W.，Prog。 能量燃烧。 科学，25:595-687，1999年。 [17] 克里梅科，A. Y。 \& 波普，S. B.，物理。 流体，15：1907-1925，2003年。 [18] Li, J., Zhao, Z., Kazakov, A., \& Dryer, F. L.，技术报告，宾夕法尼亚州立大学燃烧研究所东部州分部秋季技术会议，宾夕法尼亚州大学公园，2003年。 [19]林德斯泰特，R. P.，卢卢迪，S. A，，\& VBos，E。 M.，Proc。 燃烧。 Inst，28：149-156，2000年。 [20]马斯里，A。 R.，曹，R.，波普，S. B.，和Goldin，G. M.，燃烧。 理论建模，8:l-22，2004年。 [21] 穆勒，M. A.，Kim，T. J.，Yetter，R. A.，\& 烘干机，F. L.，Int. J。 化学。 Kznet.，31：113-125，1999年。 [22] Peters，N.，湍流 燃烧。 剑桥大学出版社，2000年。 [23] Pitsch，H. \& Steiner，H.，物理。 流体，12:2541-2554，2000年。 [24] Pope，S. B.，Prog. 能量燃烧。 科学，11:119-192，1985年。 [25] 教皇，S. B.，Proc。 燃烧。 Inst，23:591-612，1990年。 [26] Pope，S. B.，Theor。 计算机。 Fluid Dyn.，2:255-270，1991年。

S.B. 教皇429 [27] 教皇，S. B.，湍流流量。 剑桥大学出版社，剑桥，2000年。 [28] 教皇，S. B.，流动，湍流和燃烧，（印刷中），2004年。 [29] Ren，Z。 \& 波普，S. B.，燃烧。 火焰，136：208-216，2004年。 [30]桑卡兰，V。 \& Menon, S., Proc. 燃烧。 Inst，28:203-209，2000 [31] Sheikhi，M。 R。 H.，Drozda，T. G.，Givi，P.，\& Pope，S. B.，物理。 流体，15:2321-2337，2003年。 [32] Subramaniam，S. \& 波普，S. B.，燃烧。 火焰，115：487-514，1998年。 [33] Subramaniam，S。 \& 波普，S. B.，燃烧。 火焰，117：732-754，1999年。 [34] Sung，C. J.，法律，C. K.，\& Chen，J.-Y.，Proc. 燃烧。 Inst，27:295-304，1998年。 [35] Tang, Q., Xu, J., \& Pope, S. B.，Proc。 燃烧。 Inst，28:133-139，2000年。 [36] Van Slooten，P. R., Jayesh, \& Pope, S. B.，物理。 流体，10:246-265，1998年。 [37] Vervisch，L. \& 庞索，T.，安。 牧师。 流体。 机械，30：655-691，1998年。 [38] Villermaux，J. \& Devillon，J. C.，第二届化学反应工程国际研讨会论文集，第1-13页，纽约，1972年。 爱思唯尔。 [39]徐，J。 \& 波普，S. B.，燃烧。 火焰，123:281-307，2000年。

此页面故意留空

\clearpage
\part{V。 教育}
第五部分教育

此页面故意留空

\clearpage
\chapter{第20章 教育物理：教学改革对空气动力学的影响 David L. 达尔莫法尔}
\section{20.1 介绍}
第20章 教育适合：教学改革对空气动力学的影响 David L. Darmofal' 20.1 导言 近年来，工程课程改革受到了行业、政府和学术团体的严重关注，因为改变工程教育的需要已经成为一个公认的问题。 在空气动力学方面，重新设计传统课程的必要性至关重要。 计算方法的开发和成熟彻底改变了空气动力学。 与此同时，科学教育研究表明，使用与标准讲座方法不同的教学方法可以显著改善学习。 这些因素严重怀疑传统的空气动力学课程和教学法仍然是下一代航空航天工程师最有效的教育。 本文描述了麻省理工学院五年来改革本科航空动力学教育的努力。 进行空气动力学教育改革的决定不仅受到上述外部力量的刺激，还受到教授该学科的个人经验的刺激。 特别是，我们发现我们的学生处理空气动力学问题的能力有限，这些问题与马萨诸塞州剑桥市麻省理工学院航空航天系02139的具体情况不同。

\section{20.2 课程概述}
课程涵盖的434教育未来。 例如，在修改我们的课程之前的学期，开发了期末考试，以评估学生的能力（1）将概念应用于与学期中遇到的不同情况，以及（2）整合概念并将其应用于一个更复杂、开放式的问题中，即他们作为执业工程师将面临的问题类型。 学生在考试中的表现很差，没有表现出任何能力。 虽然我们认为我们的学生通过我们的教学实现了深刻的概念理解，但他们并没有。 因此，在期末考试中，我们评估了学生没有通过学科教学法发展的好机会的技能。 由于我们强烈认为概念理解和整合概念以解决复杂问题的能力是我们学科的主要目标，我们决心改变我们的教学。 课程改革主要集中在三个问题上，具体来说：1. 应用主动学习来增强概念理解。 2. 将理论、实验和计算技术整合到现代空气动力学课程中。 3. 使用一个学期的空气动力学设计专案来提供教育动机和真实的学习体验。 Darmofal et aZ描述了这项工作的前两年。[ 3]。 本文描述了自最初报告以来已经完善的当前教学法，更重要的是，它包括各种数据，展示了新教学法带来的改进。 20.2 课程概述 空气动力学（M.I.T.科目编号16.100）是本科生可以选择使用的一组高级科目之一，以完成他们的学位要求。 该课程每年秋季学期提供一次，在过去五年中，约有40名学生入学。 在本课程之前，学生会接触一些基本的流体动力学，包括守恒原理、势流动和一些不可压缩的空气动力学，包括薄翼型理论和提升线。 本课程的主要目标是让学生掌握制定和应用适当的空气动力学模型的能力，以估计现实三维配置上的力。 涵盖的主要主题是：2-D/3-D势流（不可压缩到超音速），包括面板和涡流晶格方法；边界层理论，包括过渡和湍流的影响；激波和膨胀扇；2-D/3-D欧拉和Navier-Stokes计算，包括一些基本的湍流建模；以及风洞测试。

\section{20.3 概念理解和主动学习}
D。 L。 Darmofal 435 20.3 概念理解和主动学习 在认知科学中，建构主义的学习模式已经变得流行，因为它解释了大量关于学习和解决问题的实验证据[l3，111。 建构主义学习模型认为，个人通过测试先前经验的概念，将这些概念应用于新情况，并将概念整合到先前的知识中，积极构建他们的知识。 当新的（可能正确的）概念与现有的（可能不正确的）概念发生冲突时，就会出现一种困难的情况。 除非学习者有充分的理由拒绝这些误解，否则这些新概念将难以适应，学习通常是肤浅和短期的。 在物理教育领域，Halloun和Hestenes调查了初步知识对第一门物理课程中学生表现的影响[8]。 为此，他们开发了一个诊断测试来评估学生对牛顿概念的了解。 诊断经过精心设计，包括学生从个人、日常运动经验中经常拥有的误解。 结果表明，初始知识（通过使用诊断的预测试进行评估）是物理课程成绩的主要指标，而特定讲师、学术专业、高中数学、性别和年龄等因素没有影响。 此外，诊断后的测试表现表明，考试的整体表现虽然比测试前要好，但相当差。 因此，虽然学生可以在决定他们成绩的常规课程考试中表现良好，但他们接受的传统指导对改变他们对力学的误解没有什么作用。 预先存在的知识和建构主义学习模式的重要性对传统教学产生了相当大的怀疑。 传统的教学使用传递方法，即假设学生在被动地听讲座的同时获得知识，从而在学生和白板之间产生了类比。 这种教学风格与学习的建构主义观点直接冲突，因为它没有考虑到学生积极面对他们的误解的必要性，这样他们就可以被更先进的理解所取代。 因此，仅仅提高讲座中概念的表述质量并不能产生更大的理解，相反，建构主义的学习模式表明，需要对教育学进行更实质性的改变来解决误解。 加强概念学习的一种策略是一套被称为主动学习的教学方法。 Bonwell \& Eison将主动学习定义为“涉及学生做事和思考他们正在做的事情”的教学策略[l]。 根据这个定义，学生被动聆听所呈现材料的传统讲座不是主动学习

436 教育未来战略。 相反，主动学习需要学生在课堂上对材料的某种参与。 利用Halloun和Hestenes开发的物理考试，Hake在一项对来自各种大学（以及一些高中）校园的六千名学生的研究中表明，主动学习方法与传统讲座相比，学习收益有统计学上显著的改善[6]。 此外，教育研究表明，这种主动学习也可以增加信心、对学科的享受和人际交往技能[l2]。 在1999年秋季学期，我们开始使用同伴指导，这是一种由Mazur[9，21）在物理学中开发的主动学习方法。 在这种方法中，概念性问题（由Mazur称为ConcepTests）在课堂上给学生，并有时间进行个人反思。 在检查了对问题的理解程度后，可以进行小组讨论。 此外]讲师通常会澄清误解，并带领学生进一步探索概念，通常会进行小型讲座。 在典型的课堂上，通常会讨论两到三个概念问题。 有几个选项可以衡量课堂理解。 在16.100中，我们发现使用手持式个人遥控器非常有效。 与举手或抽认卡相比，个人遥控器具有多个优势，包括学生回复的匿名性和高效生成评估数据以分析总体性能。 Hall等人还讨论了在一组大二航空航天工程课程中使用同伴指导的问题。[ 7]为了说明一个典型的概念问题，考虑电梯的生成。 由于关于飞机如何飞行的（通常不准确的）民间传说，以及在以前课程中获得的知识进一步复杂化，翼型的升力发电充满了许多误解。 在讨论升力生成时，使用了一系列概念问题，集中于通过动量变化、流式曲率和反应力来理解升力生成。 第一个问题涉及水射流碰撞在气缸上，如图20.1所示。 尽管许多学生认为喷气式飞机会导致气缸被推离溪流，但实际上]物体会旋转到溪流中。 一个简单的动量平衡直接导致力（升力）产生和动量变化之间的联系。 当我们使用这个问题时，我们包括一个课堂演示，清楚地展示了圆柱体被拉入流中。 这个问题之后是一系列问题，将流动转动的概念与力产生联系起来]并扩展了想法，以理解当翼型不再有效地转动流动时失速时的升力损失。 我们对概念问题的经验表明，学生在上课前必须对材料有一些经验。 否则]讨论概念和误解几乎是不可能的，因为学生在课程之前不太可能接触到很多材料。 为了解决这个问题，在dis-之前给出了阅读作业和分级作业

D。 L。 Damofal水射流 水射流将附着在钟摆上的圆柱体上。 钟摆会向哪个方向摆动？ 1. 进入水中（顺时针方向）2。 远离水（逆时针）3。 信息不足图20.1：在课堂上骂材料的概念问题示例。 课前作业的使用与传统工程教学法相比是一个重大的转变，传统工程教学法是分配作业，只有在课堂上讨论材料后才到期。 课前作业不仅对课堂主动学习的成功至关重要，而且还鼓励学生自学。 此外，通过扫描家庭作业，可以立即发现学生的误解和常见困难，而不仅仅是在讨论材料后几周。 随着学生准备工作的改善，随着教师对学生和学生对学生对学科概念的讨论的增加，教室成为一个更加活跃的环境。 除了改变课堂教学法外，我们还将考试从笔试改为口试。 虽然笔试只能分析纸上显示的信息，即学生思维过程的最终结果，但口试是一种积极的评估，可以更深入地了解学生如何理解和关联概念。 此外，口试适应每个学生。 如果学生卡住或误解了一个问题，教师可以帮助个人。 与浪费的评估机会相反，口试的动态适应性提高了有效评估的可能性。 最后，执业工程师每天都面临着基于基本原则应用理性论证的实时需求。 通过使用口试，这种能力可以直接评估。

\section{20.4 理论、计算和实验的整合}
\section{20.5 基于项目的学习}
438 教育未来 20.4 理论、计算和实验的整合 我们过去的空气动力学课程非常注重理论空气动力学，对实验和计算空气动力学的接触有限。 然而，正如Murman和Rizzi[lO]所建议的那样，“今天的空气动力学工程师需要精通现代CFD方法和工具，并且必须知道如何将它们与理论和实验结合使用，进行空气动力学分析和设计。” 困难在于如何最好地将计算方法整合到主流航空航天课程中。 我们设想，绝大多数航空航天工程师只对CFD计算结果有经验，一些工程师将是CFD的最终用户，极少数将参与CFD开发的某些方面。 因此，我们将计算空气动力学整合到本科课程的一般哲学是，现代空气动力学家必须很好地理解计算工具所体现的基本空气动力学近似，然而，数值方法的细节不太重要。 例如，我们希望学生了解三维、可压缩的欧拉计算可以对激波进行建模，但由于非黏性，当黏性效应可能至关重要时，不适合。 相比之下，我们不希望学生了解什么是二差人工耗散算子，或者通量差分分裂与通量向量分割有何不同。设计项目的使用（在第20.5节中进行了更详细的描述）。 该项目要求学生团队根据所需的操作条件开发经过验证的空气动力学模型。 为此，学生进行计算模拟和风洞测试。 此外，由于学生团队需要减少和更正原始风洞数据，他们开始欣赏风洞测试与纯粹的理论或计算技术一样是一个模型。 在这个过程中，学生们很快了解到，无论是计算方法还是实验都无法为所有应用提供可靠的预测，了解模拟和实验之间的一致性与否是空气动力学家的关键角色。 计算和实验方法通过20.5基于项目的学习整合到课程中。通常，空气动力学和其他高级工程主题的教学重点是理论，但几乎没有机会应用理论，特别是接近现代飞机设计中面临的复杂性的问题。 因此，学生们认为他们正在学习材料，以防他们以后在职业生涯中可能需要它。 在16.100中使用的基于项目的方法中，

\subsection{20.5.1 军用飞机设计项目}
D。 L。 Darmofal 439 知识正在立即应用。 包括基于项目的方法的另一个好处是，它增加了课程中教学技术的丰富性。 这种多样的学习经验已被公认为有效教学的关键原则[4，51。 在过去的五年里，开发了两个设计项目：一个基于军用战斗机，另一个基于混合翼体商用运输机。 这两个项目都有一个初始建模阶段，学生团队为基线配置开发和验证空气动力学模型，然后是设计阶段，模型用于提高空气动力学性能。 我们设计项目实施的一个关键特征是每周的项目工作会议。 这些项目会议的目标是提供一个预定的时间段，课程工作人员（通常是一名教师和一名教学助理）可以在团队开始处理项目时与团队互动。 这些两个小时的课程在一个大型电子教室里举行，大约有25台电脑，或者大约每两名学生一台电脑。 我们发现，这种计算机与学生的比例在促进合作方面是有效的。 在学期开始时，这些项目会议通常用于向学生提供有关项目的信息，澄清要求，并介绍各种计算工具和实验设施。 然而，在本学期的后期，员工的角色更倾向于辅导和故障排除。 学生团队由大约四名学生组成。 每个团队都会提交一份临时和最终的书面报告，这是他们成绩的基础。 对于中期报告，该报告大约在本学期的2/3到期，团队需要充分描述他们开发的所有空气动力学模型，包括他们的验证研究。 最终报告侧重于使用这些经过验证的模型进行设计（除了更正临时报告中发现的任何错误外）。 我们在项目设计阶段发现的一个最佳做法是，在进行任何重新设计之前，要求团队根据他们对空气动力学性能的概念理解，假设哪些设计更改可能会提高他们满足设计要求的能力。 然后，最终设计阶段成为对拟议设计修改是否具有预期效果的研究；如果没有，学生需要解释为什么他们的初始假设不正确。 20.5.1 军用飞机设计项目 1999年夏天，我们联系了几位行业和政府代表，要求一个设计项目，作为我们秋季晚些时候空气动力学课程的整个学期主题。 洛克希德·马丁航空公司（LMAC）提出了一个基于军用飞机行业典型的重新设计场景的项目。 具体来说，学生团队

\subsection{20.5.2 混合翼车身设计项目}
\section{20.6 结果}
\subsection{20.6.1 教育学的有效性}
44 0 教育未来是在几个关键操作条件下开发类似F-16的机翼体几何的空气动力学模型，然后在机翼设计贸易研究中使用这些空气动力学模型。 该项目在1999年秋季、2000年和2001年三个学期使用。 该项目包括从低亚音速到超音速的飞行模式，包括一些高攻角的飞行系统。 感兴趣的绩效指标是：1. 假设攻角在海平面条件下的起飞距离最多限制在25度，以避免尾部撞击地面。 2. 在10K英尺3下巡航0.9马赫的动作半径（即范围）。 在10K英尺的速度下，从0.9马赫到1.2马赫的冲刺时间估计。对于亚音速（即起飞）制度，在M.I.T.的低速隧道中建造并测试了一个1/9比例的风洞模型。 在高速下，实验数据可以从以前的LMAC测试中获得。 该项目的设计阶段侧重于通过引入前缘和后缘襟翼，以及机翼扫荡和跨度的变化来提高起飞、巡航和冲刺性能。 20.5.2 混合翼体设计项目 2002年秋季学期，与波音公司合作开发了一个基于波音混合翼体飞机设计的新设计项目。 该项目的目标是重新设计基线配置，以提高静态稳定性，同时最大限度地减少阻力并保持平衡。 具体来说，考虑了两种飞行条件：超音速巡航和低速进场。 在进近时，前缘和后缘设备被允许处于活动状态，而在巡航中，飞机必须保持清洁。 与战斗机项目一样，进行了低速风洞测试，为空气动力学模型提供验证数据。 20.6 结果 量化教学变革对学习的影响是困难的。 我们的方法是从各种来源获取数据，并从汇总中得出结论。 虽然任何单一来源都是可疑的，但加在一起，结果更确凿。 20.6.1 教学法的有效性 在过去三年（2001-2003年秋季）中，上述教学法几乎保持不变，只是略有调整。 学生评分

D。 L。 该死的阅读和家庭作业在60 5 30 0 20 a无效有效非常有效的讲座项目z a无效有效非常有效图20.2：2000年和2000年后（2001-2003）学期学生对阅读/家庭作业、讲座和项目有效性的评价比较

\subsection{20.6.2 课前作业的影响}
442 教育教育学有效性的未来（特别是阅读/家庭作业、讲座和项目）如图20.2所示。 在2000年后的学期，大多数学生将教学法的各个方面评为非常有效，尽管项目有效性的评价略低于阅读/家庭作业和讲座（对项目有效性的观察也与第20.6.3节中讨论的学生评论一致）。 20.6.2 课前作业的影响 发现使用具有挑战性的课前作业可以显著提高讲座的有效性。 在2000年秋季学期，虽然教学法如上所述，但课前作业旨在鼓励阅读，但不需要对材料进行重大参与。 因此，学生们没有为课堂上的主动学习做好充分的准备。 事实上，2000年秋季学期的学生反馈直接导致了增加家庭作业难度的决定。 家庭作业难度增加的结果是，学生们不仅发现阅读/家庭作业更有效，而且讲座也更有效。 例如，如图20.2所示，从2000年到2000年后学期，将阅读/作业和讲座评为非常有效的学生百分比在统计学上显著增加。 使用具有挑战性的课前家庭作业也对学生的考试成绩产生了有利的影响。 在2000年秋季和2003年学期，进行了笔试。 两个期末考试都由五个问题组成，其中三个是相同的。 剩下的两个问题不同，但难度相似。 三个相同的问题评估了不同的技能，特别是概念、综合和定量能力。 图20.3显示了学生在这三个问题上的表现。 0 概念问题：这个问题侧重于使用不同类型的流动模型（例如 2-D 势流、3-D 势流、3-D欧拉和3-D Navier-Stokes）。 给学生提供了不同的阻力极性曲线和升力曲线，并被要求识别用于生成每个曲线的模型。 如图20.3所示，这两个年对这个问题的表现几乎相同。 0 综合问题：这个问题侧重于对油轮加油臂上的气动力学力进行建模，不仅需要识别重要的物理效应，还需要对这些效应进行定量建模的能力。 虽然最高分的百分比相似，但最低分的百分比（即 0-60\%）从2000年到2003年从40\%左右提高到不到20\%。 0 定量问题：这个问题侧重于使用积分边界层方法来估计管道流中的边界层增长。 这些结果的差异表明，与2000年相比有了实质性的改善

D。 L.. Damnofa1 443 80-100 综合问题 定量问题 70 1 I 2000 图 20.3：2000年和2003年评估概念、综合和定量技能问题的考试表现比较。

\subsection{20.6.3 学生评论}
444 教育未来到2003年。 然而，这种进步并不令人惊讶，因为在2000年，学生进行这种更详细、定量分析的机会较少。 因此，通过有效地将家庭作业（或类似的应用活动）与基于概念的主动学习相结合，学生可以在从定量应用到综合多个概念的技能方面取得高水平的表现。 20.6.3 学生评论 除了教学法的有效性评级外，学生还得到了开放式回答问题，问“课程最好的部分是什么？” 以及“课程如何改进？” 在本节中，我们介绍了其中一些评论，并总结了主要结论。 开放式回答问题表明，学生最初往往对课前作业犹豫不决，但到了学期结束时，他们认识到了这种技术的好处。 一些评论包括：在讲座前做作业很好......使在讲座中实际学习成为可能。 0 教授 Darmofal强迫你通过分配他在讲座中没有涵盖的作业来学习主题材料，因此我必须强迫自己阅读文本并去办公时间。 当他在Pset到期后进行讲座时，我确实更好地吸收了材料。 教学方法非常出色......让我们在p-set之前阅读是一种很好的形式。 0 我最初反对这样的想法，即在我们涵盖这些科目之前，我必须做阅读和做作业。 一旦我过渡了，我意识到这让学习变得容易多了！！ 起初，我对[概念问题]等新技术持怀疑态度，这是关于讲座中没有学过的材料的家庭作业。 最后，它非常顺利。 这是一门课程，我真的觉得我物有所值。 这些评论还加强了课前作业对讲座效果的影响。 开放式回答问题的另一个共同主题是学生对口试的满意。 许多学生发现，与更传统的笔试相比，口试更准确地表达了他们的理解。 事实上，一些学生说口试是课程中最好的部分。 在公开回应评估中对口试的21条评论中，有19条评论是赞成的，只有1条评论是负面的，1条评论建议修改实施。 一些典型的评论是：口试是一种不同的学习评估方法，我非常喜欢。 0 口试是衡量理解的绝佳标准。

\section{20.7 展望}
D。 L。 Damnofal 44 5 口试[是该科目最好的部分]，我认为这些提供了展示你理解内容的好机会。 我真的很喜欢口头决赛的想法。 尽管它很可怕，但它确实表明了你对这门学科的了解，比任何考试都好。 新教学法最具挑战性的方面之一是团队项目的实施。 首先，该项目有多个方面（特别是风洞实验和计算模拟），必须成功管理。 此外，保持十个或更多由四名学生组成的团队有效运作对教师和学生来说可能非常耗时。 过去三年的开放式问题清楚地显示了团队项目的好处和困难。 在此期间，对该项目发表了31条正面评论，只有2条负面评论；然而，29名学生建议需要改进实施。 典型的评论包括：在一个项目中，你必须把你学到的东西直接应用于某事。 这比问题集更有效，因为它的规模更大--虽然在问题集上，你可能只执行一次计算，但一个项目会让你做很多次。 你开始理解为什么以及何时你正在做的事情适用于更深层次、更亲密的层面。 我认为团队项目真的很好。 有一些小点需要解决，并可能更早地解释，但它们真的让我们了解了将理论融入设计中需要哪些要素。 这些项目非常有趣。 学习如何使用计算工具，并看到所有理论和测试如何结合使用以获得准确的结果，非常有用和令人愉快。 我的小组在这个项目上挣扎了一会儿。 最后，我们把所有事情都放在一起了，但很难通过。 现在我回想起来，我不确定我是否会以任何其他方式想要它。 当我在材料上挣扎一段时间时，我学得最好，前提是我有足够的时间最终理解它。 我有足够的时间做这个项目。 学生们已经认识到将他们在课堂上学到的材料应用于一个复杂的问题的教育益处；此外，一些学生（包括此处未显示的其他评论）注意到，该项目使他们能够更好地理解理论和计算如何赞美空气动力学设计中的实验。 然而，有效利用项目仍然是一个具有挑战性的问题。 20.7 展望 为了应对外部和内部的改革力量，我们重新设计了本科空气动力学课程。 改革后的教学法的关键要素包括：

\section{20.8 致谢}
\section{20.9 参考书目}
446 教育未来主动学习与课前作业相结合，以增加概念理解，口试以评估概念理解，0一个为期一个学期的团队项目，强调实验、理论和计算在现代空气动力学设计中的作用。 由学生评估和表现衡量的结果表明，教学效果和学习得到了改善。 此外，充分的学生准备对主动学习的有效性发挥了关键作用。 虽然学生通常对教学法不太传统的方面表示最初的不期待，但他们最终发现这些方法非常有效。 20.8 致谢 过去五年的这项工作涉及麻省理工学院内部和外部的重大互动。 课程的最初重新设计始于与教授的合作。 Earl1 Murman。 军用飞机项目是由迈克·洛夫和丹尼斯·芬利与洛克希德·马丁公司共同开发的。 混合翼体是由鲍勃·利贝克（波音员工和麻省理工学院实践教授）共同开发的。 在2000年秋季学期，16.100与从佐治亚理工学院休假的Steve Ruffin教授共同授课。 这项工作得到了国家科学基金会、麻省理工学院教育创新校友基金和麻省理工学院MacVicar的教师研究员计划的部分支持。 20.9 参考书目[l] Bonwell，C.C. \& Eison，J.A. 主动学习：在课堂中营造兴奋。 ASHE-ERIC高等教育报告，编号 1991年1月。[a] Crouch，C.H. 和马祖尔，E。 同伴指导：十年的经验和成果。 美国物理学杂志，69，2001年。 [3] Darmofal, D.L., Murman, E.M., \& Love, M. 重新设计空气动力学教育，AIAA论文2001-0870，2001年1月。 [4] 弗尔德，R.M. 和西尔弗曼，L.K. 工程教育中的学习和教学风格。 工程教育，78（7），1988年。 [5] Felder，R.M.，风格问题。 ASEE棱镜，6（4），1996年。

D。 L。 达莫法尔 44 7 [6] 哈克，R. R。 互动参与11s。传统方法：六千名学生对介绍物理课程的力学测试数据的调查。 美国物理学杂志，66，1998年。 [7] Hall, S.R., Waitz, I.A., Brodeur, D.R., Soderholm, D.H. and Nasr, R. 在以讲座为基础的工程课上采用主动学习。 第32届ASEE/IEEE教育前沿会议，马萨诸塞州波士顿，2002年。 [8] Halloun，I.A. 和赫斯滕斯，D。 大学物理学生的初始知识状态。 美国物理学杂志，53，1985年。 [9] Mazur，E。 同行指导：用户手册。 新泽西州上马鞍河：普伦蒂斯大厅，1997年。[lo] Murman，E.M. \& Rizzi，A。 将CFD整合到空气动力学教育中。 计算流体动力学 - 2000的前沿。 编辑：D.A. Caughey和M.M. 哈菲兹。 2000年。[ll]国家研究委员会。 人们如何学习：大脑、思想、经验和学校。 国家科学院出版社，华盛顿特区 2000年。 [12]王子，M。 主动学习有用吗？ 研究回顾。 工程教育杂志，93，2004年。 [13] Wankat，P.C. \& Oreovicz，F.S. 教学工程。 纽约，麦格劳-希尔。 可访问http://www.asee.org/publications/teaching.cfm。 1993年。

\backmatter
\chapter*{翻译与版式说明}
全文按原书章节顺序生成，并采用连续的纯中文书籍版式。
\end{document}
